%% titre: Exercices d'algèbre 1
%% matiere: Algèbre 1
%% niveau: Semestre 1/ Licence 1
%% annee: 2026-2027
%% ordre: 2
%% description: Exercices complementaires.
\documentclass[11pt,a4paper]{article}
\usepackage[utf8]{inputenc}
\usepackage[T1]{fontenc}
\usepackage{lmodern}
\usepackage[french]{babel}
\usepackage{amsmath,amssymb,amsthm}
\usepackage{enumitem}
\usepackage{xcolor}
\usepackage{fancyhdr}
\usepackage{comment}
\usepackage[most]{tcolorbox}
\usepackage[a4paper,margin=1.5cm]{geometry}

% ---------- Macros issues des sources (Bibm@th) ----------
\newcommand{\mtn}{\mathbb{N}}\newcommand{\mtns}{\mathbb{N}^*}\newcommand{\mtz}{\mathbb{Z}}\newcommand{\mtr}{\mathbb{R}}\newcommand{\mtk}{\mathbb{K}}\newcommand{\mtq}{\mathbb{Q}}\newcommand{\mtc}{\mathbb{C}}\newcommand{\mch}{\mathcal{H}}\newcommand{\mcp}{\mathcal{P}}\newcommand{\mcb}{\mathcal{B}}\newcommand{\mcl}{\mathcal{L}}
\newcommand{\mcm}{\mathcal{M}}\newcommand{\mcc}{\mathcal{C}}
\newcommand{\mcmn}{\mathcal{M}}\newcommand{\mcmnr}{\mathcal{M}_n(\mtr)}
\newcommand{\mcmnk}{\mathcal{M}_n(\mtk)}\newcommand{\mcsn}{\mathcal{S}_n}
\newcommand{\mcs}{\mathcal{S}}\newcommand{\mcd}{\mathcal{D}}
\newcommand{\mcsns}{\mathcal{S}_n^{++}}\newcommand{\glnk}{GL_n(\mtk)}
\newcommand{\mnr}{\mathcal{M}_n(\mtr)}\DeclareMathOperator{\ch}{ch}
\DeclareMathOperator{\sh}{sh}\let\th\relax\DeclareMathOperator{\th}{th}
\DeclareMathOperator{\vect}{vect}\DeclareMathOperator{\card}{card}
\DeclareMathOperator{\comat}{comat}\DeclareMathOperator{\imv}{Im}
\DeclareMathOperator{\rang}{rg}\DeclareMathOperator{\Fr}{Fr}
\DeclareMathOperator{\diam}{diam}\DeclareMathOperator{\supp}{supp}
\newcommand{\veps}{\varepsilon}\newcommand{\mcu}{\mathcal{U}}
\newcommand{\mcun}{\mcu_n}\newcommand{\dis}{\displaystyle}
\newcommand{\croouv}{[\![}\newcommand{\crofer}{]\!]}
\newcommand{\rab}{\mathcal{R}(a,b)}\newcommand{\pss}[2]{\langle #1,#2\rangle}

% ---------- Macros complémentaires ----------
\newcommand{\N}{\mathbb{N}}\newcommand{\Z}{\mathbb{Z}}\newcommand{\Q}{\mathbb{Q}}
\newcommand{\R}{\mathbb{R}}\newcommand{\C}{\mathbb{C}}
\newcommand{\id}{\mathrm{id}}\newcommand{\pgcd}{\mathrm{pgcd}}\newcommand{\ppcm}{\mathrm{ppcm}}
\newcommand{\Ker}{\mathrm{Ker}\,}

% ---------- Mise en page ----------
\setlist[enumerate,1]{label=\arabic*.,leftmargin=*,itemsep=2pt,topsep=3pt}
\setlist[enumerate,2]{label=\alph*),leftmargin=1.6em,itemsep=1pt,topsep=2pt}
\setlist[enumerate,3]{label=(\roman*),leftmargin=1.6em}
\pagestyle{fancy}\fancyhf{}
\lhead{Fiche de TD --- Bases d'algèbre}\rhead{\thepage}
\setlength{\parindent}{0pt}
% Pour une fiche SANS corrigés : remplacer \corrigestrue par \corrigesfalse
\newif\ifcorriges\corrigestrue
\ifcorriges
\newtcolorbox{corrige}{breakable,enhanced jigsaw,colback=green!4!white,colframe=green!45!black,
	boxrule=0.4pt,left=4pt,right=4pt,top=2pt,bottom=2pt,title={Corrigé},fonttitle=\bfseries\small,
	attach boxed title to top left={yshift=-2mm},boxed title style={size=small,colback=green!45!black}}
\else\excludecomment{corrige}\fi
\newcounter{exo}[section]
\renewcommand{\theexo}{\thesection.\arabic{exo}}
\newenvironment{exo}[1][]{\refstepcounter{exo}\medskip\par\noindent
	\textbf{\textcolor{blue!60!black}{Exercice \theexo}}%
	\if\relax\detokenize{#1}\relax\else\ \textbf{(#1)}\fi\par\nopagebreak\smallskip}{\par}

\title{\textbf{Fiche de TD}\\[2pt] Bases d'algèbre}
\author{}
\date{Année 2026--2027}

\begin{document}
	\maketitle
	\tableofcontents
	\newpage
	\section{Logique et raisonnements}
	
	\subsection*{Quantificateurs}
	
	\begin{exo}[{Vraies ou fausses}]
		Déterminer parmi les propositions suivantes lesquelles sont vraies :
		
		\begin{enumerate}
			
			\item  136 est un multiple de 17 et 2 divise 167.
			
			\item  136 est un multiple de 17 ou 2 divise 167.
			
			\item  $\exists x\in \mathbb R,\ (x+1=0\ \textrm{ et }x+2=0)$.
			
			\item  $(\exists x\in\mathbb R,\ x+1=0)\textrm{ et }(\exists x\in\mathbb R,\ x+2=0)$.
			
			\item  $\forall x\in\mathbb R,\ (x+1\neq 0\textrm{ ou }x+2\neq 0)$.
			
			\item  $\exists x\in\mathbb R^*,\ \forall y\in\mathbb R^*,\ \forall z\in\mathbb R^*,\ z-xy=0$;
			
			\item  $\forall y\in\mathbb R^*,\exists x\in\mathbb R^*,\ \forall z\in\mathbb R^*,\ z-xy=0$;
			
			\item  $\forall y\in\mathbb R^*,\forall z\in\mathbb R^*,\ \exists x\in\mathbb R^*,\ z-xy=0$;
			
			\item  $\exists a\in\mathbb R,\ \forall \veps>0,\ |a|<\veps$;
			
			\item  $\forall \veps>0,\ \exists a\in\mathbb R,\ |a|<\veps$.
			
		\end{enumerate}
	\end{exo}
	\begin{corrige}
		\begin{enumerate}
			
			\item  Cette propositions est fausse, car 2 ne divise pas 167.
			
			\item  Cette proposition est vraie, car 136 est un multiple de 17.
			
			\item  Cette proposition est fausse, car $x$ devrait être simultanément égal à -1 et à -2.
			
			\item  Cette proposition est vraie car $(\exists x\in\mathbb R,\ x+1=0)$ est vraie (il suffit de prendre $x=-1$) et de la même façon $(\exists x\in\mathbb R,\ x+2=0)$ est vraie (il suffit de prendre $x=-2$).
			
			\item  Cette proposition est vraie, par exemple car il s'agit de la négation de la proposition 3, qui est fausse.
			
			\item  Cette assertion est fausse. Si on considère $x$ n'importe quel réel non nul, alors le choix de $y=1$ et de $z=2x$
			fait que $z$ est différent de $xy$.
			
			\item  Cette assertion est fausse. Prenons n'importe quel $y$ dans $\mathbb R^*$. 
			On voudrait trouver $x$ dans $\mathbb R^*$ tel que, pour tout $z$ dans $\mathbb R^*$, on ait $z=xy$. 
			Bien sûr, ce n'est pas possible, car le $x$ que l'on choisit devrait convenir à toute valeur de $z$,
			ce qui n'est pas possible car il suffit de considérer un $z$ différent de $xy$.
			
			\item  Cette assertion est vraie, car on peut choisir $x$ une fois $y$ et $z$ fixés. On choisit alors $x=z/y$.
			
			\item  L'assertion est vraie, il suffit de prendre $a=0$ (convient pour toute valeur de $\veps>0$).
			
			\item  Cette assertion est "évidemment" vraie car elle est plus faible que la précédente (on peut choisir
			$a$ après $\veps>0$).
			
		\end{enumerate}
	\end{corrige}
	
	\begin{exo}[{Nier des assertions avec quantificateurs}]
		Soit $f:\mathbb R\to\mathbb R$ une fonction. Nier les assertions suivantes :
		
		\begin{enumerate}
			
			\item  $\forall x\in \mathbb R,\ f(x)\neq 0$.
			
			\item  $\forall M>0,\ \exists A>0,\ \forall x\geq A,\ f(x)>M$.
			
			\item  $\forall x\in \mathbb R,\ (f(x)>0\implies x\leq 0)$.
			
			\item  $\forall \veps>0,\ \exists \eta>0, \forall (x,y)\in I^2,\ \big(|x-y|\leq \eta\implies |f(x)-f(y)|\leq\veps\big).$
			
		\end{enumerate}
	\end{exo}
	\begin{corrige}
		\begin{enumerate}
			
			\item  $\exists x\in\mathbb R,\ f(x)=0$.
			
			\item   $\exists M>0,\ \forall A>0,\ \exists x\geq A,\ f(x)\leq M$.
			
			\item  $\exists x\in\mathbb R,\ f(x)>0\textrm{ et }x>0$.
			
			\item  $\exists \veps>0,\ \forall \eta>0,\ \exists (x,y)\in I^2, |x-y|\leq \eta\textrm{ et } |f(x)-f(y)|>\veps$.
			
		\end{enumerate}
	\end{corrige}
	
	\begin{exo}[{Vraies ou fausses}]
		Soit $n$ un entier naturel non nul. On note $C_n$ la courbe d'équation $y=(1+x)^n$ et $D_n$ la droite d'équation $y=1+nx$.
		
		\begin{enumerate}
			
			\item  Rappeler l'équation de la tangente à $C_n$ au point $A$ de $C_ n$ d'abscisse 0.
			
			\item  Tracer (par exemple à l'aide d'un logiciel) $C_n$ et $D_n$ lorsque $n=2,3$.
			
			\item  En vous aidant du graphique pour obtenir une conjecture, démontrer si les propositions suivantes sont vraies ou fausses.
			
			\begin{enumerate}
				
				\item  $\forall n\in\mathbb N^*,\ \forall x\in\mathbb R,\ (1+x)^n\geq 1+nx$;
				
				\item  $\forall n\in\mathbb N^*,\ \forall x\in\mathbb R_+,\ (1+x)^n \geq 1+nx$;
				
				\item  $\exists n\in\mathbb N^*,\ \forall x\in\mathbb R,\ (1+x)^n =1+nx$;
				
				\item  $\forall n\in\mathbb N^*,\ \exists x\in\mathbb R,\ (1+x)^n=1+nx$;
				
				\item  $\exists n\in\mathbb N^*,\ \forall x\in\mathbb R^*,\ (1+x)^n>1+nx$.
				
			\end{enumerate}
			
			
		\end{enumerate}
	\end{exo}
	\begin{corrige}
		\begin{enumerate}
			
			\item  D'après le cours de lycée, l'équation de la tangente au point $A$ de coordonnée $(0,1)$ est $y=1+nx$.
			
			\item  
			
			\item  
			
			\begin{enumerate}
				
				\item  C'est faux, par exemple si $n=3$ et $x=-4$.
				
				\item  C'est vrai. Il y a plusieurs façons de démontrer cela. La plus facile est d'étudier la fonction 
				$f(x)=(1+x)^n-(1+nx)$ sur $[0,+\infty[$. On peut aussi utiliser la formule du binôme de Newton.
				
				\item  C'est vrai : $n=1$ convient.
				
				\item  C'est vrai : prendre $x=0$.
				
				\item  C'est vrai : prendre $n=2$, car $(1+x)^2=1+2x+x^2>1+2x$ si $x\neq 0$.
				
			\end{enumerate}
			
			
		\end{enumerate}
	\end{corrige}
	
	\begin{exo}[{Du texte aux quantificateurs}]
		Soit  $f:\mathbb R\to\mathbb R$ une fonction. 
		Exprimer à l'aide de quantificateurs les assertions suivantes :
		
		\begin{enumerate}
			
			\item  $f$ est constante;
			
			\item  $f$ n'est pas constante;
			
			\item  $f$ s'annule;
			
			\item  $f$ est périodique.
			
		\end{enumerate}
	\end{exo}
	\begin{corrige}
		\begin{enumerate}
			
			\item  On peut l'écrire sous la forme : 
			$\exists C\in\mathbb R$, $\forall x\in\mathbb R$, $f(x)=C$; une autre écriture possible est $\forall x,y\in\mathbb R,\ f(x)=f(y)$.
			
			\item  Si on nie l'assertion précédente, on trouve $\forall C\in\mathbb R,$ $ \exists x\in\mathbb R$, $f(x)\neq C$. Si on nie la seconde, on trouve $\exists x,y\in\mathbb R,\ f(x)\neq f(y)$.
			
			\item  $\exists x\in\mathbb R,\ f(x)=0$;
			
			\item  $\exists T\in\mathbb R^*$, $\forall x\in\mathbb R$, $f(x+T)=f(x)$.
			
		\end{enumerate}
	\end{corrige}
	
	\begin{exo}[{Du quantificateur au texte}]
		Soit $f:\mathbb R\to\mathbb R$ une fonction. Énoncer en langage courant les assertions suivantes écrites à l'aide de quantificateurs. Peut-on trouver
		une fonction qui satisfait cette assertion? Qui ne la satisfait pas?
		
		\begin{enumerate}
			
			\item  $\forall x\in \mathbb R,\ \exists y\in \mathbb R,\ f(x)< f(y);$
			
			\item  $\forall x\in\mathbb R,\ \exists T\in\mathbb R,\ f(x)=f(x+T);$
			
			\item  $\forall x\in\mathbb R,\ \exists T\in\mathbb R^*,\ f(x)=f(x+T);$
			
			\item  $\exists x\in\mathbb R,\ \forall y\in\mathbb R,\ y=f(x).$
			
		\end{enumerate}
	\end{exo}
	\begin{corrige}
		\begin{enumerate}
			
			\item  $f$ n'admet pas de maximum. Par exemple, $f(x)=x$ satisfait cette assertion.
			$f(x)=\sin x$ ne satisfait pas cette assertion.
			
			\item  Toute fonction vérifie cette assertion, il suffit de prendre $T=0$.
			
			\item  Toutes les valeurs de $f$ sont prises au moins deux fois. $f(x)=\sin x$ est un exemple de fonctions vérifiant
			cette assertion, $f(x)=x$ un contre-exemple.
			
			\item  Cette assertion est absurde, car elle signifierait qu'il existe un réel $x$ tel que $f(x)$ prenne toutes les valeurs réelles possibles.
			Aucune fonction ne satisfait cette assertion.
			
		\end{enumerate}
	\end{corrige}
	
	\begin{exo}[{Limites de validité d'une proposition}]
		Déterminer les réels $x$ pour lesquels l'assertion suivante est vraie :
		$$\forall y\in[0,1],\ x\geq y\implies x\geq 2y.$$
	\end{exo}
	\begin{corrige}
		On sépare en 4 cas :
		
		\begin{enumerate}
			
			\item  $x\geq 2$. Alors, pour tout $y\in[0,1]$, les propositions $x\geq y$ et $x\geq 2y$ sont vraies.
			D'après la table de vérité de l'implication, l'assertion est vraie.
			
			\item  $x<0$. Alors, pour tout $y\in [0,1]$, les propositions $x\geq y$ et $x\geq 2y$ sont fausses.
			D'après la table de vérité de l'implication, l'assertion est vraie.
			
			\item  $x\in]0,2[$. On peut alors trouver un réel $y\in[0,1]$ tel que $2y>x$ et $x\geq y$. En effet,
			$x/2\in ]0,1[$ et il suffit de choisir $x/2<y<\min(1,x)$. Dans ce cas, $x\geq y$ est vraie et $x\geq 2y$ est fausse.
			L'assertion est fausse.
			
			\item  Si $x=0$, alors les assertions $x\geq y$ et $x\geq 2y$ sont ou bien simultanément fausses
			(lorsque $y\in]0,1]$), ou bien simultanément vraies. L'assertion est donc vraie.
			
		\end{enumerate}
	\end{corrige}
	
	\begin{exo}[{Une proposition sur les fonctions}]
		Soit $f:\mathbb R\to\mathbb R$ une fonction. On considère la proposition $p$ suivante :
		$$p=(\exists t\in\mathbb R,\ \forall x\in\mathbb R,\ f(x)<t).$$
		
		\begin{enumerate}
			
			\item  Écrire la négation de $p$.
			
			\item  Donner un exemple de fonction $f$ qui vérifie $p$; un exemple qui ne vérifie pas $p$.
			
			\item  Parmi les propositions ci-dessous, déterminer celles qui sont équivalentes à $p$, celles qui sont toujours vraies, celles qui sont toujours fausses, et celles pour lesquelles on ne peut rien dire.
			
			\begin{enumerate}
				
				\item  $p_1=(\exists x\in\mathbb R,\ \forall t\in\mathbb R,\ f(t)<x);$
				
				\item  $p_2=(\exists t\in\mathbb R,\ \forall x\in\mathbb R,\ f(t)<x);$
				
				\item  $p_3=(\forall t\in\mathbb R,\ \exists x\in\mathbb R,\ f(x)<t);$
				
				\item  $p_4=(\forall t\in\mathbb R,\ \exists x\in\mathbb R,\ f(t)<x).$
				
			\end{enumerate}
			
			
		\end{enumerate}
	\end{exo}
	\begin{corrige}
		\begin{enumerate}
			
			\item  On a 
			$$\textrm{non }p=(\forall t\in\mathbb R,\ \exists x\in\mathbb R,\ f(x)\geq t).$$
			
			\item  La proposition $p$ signifie que la fonction $f$ est majorée. Elle est vérifiée par exemple si $f(x)=\sin(x)$. Il suffit de choisir $t=2$. Elle n'est pas vérifiée par la fonction $f(x)=x$.
			Sinon, il existerait $t\in\mathbb R$ tel que, pour tout $x\in\mathbb R$, $x<t$. Cela ne fonctionne pas avec $x=t+1$.
			
			\item  
			
			\begin{enumerate}
				
				\item  $p_1$ est équivalente à $p$ : en effet, les variables $x$ et $t$ sont muettes, on peut les échanger!
				
				\item  $p_2$ est toujours fausse : pour n'importe quel $t\in\mathbb R$, si on choisit $x=f(t)-1$, on n'a pas $f(t)<x$.
				
				\item  $p_3$ n'est ni toujours vraie, ni toujours fausse. Elle est vraie pour $f(x)=x$ : pour tout $t\in\mathbb R$, si on choisit $x=t-1$, on a bien $f(x)<t$. 
				Elle est fausse pour $f(x)=x^2$ : si on choisit $t=-1$, il est impossible de trouver $x\in\mathbb R$ avec $x^2< -1$. De plus, $p_3$ n'est pas équivalente à $p$ car la fonction $f(x)=x$ vérifie $p_3$ mais pas $p$.
				
				\item  $p_4$ est toujours vraie : pour n'importe quel $t\in\mathbb R$, si je choisis $x=f(t)+1$, j'ai bien $f(t)<x$.
				
			\end{enumerate}
			
			
		\end{enumerate}
	\end{corrige}
	
	\subsection*{Raisonnement par l'absurde}
	
	\begin{exo}[{Corps de nombres}]
		On rappelle que $\sqrt 2$ est un nombre irrationnel.
		
		\begin{enumerate}
			
			\item  Démontrer que si $a$ et $b$ sont deux entiers relatifs tels que $a+b\sqrt 2=0$, alors $a=b=0$.
			
			\item  En déduire que si $m,n,p$ et $q$ sont des entiers relatifs, alors 
			$$m+n\sqrt 2=p+q\sqrt 2\iff (m=p\textrm{ et }n=q).$$
			
		\end{enumerate}
	\end{exo}
	\begin{corrige}
		\begin{enumerate}
			
			\item  Raisonnons par l'absurde. Supposons que $a+b\sqrt 2=0$ sans que $a=b=0$. Alors, nécessairement $b\neq 0$ car si $b=0$ alors on devrait aussi avoir $a=0$, ce qui est contraire à l'hypothèse $(a,b)\neq (0,0)$. Mais alors, on a $\sqrt 2=\frac {-a}b\in\mathbb Q$ ce qui est faux. L'hypothèse de départ est donc fausse, et on a $a=b=0$.
			
			\item  D'une part, si $m=p$ et $n=q$, alors $m+n\sqrt 2=p+q\sqrt 2$. Réciproquement, supposons que $m+n\sqrt 2=p+q\sqrt 2$. Alors on a aussi
			$$(m-p)+(n-q)\sqrt 2=0.$$
			Le résultat de la question précédent nous dit que $m-p=0$ et que $n-q=0$, soit encore $m=p$ et $n=q$.
			
		\end{enumerate}
	\end{corrige}
	
	\begin{exo}[{Principe des tiroirs}]
		Démontrer que si vous rangez $(n+1)$ paires de chaussettes dans $n$ tiroirs distincts, alors il
		y a au moins un tiroir contenant au moins $2$ paires de chaussettes.
	\end{exo}
	\begin{corrige}
		On va faire un raisonnement par l'absurde. Supposons que tous les tiroirs contiennent au plus une paire de chaussettes. Alors il y aura au plus 
		$1+1+\dots+1=n$ paires de chaussettes, ce qui contredit qu'il y en a $(n+1)$. Donc un tiroir contient au moins deux paires de chaussettes.
	\end{corrige}
	
	\begin{exo}[{Nombres dans un intervalle}]
		Soit $n\geq 1$ un entier naturel. On se donne $n+1$ réels $x_0,x_1,\dots,x_n$ de $[0,1]$
		vérifiant $0\leq x_0\leq x_1\leq\dots\leq x_n\leq 1$. On veut démontrer par l'absurde la
		propriété suivante : il y a deux de ces réels dont la distance est inférieure ou égale à $1/n$.
		
		\begin{enumerate}
			
			\item  Ecrire à l'aide de quantificateurs et des valeurs $x_i-x_{i-1}$ une formule logique équivalente
			à la propriété.
			
			\item  Ecrire la négation de cette formule logique.
			
			\item  Rédiger une démonstration par l'absurde de la propriété (on pourra montrer que
			$x_n-x_0>1$).
			
			\item  Donnez-en une preuve en utilisant le principe des tiroirs.
			
		\end{enumerate}
	\end{exo}
	\begin{corrige}
		\begin{enumerate}
			
			\item  La traduction immédiate de la propriété est :
			$$\exists(i,j)\in\{0,\dots,n\}^2,\ (i<j)\textrm{ et } (x_j-x_i\leq 1/n).$$
			Pour utiliser simplement les valeurs $x_i-x_{i-1}$, il suffit
			de remarquer que si deux nombres sont distants de moins de $1/n$, alors
			il y aura deux nombres consécutifs qui seront distants de moins de $1/n$.
			Ceci signifie :
			$$\exists i\in\{1,\dots,n\},\ x_i-x_{i-1}\leq 1/n.$$
			
			\item  Le contraire de cette assertion est donc :
			$$\forall i\in\{1,\dots,n\},\ x_i-x_{i-1}>1/n.$$
			
			\item  Supposons la propriété de l'exercice fausse, c'est-à-dire que
			$$\forall i\in\{1,\dots,n\},\ x_i-x_{i-1}>1/n.$$
			En écrivant 
			$$x_n-x_0=(x_n-x_{n-1})+(x_{n-1}-x_{n-2})+\dots+(x_2-x_1)+(x_1-x_0)$$
			on a donc, en utilisant la propriété précédente :
			$$x_n-x_0>n\times \frac{1}{n}=1.$$
			C'est absurde. La propriété initiale est donc vraie.
			
			\item  On considère les $n$ intervalles $I_1=[0,1/n[$,
			$I_2=[1/n,2/n[$,...,$I_n=[(n-1)/n,1]$ (qui correspondent si
			l'on veut à $n$ tiroirs). On veut y placer $(n+1)$ points
			(les chaussettes!). Il faut obligatoirement qu'il y ait deux points
			dans le même intervalle. Ces deux points seront distants de moins de $1/n$.
			Bien sûr, le principe des tiroirs se démontre aussi à l'aide d'un raisonnement par l'absurde.
			
		\end{enumerate}
	\end{corrige}
	
	\subsection*{Raisonnement par contraposée}
	
	\begin{exo}[{Somme de rationnels et d'irrationnels}]
		Soit $a$ et $b$ deux réels. On considère la proposition suivante : si $a+b$ est irrationnel, alors $a$ ou $b$ sont irrationnels.
		
		\begin{enumerate}
			
			\item  Quelle est la contraposée de cette proposition?
			
			\item  Démontrer la proposition.
			
			\item  Est-ce que la réciproque de cette proposition est toujours vraie?
			
		\end{enumerate}
	\end{exo}
	\begin{corrige}
		\begin{enumerate}
			
			\item  La contraposée de cette proposition est : si $a$ et $b$ sont rationnels, alors $a+b$ est rationnel.
			
			\item  On va démontrer la contraposée, ce qui va prouver la proposition. On suppose donc que $a$ et $b$ sont rationnels. Ils s'écrivent alors $a=\frac pq$ et $b=\frac{p'}{q'}$, avec $p,p',q,q'$ des entiers relatifs, et $q,q'$ non nuls. On calcule $a+b$ en réduisant au même dénominateur : 
			$$a+b=\frac pq+\frac{p'}{q'}=\frac{pq'+p'q}{qq'}.$$
			Le nombre $a+b$ est le quotient de deux entiers, c'est un rationnel. 
			
			\item  La réciproque de la proposition est : si $a$ ou $b$ sont irrationnels, alors $a+b$ est irrationnel. Ce n'est pas toujours le cas. Par exemple, si $a=\sqrt 2$ et $b=-\sqrt 2$, alors $a$ ou $b$ sont irrationnels (en réalité, on a même $a$ et $b$ qui sont irrationnels), et pourtant $a+b$ n'est pas irrationnel. On prendra bien garde ici à la signification du "ou" en mathématiques.
			
		\end{enumerate}
	\end{corrige}
	
	\begin{exo}[{Pair/impair}]
		Soit $n$ un entier. Énoncer et démontrer la contraposée de la proposition suivante :
		
		\begin{quote}
			
			Si $n^2$ est impair, alors $n$ est impair.
			
		\end{quote}
		
		A-t-on démontré la proposition initiale?
	\end{exo}
	\begin{corrige}
		La contraposée de la proposition est : si $n$ est pair, alors $n^2$ est pair.
		Démontrons cela. Si $n$ est pair, alors il s'écrit $2k$ où $k$ est un autre entier.
		Mais alors $n^2$ s'écrit $(2k)^2=4k^2=2(2k^2)$ et est donc pair. Par le principe
		de contraposition, on a démontré la proposition de l'énoncé.
	\end{corrige}
	
	\begin{exo}[{Divisibilité par 8}]
		Le but de cet exercice est de démontrer par contraposition la propriété suivante, pour $n\in\mtn^*$ :
		
		\begin{quote}
			
			Si l'entier $(n^2-1)$ n'est pas divisible par 8, alors l'entier $n$ est pair.
			
		\end{quote}
		
		
		\begin{enumerate}
			
			\item  Ecrire la contraposée de la proposition précédente.
			
			\item  En remarquant qu'un entier impair $n$ s'écrit sous la forme $n=4k+r$ avec $k\in\mtn$ et $r\in\{1,3\}$
			(à justifier), prouver la contraposée.
			
			\item  A-t-on démontré la propriété de l'énoncé?
			
		\end{enumerate}
	\end{exo}
	\begin{corrige}
		\begin{enumerate}
			
			\item  Si $n$ est impair, alors $n^2-1$ est divisible par 8.
			
			\item  Prenons $n$ un entier impair. $n$ s'écrit donc $2l+1$
			où $l$ est un entier. Si $l$ est pair, $l=2k$ et donc $n=4k+1$.
			Si $l$ est impair, $l=2k+1$ et donc $n=4k+3$. Dans tous les cas,
			on a donc $n=4k+r$ avec $k\in\mtn$ et $r\in\{1,3\}$. On passe au carré :
			$$n^2-1=(4k+r)^2-1=16k^2+8kr+r^2-1=8(2k^2+kr)+r^2-1.$$
			Or, si $r=1$, $r^2-1=0$ et $n^2-1$ est divisible par 8.
			Si $r=3$, $r^2-1=8$ et $n^2-1$ est aussi divisible par 8.
			
			\item  Par le principe de contraposition, oui!
			
		\end{enumerate}
	\end{corrige}
	
	\subsection*{Raisonnement par récurrence}
	
	\begin{exo}[{Pour se mettre en confiance...}]
		Démontrer que, pour tout $n\in\mathbb N^*$, on a $2^{n-1}\leq n!\leq n^n$.
	\end{exo}
	\begin{corrige}
		On va procéder par récurrence sur $n\in\mathbb N^*$.
		On considère la propriété $\mathcal P(n)="2^{n-1}\leq n!\leq n^n"$. $\mathcal P(1)$
		est vérifiée, puisque $2^0=1!=1^1=1$. Supposons maintenant que $\mathcal P(n)$ est vraie
		pour un certain $n\in\mathbb N^*$ et prouvons que $\mathcal P(n+1)$ est vraie. 
		On commence par prouver l'inégalité de gauche. On remarque d'abord que, puisque 
		$n\geq 1$, on a $2\leq n+1$ d'où l'on déduit
		$$2^n=2\times 2^{n-1}\leq 2\times n!\leq (n+1)\times n!\leq (n+1)!$$
		Pour l'inégalité de droite, on part de l'hypothèse de récurrence et on multiplie par $(n+1)$, pour obtenir 
		$$(n+1)!\leq (n+1)n^n.$$
		Or, $n\leq n+1$ et donc $n^n\leq (n+1)^n$. On en déduit que 
		$$(n+1)!\leq (n+1)\times (n+1)^n=(n+1)^{n+1}.$$
		$\mathcal P(n+1)$ est donc vérifiée, ce qui prouve, par récurrence, l'inégalité voulue pour tout $n\in\mathbb N^*$.
	\end{corrige}
	
	\begin{exo}[{Limite de validité}]
		Pour $n\in\mtn$, on considère la propriété suivante :
		$$P_n:\ 2^n>n^2.$$
		
		\begin{enumerate}
			
			\item  Montrer que l'implication $P_n\implies P_{n+1}$ est vraie pour $n\geq 3$.
			
			\item  Pour quelles valeurs de $n$ la propriété $P_n$ est vraie?
			
		\end{enumerate}
	\end{exo}
	\begin{corrige}
		\begin{enumerate}
			
			\item  Soit $n\geq 3$ tel que $P_n$ est vraie. On a alors 
			$$2^{n+1}=2\times 2^n> 2n^2.$$
			Il suffit donc de montrer que $2n^2\geq (n+1)^2\iff n^2-2n-1\geq 0$.
			On calcule le discriminant de ce dernier polynôme, on fait son tableau
			de signes et on s'aperçoit qu'il est toujours positif pour $n\geq 3$.
			Ceci montre bien $P_{n+1}$.
			
			\item  Il faut faire attention, car $P_3$ est fausse!
			On ne peut donc pas en déduire que $P_n$ est vraie pour $n\geq 3$.
			D'ailleurs, $P_4$ est fausse, mais $P_5$ est vraie. Par le principe de récurrence,
			$P_n$ est vraie pour $n\geq 5$. Pour les premières valeurs de $n$, on constate
			également que $P_0$ et $P_1$ sont vraies.
			
		\end{enumerate}
	\end{corrige}
	
	\begin{exo}[{Plusieurs paramètres?}]
		On souhaite démontrer par récurrence que pour tout entier naturel $n$ et pour tout réel $x>-1$,
		on a $(1+x)^n\geq 1+nx$. 
		
		\begin{enumerate}
			
			\item  La récurrence porte-t-elle sur $n$? Sur $x$? Sur les deux?
			
			\item  Énoncer l'hypothèse de récurrence.
			
			\item  Vérifier que $(1+nx)(1+x)=1+(n+1)x+nx^2$.
			
			\item  Rédiger la démonstration.
			
		\end{enumerate}
	\end{exo}
	\begin{corrige}
		\begin{enumerate}
			
			\item  La récurrence ne peut porter que sur un seul paramètre, entier de surcroît. Elle porte donc
			sur $n$.
			
			\item  $P_n :$ Pour tout $x>-1$, $(1+x)^n\geq 1+nx$.
			
			\item  Vérification immédiate en développant.
			
			\item  Initialisation : Montrons que $P_0$ est vraie. Soit $x>-1$. On a 
			$$(1+x)^0=1 \geq 1+0x$$
			ce qui prouve bien $P_0$.\\
			
			Hérédité : Soit $n\in\mathbb N$ telle que $P_n$ vraie et prouvons $P_{n+1}$. Soit $x>-1$.
			On a
			$$(1+x)^{n+1}=(1+x)^n(1+x)\geq (1+nx)(1+x)$$
			cette dernière inégalité provenant du fait que $P_n$ est vraie et que $1+x>0$ puisque $x>-1$. On en déduit que 
			$$(1+x)^{n+1}\geq 1+(n+1)x+nx^2\geq 1+(n+1)x$$
			puisque $nx^2\geq 0$.\\
			
			On peut en fait rédiger cette récurrence de deux façons différentes : comme on l'a fait ci-dessus, en mettant le quantificateur portant sur $x$ dans l'hypothèse de récurrence $P_n$. Il faut alors bien prendre garde que dans la phase d'initialisation et celle d'héridité, il faut commencer par "soit $x>-1$". On peut aussi choisir de fixer $x>-1$ avant la récurrence, et de travailler avec ce $x>-1$ fixé. L'hypothèse de récurrence devient $Q_n: (1+x)^n\geq 1+nx$, et dans les phases d'initialisation et d'hérédité, il faut supprimer la phrase "Soit $x>-1$."
			
		\end{enumerate}
	\end{corrige}
	
	\begin{exo}[{Décomposition du nombre 1}]
		Démontrer que, pour tout entier $n\geq 3$, il existe $n$ entiers strictement positifs $x_1,\dots,x_n$, deux à deux distincts,
		tels que 
		$$\frac1{x_1}+\cdots+\frac1{x_n}=1.$$
	\end{exo}
	\begin{corrige}
		Pour $n\geq 3$, notons $\mathcal P(n)$ la propriété suivante : ``Il existe $n$ entiers strictement positifs $x_1,\dots,x_n$, deux à deux distincts,
		tels que $\frac1{x_1}+\cdots+\frac1{x_n}=1$''. Prouvons par récurrence que pour tout $n\geq 3$, $\mathcal P(n)$ est vraie.\\
		
		Initialisation : On a 
		$$1=\frac12+\frac 13+\frac 16$$
		donc $\mathcal P(3)$ est vraie.\\
		
		Hérédité : Soit $n\geq 3$ tel que $\mathcal P(n)$ est vraie et prouvons $\mathcal P(n+1)$. Il existe des entiers strictement positifs et deux à deux distincts
		$x_1,\dots,x_n$ tels que 
		$$\frac1{x_1}+\cdots+\frac1{x_n}=1.$$
		On peut toujours supposer que $x_1<x_2<\dots<x_n$. Il est alors facile de remarquer que l'on a forcément $x_1\geq 2$ (puisque $\frac{1}{x_2}+\cdots+\frac1{x_n}>0$, $x_1$ ne peut pas
		être égal à $1$). On écrit alors
		\begin{eqnarray*}
			1&=&\frac 12+\frac 12\\
			&=&\frac 12+\frac12\times\left(\frac1{x_1}+\cdots+\frac1{x_n}\right)\\
			&=&\frac 12+\frac{1}{2x_1}+\cdots+\frac{1}{2x_n}\\
			&=&\frac 1{y_1}+\frac{1}{y_2}+\cdots+\frac{1}{y_{n+1}}
		\end{eqnarray*}
		en ayant posé $y_1=2$ et $y_k=2x_{k-1}$ pour $k=2,\dots,n+1$. On a alors
		$$y_1=2<4\leq 2x_1=y_2<y_3<\cdots<y_{n+1}.$$
		La propriété est donc démontrée au rang $n+1$.\\
		
		Conclusion : Par le principe de récurrence, pour tout $n\geq 3$, $\mathcal P(n)$ est vraie.
	\end{corrige}
	
	\begin{exo}[{Une récurrence double}]
		Soit $(u_n)_{n\in\mathbb N}$ la suite définie par $u_0=2$, $u_1=3$ et, pour tout $n\in\mathbb N$, $u_{n+2}=3u_{n+1}-2u_n$. Démontrer que, pour tout $n\in\mathbb N$, $u_{n}=1+2^n$.
	\end{exo}
	\begin{corrige}
		On va démontrer ce résultat par récurrence. Comme la formule de récurrence porte sur deux termes de la suite, nous allons effectuer une récurrence double. Pour $n\in\mathbb N$, on note $\mathcal P(n)$ la propriété $"u_n=1+2^n"$.\\
		
		Initialisation : On a $u_0=2=1+2^0$, donc $\mathcal P(0)$ est vraie. On a aussi $u_1=3=1+2^1$ donc $\mathcal P(1)$ est vraie.\\
		
		Hérédité : Soit $n\in\mathbb N$ tel que $\mathcal P(n)$ et $\mathcal P(n+1)$ sont vraies. Alors on a 
		$$u_{n+2}=3u_{n+1}-2u_n.$$
		Puisque $\mathcal P(n+1)$ et $\mathcal P(n)$ sont vraies, on a
		\begin{align*}
			u_{n+2}&=3(1+2^{n+1})-2(1+2^n)\\
			&=3+3\cdot 2^{n+1}-2-2\cdot 2^n\\
			&=1+6\cdot 2^n-2\cdot 2^n\\
			&=1+4\cdot 2^n\\
			&=1+2^{n+2}.
		\end{align*}
		Donc $\mathcal P(n+2)$ est vraie.\\
		
		En conclusion, par le principe de récurrence double, $\mathcal P(n)$ est vraie pour tout entier naturel $n$.
	\end{corrige}
	
	\begin{exo}[{Simple ou double?}]
		On considère la suite $(u_n)$ (suite de Fibonacci) définie par $u_0=u_1=1$ et, pour tout $n\geq 0$, $u_{n+2}=u_n+u_{n+1}$. Démontrer que la suite $(u_n)$ vérifie les propriétés suivantes :
		
		\begin{enumerate}
			
			\item  pour tout $n\in\mathbb N$, $u_n\geq n$;
			
			\item  pour tout $n\in\mathbb N$, $u_n u_{n+2}-u_{n+1}^2=(-1)^n$.
			
		\end{enumerate}
		
		Avez-vous utilisé une récurrence simple ou une récurrence double?
	\end{exo}
	\begin{corrige}
		\begin{enumerate}
			
			\item  On va procéder par récurrence double. Pour $n\in\mathbb N$, on note $\mathcal P(n)$ la propriété "$u_n\geq n$".\\
			
			Initialisation : $u_0=1\geq 0$ et $u_1=1\geq 1$ : les propriétés $\mathcal P(0)$ et $\mathcal P(1)$ sont vérifiées.\\
			
			Hérédité : Soit $n\in\mathbb N$ et supposons que $\mathcal P(n)$ et $\mathcal P(n+1)$ sont vérifiées. Prouvons alors que $\mathcal P(n+2)$ est vraie. Pour cela, on remarque que 
			$$u_{n+2}=u_{n+1}+u_n\geq n+1+n\geq n+1+1=n+2$$
			à condition que $n\geq 1$ et non $n\in\mathbb N$!!! Pour ne pas faire d'erreurs, il faut donc ici initialiser la propriété jusque $\mathcal P(2)$, puis faire l'étape d'hérédité à partir de $n\geq 1$. Reprenons donc une rédaction correcte.\\
			
			Initialisation : $u_0=1\geq 0$, $u_1=1\geq 1$, $u_2=1+1=2\geq 2$ : les propriétés $\mathcal P(0)$, $\mathcal P(1)$ et $\mathcal P(2)$ sont vérifiées.\\
			
			Hérédité : Soit $n\geq 1$ et supposons que $\mathcal P(n)$ et $\mathcal P(n+1)$ sont vérifiées. Prouvons alors que $\mathcal P(n+2)$ est vraie. Pour cela, on remarque que 
			$$u_{n+2}=u_{n+1}+u_n\geq n+1+n\geq n+1+1=n+2.$$
			Ainsi, $\mathcal P(n+2)$ est vérifiée.\\
			
			Conclusion : par le principe de récurrence double, $\mathcal P(n)$ est vraie pour tout entier $n\in\mathbb N$.
			
			\item  Ici, une récurrence simple suffit! Notons, pour $n\in\mathbb N$, 
			$$\mathcal Q(n):"u_n u_{n+2}-u_{n+1}^2=(-1)^n".$$
			Initialisation : 
			$$u_0u_2-u_1^2=2-1=1=(-1)^0$$
			donc $\mathcal Q(0)$ est vraie.\\
			
			Hérédité : Soit $n\in\mathbb N$ tel que $\mathcal Q(n)$ est vraie,et prouvons $\mathcal Q(n+1)$. On a, en remplaçant $u_{n+3}$ et un des termes de $u_{n+2}$ par leur expression utilisant la formule de récurrence de $(u_n)$,
			\begin{eqnarray*}
				u_{n+1}u_{n+3}-u_{n+2}^2&=&u_{n+1}(u_{n+2}+u_{n+1})-u_{n+2}(u_{n+1}+u_n)\\
				&=&u_{n+1}^2-u_nu_{n+2}\\
				&=&-(-1)^n=(-1)^{n+1}
			\end{eqnarray*}
			La propriété est donc vraie au rang $n+1$.\\
			
			Par le principe de récurrence, $\mathcal Q(n)$ est vraie pour tout entier $n\in\mathbb N$.
			
		\end{enumerate}
	\end{corrige}
	
	\begin{exo}[{Une décomposition des entiers}]
		Démontrer que tout entier $n\in\mathbb N^*$ peut s'écrire de façon unique 
		sous la forme $n=2^p(2q+1)$ où $(p,q)\in\mathbb N^2$.
	\end{exo}
	\begin{corrige}
		On va prouver l'existence par récurrence forte. Le résultat est en effet vrai pour $n=1$,
		il suffit de prendre $p=q=0$. Supposons maintenant que tout entier inférieur ou égal à $n-1$ vérifie cette propriété,
		et prouvons qu'elle est encore vérifiée pour $n$. On distingue alors deux cas :
		
		\begin{enumerate}
			
			\item  ou bien $n$ est pair : dans ce cas, $n=2m$, et $m\in\{1,\dots,n-1\}$. On peut lui appliquer l'hypothèse de
			récurrence, et donc $m$ s'écrit $m=2^a(2b+1)$, avec $a,b\in\mathbb N$. Mais alors $n=2^{a+1}(2b+1)$, et
			donc l'existence est démontrée avec $p=a+1$ et $q=b$. 
			
			\item  ou bien $n$ est impair : alors $n=2m+1=2^0(2m+1)$, et la proposition est démontrée pour $p=0$ et $q=m$.
			
		\end{enumerate}
		
		Ainsi, l'existence de la décomposition est établie pour tout entier $n\in\mathbb N^*$. \\
		
		Prouvons désormais l'unicité. Supposons que $n=2^p(2q+1)=2^a(2b+1)$. Si $p\geq a$, alors 
		$2^{p-a}(2q+1)=2b+1$ et le terme de droite est impair. Celui de gauche l'est aussi, et donc $p=a$.
		On en déduit donc que $q=b$ également, ce qui prouve l'unicité.
	\end{corrige}
	
	\begin{exo}[{Une décomposition des entiers}]
		Démontrer que tout entier $n\geq 1$ peut s'écrire comme somme de puissances de 2 toutes distinctes.
	\end{exo}
	\begin{corrige}
		On va procéder par récurrence forte sur $n$. La propriété est vraie pour $n=1$, puisqu'on peut alors écrire $1=2^0$.
		Soit $n\geq 2$ et supposons la propriété vraie pour tous les entiers inférieurs stricts à $n$. Si $n$ est pair, alors $n=2m$ et 
		$m$ s'écrit comme somme de puissances de 2 toutes distinctes, $m=2^{p_1}+\dots+2^{p_r}$. Alors 
		$n=2^{p_1+1}+\dots+2^{p_r+1}$ s'écrit aussi comme somme de puissances de 2 toutes distinctes. Si $n$ est impair,
		alors $n-1=2m$ est pair. $m$ s'écrit à nouveau $m=2^{p_1}+\dots+2^{p_r}$, où les $p_i$ sont tous différents. Mais alors,
		$$n=1+2m=2^0+2^{p_1+1}+\dots+2^{p_r+1}$$
		s'écrit bien comme somme de puissances de 2 toutes distinctes.
	\end{corrige}
	
	\subsection*{Raisonnement par disjonction de cas}
	
	\begin{exo}[{Une inégalité}]
		Démontrer que, pour tout $x\in\mathbb R$, $|x-1|\leq x^2-x+1$.
	\end{exo}
	\begin{corrige}
		La difficulté de l'exercice vient du fait que la valeur absolue a deux expressions distinctes suivant le signe de la quantité à l'intérieur. Ceci nous incite à raisonner par disjonction de cas. 
		
		\begin{itemize}
			
			\item  Si $x\leq 1$, alors $x-1\leq 0$ et $|x-1|=1-x$. On doit alors démontrer que, si $x\leq 1$, on a $1-x\leq x^2-x+1$. Étudions donc le signe du trinôme 
			$$P(x)=x^2-x+1-(1-x)=x^2\geq 0.$$
			La propriété est vraie si $x\leq 1$.
			
			\item  Si $x\geq 1$, alors $x-1\geq 0$ et $|x-1|=x-1$. On doit alors démontrer que, si $x\geq 1$, on a $x-1\leq x^2-x+1$. Étudions donc le signe du trinôme
			$$Q(x)=x^2-x+1-(x-1)=x^2-2x+2=(x-1)^2+1\geq 0.$$
			La propriété est donc vraie si $x\geq 1$.
			
		\end{itemize}
		
		En conclusion, la propriété est vraie pour tout $x\in\mathbb R$.
	\end{corrige}
	
	\begin{exo}[{Produit de nombres qui ne sont pas divisibles par 3}]
		Le but de l'exercice est de démontrer que le produit de deux nombres entiers qui ne sont pas divisibles par 3 n'est pas divisible par 3.
		
		\begin{enumerate}
			
			\item  Soit $n$ un entier. Quels sont les restes possibles dans la division euclidienne de $n$ par $3$?
			
			\item  En déduire que si $n$ n'est pas divisible par 3, alors $n$ s'écrit $3k+1$ ou $3k+2$, avec $k$ un entier. La réciproque est-elle vraie?
			
			\item  Soit $n$ un entier s'écrivant $3k+1$ et $m$ un entier s'écrivant $3l+1$. Vérifier que 
			$$n\times m=3(3kl+k+l)+1.$$
			En déduire que $n\times m$ n'est pas divisible par $3$.
			
			\item  Démontrer la propriété annoncée par l'exercice.
			
		\end{enumerate}
	\end{exo}
	\begin{corrige}
		\begin{enumerate}
			
			\item  Les restes possibles sont $0$, $1$ ou $2$. 
			
			\item  Écrivons la division euclidienne de $n$ par $3$ : $n=3k+a$, avec $k$ un entier, et $a\in\{0,1,2\}$. Puisque $n$ n'est pas divisible par $3$, $a$ est non nul et donc $a=1$ ou $a=2$. Réciproquement, si $n=3k+1$ ou $n=3k+2$, on a écrit la division euclidienne de $n$ par 3. Comme le reste est non nul, c'est bien que $n$ n'est pas divisible par $3$.
			
			\item  Un simple calcul donne 
			$$n\times m=(3k+1)(3l+1)=9kl+3k+3l+1=3(3kl+k+l)+1.$$
			Par la propriété réciproque démontrée à la question précédente, $n\times m$ n'est pas divisible par 3.
			
			\item  Considérons deux entiers $n$ et $m$ qui ne sont pas divisibles par $3$. Alors il existe deux entiers $k$ et $l$ tels que $n$ s'écrit $n=3k+1$ ou $n=3k+2$, et $m$ s'écrit $3l+1$ ou $3l+2$. On raisonne ensuite par disjonction de cas : 
			
			\begin{itemize}
				
				\item  Si $n=3k+1$ et $m=3l+1$, alors on vient de prouver que $n\times m$ n'est pas divisible par $3$.
				
				\item  Si $n=3k+1$ et $m=3l+2$, alors
				$$n\times m=(3k+1)(3l+2)=9kl+6k+3l+2=3(3kl+2k+l)+2.$$
				Ainsi, dans ce cas $n\times m$ n'est pas divisible par $3$.
				
				\item  Si $n=3k+2$ et $m=3l+1$, alors
				$$n\times m=(3k+2)(3l+1)=9kl+3k+6l+2=3(3kl+k+2l)+2.$$
				Ainsi, dans ce cas $n\times m$ n'est pas divisible par $3$.
				
				\item  Si $n=3k+2$ et $m=3l+2$, c'est un petit peu plus difficile... L'astuce consiste à écrire que $4=3+1$ :
				\begin{align*}
					n\times m&=(3k+2)(3l+2)\\
					&=9kl+6k+6l+4\\
					&=9kl+6k+6l+3+1\\
					&=3(3kl+2k+2l+1)+1.
				\end{align*}
				Dans ce cas, $n\times m$ n'est pas divisible par $3$.
				
			\end{itemize}
			
			On a traité tous les cas, et on a donc prouvé que le produit de deux nombres qui ne sont pas divisibles par 3 n'est jamais divisible par 3.
			
		\end{enumerate}
	\end{corrige}
	
	\begin{exo}[{Reste dans la division par quatre d'une somme de deux carrés}]
		Démontrer que si $n$ est la somme de deux carrés, alors le reste de la division euclidienne de $n$ par 4 est toujours différent de $3$.
	\end{exo}
	\begin{corrige}
		On écrit $n=a^2+b^2$, et on raisonne suivant la parité de $a$ et de $b$ :
		
		\begin{itemize}
			
			\item  Si $a$ et $b$ sont tous les deux pairs, alors $a=2k$ et $b=2\ell$, $n=4(k^2+\ell^2)$ est divisible par 4 : la propriété est vraie.
			
			\item  Si $a$ est pair et $b$ est impair, alors $a=2k$ et $b=2\ell+1$. Alors $n=4k^2+4\ell^2+4l+1=4(k^2+\ell^2+\ell)+1$. Le reste de la division de $n$ par $4$ sera égal à 1. La propriété est vraie dans ce cas. Le cas $a$ impair et $b$ pair est symétrique.
			
			\item  Si $a$ est impair et $b$ est impair, alors $a=2k+1$, $b=2\ell+1$ et $n=4k^2+4k+1+4\ell^2+4\ell+1=4(k^2+k+\ell^2+\ell)+2$. Le reste de la division de $n$ par $4$ sera égal à 2. La propriété est vraie dans ce cas.
			
		\end{itemize}
		
		Ainsi, dans tous les cas, la propriété est vraie. Ce raisonnement par disjonction de cas est particulièrement utile en arithmétique, en particulier lorsqu'on utilise l'outil des congruences.
	\end{corrige}
	
	\subsection*{Raisonnement par analyse/synthèse}
	
	\begin{exo}[{Une équation avec des racines carrées}]
		Déterminer les réels $x$ tels que $\sqrt{2-x}=x$.
	\end{exo}
	\begin{corrige}
		On va raisonner par analyse-synthèse. \\
		
		Analyse : Imaginons que $x$ soit une solution de cette équation. Alors il est déjà clair que $x\in ]-\infty,2]$, sinon la racine carrée n'aurait pas de sens. On doit aussi avoir $x\geq 0$, car la racine carrée est positive et donc $x\in [0,2]$. Élevons ensuite l'équation au carré. Si $x$ est solution, alors $2-x=x^2$ (on a bien ici simplement une implication, pas une condition nécessaire et suffisante!), c'est-à-dire $x^2+x-2=0$. La résolution de cette équation du second degré donne $x_1=1$ et $x_2=-2$. Seul $x_1$ est dans l'intervalle $[0,2].$ Donc la seule solution possible est $1.$\\
		
		Synthèse : Prouvons que $x=1$ est solution de l'équation. C'est presque évident, car $\sqrt{2-1}=1$.\\
		
		Conclusion : la seule solution de l'équation est $1$.
	\end{corrige}
	
	\begin{exo}[{Une équation fonctionnelle}]
		Dans cet exercice, on souhaite déterminer toutes les fonctions $f:\mathbb R\to\mathbb R$ vérifiant la relation suivante :
		\begin{equation}
			\forall x\in\mathbb R,\ f(x)+xf(1-x)=1+x.
		\end{equation}
		
		\begin{enumerate}
			
			\item  On considère $f$ une fonction satisfaisant la relation précédente. Que vaut $f(0)$? $f(1)$?
			
			\item  Soit $x\in\mathbb R$. En substituant $x$ par $1-x$ dans la relation, déterminer $f(x)$.
			
			\item  Quelles sont les fonctions $f$ solution du problème?
			
		\end{enumerate}
	\end{exo}
	\begin{corrige}
		\begin{enumerate}
			
			\item  Choisissons $x=0$ dans la relation. On a alors $f(0)=1$. Choisissons ensuite $x=1$ dans la relation. On a alors $f(1)+1f(0)=2$ qui donne $f(1)=1$.
			
			\item  On sait que pour tout $x\in \mathbb R$, on a 
			$$f(x)+xf(1-x)=1+x.$$
			Si on substitue $1-x$ à $x$ dans cette relation, on trouve que 
			$$f(1-x)+(1-x)f(x)=2-x.$$
			On a donc un système de deux équations à deux inconnues qui va nous permettre de déterminer $f(x)$ et $f(1-x)$. Si on multiplie la deuxième équation par $x$, on trouve :
			$$x(1-x)f(x)+xf(1-x)=2x-x^2.$$
			On lui retranche la première équation et on trouve :
			$$\big(x(1-x)-1\big)f(x)=-x^2+x-1\implies (-x^2+x-1)f(x)=-x^2+x-1.$$
			Or, pour tout $x\in\mathbb R$, on a $-x^2+x-1\neq 0$. En effet, le discriminant de l'équation $-x^2+x-1=0$ vaut $\Delta=-3<0$. Ainsi, on peut simplifier par $-x^2+x-1$ et on trouve que $f(x)=1$.
			
			\item  On a démontré que si $f$ est solution, alors $f$ est la fonction constante égale à $1$. Réciproquement, la fonction constante égale à $1$ est solution. On en déduit que l'ensemble des solutions est constitué par cette fonction.
			
		\end{enumerate}
	\end{corrige}
	
	\begin{exo}[{Une équation fonctionnelle}]
		Déterminer toutes les fonctions $f:\mathbb R\to\mathbb R$ telles que, pour tous $x,y\in\mathbb R$, 
		$$f(x)\times f(y)-f(x\times y)=x+y.$$
	\end{exo}
	\begin{corrige}
		On va raisonner par analyse-synthèse en déterminant des conditions nécessaires sur $f$ pour vérifier cette équation, puis en vérifiant que ces conditions sont suffisantes.
		
		\begin{itemize}
			
			\item  Analyse : soit $f$ une fonction vérifiant cette équation. Alors, choisissant $x=y=0$, on obtient $f(0)^2-f(0)=0$, ce qui entraîne $f(0)=0$ ou $f(0)=1$. Choisissant ensuite $x=1$ et $y=0$, on a $f(1)f(0)-f(0)=1$, ce qui interdit $f(0)=0$. On a donc $f(0)=1$ (et aussi $f(1)=2$). Enfin, fixant $y=0$, on obtient que pour tout $x\in\mathbb R$, on a $f(x)=x+1$. On a donc prouvé que si $f$ est solution de l'équation, alors, pour tout $x\in\mathbb R$, $f(x)=x+1$. 
			
			\item  Synthèse : réciproquement, considérons la fonction $f$ définie par $f(x)=x+1$. Alors, pour tous $x,y\in\mathbb R$, on a
			$$f(x)\times f(y)=(x+1)(y+1)=xy+x+y+1$$
			et 
			$$f(xy)=xy+1$$
			ce qui prouve bien que $f$ est solution de l'équation.
			
		\end{itemize}
		
		En conclusion, l'équation admet une solution unique : la fonction $f$ donné par $f(x)=x+1$.
	\end{corrige}
	
	\subsection*{Exercices complémentaires}
	
	\begin{exo}[{Quantificateurs et négation}]\ \textcolor{red!70!black}{$\star$}
		Soit $f:\R\to\R$. Traduire à l'aide de quantificateurs, puis écrire la négation de chacune des assertions suivantes :
		\begin{enumerate}
			\item $f$ est croissante ;
			\item $f$ est bornée ;
			\item $f$ s'annule au moins une fois ;
			\item $f$ est continue en $a\in\R$ ;
			\item $f$ n'est pas la fonction nulle.
		\end{enumerate}
	\end{exo}
	\begin{corrige}
		\begin{enumerate}
			\item Croissante : $\forall x,y\in\R,\ x\leq y\Rightarrow f(x)\leq f(y)$. Négation : $\exists x,y\in\R,\ x\leq y$ et $f(x)>f(y)$.
			\item Bornée : $\exists M\in\R,\ \forall x\in\R,\ |f(x)|\leq M$. Négation : $\forall M\in\R,\ \exists x\in\R,\ |f(x)|>M$.
			\item S'annule : $\exists x\in\R,\ f(x)=0$. Négation : $\forall x\in\R,\ f(x)\neq 0$.
			\item Continue en $a$ : $\forall\veps>0,\ \exists\delta>0,\ \forall x\in\R,\ |x-a|\leq\delta\Rightarrow|f(x)-f(a)|\leq\veps$. Négation : $\exists\veps>0,\ \forall\delta>0,\ \exists x\in\R,\ |x-a|\leq\delta$ et $|f(x)-f(a)|>\veps$.
			\item « $f$ n'est pas nulle » s'écrit $\exists x\in\R,\ f(x)\neq 0$ ; sa négation est $\forall x\in\R,\ f(x)=0$.
		\end{enumerate}
	\end{corrige}
	
	\begin{exo}[{Irrationalité par l'absurde}]\ \textcolor{red!70!black}{$\star$}
		\begin{enumerate}
			\item Montrer par l'absurde que $\sqrt{2}\notin\Q$.
			\item Montrer que si $a\in\Q$ et $b\notin\Q$, alors $a+b\notin\Q$. En déduire que $1+\sqrt 2\notin\Q$.
			\item Montrer par l'absurde qu'il n'existe pas de plus petit réel strictement positif.
		\end{enumerate}
	\end{exo}
	\begin{corrige}
		\begin{enumerate}
			\item Supposons $\sqrt2=p/q$ avec $p,q\in\N^*$ premiers entre eux. Alors $p^2=2q^2$, donc $p^2$ est pair, donc $p$ est pair : $p=2k$. Alors $4k^2=2q^2$, soit $q^2=2k^2$, donc $q$ est pair. Ainsi $2$ divise $p$ et $q$ : contradiction.
			\item Si $a+b=r\in\Q$, alors $b=r-a\in\Q$, ce qui est faux. Avec $a=1$ et $b=\sqrt2\notin\Q$, on obtient $1+\sqrt2\notin\Q$.
			\item Si $m>0$ était le plus petit réel strictement positif, alors $m/2>0$ et $m/2<m$ : contradiction.
		\end{enumerate}
	\end{corrige}
	
	\begin{exo}[{Récurrences}]\ \textcolor{red!70!black}{$\star$}
		\begin{enumerate}
			\item Montrer que pour tout $n\in\N^*$, $\dis\sum_{k=1}^n k^3=\left(\frac{n(n+1)}{2}\right)^2$.
			\item Montrer que pour tout $n\geq 4$, $2^n\geq n^2$.
			\item Montrer par récurrence forte que tout entier $n\geq 2$ est un produit de nombres premiers.
		\end{enumerate}
	\end{exo}
	\begin{corrige}
		\begin{enumerate}
			\item Pour $n=1$ : $1=1^2$. Si la formule est vraie au rang $n$, alors
			$\dis\sum_{k=1}^{n+1}k^3=\frac{n^2(n+1)^2}{4}+(n+1)^3=(n+1)^2\frac{n^2+4n+4}{4}=\left(\frac{(n+1)(n+2)}{2}\right)^2.$
			\item Pour $n=4$ : $16\geq 16$. Si $2^n\geq n^2$ avec $n\geq 4$, alors $2^{n+1}\geq 2n^2\geq(n+1)^2$ car $n^2-2n-1=(n-1)^2-2>0$.
			\item $2$ est premier. Soit $n\geq 3$ et supposons le résultat vrai pour tous les entiers de $2$ à $n-1$. Si $n$ est premier, c'est fini. Sinon $n=ab$ avec $2\leq a,b\leq n-1$ ; $a$ et $b$ sont produits de premiers, donc $n$ aussi.
		\end{enumerate}
	\end{corrige}
	
	\begin{exo}[{Disjonction de cas}]\ \textcolor{red!70!black}{$\star$}
		Montrer que pour tout $n\in\N$, l'entier $n(n+1)(n+2)$ est divisible par $3$, puis par $6$. Montrer ensuite que pour tous $x,y\in\R$, $|x+y|\leq |x|+|y|$ en distinguant les cas selon les signes de $x$ et $y$.
	\end{exo}
	\begin{corrige}
		Selon le reste de $n$ modulo $3$ : si $n\equiv0$, $3\mid n$ ; si $n\equiv1$, $3\mid n+2$ ; si $n\equiv2$, $3\mid n+1$. Dans tous les cas $3\mid n(n+1)(n+2)$. Par ailleurs $n$ ou $n+1$ est pair, donc $2$ divise le produit. Comme $2$ et $3$ sont premiers entre eux, $6$ divise $n(n+1)(n+2)$.
		
		Inégalité triangulaire : si $x,y\geq0$, $|x+y|=x+y=|x|+|y|$ ; si $x,y\leq0$, $|x+y|=-x-y=|x|+|y|$ ; si $x\geq0\geq y$, alors $x+y\leq x\leq|x|+|y|$ et $-(x+y)=-y-x\leq -y\leq|x|+|y|$, donc $|x+y|\leq|x|+|y|$ ; le cas $y\geq0\geq x$ est symétrique.
	\end{corrige}
	\section{Ensembles, applications, relations}
	
	\subsection*{Pour comprendre le cours}
	
	\begin{exo}[{Écriture en extension}]
		Écrire en extension (c'est-à-dire en donnant tous leurs éléments) les ensembles suivants :
		$$A=\left\{\textrm{nombres entiers compris entre $\sqrt{2}$ et $2\pi$}\right\}.$$
		$$B=\left\{x\in\mtq;\ \exists(n,p)\in\mtn^*\times\mtn,\ x=\frac{p}{n}\textrm{ et }1\leq p\leq 2n\leq 7\right\}.$$
	\end{exo}
	\begin{corrige}
		On a :
		$$A=\left\{2,3,4,5,6\right\}.$$
		Pour écrire $B$, on remarque que $1/2\leq n\leq 7/2\implies n=1,2$ ou $3$. Pour chaque valeur possible de $n$, on détermine ensuite les valeurs possibles de $p$ : 
		
		\begin{itemize}
			
			\item  Pour $n=1$, on a $p\in\{1,2\}$.
			
			\item  Pour $n=2$, on a $p\in\{1,2,3,4\}$.
			
			\item  Pour $n=3$, on a $p\in\{1,2,3,4,5,6\}$.
			
		\end{itemize}
		
		Finalement, on trouve
		$$B=\left\{1,2,\frac{1}{2},\frac{3}{2},\frac{1}{3},\frac{2}{3},\frac{4}{3},\frac{5}{3}\right\}.$$
		On n'a pas écrit plusieurs fois 1, qui s'obtient aussi avec 2/2 et 3/3, ni plusieurs fois 2, qu'on obtient avec 2/1, 4/2 et 6/3.
	\end{corrige}
	
	\begin{exo}[{Ensemble des parties}]
		Écrire l'ensemble des parties de $E=\left\{a,b,c,d\right\}$.
	\end{exo}
	\begin{corrige}
		On classe les parties suivant leur nombre d'éléments :
		
		\begin{enumerate}
			
			\item  0 éléments : Il n'y a que l'ensemble vide : $\varnothing$.
			
			\item  1 élément : Il y a les 4 singletons : $\{a\},\ \{b\},\ \{c\},\ \{d\}$.
			
			\item  2 éléments : Il y a 6 parties à 2 éléments :
			$$\{a,b\},\ \{a,c\},\ \{a,d\},\ \{b,c\},\ \{b,d\},\ \{c,d\}.$$
			
			\item  3 éléments : Il y a 4 parties à 3 éléments :
			$$\{a,b,c\},\ \{a,b,d\},\ \{a,c,d\},\ \{b,c,d\}.$$
			
			\item  4 éléments : Il n'y a qu'une partie à 4 éléments : l'ensemble $E$ lui-même.
			
		\end{enumerate}
		
		L'ensemble des parties de $E$ comporte donc $16=2^4$ éléments.
	\end{corrige}
	
	\begin{exo}[{Encore des exemples}]
		Les applications suivantes sont-elles injectives, surjectives, bijectives?
		
		\begin{enumerate}
			
			\item  $f:\mtn\to\mtn,\ n\mapsto n+1$.
			
			\item  $g:\mtz\to\mtz,\ n\mapsto n+1$.
			
			\item  $h:\mtr^2\to\mtr^2,\ (x,y)\mapsto (x+y,x-y)$.
			
		\end{enumerate}
	\end{exo}
	\begin{corrige}
		\begin{enumerate}
			
			\item  $f$ est clairement injective, mais n'est pas surjective car $0$ n'a pas d'antécédent.
			
			\item  $g$ est bijective : l'équation $n+1=k$, avec $k\in\mtz$ admet une unique solution $n\in\mtz$ qui vaut $n=k-1$.
			
			\item  $h$ est bijective : prenons en effet un couple $(x_1,y_1)$ de $\mtr^2$, et essayons de résoudre le système :
			$$\left\{
			\begin{array}{rcl}
				x+y&=&x_1\\
				x-y&=&y_1
			\end{array}
			\right.$$
			Ce système possède une unique solution, donnée par $x=(x_1+y_1)/2$ et $y=(x_1-y_1)/2$. L'application est bijective.
			
		\end{enumerate}
	\end{corrige}
	
	\begin{exo}[{Exemples d'image directe et d'image réciproque}]
		\begin{enumerate}
			
			\item  Soit $f:\mathbb R\to \mathbb R$, $x\mapsto x^2$, et soit $A=[-1,4]$.
			Déterminer
			
			\begin{enumerate}
				
				\item  l'image directe de $A$ par $f$;
				
				\item  l'image réciproque de $A$ par $f$.
				
			\end{enumerate}
			
			
			\item  On considère la fonction $\sin:\mathbb R\to \mathbb R$. Quelle est l'image directe, par $\sin$, de $\mathbb R$? De $[0,2\pi]$? de $[0,\pi/2]$?
			Quelle est l'image réciproque, par $\sin$, de $[0,1]$? de $[3,4]$? de $[1,2]$?
			
		\end{enumerate}
	\end{exo}
	\begin{corrige}
		\begin{enumerate}
			
			\item  
			\begin{enumerate}
				
				\item  On cherche toutes les valeurs prises par $x^2$ lorsque $x$ parcourt $[-1,4]$. Entre $-1$ et $0$, ce sont toutes les valeurs de 0 à 1
				qui sont prises, et entre $0$ et $4$, toutes les valeurs entre 0 et 16. On a donc $f(A)=[0,16]$.
				
				\item  On a $x\in f^{-1}(A)$ si et seulement si $x^2\in [-1,4]$. Bien sûr, les valeurs négatives sont exclues, et pour que $x^2$ soit dans $[0,4]$, il est
				nécessaire et suffisant que $x\in[-2,2]$. On a donc $f^{-1}(A)=[-2,2]$.
				
			\end{enumerate}
			
			
			\item   L'image directe de $\mathbb R$ comme de $[0,2\pi]$ est $[-1,1]$. L'image directe de $[0,\pi/2]$ est $[0,1]$. 
			Pour déterminer l'image réciproque de $[0,1]$, on cherche les réels $x$ tels que $\sin(x)\in[0,1]$. Ce sont tous les réels
			qui peuvent s'écrire $u+k2\pi$, avec $u\in[0,\pi]$ et $k\in\mathbb Z$. On peut encore écrire cet ensemble
			$$\bigcup_{k\in\mathbb Z}[2k\pi,(2k+1)\pi].$$ Aucun réel n'a son sinus dans $[3,4]$. L'image réciproque de $[3,4]$ est donc l'ensemble vide.
			Enfin, l'image réciproque de $[1,2]$ est identique à l'image réciproque de $\{1\}$, et elle est égale à $\{\pi/2+2k\pi;\ k\in\mathbb Z\}$.
			
		\end{enumerate}
	\end{corrige}
	
	\begin{exo}[{Composition non commutative}]
		\begin{enumerate}
			
			\item  Soient $f$ et $g$ les deux fonctions de $\mathbb R$ dans $\mathbb R$ définies par
			$$f(x)=3x+1\textrm{ et }g(x)=x^2-1.$$
			Calculer $f\circ g$ et $g\circ f$.
			
			\item  Dans les exemples suivants, déterminer deux fonctions $u$ et $v$ telles que $h=u\circ v$ :
			$$h_1(x)=\sqrt{3x-1}\quad h_2(x)=\sin\left(x+\frac \pi 2\right)\quad h_3(x)=\frac 1{x+7}.$$
			
		\end{enumerate}
	\end{exo}
	\begin{corrige}
		\begin{enumerate}
			
			\item  Soit $x\in\mathbb R$. On a
			$$f\circ g(x)=f(x^2-1)=3(x^2-1)+1=3x^2-2.$$
			D'autre part, on a 
			$$g\circ f(x)=g(3x+1)=(3x+1)^2-1=9x^2+6x.$$
			En particulier, on a $f\circ g\neq g\circ f$.
			
			\item  Pour chacun des cas, on peut poser :
			
			\begin{enumerate}
				
				\item  $u(x)=\sqrt x$, $v(x)=3x-1$;
				
				\item  $u(x)=\sin x$, $v(x)=x+\frac\pi 2$;
				
				\item  $u(x)=\frac 1x$, $v(x)=x+7$.
				
			\end{enumerate}
			
			
		\end{enumerate}
	\end{corrige}
	
	\begin{exo}[{Composition et injectivité}]
		Soient $f$ et $g$ les applications de $\mathbb N$ dans $\mathbb N$ définies par $f(x)=2x$ et 
		$$g(x)=\left\{
		\begin{array}{ll}
			\frac x2&\textrm{ si $x$ est pair}\\
			0&\textrm{ si $x$ est impair.}
		\end{array}
		\right.$$
		Déterminer $g\circ f$ et $f\circ g$. Les fonctions $f$ et $g$ sont-elles injectives? surjectives? bijectives?
	\end{exo}
	\begin{corrige}
		On a $g\circ f(x)=g(2x)$. Mais $2x$ est pair, et donc $g(2x)=(2x)/2=x$. Ainsi, $g\circ f(x)=x$.\\
		
		D'autre part, si $x$ est pair, on a $f\circ g(x)=f(x/2)=x$. Si $x$ est impair, $f\circ g(x)=f(0)=0$. En particulier, on a $f\circ g\neq g\circ f$ puisque $f\circ g(1)=0$ alors que $g\circ f(1)=1$.\\
		
		$f$ n'est pas surjective, car les nombres impairs ne sont pas des valeurs prises par $f$. En revanche, $f$ est injective car si $f(x)=f(y)$, on a $2x=2y$ et donc $x=y$.\\
		
		$g$ n'est pas injective, car $g(1)=g(3)=0$ alors que $1\neq 3$. En revanche, $g$ est surjective. Prenons en effet $y$ n'importe quel entier naturel. Alors, $2y$ est pair et $g(2y)=(2y)/2=y$.\\
		
		Des deux études précédentes, on déduit par définition que ni $f$ ni $g$ ne sont bijectives.
	\end{corrige}
	
	\begin{exo}[{Calcul de la réciproque}]
		Démontrer que la fonction $f:\mathbb R\to\mathbb R_+^*$ définie par 
		$$f(x)=\frac{e^x+2}{e^{-x}}$$
		est bijective. Calculer sa bijection réciproque $f^{-1}$.
	\end{exo}
	\begin{corrige}
		La première chose à remarquer est que $f$ s'écrit plus facilement 
		$$f(x)=e^{2x}+2e^x.$$
		Soit $y>0$. Alors on a
		$$y=f(x)\iff e^{2x}+2e^x-y=0.$$
		On pose alors $X=e^x$ et l'équation est équivalente à l'équation du second degré 
		$$X^2+2X-y=0$$
		dont les racines sont
		$$X_1=-1-\sqrt{1+y}\textrm{ et }X_2=-1+\sqrt{1+y}.$$
		Mais $e^x>0$, et donc on a 
		$$y=f(x)\iff  e^x=-1+\sqrt{1+y}\iff x=\ln\left(-1+\sqrt{1+y}\right).$$
		L'équation $y=f(x)$ admet donc toujours une unique solution. La fonction $f$ est bijective et 
		$$f^{-1}(y)=\ln\left(-1+\sqrt{1+y}\right).$$
	\end{corrige}
	
	\begin{exo}[{Avec des nombres complexes}]
		Démontrer que la fonction $$
		\begin{array}{rcl}
			f:\mathbb C\backslash \{-3\}&\to&\mathbb C\\
			z&\mapsto&\frac{iz-i}{z+3}
		\end{array}$$
		est une bijection de $\mathbb C\backslash\{-3\}$ sur $\mathbb C\backslash\{i\}.$ Déterminer sa bijection réciproque.
	\end{exo}
	\begin{corrige}
		On va démontrer directement que $f$ est bijective en prouvant que, pour tout $w\in\mathbb C\backslash\{i\}$, l'équation 
		$f(z)=w$ admet une unique solution $z\in\mathbb C\backslash\{-3\}$. Pour cela, on remarque que
		\begin{eqnarray*}
			\frac{iz-i}{z+3}=w&\iff&iz-i=wz+3w\\
			&\iff&(i-w)z=3w+i.
		\end{eqnarray*}
		Puisque $w\neq i$, ceci est encore équivalent à 
		$$z=\frac{3w+i}{i-w}.$$
		L'équation admet une unique solution : $f$ est bijective et sa bijection réciproque est donnée par
		$$f^{-1}(w)=\frac{3w+i}{i-w}.$$
	\end{corrige}
	
	\begin{exo}[{Nature des relations}]
		Dire si les relations suivantes sont réflexives, symétriques, antisymétriques, transitives :
		
		\begin{enumerate}
			
			\item  $E=\mathbb Z$ et $x\mathcal R y\iff x=-y$;
			
			\item  $E=\mathbb R$ et $x\mathcal R y\iff \cos^2 x+\sin^2 y=1$;
			
			\item  $E=\mathbb N$ et $x\mathcal R y\iff \exists p,q\geq 1,\ y=px^q$ ($p$ et $q$ sont des entiers).
			
		\end{enumerate}
		
		Quelles sont parmi les exemples précédents les relations d'ordre et les relations d'équivalence?
	\end{exo}
	\begin{corrige}
		\begin{enumerate}
			
			\item  La relation n'est pas réflexive, car $1$ n'est pas en relation avec lui-même. En effet, $1\neq -1$.
			La relation est symétrique, car $x=-y\iff y=-x$. Elle n'est pas antisymétrique, car $1\mathcal R -1$ et
			$-1\mathcal R 1$, alors que $1\neq -1$. Elle n'est pas transitive. En effet, on a $1\mathcal R -1$, $-1\mathcal R 1$ et $1$ et $1$ ne sont pas en relation.
			Cette relation n'est ni une relation d'équivalence, ni une relation d'ordre.
			
			\item  De la formule $\cos^2 x+\sin^2 x=1$, on déduit que la relation est réflexive. Elle est aussi symétrique.
			En effet, si $x\mathcal R y$, ie $\cos^2 x+\sin^2 y=1$, alors on a
			$$\sin^2 x+\cos ^2 x+\cos ^2 y+\sin^2 y=(\cos^2 x+\sin^2 y)+(\cos^2 y+\sin^2 x)=1+(\cos^2 y+\sin^2 x)$$
			d'une part, et
			$$\sin^2 x+\cos ^2 x+\cos ^2 y+\sin^2 y=1+1=2$$
			d'autre part, ce qui entraîne bien 
			$$\cos^2 y+\sin^2 x=1$$
			et donc la relation est symétrique.\\
			
			Elle n'est pas antisymétrique, car $0\mathcal R 2\pi$ et $2\pi\mathcal R 0$ alors que $0\neq 2\pi$.\\
			
			Elle est transitive. Si $x\mathcal R y$ et $y\mathcal R z$, on a 
			$$\cos^2 x+\sin^2 y=1\textrm{ et }\cos^2 y+\sin^2 z=1$$
			soit en sommant
			$$\cos^2 x+(\cos^2 y+\sin^2 y)+\sin^2 z=2$$
			ce qui implique
			$$\cos^2 x+\sin^2 z=1.$$
			On a donc affaire à une relation d'équivalence.
			
			\item  La relation est réflexive (prendre $p=q=1$), elle n'est pas symétrique car si $x\mathcal R y$, on a 
			forcément $x\leq y$. Ainsi, on a $2\mathcal R 4$ (prendre $p=2$, $q=1$), alors qu'on n'a pas $4\mathcal R 2$.
			La relation est antisymétrique : si $x\mathcal R y$ et $y\mathcal R x$, alors on a
			$x\leq y$ et $y\leq x$ et donc $x=y$. Enfin, la relation est transitive. Si $x\mathcal R y$
			et $y\mathcal R z$, alors il existe des entiers $p,q,a,b\geq 1$ tels que
			$$y=px^q\textrm{ et }z=ay^b.$$
			On en déduit
			$$z=a(px^q)^b=(ap^b)x^{bq}$$
			et donc $x\mathcal R z$. La relation est une relation d'ordre.
			
		\end{enumerate}
	\end{corrige}
	
	\begin{exo}[{Classe d'équivalence sur $\mathbb R^2$}]
		Sur $\mathbb R^2$, on définit la relation d'équivalence $\mathcal R$ par 
		$$(x,y)\mathcal R (x',y')\iff x=x'.$$
		Démontrer que $\mathcal R$ est une relation d'équivalence, puis déterminer la classe d'équivalence d'un élément $(x_0,y_0)\in\mathbb R^2$.
	\end{exo}
	\begin{corrige}
		La relation $\mathcal R$ est une relation d'équivalence. En effet, elle est
		
		\begin{itemize}
			
			\item  réflexive : $(x,y)\mathcal R (x,y)$ car $x=x$;
			
			\item  symétrique : si $(x,y)\mathcal R(x',y')$, alors $x=x'$ ce qui s'écrit aussi $x'=x$ et qui est équivalent à $(x',y')\mathcal R(x,y)$.
			
			\item  transitive : si $(x,y)\mathcal R(x',y')$ et si $(x',y')\mathcal R (x'',y'')$, alors on a $x=x'$ d'une part et $x'=x''$ d'autre part, donc $x=x''$ ce qui entraine $(x,y)\mathcal R (x'',y'')$.
			
		\end{itemize}
		
		Chercher la classe d'équivalence de $(x_0,y_0)$, c'est déterminer tous les couples $(x,y)$ tels que $(x,y)\mathcal R  (x_0,y_0)$. Mais, 
		$$(x,y)\mathcal R(x_0,y_0)\iff x=x_0.$$
		Autrement dit, $x$ doit être égal à $x_0$ et $y$ peut être quelconque. On en déduit que la classe d'équivalence de $(x_0,y_0)$ pour la relation d'équivalence $\mathcal R$ est $\{(x_0,y);\ y\in\mathbb R\}$.
	\end{corrige}
	
	\subsection*{Ensembles}
	
	\begin{exo}[{Partie d'une union}]
		Est-ce que $C\subset A\cup B$ entraîne $C\subset A$ ou $C\subset B$?
	\end{exo}
	\begin{corrige}
		Non! Prendre par exemple $A=\{1,2\}$, $B=\{3,4\}$ et $C=\{2,3\}$.
	\end{corrige}
	
	\begin{exo}[{Trois ensembles}]
		Soient $A,B,C$ trois ensembles tels que $A\cup B=B\cap C$.	
		Montrer que $A\subset B\subset C$.
	\end{exo}
	\begin{corrige}
		Prenons $x\in A$. Alors $x\in A\cup B$, et donc $x\in B\cap C$.
		En particulier, $x\in B$, et donc $A\subset B$.\\
		
		Prenons maintenant $x\in B$. Alors $x\in A\cup B$, et donc $x\in B\cap C$.
		En particulier, $x\in C$, et donc $B\subset C$.
	\end{corrige}
	
	\begin{exo}[{Deux descriptions d'un même ensemble}]
		Soit $A=\{(x,y)\in\mathbb R^2;\ 4x-y=1\}$ et $C=\{(t+1,4t+3);\ t\in\mathbb R\}$. Démontrer que $A=C$.
	\end{exo}
	\begin{corrige}
		On va procéder par double inclusion. Le plus facile est de prouver que $C\subset A$. En effet, prenons un couple $(x,y)\in C$. Alors on sait qu'il existe un $t\in\mathbb R$ tel que $x=t+1$ et $y=4t+3$. Mais alors,
		$$4x-y=4t+4-4t-3=1$$
		et donc on a bien $(x,y)\in A$.\\
		
		Réciproquement, prenons $(x,y)\in A$ et prouvons que $(x,y)\in C$. C'est plus difficile, car il faut construire un réel $t$. On va procéder par analyse-synthèse. Si un tel $t$ existe, alors nécessairement on doit avoir $t=x-1$. Posons donc $t=x-1$. Alors, $y=4x-1=4(t+1)-1=4t+3$. On a donc bien $(x,y)=(t+1,4t+3)$ et $(x,y)\in C$.
	\end{corrige}
	
	\begin{exo}[{Lois de De Morgan}]
		Soient $A$, $B$ et $C$ trois parties d'un ensemble $E$. Pour $X\subset E$, on note $X^c$ le complémentaire de
		$X$ dans $E$. Démontrer les lois de Morgan suivantes :
		$$\begin{array}{lll}
			\mathbf{1.}\ (A\cap B)\cup C=(A\cup C)\cap (B\cup C)&&\mathbf{2.}\ (A^c)^c=A\\
			\mathbf{3.}\ (A\cap B)^c=A^c\cup B^c&&\mathbf{4.}\ (A\cup B)^c=A^c\cap B^c.\\
		\end{array}$$
	\end{exo}
	\begin{corrige}
		On raisonne à chaque fois par double inclusion.
		
		\begin{enumerate}
			
			\item  Soit $x\in (A\cap B)\cup C$. Alors ($x\in A$ et $x\in B$) ou $x\in C$. Si $x\in A$ et $x\in B$, alors $x\in A\cup C$ et $x\in B\cup C$, et l'inclusion est prouvée. Sinon, c'est que $x\in C$, et dans ce cas on a aussi $x\in A\cup C$ et $x\in B\cup C$.\\
			
			Réciproquement, si $x\in A\cup C$ et $x\in B\cup C$, on distingue deux cas :
			
			\begin{itemize}
				
				\item  Si $x\in C$, alors $x\in (A\cap B)$ ou $x\in C$ et donc $x\in (A\cap B)\cup C$.
				
				\item  Sinon, $x\notin C$. Mais alors, puisque $x\in A\cup C$, on a $x\in A$. De même, puisque $x\in B\cup C$, on a $x\in B$. Ceci prouve que $x\in A\cap B$ et donc $x\in (A\cap B)\cup C$.
				
			\end{itemize}
			
			
			\item  On suppose que $x\in (A^c)^c$. Alors $x\notin A^c$, et donc $x\in A$. Réciproquement, si $x\in A$, alors $x\notin
			A^c$ et donc $x\in (A^c)^c$.
			
			\item  Soit $x\in(A\cap B)^c$. Alors $x\notin A\cap B$. On a donc $x\notin A$ ou $x\notin B$, c'est-à-dire
			$x\in A^c$ ou $x\in B^c$. On en déduit que $x\in A^c\cup B^c$. Réciproquement, soit $x\in A^c\cup B^c$. Alors
			$x\in A^c$ ou $x\in B^c$, c'est-à-dire $x\notin A$ ou $x\notin B$. En particulier, $x\notin A\cap B$ et
			donc $x\in(A\cap B)^c$.
			
			\item  On peut présenter aussi les raisonnements précédents sous forme d'équivalence. C'est ce que l'on fait pour ce
			dernier exemple :
			\begin{eqnarray*}
				x\in (A\cup B)^c&\iff&x\notin A\cup B\\
				&\iff&x\notin A\textrm{ et }x\notin B\\
				&\iff&x\in A^c\textrm{ et }x\in B^c\\
				&\iff&x\in A^c\cap B^c.
			\end{eqnarray*}
			
		\end{enumerate}
	\end{corrige}
	
	\begin{exo}[{Produit cartésien}]
		Soit $D=\{(x,y)\in\mathbb R^2;\ x^2+y^2\leq 1\}$. Démontrer que $D$ ne peut pas s'écrire comme le produit cartésien de deux parties
		de $\mathbb R$.
	\end{exo}
	\begin{corrige}
		Procédons par l'absurde et supposons qu'il existe $A,B\subset\mathbb R$ tels que $D=A\times B$.
		Alors le point $(1,0)\in D$, et donc $1\in A$. De même, $(0,1)\in D$, et donc $1\in B$. On en déduit que $(1,1)\in A\times B=D$,
		ce qui n'est pas le cas. $D$ n'est donc pas le produit cartésien de deux parties de $\mathbb R$.
	\end{corrige}
	
	\begin{exo}[{Réunion et intersection égales}]
		Soit $E$ un ensemble et $A,B,C$ trois éléments de $\mathcal P(E)$.
		
		\begin{enumerate}
			
			\item  Démontrer que, si $A\cap B=A\cup B$, alors $A=B$.
			
			\item  Démontrer que, si $A\cap B=A\cap C$ et $A\cup B=A\cup C$, alors $B=C$.
			Une seule des deux conditions suffit-elle?
			
		\end{enumerate}
	\end{exo}
	\begin{corrige}
		\begin{enumerate}
			
			\item  Par symétrie du problème en $A$ et $B$, il suffit de démontrer que $A\subset B$. Prenons $x\in A$ et
			supposons que $x\notin B$. Alors $x\in A\cup B$ mais $x\notin A\cap B$ et donc les ensembles $A\cap B$ et $A\cup B$ sont
			différents, une contradiction. C'est donc que $x\in B$.
			
			\item  Par symétrie du problème en $B$ et $C$, il suffit de démontrer l'inclusion $B\subset C$.
			Soit donc $x\in B$. On distingue deux cas :
			
			\begin{itemize}
				
				\item  ou bien $x\in A$. Dans ce cas, $x\in A\cap B=A\cap C$, et donc $x\in C$.
				
				\item  ou bien $x\notin A$. Dans ce cas, $x\in A\cup B=A\cup C$ et donc $x\in A$ ou $x\in C$.
				Puisqu'on est dans le cas $x\notin A$, on en déduit que $x\in C$.
				
			\end{itemize}
			
			Dans tous les cas, on a démontré $x\in C$, et donc $B\subset C$. Une seule des deux conditions n'est pas suffisante :
			
			\begin{itemize}
				
				\item  Si on suppose seulement que $A\cup B\subset A\cup C$, il suffit de prendre $A=\{1,2\}$, $B=\{1\}$ et $C=\{2\}$. On a bien $A\cup B\subset A\cup C$, mais on n'a pas $B\subset C$.
				
				\item  Si on suppose seulement que $A\cap B\subset A\cap C$, il suffit de prendre $A=C=\{1\}$ et $B=\{1,2\}$.
				
			\end{itemize}
			
			
		\end{enumerate}
	\end{corrige}
	
	\begin{exo}[{Différence symétrique}]
		Soit $E$ un ensemble, et $A,B$ deux sous-ensembles de $E$. On appelle différence symétrique de $A$ et $B$, notée $A\Delta B$, le sous-ensemble de $E$ :
		$$A\Delta B=\{x\in A\cup B;\ x\notin A\cap B\}.$$
		
		\begin{enumerate}
			
			\item  Interpréter les éléments de $A\Delta B$.
			
			\item  Montrer que $A\Delta B=(A\cap C_EB)\cup (B\cap C_EA)$ ($C_EA$ désigne le complémentaire de $A$ dans $E$).
			
			\item  Calculer $A\Delta A$, $A\Delta \varnothing$, $A\Delta E$, $A\Delta C_E A$.
			
			\item  Démontrer que pour tous $A,B,C$ sous-ensembles de $E$, on a :
			$$(A\Delta B)\cap C=(A\cap C)\Delta (B\cap C).$$
			
		\end{enumerate}
	\end{exo}
	\begin{corrige}
		\begin{enumerate}
			
			\item  Les éléments de $A\Delta B$ sont les éléments qui appartiennent à $A$ ou à $B$, mais qui n'appartiennent pas simultanément aux deux (dans un langage informatique, on pourrait parler de "ou exclusif").
			
			\item  Soit $x\in A\Delta B$. Par symétrie du problème, on peut toujours supposer que $x\in A$. Nécessairement, $x\notin B$. On en déduit que $x\in A$ et
			$x\in C_E B$. Ceci donne $x\in A\cap C_EB$.\\
			
			Réciproquement, si par exemple $x\in A\cap C_EB$, $x\in A$ et $x\notin B$, et donc $x\in A\cup B$ et $x\notin A\cap B$. L'autre possibilité se traite exactement de la même façon.
			
			\item  En utilisant ou la définition (c'est plus clair avec l'interprétation de la première question), ou le résultat précédent, on a :
			$$A\Delta A=\varnothing.$$
			$$A\Delta \varnothing=A.$$
			$$A\Delta E=C_E A.$$
			$$A\Delta C_E A=E.$$
			
			\item  Comme toujours, on raisonne par double inclusion. Si $x\in (A\Delta B)\cap C$, alors $x\in A\Delta B$ et $x\in C$. Si on a $x\in A$, alors $x\notin B$, et $x\in C$, ce qui donne encore : $x\in A\cap C$ et $x\notin B\cap C$. On a bien $x\in (A\cap C)\Delta (B\cap C)$. Le cas où $x\in B$ se traite exactement de la même façon (par symétrie).\\
			
			Réciproquement, si $x\in (A\cap C)\Delta (B\cap C)$, supposons par exemple que $x\in (A\cap C)$. Alors $x\notin B\cap C$, ce qui implique $x\notin B$ ou $x\notin C$. Mais, $x\in (A\cap C)\implies x\in A\textrm{ et }x\in C$. On en déduit que $x\notin B$. D'où $x\in A\Delta B$, et $x\in C$, ce qui est le résultat que nous voulions prouver. L'autre cas se traite également par symétrie.\\
			
			
		\end{enumerate}
	\end{corrige}
	
	\begin{exo}[{Retour sur la différence symétrique}]
		Soit $E$ un ensemble et soient $A,B$ deux parties de $E$. On rappelle que la différence symétrique de
		$A$ et $B$ est définie par 
		$$A \Delta  B =  (A\cap \bar{B})\cup \left(\bar{A}\cap B\right)$$
		où $\bar A$ (resp. $\bar B$) désigne le complémentaire de $A$ (resp. de $B$) dans $E$.
		Démontrer que $A\Delta B=B$ si et seulement si $A=\varnothing$.
	\end{exo}
	\begin{corrige}
		Il y a d'abord un sens qui est facile : si $A=\varnothing$, alors par définition de la différence
		symétrique, on a bien $A\Delta B=B$ car $A=\varnothing$ et $\bar A\cap B=B$.
		Réciproquement, si $A\Delta B=B$, il faut prouver que $A=\varnothing$.
		On va découper la preuve en deux parties :
		
		\begin{enumerate}
			
			\item  on prouve que $A\cap B=\varnothing$.
			En effet, prenons $x\in  B$. Alors, en particulier $x\in A\Delta B$, et donc $x\in A\cap\bar B$ ou $x\in \bar A\cap B$.
			La première éventualité est impossible (car $x\in B$) et donc on a $x\in\bar A\cap B$. Ainsi, tout élément de $B$ est aussi dans $\bar A$,
			et donc $A\cap B=\varnothing$.
			
			\item  on va aussi prouver que $A\cap \bar B=\varnothing$. En effet, imaginons qu'on puisse trouve un élément dans $A\cap\bar B$.
			Alors cet élément serait aussi dans $A\Delta B=B$, ce qui est impossible puisqu'il serait simultanément dans $B$ et dans $\bar B$.
			
		\end{enumerate}
		
		La confrontation des deux propriétés précédentes entraîne alors que $A=\varnothing$.
	\end{corrige}
	
	\begin{exo}[{Fonction caractéristique}]
		Soit $A$ une partie d'un ensemble  $E$. On appelle fonction caractéristique de $A$ l’application $f$ de $E$ dans l’ensemble à deux
		éléments $\{0, 1\}$ telle que : 
		$$f(x)=\left\{
		\begin{array}{ll}
			1&\textrm{ si }x\in A\\
			0&\textrm{ si }x\notin A
		\end{array}\right.$$
		Soient $A$ et $B$ deux parties de $E$, $f$ et $g$ leurs fonctions caractéristiques. Montrer que les fonctions suivantes sont les fonctions caractéristiques d’ensembles que l’on déterminera : 
		
		\begin{enumerate}
			
			\item  $1-f$;
			
			\item  $fg$;
			
			\item  $f+g-fg$.
			
		\end{enumerate}
	\end{exo}
	\begin{corrige}
		\begin{enumerate}
			
			\item  $h=1-f$ est la fonction caractéristique du complémentaire de $A$, $\bar A$. En effet, on a 
			$h(x)=0$ si $f(x)=1$, c'est-à-dire si $x\in A$, c'est-à-dire si $x\notin \bar A$, et $h(x)=1$ si $f(x)=0$, c'est-à-dire si $x\notin A$, c'est-à-dire si $x\in\bar A$.
			
			\item  $h=fg$ est la fonction caractéristique de $A\cap B$. En effet, si $x\in A\cap B$, alors $f(x)=g(x)=1$, et donc $h(x)=1$. Si $x\notin A\cap B$, alors ou bien $x\notin A$ et $f(x)=0$ ou bien $x\notin B$ et $g(x)=0$. Dans tous les cas, $h(x)=0$.
			
			\item  $h=f+g-fg$ est la fonction caractéristique de $A\cup B$. En effet, si $x\notin A\cup B$, alors $f(x)=g(x)=0$, et donc $h(x)=0$. Si $x\in A\cup B$, alors on peut distinguer trois cas :
			
			\begin{itemize}
				
				\item  $x\in A$ et $x\in B$ : on a alors $f(x)=g(x)=1$, et $h(x)=1+1-1=1$;
				
				\item  $x\in A$ et $x\notin B$ : on a alors $f(x)=1$ et $g(x)=0$, soit $h(x)=1+0-0=1$;
				
				\item  $x\notin A$ et $x\in B$ : on a alors $f(x)=0$ et $g(x)=1$, soit $h(x)=0+1-0=1$.
				
			\end{itemize}
			
			
		\end{enumerate}
	\end{corrige}
	
	\subsection*{Applications}
	
	\begin{exo}[{Un exemple avec des fonctions}]
		Soit $f:\mtr\to\mtr$ définie par $f(x)=2x/(1+x^2)$.
		
		\begin{enumerate}
			
			\item  $f$ est-elle injective? surjective?
			
			\item  Montrer que $f(\mtr)=[-1,1]$.
			
			\item  Montrer que la restriction $g:[-1,1]\to[-1,1]$, $g(x)=f(x)$ est une bijection.
			
		\end{enumerate}
	\end{exo}
	\begin{corrige}
		\begin{enumerate}
			
			\item   $f$ n'est pas injective, car $f(2)=f(1/2)=4/5$. $f$ n'est pas surjective, car $y=2$ n'a pas d'antécédent :
			en effet, l'équation $f(x)=2$ devient $2x=2(1+x^2)$ soit $x^2-x+1=0$ qui n'a pas de solutions dans $\mtr$.
			
			\item  L'équation $f(x)=y$ est équivalent à l'équation $yx^2-2x+y=0$. Cette équation a des solutions réelles
			si et seulement si $\Delta=4-4y^2\geq 0$, donc il y a des solutions si et seulement si $y\in[-1,1]$. Ainsi,
			on a exactement $f(\mtr)=[-1,1]$.
			
			\item  Soit $y\in]-1,1[\backslash \{0\}$. Les solutions possibles de l'équation $g(x)=y$ sont $x=\frac{1-\sqrt{1-y^2}}{y}$ ou 
			$x=\frac{1+\sqrt{1-y^2}}{y}$. La deuxième solution n'appartient pas à $[-1,1]$ (elle est strictement supérieure à $1$ si $y>0$, et strictement inférieure à $-1$ si $y<0$). D'autre part,
			$x=\frac{1-\sqrt{1-y^2}}{y}=\frac{y}{1+\sqrt{1-y^2}}$ est dans $[-1,1]$. En effet, 
			$$1\leq 1+\sqrt{1-y^2}\implies 0<\frac 1{1+\sqrt{1-y^2}}\leq 1$$
			tandis que 
			$$-1<y<1.$$
			On en déduit que 
			$$-1\leq \frac{-1}{1+\sqrt{1-y^2}}\leq x\leq \frac 1{1+\sqrt{1-y^2}}\leq 1.$$
			D'autre part, si $y=1$, l'équation $g(x)=1$ a pour unique solution $x=1$ tandis que si $y=-1$, l'équation $g(x)=-1$ a pour unique solution $x=-1$. Enfin, si $y=0$, l'équation $g(x)=0$ admet pour unique solution $x=0$.\\
			
			Dans tous les cas, on a prouvé que pour tout $y\in [-1,1]$, l'équation $g(x)=y$ admet  
			une unique solution avec $x\in[-1,1]$. Nous avons bien prouvé que $g$ est une bijection. Bien sûr, une méthode d'analyse, en utilisant la continuité et la stricte monotonie de $g$, serait plus facile!
			
		\end{enumerate}
	\end{corrige}
	
	\begin{exo}[{Composition itérée}]
		Soit pour $x\in\mathbb R_+,$ $f(x)=\frac{x}{x+1}$. Déterminer $f\circ f\circ\dots\circ f(x)$ (où le symbole $f$ apparaît $n$ fois) en fonction de $n\in\mathbb N^*$ et de $x\in\mathbb R_+$.
	\end{exo}
	\begin{corrige}
		On va commencer par calculer $f\circ f(x)$ pour émettre une conjecture.
		On a 
		$$f\circ f(x)=\frac{\frac x{x+1}}{\frac x{x+1}+1}=\frac{x}{2x+1}.$$
		Ceci nous conduit à conjecturer que 
		$$f\circ f\circ\dots\circ f(x)=\frac{x}{nx+1}.$$
		Démontrons ce résultat par récurrence. Il est vrai au rang $n=1$ et $n=2$. Supposons le résultat démontré au rang $n-1$. Alors 
		$$f\circ f\circ\dots\circ f(x)=\frac{\frac x{x+1}}{(n-1)\frac x{x+1}+1}=\frac{ \frac x{x+1}}{\frac {(n-1)x+x+1}{x+1}}=\frac{x}{nx+1}.$$
		Le résultat est donc encore vrai au rang $n$. Par le principe de récurrence, il est vrai pour tout entier $n$.\\
		
		On peut se poser la question subsidiaire suivante : pour quelles valeurs de $x$ peut-on définir $f\circ f\circ\dots\circ f(x)$?
	\end{corrige}
	
	\begin{exo}[{Une bijection de $\mathbb N^2$ dans $\mathbb N$}]
		Soit $f:\mathbb N^2\to\mathbb N^*$, $(n,p)\mapsto 2^n(2p+1)$. Démontrer que $f$ est une bijection. En déduire une
		bijection de $\mathbb N^2$ sur $\mathbb N$.
	\end{exo}
	\begin{corrige}
		Remarquons d'abord que $f$ est bien à valeurs dans $\mathbb N^*$. Prouvons ensuite que $f$ est injective.
		Soient $(n,p)$ et $(m,q)$ deux éléments de $\mathbb N^2$ tels que $f(n,p)=f(m,q)$. Alors on a 
		$2^n(2p+1)=2^m(2q+1).$
		Supposons par exemple $n\geq m$. Alors on obtient
		$$2^{n-m}(2p+1)=2q+1.$$
		Si $n\neq m$, le terme de gauche est pair et celui de droite est impair, une contradiction. Donc $n=m$. On obtient alors
		$2p+1=2q+1$, soit aussi $p=q$. Finalement, on a bien $(n,p)=(m,q)$ et donc $f$ est injective. Prouvons ensuite que $f$
		est surjective. Soit $l\in\mathbb N^*$. Alors $l$ se décompose en produit de facteurs premiers
		$$l=2^n p_2^{\alpha_2}\dots p_r^{\alpha_r},$$
		avec $p_i$, $i\geq 2$, des nombres premiers impairs. Le produit de nombres impairs étant impair, on sait que $p_2^{\alpha_2}\dots p_r^{\alpha_r}$
		est impair et s'écrit donc $2p+1$, avec $p\in\mathbb N$. On a donc $l=2^n (2p+1)=f(n,p)$ et $f$ est surjective.
		Pour trouver une bijection de $\mathbb N^2$ dans $\mathbb N$, il suffit de composer avec une bijection de $\mathbb N^*$
		dans $\mathbb N$, par exemple l'application $n\mapsto n-1$. Ainsi, l'application $g:(n,p)\mapsto 2^n(2p+1)-1$ est une
		bijection de $\mathbb N^2$ dans $\mathbb N$.
	\end{corrige}
	
	\begin{exo}[{Un exemple avec de l'arithmétique}]
		Soit $f:\mathbb Z\times\mathbb N^*\to\mathbb Q$, $(p,q)\mapsto p+\frac 1q$. $f$ est-elle injective, surjective?
	\end{exo}
	\begin{corrige}
		On va commencer par prouver que $f$ est injective. Prenons $(p,q)$ et $(p',q')$ tels que $f(p,q)=f(p',q')$. On souhaite prouver
		que $(p,q)=(p',q')$. 
		$$f(p,q)=f(p',q')\iff p'-p=\frac 1{q}-\frac 1{q'}.$$
		Mais $0<\frac 1{q}\leq 1$ et $-1\leq -\frac 1{q'}<0$. On en déduit que
		$$-1<\frac 1{q}-\frac1{q'}<1$$
		On a donc $|p-p'|<1$, alors que $p-p'$ est un entier. On en déduit que $p=p'$, puis, utilisant de nouveau que $f(p,q)=f(p',q')$, que $q=q'$.\\
		
		On va ensuite prouver que $f$ n'est pas surjective. En effet, $2/3$ n'a pas d'antécédents. Sinon, il s'écrirait 
		$p+\frac 1q$. Puisque $p< p+1/q=2/3\leq p+1$, il faudrait nécessairement que $p=0$. Et donc on aurait
		$\frac 23=\frac 1q$ soit $q=3/2$, ce qui n'est pas un entier. Donc $f$ n'est pas surjective.
	\end{corrige}
	
	\begin{exo}[{Devinettes...}]
		\begin{enumerate}
			
			\item  Déterminer une bijection de $\mtn\to \mtn^*$.
			
			\item  Déterminer une bijection de $\{1/n;\ n\geq 1\}$ dans $\{1/n;\ n\geq 2\}$.
			
			\item  Déduire de la question précédente une bijection de $[0,1]$ dans $[0,1[$.
			
			\item  Déterminer une bijection de $\mtn\to\mtz$
			
		\end{enumerate}
	\end{exo}
	\begin{corrige}
		\begin{enumerate}
			
			\item  Posons $f:\mtn\to \mtn^*$, définie par $f(n)=n+1$. Remarquons $f$ est bien à image dans $\mtn^*$. Il reste à prouver
			que $f$ est bijective, ce qui est très facile avec la définition : si $n\in\mtn^*$, on a $f(k)=n\iff k+1=n\iff k=n-1$, l'équation
			$f(k)=n$ admet une unique solution dans $\mtn$, ce qui dit bien que $f$ est bijective.
			
			\item  Posons $g:\{1/n;\ n\geq 1\}\to \{1/n;\ n\geq 2\}$, définie par $g(1/n)=1/(n+1)$. Remarquons là aussi que l'ensemble
			d'arrivée est bien cohérent avec l'ensemble de départ. D'autre part, $g$ est bijective.
			
			\item  C'est plus compliqué!!! Ecrivons $[0,1]=\left\{1/n;\ n\geq 1\right\}\cup A$, où $A$ est le complémentaire de $\left\{1/n;\ n\geq 1\right\}$
			dans $[0,1]$. On définit $h$ de la façon suivante :
			
			\begin{itemize}
				
				\item  Si $x=1/n$, alors $h(x)=1/(n+1)$.
				
				\item  Sinon, c'est-à-dire si $x\in A$, $h(x)=x$.
				
			\end{itemize}
			
			Alors $h$ est bijective! Prouvons d'abord qu'elle est injective : si $h(x)=h(x')$, on distingue 3 cas :
			
			\begin{itemize}
				
				\item  Si $x\in A$ et $x'\in A$, alors $h(x)=x$ et $h(x')=x'$ ce qui entraîne $x=x'$.
				
				\item  Si $x\in A$ et $x'\notin A$, écrivant $x'=1/k$, on a $x=h(x)=h(x')=1/(k+1)$, ce qui implique $x\notin A$, ce qui est impossible.
				
				\item  Si $x\notin A$ et $x'\notin A$, écrivant $x=1/k$ et $x'=1/n$, on a $1/(k+1)=h(x)=h(x')=1/(n+1)$ ce qui entraîne $k+1=n+1$ et par suite $x=x'$.
				
			\end{itemize}
			
			Dans tous les cas possibles, on trouve $x=x'$, et $h$ est injective. Prouvons maintenant que $h$ est surjective, et choisissons $y\in[0,1[$. Si $y\in A$,
			en particulier $y\neq 1$, et on a $h(y)=y$. Si $y\notin A$, $y=1/n$, où $n$ est entier strictement plus grand que 1 puisque $y\neq 1$.
			On a alors $h(1/(n-1))=y$. Dans tous les cas, $y$ possède un antécédent, ce qui prouve que $h$ est surjective. 
			
			\item  Rappelons que tout entier naturer peut s'écrire $2k$ s'il est pair, et $2k+1$ s'il est impair, avec $k\in\mathbb N.$ Posons $f(2k)=k$, et $f(2k+1)=-k-1$. Reste à vérifier que
			$f$ est bijective, ce qui est laissé au lecteur!
			
		\end{enumerate}
	\end{corrige}
	
	\begin{exo}[{Composition, injectivité et surjectivité}]
		On considère 4 ensembles $A,B,C$ et $D$, et des applications $f:A\to B$, $g:B\to C$ et $h:C\to D$. Montrer que 
		$$g\circ f\textrm{ injective}\implies f\textrm{ injective,}$$
		$$g\circ f\textrm{ surjective}\implies g\textrm{ surjective.}$$
		Montrer que :
		$$\big(g\circ f\textrm{ et }h\circ g\textrm{ sont bijectives }\big)\iff \big(f,g\textrm{ et }h\textrm{ sont bijectives}\big).$$
	\end{exo}
	\begin{corrige}
		Première implication : si $f(a)=f(b)$, alors $g\circ f(a)=g\circ f(b)$, et puisque $g\circ f$ est injective, on en déduit $a=b$.\\
		
		Deuxième implication : soit $y\in C$. Puisque $g\circ f$ est surjective, il existe $a\in A$ avec $g\circ f(a)=y$. Posons 
		$b=f(a)$. On a alors $g(b)=y$, ce qui prouve que $g$ est surjective.\\
		
		Equivalence : d'abord, si $f$, $g$ et $h$ sont bijectives, la composée d'applications bijectives étant bijective,
		on en déduit que $g\circ f$ et $h\circ g$ sont bijectives. Réciproquement, puisque $g\circ f$ est bijective, elle
		est surjective, et on trouve que $g$ est surjective. D'autre part, puisque $h\circ g$ est bijective, elle est injective
		et donc $g$ est injective. On trouve donc que $g$ est bijective. Puisque $g\circ f$ est bijective, composant par $g^{-1}$ à gauche
		qui est bijective, $f$ est bijective. De même, puisque $h\circ g$ est bijective et que $g$ est bijective,
		$(h\circ g)\circ g^{-1}$ est bijective, et ceci est égal à $h$.
	\end{corrige}
	
	\begin{exo}[{Image directe de l'image réciproque...et vice-versa!}]
		Soient $E$ et $F$ deux ensembles et $f:E\to F$. Démontrer que 
		
		\begin{enumerate}
			
			\item  $\forall A \in {\mathcal{P}}(E),A \subset f^{ - 1} (f(A))$;
			
			\item  $\forall B \in {\mathcal{P}}(F),f(f^{ - 1} (B)) \subset B$.
			
			\item  Question subsidiaire (plus difficile) : a-t-on égalité en général?
			
		\end{enumerate}
	\end{exo}
	\begin{corrige}
		\begin{enumerate}
			
			\item  Soit $x\in A$. On a $f(x)\in f(A)$, et donc $x\in f^{-1}(f(A))$. Ainsi, on a prouvé l'inclusion $A\subset f^{-1}(f(A))$.
			
			\item  Soit $y\in f(f^{-1}(B))$. Alors il existe $x\in f^{-1}(B)$ tel que $y=f(x)$. Mais puisque $x\in f^{-1}(B)$, 
			alors $f(x)\in B$. Et donc $y=f(x)\in B$. On a bien prouvé l'inclusion $f(f^{ - 1} (B)) \subset B$.
			
			\item  Les inclusions peuvent être strictes. Prenons $f:\{1,2\}\to \{1,2\}$, $1\mapsto 1$ et $2\mapsto 1$
			et $B=\{1,2\}$. Alors $f^{-1}(B)=\{1,2\}$ et $f(f^{-1}(B))=\{1\}$ qui est différent de $B$.\\
			
			Pour l'autre exemple, prenons $A=\{1\}$. Alors $f(A)=\{1\}$ et $f^{-1}(f(A))=\{1,2\}$ qui est 
			différent de $A$.
			
		\end{enumerate}
	\end{corrige}
	
	\begin{exo}[{Ensembles et images réciproques}]
		Soient $E$ et $F$ deux ensembles et soit $f:E\to F$. Soient également $A$ et $B$ deux parties de $F$.
		
		\begin{enumerate}
			
			\item  Démontrer que $A\subset B\implies f^{-1}(A)\subset f^{-1}(B)$. La réciproque est-elle vraie?
			
			\item  Démontrer que $f^{-1}(A\cap B)=f^{-1}(A)\cap f^{-1}(B)$.
			
			\item  Démontrer que $f^{-1}(A\cup B)=f^{-1}(A)\cup f^{-1}(B)$.
			
		\end{enumerate}
	\end{exo}
	\begin{corrige}
		\begin{enumerate}
			
			\item   Prenons $x\in f^{-1}(A)$. Alors $f(x)\in A$ et donc $f(x)\in B$, et donc $x\in f^{-1}(B)$, ce
			qui prouve l'inclusion $f^{-1}(A)\subset f^{-1}(B)$. La réciproque n'est pas toujours vraie.
			Prenons $E=\{1\}$, $F=\{1,2\}$, $f:E\to F$ définie par $f(1)=1$, $A=\{2\}$ et $B=\{1\}$. 
			Alors $f^{-1}(A)=\varnothing\subset f^{-1}(B)=\{1\}$. Et pourtant, $A$ n'est pas inclus dans $B$.
			
			\item  On a $A\cap B\subset A$, et donc $f^{-1}(A\cap B)\subset f^{-1}(A)$. De même,
			$f^{-1}(A\cap B)\subset f^{-1}(B)$, et donc $f^{-1}(A\cap B)\subset f^{-1}(A)\cap f^{-1}(B)$.\\
			
			Réciproquement, soit $x\in f^{-1}(A)\cap f^{-1}(B)$. Alors on a à la fois $f(x)\in A$ et $f(x)\in B$, ce qui entraîne
			$f(x)\in A\cap B$. Ainsi, $x\in f^{-1}(A\cap B)$, et donc $f^{-1}(A)\cap f^{-1}(B)\subset f^{-1}(A\cap B)$.
			
			\item  On a $A\subset A\cup B$ et donc $f^{-1}(A)\subset f^{-1}(A\cup B)$. De même, $f^{-1}(B)\subset f^{-1}(A\cup B)$
			et donc $f^{-1}(A)\cup f^{-1}(B)\subset f^{-1}(A\cup B)$.\\
			
			Réciproquement, si $x\in f^{-1}(A\cup B)$, alors $f(x)\in A\cup B$. Ainsi, ou bien $f(x)\in A$, ou bien $f(x)\in B$. 
			Mais si $f(x)\in A$, on a $x\in f^{-1}(A)\subset f^{-1}(A)\cup f^{-1}(B)$ et de même, si $f(x)\in B$, on a $x\in f^{-1}(B)\subset f^{-1}(A)\cup f^{-1}(B)$.
			Dans tous les cas, on a prouvé que $x\in f^{-1}(A)\cup f^{-1}(B)$ et donc l'inclusion $f^{-1}(A\cup B)\subset f^{-1}(A)\cup f^{-1}(B)$.
			
		\end{enumerate}
	\end{corrige}
	
	\begin{exo}[{Ensembles et images directes}]
		Soient $E$ et $F$ deux ensembles et soit $f:E\to F$. Soient également $A$ et $B$ deux parties de $E$.
		
		\begin{enumerate}
			
			\item  Démontrer que $A\subset B\implies f(A)\subset f(B)$. La réciproque est-elle vraie?
			
			\item  Démontrer que $f(A\cap B)\subset f(A)\cap f(B)$. L'inclusion réciproque est-elle vraie?
			
			\item  Démontrer que $f(A\cup B)=f(A)\cup f(B)$.
			
		\end{enumerate}
	\end{exo}
	\begin{corrige}
		\begin{enumerate}
			
			\item   Prenons $y\in f(A)$. Alors il existe $x\in A$ tel que $y=f(x)$. Mais alors, $x\in B$ et donc $y\in f(B)$.
			Ceci prouve l'inclusion $f(A)\subset f(B)$. La réciproque n'est pas toujours vraie.
			Prenons $E=\{1,2\}$, $F=\{1\}$, $f:E\to F$ définie par $f(1)=1$, $f(2)=1$, $A=\{1\}$ et $B=\{2\}$. 
			Alors $f(A)=f(B)=\{1\}$ alors que pourtant $A$ n'est pas inclus dans $B$.
			
			\item  On a $A\cap B\subset A$, et donc $f(A\cap B)\subset f(A)$. De même,
			$f(A\cap B)\subset f(B)$, et donc $f(A\cap B)\subset f(A)\cap f(B)$.\\
			
			L'inclusion réciproque est fausse, ce que l'on constate en prenant exactement le même exemple.
			
			\item  On a $A\subset A\cup B$ et donc $f(A)\subset f(A\cup B)$. De même, $f(B)\subset f(A\cup B)$
			et donc $f(A)\cup f(B)\subset f(A\cup B)$.\\
			
			Réciproquement, si $y\in f(A\cup B)$, alors il existe $x\in A\cup B$ tel que $y=f(x)$.
			Mais si $x\in A$, on a $y\in f(A)\subset f(A)\cup f(B)$ et de même, si $x\in B$, on a $y\in f(B)\subset f(A)\cup f(B)$.
			Dans tous les cas, on a prouvé que $y\in f(A)\cup f(B)$ et donc l'inclusion $f(A\cup B)\subset f(A)\cup f(B)$.
			
		\end{enumerate}
	\end{corrige}
	
	\begin{exo}[{Image directe et injectivité}]
		Soit $f:X\to Y$. Montrer que les conditions suivantes sont équivalentes :
		
		\begin{enumerate}
			
			\item  $f$ est injective.
			
			\item  Pour toutes parties $A,B$ de $X$, on a $f(A\cap B)=f(A)\cap f(B)$.
			
		\end{enumerate}
	\end{exo}
	\begin{corrige}
		\begin{itemize}
			
			\item  $1\implies 2$ : D'abord, une inclusion est toujours vérifiée : prenons en effet $y\in f(A\cap B)$. Alors
			il existe $x\in A\cap B$ tel que $y=f(x)$. Mais alors, $y\in f(A)$ puisque $y=f(x)$ avec $x\in A$. De même,
			$y\in f(B)$. On en déduit que $y\in f(A)\cap f(B)$. Réciproquement, 
			si $y\in f(A)\cap f(B)$, alors il existe $a\in A$ tel que $y=f(a)$ et $b\in B$ tel que $y=f(b)$. Mais puisque
			$f$ est injective et que $f(a)=f(b)$, on a $a=b$, et donc $a\in A\cap B$. On en déduit que $y\in f(A\cap B)$.
			
			\item  $2\implies 1$ : Soit $a,b\in X$ tels que $f(a)=f(b)$. Notons $y=f(a)=f(b)$  et prenons $A=\{a\}$ et $B=\{b\}$. Remarquons que 
			$f(A)=f(B)=\{y\}$. Alors, puisque 2 est vrai, on a $f(A\cap B)=\{y\}\neq f(\varnothing)=\varnothing$. En particulier, $A\cap B\neq \varnothing$, et donc $a=b$.
			
		\end{itemize}
	\end{corrige}
	
	\begin{exo}[{Caractérisations}]
		Soient $X,Y$ deux ensembles et $f:X\to Y$ une application.
		
		\begin{enumerate}
			
			\item  Montrer que $f$ est injective si et seulement si, pour tout ensemble $Z$, pour tout $g:Z\to X$ et tout $h:Z\to X$, on a
			$f\circ g=f\circ h\implies g=h$.
			
			\item  Montrer que $f$ est surjective si et seulement si, pour tout ensemble $Z$, pour tout $g:Y\to Z$ et tout $h:Y\to Z$, on 
			a $g\circ f=h\circ f\implies g=h$.
			
		\end{enumerate}
	\end{exo}
	\begin{corrige}
		\begin{enumerate}
			
			\item  Supposons d'abord $f$ injective et soient $g:Z\to X$ et $h:Z\to X$ telles que
			$f\circ g=f\circ h$. Alors, pour tout $z$ de $Z$, on a $f(g(z))=f(h(z))\implies g(z)=h(z)$
			puisque $f$ est injective. On a donc bien $g=h$.\\
			
			Pour montrer l'implication réciproque, on procède par contraposée en supposant que $f$ n'est pas injective.
			Soient $x\neq y$ tel que $f(x)=f(y)$. Posons $Z=\{0\}$, $g(0)=x$ et $h(0)=y$. Alors on a $f\circ g(0)=f\circ h(0)=f(x)=f(y)$ alors que $g\neq h$.
			
			\item  Supposons d'abord $f$ surjective et soient $g:Y\to Z$ et $h:Y\to Z$ telles que
			$g\circ f=h\circ f$. Soit $y\in Y$. Il existe $x$ de $X$ tel que $y=f(x)$. On en déduit
			$g(y)=g\circ f(x)=h\circ f(x)=h(y)$, ce qui prouve $g=h$.\\
			
			Pour montrer l'implication réciproque, on procède par contraposée en supposant que $f$ n'est pas surjective.
			Il existe donc un point $y_0$ de $Y$ qui n'est pas dans $f(X)$. On considère alors $Z=\{0,1\}$, $g$ défini sur $Y$ par $g(y_0)=1$ et $g(y)=0$ sinon,
			$h$ défini sur $Y$ par $h(y)=0$ pour tout y. Alors on a bien $g\circ f=h\circ f$ (car $f(x)\neq y_0$ pour tout $x$ de $X$)
			et $h\neq g$.
			
		\end{enumerate}
	\end{corrige}
	
	\begin{exo}[{Bijectivité et passage au complémentaire}]
		Soit $f:E\to F$. Montrer que $f$ est bijective si et seulement si, pour tout $A$ de $\mathcal P(E)$,
		on a $f(\overline A)=\overline{f(A)}$ ($\overline A$ désigne le complémentaire de $A$).
	\end{exo}
	\begin{corrige}
		Pour l'implication directe, on suppose que $f$ est bijective, et prenons $A$ un élément de $\mathcal{P}(E)$. On doit montrer une double inclusion. Soit d'abord $x$ dans $f(\overline A)$. Alors
		$x=f(y)$ où $y\in\overline A$. Supposons que $x\in f(A)$. Alors $x=f(z)$ où $z\in A$. Mais alors, on a $f(y)=f(z)$ et par injectivité de $f$, on a $y=z$.
		Comme $y$ est élément de $A$ et $z$ est élément de son complémentaire, ceci est impossible et donc $x\notin f(A)$,
		c'est-à-dire qu'on a prouvé que $f(\overline A)\subset \overline{f(A)}$. \\
		
		Prouvons maintenant l'autre inclusion. Soit $x\in \overline{f(A)}$. Alors, puisque $f$ est surjective, il existe $y$ élément de $E$ tel que $x=f(y)$. Mais $y$ ne peut pas être élément de $A$ sinon $x$ serait élément de $f(A)$ ce qui n'est pas. Et donc $y$
		est élément de $\overline A$ et $x\in f(\overline A)$.\\
		
		Etudions maintenant l'implication réciproque, c'est-à-dire qu'on suppose que pour tout $A$
		de $\mathcal{P}(E)$, on a $f(\overline A)=\overline{f(A)}$. Prouvons d'abord que ceci entraine que $f$ est injective. En
		effet, pour $x,y$ de $E$ tels que $f(x)=f(y)$, supposons $x\neq y$. Posons $A=\{x\}$. On a $y\in \overline{A}$ et donc $f(y)\in
		f(\overline{A}) \subset \overline{f(A)}$. Or, $f(A)=\{f(x)\}$, et donc $f(y)\neq f(x)$, une contradiction.\\
		
		Prouvons enfin que $f$ est surjective. Par hypothèse appliquée à $A=E$, on sait que
		$f(\overline E)=\overline{f(E)}$. Mais $f(\overline E)=f(\emptyset)=\emptyset$,
		et donc $\overline{f(E)}=\emptyset$ ce qui, en prenant le complémentaire, se traduit en
		$f(E)=F$, c'est-à-dire que $f$ est surjective.
	\end{corrige}
	
	\begin{exo}[{Fonction définie sur l'ensemble des parties}]
		Soient $E$ un ensemble, $\mathcal P(E)$ l'ensemble de ses parties, et $A$ et $B$
		deux parties de $E$. On définit
		$$\begin{array}{cccc}
			f:&\mathcal P(E)&\to&\mathcal P(A)\times\mathcal P(B)\\
			&X&\mapsto &(X\cap A,X\cap B).
		\end{array}
		$$
		
		\begin{enumerate}
			
			\item  Montrer que $f$ est injective si et seulement si $A\cup B=E$.
			
			\item  Montrer que $f$ est surjective si et seulement si $A\cap B=\varnothing$.
			
			\item  Donner une condition nécessaire et suffisante sur $A$ et $B$ pour que $f$ soit bijective.
			Donner dans ce cas la bijection réciproque.
			
		\end{enumerate}
	\end{exo}
	\begin{corrige}
		\begin{enumerate}
			
			\item  Pour démontrer le sens direct, on raisonne par contraposée: si $A\cup B\neq E$, on prend
			$x\in E\backslash (A\cup B)$ et $X=\{x\}$. Alors $f(X)=(X\cap A,X\cap B)=(\varnothing,\varnothing)$ 
			car $x$ n'appartient ni à $A$ ni à $B$. D'autre part, $f(\varnothing)=(\varnothing,\varnothing)$. Donc
			$f(X)=f(\varnothing)$ alors que $X\neq\varnothing$: $f$ n'est pas injective.\\
			
			Pour le sens réciproque, remarquons que pour tout $X\subset E$, puisque $A\cup B=E$, on a 
			$$X=X\cap E=X\cap(A\cup B)=(X\cap A)\cup(X\cap B).$$
			Ainsi, si $X,X'\subset E$ sont tels que $f(X)=f(X')$, c'est-à-dire $X\cap A=X'\cap A$
			et $X\cap B=X'\cap B$, on a
			$$X=(X\cap A)\cup(X\cap B)=(X'\cap A)\cup(X'\cap B)=X'.$$
			Ainsi, $f$ est injective.
			
			\item  Supposons d'abord que $f$ est surjective et prenons $x\in A$.
			Alors il existe $X\subset E$ tel que $f(X)=(\{x\},\varnothing)$. 
			Alors, on a $X\cap B=\varnothing$ et $x\in X\cap A$. Ainsi, $x\in X$
			et donc $x\notin B$. Ainsi, on a $A\cap B=\varnothing$.\\
			
			Réciproquement, si $A\cap B=\varnothing$, et prenons $A'\subset A$ et
			$B'\subset B$. Alors, posons $X=A'\cup B'$. Puisque $A\cap B=\varnothing$,
			on a $X\cap A=A'$ et $X\cap B=B'$ et donc $f(X)=(A',B')$: $f$ est surjective.
			
			\item  D'après les questions précédentes, 
			on a $f$ bijective si et seulement si $A\cup B=E$ et $A\cap B=\varnothing$,
			ie si $(A,B)$ est une partition de $E$. La bijection réciproque
			a été établie à la question précédente et est donnée par
			$(A',B')\mapsto A'\cup B'$.
			
		\end{enumerate}
	\end{corrige}
	
	\begin{exo}[{Fonctions et fonctions d'ensemble}]
		Soient $E$ et $F$ deux ensembles, et $f:E\to F$. On définit deux applications $f^\sharp$ et $f_\sharp$ par :
		\begin{eqnarray*}
			f^\sharp:\mathcal P(E)\to\mathcal P(F),&&f^\sharp(A)=f(A)\\
			f_\sharp:\mathcal P(F)\to\mathcal P(E),&&f_\sharp(A)=f^{-1}(A).
		\end{eqnarray*}
		Démontrer que 
		
		\begin{enumerate}
			
			\item  $f^\sharp$ est injective si et seulement si $f$ est injective.
			
			\item  $f_\sharp$ est injective si et seulement si $f$ est surjective.
			
		\end{enumerate}
	\end{exo}
	\begin{corrige}
		\begin{enumerate}
			
			\item  Supposons d'abord que $f^\sharp$ est injective, et prouvons que $f$ l'est. Soient $x,y\in E$ avec $f(x)=f(y)$.
			Posons $A=\{x\}$ et $B=\{y\}$. Alors $f^\sharp(A)=f^\sharp(B)=\{f(x)\}$. Ainsi, par injectivité de $f^\sharp$, $A=B$ et donc $x=y$.
			$f$ est injective.\\
			
			Réciproquement supposons que $f$ est injective, et prouvons que $f^\sharp$ l'est aussi. Soient $A,B\in\mathcal P(E)$
			tels que $f^\sharp(A)=f^\sharp(B)$. Prenons ensuite $x\in A$ et montrons que $x\in B$. On a $f(x)\in f(A)=f(B)$, donc il existe
			$y\in B$ tel que $f(x)=f(y)$. Maintenant, puisque $f$ est injective, ceci entraîne que $x=y$. Ainsi, $x\in B$ et on a prouvé
			que $A\subset B$. Bien entendu, par symétrie du rôle joué par $A$ et $B$, on a $A=B$ et $f^\sharp$ est injective.
			
			\item  Supposons d'abord que $f_\sharp$ est injective, et prouvons que $f$ est surjective. Soit $y\in F$. Puisque $f_\sharp(\varnothing)=\varnothing$,
			on a $f_\sharp(\{y\})\neq\varnothing$, et donc il existe $x\in f_\sharp(\{y\})=f^{-1}(\{y\})$. Ainsi, $y=f(x)$ et $f$ est surjective.\\
			
			Supposons maintenant que $f$ est surjective, et prouvons que $f_\sharp$ est injective. Soient $A,B\in\mathcal P(F)$ tels que
			$f_\sharp(A)=f_\sharp(B)$. Considérons $y\in A$. Alors, puisque $f$ est surjective, il existe $x\in E$ tel que $y=f(x)$. 
			Mais $x\in f^{-1}(A)$ et donc $x\in f^{-1}(B)$ ce qui signifie que $f(x)\in B$. Mais $y=f(x)$, et donc $y\in B$. On a donc prouvé
			que $A\subset B$, et, toujours par symétrie du rôle joué par $A$ et $B$, on en déduit que $A=B$, c'est-à-dire que 
			$f_\sharp$ est injective. 
			
		\end{enumerate}
	\end{corrige}
	
	\subsection*{Relations}
	
	\begin{exo}[{Relation d'équivalence et fonction}]
		On définit sur $\mathbb R$ la relation $x\mathcal R y$ si et seulement si
		$x^2-y^2=x-y$.
		
		\begin{enumerate}
			
			\item  Montrer que $\mathcal R$ est une relation d'équivalence.
			
			\item  Calculer la classe d'équivalence d'un élément $x$ de $\mathbb R$.
			Combien y-a-t-il d'éléments dans cette classe?
			
		\end{enumerate}
	\end{exo}
	\begin{corrige}
		\begin{enumerate}
			
			\item  Il suffit de remarquer que $x\mathcal R y\iff x^2-x=y^2-y\iff f(x)=f(y)$
			avec $f:x\mapsto x^2-x$. Il est alors aisé de vérifier en appliquant la définition
			que $\mathcal R$ est une relation d'équivalence, c'est-à-dire qu'elle est réflexive, symétrique et transitive.
			
			\item  Soit $x\in\mathbb R$. On cherche les éléments $y$ de $\mathbb R$ tels que
			$x\mathcal R y$. On doit donc résoudre l'équation (en $y$) $x^2-y^2=x-y$. Elle se factorise
			en
			$$(x-y)(x+y)-(x-y)=0\iff (x-y)\times (x+y-1)=0.$$
			Ses solutions sont $y=x$ et $y=1-x$.
			La classe de $x$ est donc égale à $\{x,1-x\}$. Elle est constituée de deux éléments,
			sauf si $x=1-x\iff x=1/2$. Dans ce cas, elle est égale à $\{1/2\}$.
			
		\end{enumerate}
	\end{corrige}
	
	\begin{exo}[{Parties égales ou égales au complémentaire}]
		Soit $E$ un ensemble. On définit sur $\mathcal P(E)$, l'ensemble des parties de $E$, la relation suivante :
		$$A\mathcal R B\textrm{ si }A=B\textrm{ ou }A=\bar B,$$
		où $\bar B$ est le complémentaire de $B$ (dans $E$).
		Démontrer que $\mathcal R$ est une relation d'équivalence.
	\end{exo}
	\begin{corrige}
		Il suffit de vérifier que la définition d'une relation d'équivalence est satisfaite. Considérons
		trois parties $A$, $B$ et $C$ de $E$. Tout d'abord, 
		on a bien $A\mathcal R A$ : en effet, on a $A=A$. D'autre part, si $A\mathcal R B$, alors on distingue deux cas.
		Ou bien $A=B$, et dans ce cas $B=A$ et $B\mathcal R A$. Ou bien $A=\bar B$, mais dans ce cas, 
		$B=\bar A$, et on a encore $B\mathcal R A$. Enfin, si $A\mathcal R B$ et $B\mathcal R C$, on peut distinguer quatre cas :
		
		\begin{enumerate}
			
			\item  Si $A=B$ et $B=C$, alors $A=C$.
			
			\item  Si $A=B$ et $B=\bar C$, alors $A=\bar C$.
			
			\item  Si $A=\bar B$ et $B=C$, alors $A=\bar C$.
			
			\item  Si $A=\bar B$ et $B=\bar C$, alors $A=C$.
			
		\end{enumerate}
		
		Dans tous les cas, on a bien $A\mathcal R C$. La relation $\mathcal R$ est symétrique, réflexive et transitive.
		Il s'agit bien d'une relation d'équivalence.
	\end{corrige}
	
	\begin{exo}[{Une relation d'équivalence}]
		On définit sur $\mathbb Z$ la relation $x\mathcal R y$ si et seulement si
		$x+y$ est pair. Montrer qu'on définit ainsi une relation d'équivalence.
		Quelles sont les classes d'équivalence de cette relation?
	\end{exo}
	\begin{corrige}
		La relation est
		
		\begin{itemize}
			
			\item  réflexive, car $x+x=2x$ est pair;
			
			\item  symétrique, car $x+y=y+x$ et donc si $x+y$ est pair, $y+x$ est pair;
			
			\item  transitive, car si $x\mathcal R y$ et $y\mathcal R z$, alors
			$x+y=2k$ et $y+z=2l$ pour des entiers $k$ et $l$. Mais alors, on effectue la somme
			des ces deux égalités et on trouve
			$$x+2y+z=2k+2l\implies x+z=2(k+l-y)$$
			et donc $x+z$ est pair.
			
		\end{itemize}
		
		Pour déterminer les classes d'équivalence de $\mathcal R$, il suffit
		de trouver une famille $(E_i)$ d'ensembles tels que :
		
		\begin{itemize}
			
			\item  la réunion des $E_i$ est $\mathbb Z$;
			
			\item  les $E_i$ sont deux à deux disjoints;
			
			\item  si $x,y$ sont dans le même $E_i$, alors $x\mathcal R y$;
			
			\item  si $x$ est dans $E_i$ et $y$ est dans $E_j$ avec $i\neq j$,
			alors $x$ n'est pas en relation avec $y$.
			
		\end{itemize}
		
		Ici, on peut constater que tous les éléments en relation avec 0 sont les entiers pairs,
		tandis que tous les entiers en relation avec 1 sont les entiers impairs. Puisque l'ensemble des
		entiers pairs et des entiers impairs forme une partition de $\mathbb Z$, on en déduit que
		ces deux ensembles sont exactement les deux classes d'équivalence de la relation.
	\end{corrige}
	
	\begin{exo}[{Une relation d'ordre sur les entiers}]
		On définit la relation $\mathcal R$ sur $\mathbb N^*$ par $p\mathcal R q\iff \exists k\in\mathbb N^*,\ q=p^k$. Montrer que $\mathcal R$ définit un ordre partiel sur $\mathbb N^*$.
		Déterminer les majorants de $\{2,3\}$ pour cet ordre.
	\end{exo}
	\begin{corrige}
		La relation est 
		
		\begin{itemize}
			
			\item  réflexive : $p\mathcal R p$ puisque $p=p^1$ pour tout $p\in\mathbb N^*$;
			
			\item  antisymétrique : si $p\mathcal R q$ et $q\mathcal R p$, alors $p=q^k$ et $q=p^j$ avec $k,j\geq 1$ et donc $p=p^{jk}$. Ceci n'est possible
			que si $p=1$, mais alors $q=1=p$ ou si $jk=1$, ce qui implique $j=k=1$ et donc $p=q$.
			
			\item  transitive : si $p\mathcal R q$ et $q\mathcal R r$, alors $q=p^k$ et $r=q^j=p^{jk}$, donc $p\mathcal R r$.
			
		\end{itemize}
		
		On définit donc ainsi un ordre sur $\mathbb N^*$ qui n'est pas total (par exemple, on ne peut pas comparer $2$ et $3$).
		Soit maintenant $p$ un majorant de $\{2,3\}$. Alors $2\mathcal R p$ et donc $p=2^k$, avec $k\geq 1$. De même, $p=3^j$ avec $j\geq 1$. Par unicité
		de la décomposition en produits de facteurs premiers, ceci est impossible. L'ensemble $\{2,3\}$ n'a pas de majorant.
	\end{corrige}
	
	\begin{exo}[{Ordre lexicographique}]
		On définir sur $\mathbb R^2$ la relation $\prec$ par 
		$$(x,y)\prec (x',y')\iff \big( (x<x')\textrm{ ou }(x=x'\textrm{ et }y\leq y')\big).$$
		Démontrer que ceci définit une relation d'ordre sur $\mathbb R^2$.
	\end{exo}
	\begin{corrige}
		La relation est
		
		\begin{itemize}
			
			\item  réflexive : $x=x$ et $y\leq y$ impliquent $(x,y)\prec (x,y)$;
			
			\item  antisymétrique : si $(x,y)\prec (x',y')$ et $(x',y')\prec (x,y)$, alors on a nécessairement que 
			$x=x'$ (si $x<x'$ par exemple, on ne peut avoir $(x',y')\prec (x,y)$). Mais alors, on a à la fois $y\leq y'$ d'après la première relation, et aussi $y'\leq y$ d'après la seconde. On en déduit que $x=x'$ et $y=y'$;
			
			\item  transitive : si $(x,y)\prec (x',y')$ et $(x',y')\prec (x'',y'')$, alors :
			
			\begin{itemize}
				
				\item  ou bien $x=x'$ et $x'=x''$ : dans ce cas, on a $y\leq y'$ et $y'\leq y''$ donc $y\leq y''$ et donc
				$(x,y)\prec (x'',y'')$;
				
				\item  ou bien $x=x'$ et $x'<x''$ : dans ce cas, on a $x<x''$ et donc $(x,y)\prec (x'',y'')$;
				
				\item  ou bien $x<x'$ et $x'=x''$ : dans ce cas, on a $x<x''$ et donc $(x,y)\prec (x'',y'')$;
				
			\end{itemize}
			
			Remarquons que ceci définit un ordre total sur $\mathbb R^2$, car deux éléments sont toujours comparables.
		\end{itemize}
	\end{corrige}
	
	\begin{exo}[{Ordre sur $\mathbb R^2$}]
		On munit $\mathbb R^2$ de la relation notée $\prec$ définie par 
		$$(x,y)\prec (x',y')\iff x\leq x'\textrm{ et }y\leq y'.$$
		
		\begin{enumerate}
			
			\item  Démontrer que $\prec$ est une relation d'ordre sur $\mathbb R^2$. L'ordre est-il total?
			
			\item  Le disque fermé de centre $O$ et de rayon 1 a-t-il des majorants? un plus grand élément? une borne supérieure?
			
		\end{enumerate}
	\end{exo}
	\begin{corrige}
		\begin{enumerate}
			
			\item  La relation $\prec$ est 
			
			\begin{itemize}
				
				\item  réflexive : pour tout $(x,y)\in\mathbb R^2$, on a $x\leq x$ et $y\leq y$.
				
				\item  transitive : si $(x_1,y_1)\prec (x_2,y_2)$ et $(x_2,y_2)\prec (x_3,y_3)$, alors 
				$$x_1\leq x_2\leq x_3\textrm{ et }y_1\leq y_2\leq y_3$$
				donc $(x_1,y_1)\prec (x_3,y_3)$.
				
				\item  antisymétrique : si $(x,y)\prec (x',y')$ et $(x',y')\prec (x,y)$, alors on a à la fois
				$x\leq x'$ et $x'\leq x$ et donc $x=x'$ et de même $y=y'$.
				
			\end{itemize}
			
			Elle définit donc bien une relation d'ordre sur $\mathbb R^2$. L'ordre n'est pas total, car on ne peut pas comparer
			$(0,1)$ et $(1,0)$.
			
			\item  Soit $(x,y)$ un majorant de ce disque noté $D$. Alors $(1,0)\prec (x,y)$ et donc $x\geq 1$. De même, 
			$(0,1)\prec (x,y)$ et donc $y\geq 1$. Ainsi, on a $x\geq 1$ et $y\geq 1$. Réciproquement, soit $(x,y)\in\mathbb R^2$
			tel que $x\geq 1$ et $y\geq 1$. Alors $(x,y)$ est clairement un majorant de $D$, puisque tout élément $(x_0,y_0)$ de $D$ vérifie
			$x_0^2+y_0^2\leq 1$, et donc $x_0\leq 1$ et $y_0\leq 1$. On en déduit que l'ensemble des majorants de $D$ est 
			$$\{(x,y)\in\mathbb R^2;\ x\geq 1\textrm{ et }y\geq 1\}.$$
			$D$ n'admet donc pas de plus grand élément, puisque les majorants de $D$ ne sont pas dans $D$. En revanche, $D$ admet une borne supérieure qui
			est $(1,1)$, le plus petit des majorants de $D$.
			
		\end{enumerate}
	\end{corrige}
	
	\begin{exo}[{Relation d'équivalence et théorie des ensembles}]
		Soit $E$ un ensemble non-vide et $\alpha\subset\mathcal P(E)$ non-vide vérifiant la propriété suivante :
		$$\forall X,Y\in\alpha,\ \exists Z\in\alpha, Z\subset (X\cap Y).$$
		On définit sur $\mathcal P(E)$ la relation $\sim$ par $A\sim B\iff \exists X\in\alpha,\ X\cap A=X\cap B$.
		Prouver que ceci définit une relation d'équivalence sur $\mathcal P(E)$. Quelles sont les classes d'équivalence de $\varnothing$ et de $E$?
	\end{exo}
	\begin{corrige}
		Vérifions les 3 propriétés d'une relation d'équivalence :
		
		\begin{itemize}
			
			\item  La relation est symétrique :
			$$A\sim B\iff \exists X\in\alpha,\ X\cap A=X\cap B\iff B\sim A.$$
			
			\item  La relation est réflexive : soit $A\in\mathcal P(E)$ quelconque, et prenons $X\in\alpha$
			(peu importe lequel). Alors on a bien $X\cap A=X\cap A$ et donc $A\sim A$.
			
			\item  La relation est transitive : prenons $A,B,C\in\mathcal P(E)$ tels que $A\sim B$ et $B\sim C$.
			Alors on a
			$$A\sim B\iff \exists X\in\alpha,\ X\cap A=X\cap B$$
			$$B\sim C\iff \exists Y\in\alpha,\ Y\cap B=Y\cap C.$$
			Soit $Z\in\alpha$ tel que $Z\subset X\cap Y$. Alors on a $Z\cap X=Z$ et $Z\cap Y=Z$. De cela, on tire
			\begin{eqnarray*}
				Z\cap A&=&Z\cap X\cap A\\
				&=&Z\cap X\cap B\\
				&=&Z\cap B\\
				&=&Z\cap Y\cap B\\
				&=&Z\cap Y\cap C\\
				&=&Z\cap C.
			\end{eqnarray*}
			Ceci prouve bien que $A\sim C$ et que $\sim$ est une relation d'équivalence.
			
		\end{itemize}
		
		Cherchons ensuite la classe d'équivalence de $\varnothing$. On a
		$$A\sim\varnothing\iff \exists X\in\alpha,\  A\cap X=\varnothing\cap X=\varnothing.$$
		La classe d'équivalence de $\varnothing$ est donc constituée des parties $A$ de $E$ disjointes d'au moins un élément de $\alpha$.
		Cherchons enfin la classe d'équivalence de $E$. On a 
		$$A\sim E\iff \exists X\in\alpha,\ A\cap X=E\cap X=X.$$
		La classe d'équivalence de $E$ est donc constituée des parties $A$ de $E$ qui contiennent au moins un élément de $\alpha$.
	\end{corrige}
	
	\subsection*{Exercices complémentaires}
	
	\begin{exo}[{Opérations sur les ensembles}]\ \textcolor{red!70!black}{$\star$}
		Soient $A,B,C$ trois parties d'un ensemble $E$. Démontrer :
		\begin{enumerate}
			\item $A\cap(B\cup C)=(A\cap B)\cup(A\cap C)$ ;
			\item $A\setminus(B\cap C)=(A\setminus B)\cup(A\setminus C)$ ;
			\item $A\cap B=A\cap C$ et $A\cup B=A\cup C$ impliquent $B=C$.
		\end{enumerate}
	\end{exo}
	\begin{corrige}
		\begin{enumerate}
			\item $x\in A\cap(B\cup C)\iff x\in A$ et ($x\in B$ ou $x\in C$) $\iff$ ($x\in A$ et $x\in B$) ou ($x\in A$ et $x\in C$) $\iff x\in(A\cap B)\cup(A\cap C)$.
			\item $x\in A\setminus(B\cap C)\iff x\in A$ et ($x\notin B$ ou $x\notin C$) $\iff x\in(A\setminus B)\cup(A\setminus C)$.
			\item Soit $x\in B$. Si $x\in A$, alors $x\in A\cap B=A\cap C$, donc $x\in C$. Si $x\notin A$, comme $x\in A\cup B=A\cup C$, on a $x\in C$. Donc $B\subset C$ ; par symétrie $C\subset B$.
		\end{enumerate}
	\end{corrige}
	
	\begin{exo}[{Images directes et réciproques}]\ \textcolor{red!70!black}{$\star$}
		Soit $f:E\to F$ une application, $A,B\subset E$ et $C,D\subset F$.
		\begin{enumerate}
			\item Montrer que $f(A\cup B)=f(A)\cup f(B)$ et $f(A\cap B)\subset f(A)\cap f(B)$.
			\item Montrer que $f^{-1}(C\cap D)=f^{-1}(C)\cap f^{-1}(D)$.
			\item Montrer que $f$ est injective si et seulement si $f(A\cap B)=f(A)\cap f(B)$ pour toutes parties $A,B$ de $E$.
		\end{enumerate}
	\end{exo}
	\begin{corrige}
		\begin{enumerate}
			\item $y\in f(A\cup B)\iff\exists x\in A\cup B,\ y=f(x)\iff y\in f(A)$ ou $y\in f(B)$. De plus $A\cap B\subset A$ et $A\cap B\subset B$ donnent $f(A\cap B)\subset f(A)\cap f(B)$.
			\item $x\in f^{-1}(C\cap D)\iff f(x)\in C$ et $f(x)\in D\iff x\in f^{-1}(C)\cap f^{-1}(D)$.
			\item ($\Rightarrow$) Si $f$ est injective et $y\in f(A)\cap f(B)$, alors $y=f(a)=f(b)$ avec $a\in A$, $b\in B$ ; par injectivité $a=b\in A\cap B$, donc $y\in f(A\cap B)$. Avec la question 1, on a l'égalité.
			($\Leftarrow$) Si $f(x)=f(x')$, posons $A=\{x\}$, $B=\{x'\}$. Alors $f(A\cap B)=f(A)\cap f(B)=\{f(x)\}\neq\emptyset$, donc $A\cap B\neq\emptyset$, c'est-à-dire $x=x'$.
		\end{enumerate}
	\end{corrige}
	
	\begin{exo}[{Injectivité, surjectivité}]\ \textcolor{red!70!black}{$\star$}
		Étudier l'injectivité et la surjectivité des applications suivantes :
		\begin{enumerate}
			\item $f:\N\to\N,\ n\mapsto 2n$ ;
			\item $g:\N\to\N,\ n\mapsto \lfloor n/2\rfloor$ ;
			\item $h:\Z\to\Z,\ n\mapsto n^2$ ;
			\item $k:\R\to\R,\ x\mapsto x^3-x$.
		\end{enumerate}
		Calculer $g\circ f$ et $f\circ g$. Que remarque-t-on ?
	\end{exo}
	\begin{corrige}
		\begin{enumerate}
			\item $f$ est injective, non surjective ($1$ n'a pas d'antécédent).
			\item $g$ n'est pas injective ($g(0)=g(1)=0$), mais elle est surjective ($g(2m)=m$).
			\item $h$ n'est ni injective ($h(1)=h(-1)$) ni surjective ($2$ n'est pas un carré).
			\item $k$ n'est pas injective ($k(0)=k(1)=0$), mais elle est surjective (continue, de limites $\pm\infty$ en $\pm\infty$, théorème des valeurs intermédiaires).
		\end{enumerate}
		On a $g\circ f(n)=\lfloor 2n/2\rfloor=n$, donc $g\circ f=\id_\N$. En revanche $f\circ g(n)=2\lfloor n/2\rfloor$ vaut $n$ si $n$ est pair et $n-1$ sinon : $f\circ g\neq\id_\N$. Ainsi $g\circ f=\id$ ne suffit pas à avoir $f\circ g=\id$ : $f$ est seulement injective et $g$ seulement surjective.
	\end{corrige}
	
	\begin{exo}[{Relations d'équivalence et d'ordre}]\ \textcolor{red!70!black}{$\star$}
		\begin{enumerate}
			\item Sur $\Z$, on définit $x\sim y\iff x^2-y^2$ est divisible par $4$. Montrer que $\sim$ est une relation d'équivalence et déterminer ses classes.
			\item Sur $\N^*$, montrer que la divisibilité est une relation d'ordre. L'ordre est-il total ?
		\end{enumerate}
	\end{exo}
	\begin{corrige}
		\begin{enumerate}
			\item Réflexivité : $x^2-x^2=0$. Symétrie : évidente. Transitivité : $x^2-z^2=(x^2-y^2)+(y^2-z^2)$. Un carré est congru à $0$ modulo $4$ si l'entier est pair, à $1$ s'il est impair ; donc $x\sim y$ si et seulement si $x$ et $y$ ont même parité. Il y a deux classes : les entiers pairs et les entiers impairs.
			\item $a\mid a$ ; si $a\mid b$ et $b\mid a$ avec $a,b>0$, alors $a=b$ ; si $a\mid b$ et $b\mid c$ alors $a\mid c$. L'ordre n'est pas total : $2$ et $3$ ne sont pas comparables.
		\end{enumerate}
	\end{corrige}
	\section{Arithmétique dans $\Z$}
	
	\subsection*{Divisibilité et division euclidienne}
	
	\begin{exo}[{Comprendre la division euclidienne}]
		\begin{enumerate}
			
			\item  Déterminer les entiers positifs $a$ et $b$ sachant que $a<4000$ et que la division euclidienne de $a$ par $b$ donne un quotient de 82 et un reste de 47.
			
			\item  Déterminer le quotient et le reste de la division euclidienne de $2^{2013}+562$ par $4$.
			
			\item  Quand on divise un nombre par 12, le reste est 8. Quand on divise ce même nombre par 10, on augmente le quotient de $1$ et le reste devient 2. Quel est ce nombre?
			
			\item  Démontrer que sur la droite $y=\frac 34x+\frac 18$, il n'y a pas de points à coordonnées entières.
			
		\end{enumerate}
	\end{exo}
	\begin{corrige}
		\begin{enumerate}
			
			\item  On sait que $a=82b+47$, et que $b>47$. Mais si $b=49$, alors $82\times 49+47>4000$. La seule possibilité est donc $b=48$, et $a=3983$.
			
			\item  On commence par remarquer que $2^{2013}=4\times 2^{2011}$. Pour l'autre partie, $562=4\times 140+2$. On en déduit donc que 
			$$2^{2013}+562=4\times(2^{2011}+140)+2.$$
			Ainsi, le quotient est $2^{2011}+140$ et le reste est $2$.
			
			\item  On écrit deux fois la division euclidienne. Notons $a$ le nombre recherché, et $q$ le premier quotient, on trouve 
			$a=12q+8$ et $a=10(q+1)+2$. Il s'agit maintenant d'un petit système à résoudre, et on trouve $a=32$.
			
			\item  Supposons qu'il existe un point $(x,y)$ à coordonnées entières sur cette droite. Alors, multipliant l'équation de la droite par $8$, on trouve que 
			$$8y=6x+1.$$
			Mais $2|8y$, $2|6x$ et donc $2|1$, ce qui n'est pas le cas.
			
		\end{enumerate}
	\end{corrige}
	
	\begin{exo}[{Divisibilité et identité remarquable}]
		Démontrer que, pour tous $x,y\in \mathbb R^*$ et tout $n\in\mathbb N\backslash\{0\}$, on a 
		$$x^n-y^n=(x-y)\sum_{k=0}^{n-1}x^k y^{n-1-k}.$$
		En déduire que $609|5^{4n}-2^{4n}.$
	\end{exo}
	\begin{corrige}
		Développons le membre de droite de l'identité que l'on nous demande de démontrer.
		\begin{eqnarray*}
			(x-y)\sum_{k=0}^{n-1}x^k y^{n-1-k}&=&\sum_{k=0}^{n-1}x^{k+1}y^{n-1-k}-\sum_{k=0}^{n-1}x^k y^{n-k}\\
			&=&\sum_{k=1}^n x^ky^{n-k}-\sum_{k=0}^{n-1}x^k y^{n-k}
		\end{eqnarray*}
		en effectuant un changement d'indice dans la première somme. Mais les deux sommes sont identiques, à deux termes près, et on a donc
		$$(x-y)\sum_{k=0}^{n-1}x^k y^{n-1-k}=x^ny^{n-n}-x^0y^{n-0}=x^n-y^n.$$
		Le fait que $609|5^{4n}-2^{4n}$ est alors une application immédiate de cette identité à $x=5^4$ et $y=2^4$, puisque dans ce cas $x-y=609$.
	\end{corrige}
	
	\begin{exo}[{Une relation de divisibilité}]
		Déterminer les entiers relatifs $n$ tels que $n-4$ divise $3n-17$.
	\end{exo}
	\begin{corrige}
		Puisque $n-4$ divise $3n-17$ et que $n-4$ divise $n-4$, $n-4$ divise également $(3n-17)-3(n-4)=-5$. 
		Or les diviseurs dans $\mathbb Z$ de $-5$ sont $-5,-1,1,5$. Les valeurs possibles de $n-4$ sont donc ces valeurs, et donc on a $n\in \{-1,3,5,9\}$.\\
		
		Réciproquement, si $n=-1$, alors $n-4=-5$ divise $3n-17=-20$. Si $n=3$, $n-4=-1$ divise $3n-17=-8$. Si $n=5$, $n-4=1$ divise $3n-17=-2$. Si $n=9$, $n-4=5$ divise $3n-17=10$. \\
		
		En conclusion, les valeurs de $n$ qui conviennent sont $-1,3,5$ et $9$.
	\end{corrige}
	
	\begin{exo}
		Soit $n\geq 1$ un entier. Déterminer le reste dans la division euclidienne par $n$ de la somme des $n$ premiers entiers strictement positifs.
	\end{exo}
	\begin{corrige}
		Soit $S_n$ la somme des $n$ premiers entiers. Elle vaut $\frac{n(n+1)}2$. Si $n$ est impair, autrement dit si $n+1$ est pair,
		alors $(n+1)/2$ est un entier, et $n|S_n$. Le reste est donc 0. Si maintenant $n=2k$ est pair, alors
		$$S_n=\frac{2k(2k+1)}2=\frac{(2k)^2}{2}+k=k\times 2k+k=k\times n+n/2$$
		Le reste est donc $n/2$.
	\end{corrige}
	
	\begin{exo}[{Division euclidienne avec de grands nombres}]
		Soient $a,b,n$ trois entiers supérieurs ou égaux à 1. On note $q$ le quotient de la division euclidienne de $a-1$ par $b$, et $r$ le reste. Déterminer le quotient et le reste de la division euclidienne de $ab^n-1$ par $b^{n+1}$.
	\end{exo}
	\begin{corrige}
		La division euclidienne s'écrit $a-1=bq+r$, avec $0\leq r\leq b-1$. On multiplie ensuite tout par $b^n$ et on trouve $ab^n-b^n=b^{n+1}q+rb^n$, puis on se ramène à 
		$$ab^n-1=b^{n+1}q+rb^n+b^n-1=b^{n+1}q+r'$$
		où $r'=rb^n+b^n-1$. On remarque que $0\leq r'\leq (b-1)b^n+b^n-1=b^{n+1}-1$. On a donc écrit la division euclidienne de $ab^n-1$ par $b^{n+1}$, le quotient est $q$ et le reste est $rb^n+b^n-1$.
	\end{corrige}
	
	\begin{exo}[{Grands nombres divisibles}]
		Montrer que pour tout entier $n\geq 1$, $40^n n!|(5n)!$.
	\end{exo}
	\begin{corrige}
		On raisonne par récurrence sur $n$. La propriété est clairement vérifiée
		si $n=1$, car $5!=120=40\times 3$. Supposons la propriété vraie au rang $n$, et prouvons la au rang $n+1$.
		On a
		\begin{align*}
			(5n+5)!&=(5n)!(5n+1)(5n+2)(5n+3)(5n+4)(5n+5)\\
			&= (5n)! 5(n+1)(5n+1)(5n+2)(5n+3)(5n+4)
		\end{align*}
		et
		$$40^{n+1} (n+1)!=40^n n!\times 40(n+1).$$
		Puisque par hypothèse de récurrence, $40^n n!|(5n)!$, il suffit de prouver que $40(n+1)|5(n+1)(5n+1)(5n+2)(5n+3)(5n+4)$, ou encore
		$8|(5n+1)(5n+2)(5n+3)(5n+4)$. Mais $5n+1,5n+2,5n+3,5n+4$ sont quatre entiers consécutifs.
		Deux sont pairs, et l'un des ces deux entiers pairs est même divisible par 4. Donc 8 divise bien leur produit.
	\end{corrige}
	
	\begin{exo}[{Équations diophantiennes du second degré résolues par factorisation}]
		Résoudre les équations d'inconnues $(x,y)\in\mathbb N^2$ :
		
		\begin{enumerate}
			
			\item  $x^2-y^2=7$;
			
			\item  $9x^2-y^2=32$.
			
		\end{enumerate}
	\end{exo}
	\begin{corrige}
		\begin{enumerate}
			
			\item  On factorise $x^2-y^2=(x-y)(x+y)$. Puisque les seuls diviseurs de $7$ sont $1$ et $7$, et puisque $x+y\geq x-y$, $x$ et $y$ sont solutions du système
			$$\left\{
			\begin{array}{rcl}
				x-y&=&1\\
				x+y&=&7
			\end{array}\right.$$
			dont la seule solution est $x=4$ et $y=3$. L'équation $x^2-y^2=7$ admet donc pour unique solution le couple $(4,3)$.
			
			\item  On procède de la même façon en remarquant que $9x^2-y^2=(3x-y)(3x+y)$. Mais cette fois, $32$ a beaucoup plus de diviseurs et on peut l'écrire sous la forme
			$$32=32\times 1=16\times 2=8\times 4.$$
			Tenant compte de $3x+y\geq 3x-y$, on trouve que $(x,y)$ est solution d'un des trois systèmes suivants :
			
			\begin{itemize}
				
				\item  $\left\{\begin{array}{rcl}
					3x+y&=&32\\
					3x-y&=&1.
				\end{array}\right.\quad$ 
				Mais ce système n'a pas de solutions entières...
				
				\item   $\left\{\begin{array}{rcl}
					3x+y&=&16\\
					3x-y&=&2.
				\end{array}\right.\quad$
				Dans ce cas, on trouve $x=3$ et $y=7$.
				
				\item   $\left\{\begin{array}{rcl}
					3x+y&=&8\\
					3x-y&=&4.
				\end{array}\right.\quad$  
				Dans ce cas, on trouve $x=2$ et $y=2$.
				
			\end{itemize}
			
			Finalement, on a prouvé que l'ensemble des solutions est $\{(3,7);(2,2)\}$.
			
		\end{enumerate}
	\end{corrige}
	
	\begin{exo}[{Écriture en base $b$}]
		Soit $b\geq 2$ un entier. On souhaite démontrer que tout entier $n\geq 1$ s'écrit uniquement
		$$n=\sum_{k=0}^p a_k b^k$$
		avec $p\geq 0$, $a_k\in\{0,\dots,b-1\}$ et $a_p\geq 1$. 
		
		\begin{enumerate}
			
			\item  Existence : démontrer l'existence en procédant par récurrence forte. Pour l'hérédité, on pourra utiliser la division euclidienne de $n$ par $b$.
			
			\item  Unicité : on suppose que $n$ admet deux décompositions distinctes 
			\[ n=\sum_{k=0}^p a_k b^k\textrm{ et }n=\sum_{k=0}^{p'}a'_k b^k. \]
			On peut supposer $p\geq p'$. Quitte à compléter la suite $a'_k$ par $a'_{p'+1}=\dots=a'_p=0$, on peut supposer que $p=p'$. Soit $\ell\in\{0,\dots,p\}$ le plus grand possible tel que $a_\ell\neq a'_\ell$.
			
			\begin{enumerate}
				
				\item  Vérifier que $(a_\ell-a'_\ell)b^\ell=\sum_{k=0}^{\ell-1}(a'_k-a_k)b^k$.
				
				\item  Démontrer que, pour toute suite finie $c_0,\dots,c_{\ell-1}$ avec $0\leq c_k\leq b-1$, on a 
				\[\sum_{k=0}^{\ell-1}c_k b^k < b^\ell. \]
				
				\item  Conclure.
				
			\end{enumerate}
			
			
			\item  Donner l'écriture de 37 (écrit en base 10) en base 2, puis en base 3.
			
		\end{enumerate}
	\end{exo}
	\begin{corrige}
		\begin{enumerate}
			
			\item  Pour $n\geq 1$, on note $H_n$ la propriété suivante : 
			$$H_n="\exists p\geq 0,\exists (a_0,\dots,a_p)\in\{0,\dots,b-1\}\textrm{ avec }a_p\geq 1\textrm{ tels que }n=\sum_{k=0}^p a_k b^k. "$$
			Initialisation : $H_1$ est vraie, car $1=1\times b^0$ (on a $p=0$ et $a_0=1$).\\
			
			Hérédite : Soit $n\geq 1$ et supposons $H_1,\dots,H_{n}$ vraies. Prouvons $H_{n+1}$. On distingue deux cas : 
			
			\begin{itemize}
				
				\item  soit $n+1<b$. Dans ce cas, on peut écrire $n+1=(n+1)b^0$, et on a le résultat avec $p=0$ et $a_0=n+1$.
				
				\item  soit $n+1\geq b$. Dans ce cas, on réalise la division euclidienne de $n+1$ par $b$. On écrit donc 
				$$n+1=bq+r,\ 0\leq r\leq b-1.$$
				On a alors $q\leq n$ et on peut appliquer l'hypothèse de récurrence à $q$ : $q$ s'écrit 
				$$q=\sum_{k=0}^p a_k b^k,\ 0\leq a_k\leq b-1,\ a_p\geq 1.$$
				On obtient alors 
				$$n+1=\sum_{k=0}^p a_k b^{k+1}+r=r+\sum_{k=1}^{p+1} a_{k-1} b^k.$$
				On obtient alors bien que $n+1$ s'écrit
				$$n+1=\sum_{k=0}^{p'} a'_k b^k, 0\leq a'_k\leq b-1,\ a'_{p'}\geq 1$$
				en posant $p'=p+1$, $a'_k=a_{k-1}$ pour $k= 1,\dots,p'$ et $a'_0=r$.
				
			\end{itemize}
			
			Dans tous les cas $H_{n+1}$ est vraie.\\
			
			Donc, par le principe de récurrence forte, $H_n$ est vraie pour tout entier $n\geq 1$.
			
			\item  
			\begin{enumerate}
				
				\item  Il suffit d'écrire que 
				$$n=\sum_{k=0}^p a_k b^k=\sum_{k=0}^p a'_k b^k$$
				ce qui s'écrit encore, par définition de $\ell$,
				$$n=\sum_{k=0}^\ell a_k b^k=\sum_{k=0}^\ell a'_k b^k.$$
				Ceci s'écrit à son tour 
				$$a_\ell b^\ell+\sum_{k=0}^{\ell-1} a_k b^k=a'_\ell b^\ell+\sum_{k=0}^{\ell-1} a'_k b^k$$
				ce qui donne immédiatement le résultat voulu.
				
				\item  On écrit simplement que 
				$$\sum_{k=0}^{\ell-1} c_k b^k\leq (b-1)\sum_{k=0}^{\ell-1}b^k=b^\ell-1<b^\ell$$
				où on a reconnu une somme géométrique.
				
				\item  Supposons par exemple que $a_\ell>a'_\ell$. On a alors, puisque $|a'_k-a_k|\leq b-1$,
				$$\left|\sum_{k=0}^{\ell-1} (a'_k-a_k)b^k\right|\leq\sum_{k=0}^{\ell-1} |a'_k-a_k| b^k<b^{\ell}\leq (a_\ell-a'_\ell)b^\ell$$
				ce qui contredit l'hypothèse faite. On a donc bien démontré l'unicité de la décomposition en base b.
				
			\end{enumerate}
			
			
			\item  On va faire des divisions euclidiennes successives :
			\begin{align*}
				37&=2\times 18+1\\
				&=2\times (2\times 9+0)+1\\
				&=2\times (2\times (2\times 4+1))+1\\
				&=2\times (2\times (2\times 2^2+1))+1\\
				&=2^5+2^2+2^0
			\end{align*}
			$37$ s'écrit donc en base $2$ : $100101$. En base 3, on a 
			\begin{align*}
				37&=3\times 12+1\\
				&=3\times ( 3\times 4)+1\\
				&=3\times ( 3\times(3+1))+1\\
				&=3^3+3^2+3^0.
			\end{align*}
			$37$ s'écrit donc en base $3$ : $1101$.
			
		\end{enumerate}
	\end{corrige}
	
	\begin{exo}[{Une équation de type Fermat}]
		Soit $k\in\mathbb N^*$. Le but de l'exercice est de déterminer les entiers $a,b\in\mathbb N^*$ solutions de $a^2+b^2=2^k.$
		
		\begin{enumerate}
			
			\item  Démontrer que si $N,a$ et $b$ sont des entiers tels que $N=a^2+b^2$ et $N$ est un multiple de $4$, alors $a$ et $b$ sont pairs.
			
			\item  En déduire que l'équation $2^{2n}=a^2+b^2$, $n\in\mathbb N$, $a,b\in\mathbb N^*$ n'admet pas de solutions.
			
			\item  Démontrer que l'équation $2^{2n+1}=a^2+b^2$, $n\in\mathbb N$, $a,b\in\mathbb N^*$ admet une unique solution que l'on précisera.
			
		\end{enumerate}
	\end{exo}
	\begin{corrige}
		\begin{enumerate}
			
			\item  On distingue plusieurs cas. Si $a$ et $b$ sont tous deux impairs, $a=2k+1$ et $b=2l+1$, alors $a^2+b^2=4(k^2+l^2+k+l)+2$ n'est pas un multiple de 4. Si $a$ est pair et $b$ est impair, $a=2k$ et $b=2l+1$, alors $a^2+b^2=4(k^2+l^2+l)+1$ n'est pas un multiple de 4. Le raisonnement est identique si $a$ est impair et $b$ est pair. Dans tous les cas, pour que $a^2+b^2$ soit un multiple de 4, il est nécessaire que $a$ et $b$ soient pairs.
			
			\item  Imaginons qu'il existe au moins une solution à l'équation, et choisissons une solution $2^{2n}=a^2+b^2$ avec $n$ le plus petit possible. Alors $n\geq 1$, car l'équation $2^0=1=a^2+b^2$ n'admet pas de solutions avec $a\geq 1$ et $b\geq 1$. Donc, $2^{2n}$ est un multiple de 4. D'après la question précédente, $a=2k$ et $b=2l$ sont des nombres pairs. Ainsi, on a encore 
			$$2^{2n}=4k^2+4l^2\iff 2^{2(n-1)}=k^2+l^2.$$
			Ainsi, on a trouvé une solution à l'équation avec $n-1<n$, ce qui contredit la minimalité de $n$.  Donc l'équation n'admet aucune solution.
			
			\item  Démontrons par récurrence sur $n$ que l'équation (portant sur $a,b\in\mathbb N^*$) $2^{2n+1}=a^2+b^2$ admet pour unique solution $a=b=2^n$. Pour $n=0$, on doit résoudre $2=a^2+b^2$ avec $a,b\in\mathbb N^*$, ce qui entraîne bien $a=b=1$. Soit $n\in\mathbb N$ et supposons que la propriété est vraie au rang $n$. Prouvons-la au rang $n+1$. Alors, 
			$2^{2n+3}$ est un multiple de 4. D'après la première question, si $a^2+b^2=2^{2(n+1)+1}$, alors $a=2k$, $b=2l$ et donc 
			$$2^{2n+1}=k^2+l^2.$$
			Par hypothèse de récurrence, $k=l=2^{n}$, et donc $a=b=2^{n+1}$. La propriété est donc vraie au rang $n+1$. Par le principe de récurrence, elle est vraie pour tout $n\geq 0$. 
			
		\end{enumerate}
	\end{corrige}
	
	\subsection*{Pgcd, ppcm et nombres premiers entre eux}
	
	\begin{exo}[{Pour bien commencer...}]
		\begin{enumerate}
			
			\item  Quels sont les diviseurs positifs communs à $390$ et $525$?
			
			\item  Calculer $(3^{123}-5)\wedge 25$.
			
			\item  Soit $n\in\mathbb Z$. Démontrer que $6|n(n+1)(n+2)$.
			
		\end{enumerate}
	\end{exo}
	\begin{corrige}
		\begin{enumerate}
			
			\item  En utilisant l'algorithme d'Euclide par exemple, on trouve que le pgcd de $390$ et $525$ est égal à $15.$ Les diviseurs communs de $390$ et $525$ sont les diviseurs de $15.$ Ces diviseurs (positifs) sont donc $1,\ 3,\ 5$ et $15$.
			
			\item  Notons $d$ le pgcd à calculer. Puisqu'il divise les deux nombres, il divise
			$25$ et $d$ est donc égal à $1$, $5$ ou $25$. S'il est égal à $5$ ou $25$,
			ceci signifie que $5$ divise $3^{123}-5$. Mais alors, puisque
			$5$ divise $-5$, ceci entrainerait encore que $5$ divise $3^{123}$, ce
			qui n'est pas vrai. Donc $d=1$.
			
			\item  $n$ et $n+1$ sont deux entiers consécutifs. L'un des deux au moins
			est pair, et donc $2|n(n+1)(n+2)$. De même, $n$, $n+1$ et $n+2$ sont trois entiers consécutifs.
			L'un au moins est un multiple de $3,$ et donc $3|n(n+1)(n+2)$. Par le théorème de Gauss,
			puisque $2$ et $3$ sont premiers entre eux,
			$6|n(n+1)(n+2)$.
			
		\end{enumerate}
	\end{corrige}
	
	\begin{exo}[{Des équations de Bezout}]
		Résoudre les équations suivantes dans $\mathbb Z^2$ :
		
		\begin{enumerate}
			
			\item  $2x+5y=3$;
			
			\item  $323x-391y=612$;
			
			\item  $162x+207y=27$;
			
			\item  $221x+247y=15$.
			
		\end{enumerate}
	\end{exo}
	\begin{corrige}
		\begin{enumerate}
			
			\item  Commençons par remarquer que $2\wedge 5=1$ et donc, d'après le théorème de Bézout, l'équation admet toujours une solution. Ici, il est facile de remarquer que $(-1,1)$ est une solution particulière de l'équation. Considérons maintenant un couple d'entiers $(x,y)\in\mathbb Z^2$ tel que $2x+5y=3$. Alors, puisque $2\times(-1)+5\times 1=3$, on en déduit en soustrayant ces deux équations que 
			$$2(x+1)+5(y-1)=0\iff 2(x+1)=-5(y-1).$$
			On a donc $2|-5(y-1)$ et donc, puisque $2\wedge 5=1$, le théorème de Gauss assure que $2|y-1$. On en déduit l'existence de $k\in\mathbb Z$ tel que $y=1+2k$. Si on reporte ceci dans l'équation $2(x+1)=-5(y-1)$, on trouve que $x=-1-5k$. Réciproquement, tout couple d'entiers de la forme $(-1-5k,1+2k)$ avec $k\in\mathbb Z$ est solution de l'équation. Ainsi, on a prouvé que l'ensemble des solutions est 
			$$\mathcal S=\{(-1-5k,1+2k);\ k\in\mathbb Z\}.$$
			
			\item  On commence par rechercher le pgcd de 323 et 391 en appliquant par exemple l'algorithme d'Euclide.
			On a :
			\begin{eqnarray*}
				391&=&323\times 1+68\\
				323&=&68\times 4+51\\
				68&=&51\times 1+17\\
				51&=&17\times 3
			\end{eqnarray*}
			Ainsi, le pgcd de 391 et 323 est 17. On vérifie que 17 divise 612, car $612=17\times 36$. 
			D'autre part, $391=17\times 23$ et $323=17\times 19$. L'équation est donc équivalente à
			$$19x-23y=36.$$
			D'après le théorème de Bezout, puisque 19 et 23 sont premiers entre eux,
			on sait qu'il existe $(x_0,y_0)$ solution de $19x-23y=1$.
			Multipliant tout par 36, $(36x_0,36y_0)$ est solution de l'équation.\\
			
			Soit $(x,y)$ une solution de l'équation. Alors on a 
			$$19(x-36x_0)-23(y-36y_0)=0.$$
			Ainsi, $19$ divise $23(y-36y_0)$. Mais puisque $19$ et $23$ sont premiers entre eux,
			$19$ divise $(y-36y_0)$ et donc il existe un entier $k$ tel que
			$y=36y_0+19k$. Reportant cela dans l'équation, on trouve que
			$x=36x_0+23k$. Réciproquement, il est facile de vérifier que
			tout couple $(36x_0+23k,36y_0+19k)$ est solution. Il
			suffit donc de déterminer un couple $(x_0,y_0)$.\\
			
			On applique une nouvelle fois l'algorithme 
			d'Euclide :
			\begin{eqnarray*}
				23&=&19+4\\
				19&=&4\times 4+3\\
				4&=&3+1
			\end{eqnarray*}
			soit, en remontant les calculs :
			$$1=4-(19-4\times 4)=5\times 4-19=5\times(23-19)-19=5\times 23-6\times 19.$$
			On peut donc prendre $x_0=-6$ et $y_0=-5$. L'ensemble des solutions de l'équation est donc
			$$\{(-216+23k,-180+19k),\ k\in\mathbb Z\}.$$
			
			\item  On procède exactement comme ci-dessus, cette fois-ci, $162\wedge 207=9$ et on peut simplifier par $9$. On trouve que l'ensemble des solutions est 
			$$\{(27-23k,-21+18k);\ k\in\mathbb Z\}.$$
			
			\item  Le pgcd de $221$ et $247$ est 13. Or, $13$ ne divise pas 15. L'équation n'admet donc pas de solutions!
			
		\end{enumerate}
	\end{corrige}
	
	\begin{exo}[{Irrationnalité d'un quotient de logarithmes}]
		Soit $a,b\geq 2$ deux entiers premiers entre eux. Démontrer que $\ln(a)/\ln(b)$ est irrationnel.
	\end{exo}
	\begin{corrige}
		Supposons que $\ln(a)/\ln(b)$ est rationnel et soit $p,q\geq 1$ avec $p\wedge q=1$ tel que $\frac{\ln(a)}{\ln(b)}=\frac pq$. Ceci entraîne que $\ln(a^q)=\ln(b^p)$ et donc, par injectivité de la fonction $\ln$, que $a^q=b^p$. 
		Ainsi, $a|b^p$. Mais ceci contredit que $a\wedge b=1$.
	\end{corrige}
	
	\begin{exo}[{Rationnels qui sont entiers}]
		Soient $a$ et $b$ deux rationnels tels que $a+b$ et $a\times b$ sont des entiers. Prouver que $a$ et $b$ sont des entiers.
	\end{exo}
	\begin{corrige}
		Remarquons d'abord que si $a$ ou $b$ est nul, c'est trivial. On suppose donc que $a\neq 0$ et $b\neq 0$, et on écrit ces réels sous forme de fraction
		irréductible, $a=\frac{p_1}{q_1}$ et $b=\frac{p_2}{q_2}$, avec $q_1,q_2>0$ et $p_1\wedge q_1=1$, $p_2\wedge q_2=1$. On a alors 
		$$a+b=\frac{p_1q_2+p_2q_1}{q_1q_2}\in\mathbb Z.$$
		On en déduit que $q_1|(p_1q_2+p_2q_1)$. Or, on sait déjà que $q_1|p_2q_1$. On a donc $q_1|p_1q_2$. De plus, puisque $q_1\wedge p_1=1$, le théorème de Gauss
		nous dit que $q_1|q_2$. Par symétrie, on a aussi $q_2|q_1$, et donc $q_1=q_2=q$. Le produit $ab$ vaut alors
		$$ab=\frac{p_1p_2}{q^2},$$
		et donc $q|p_1p_2$. Ceci implique (toujours par le lemme de Gauss) que $q|p_1$ et comme $q\wedge p_1=1$, on a $q\wedge p_1=q=1$.
		Ainsi, $a$ et $b$ sont des entiers.
	\end{corrige}
	
	\begin{exo}[{Points à coordonnées entières sur une droite rationnelle}]
		Le plan est muni d'un repère $(O,\vec i,\vec j)$. On considère 4 entiers relatifs non nuls $m,n,p$ et $q$ tels que $\textrm{pgcd}(m,n)=\textrm{pgcd}(p,q)=1$. Le but de l'exercice est de déterminer une condition nécessaire et suffisante portant sur $m,n,p,q$ pour que la droite $\Delta$ d'équation $y=\frac mn x-\frac pq$ comporte au moins un point dont les coordonnées sont deux entiers relatifs.
		
		\begin{enumerate}
			
			\item  On suppose dans cette question que la droite $\Delta$ comporte un point de coordonnées $(x_0,y_0)$ où $x_0$, $y_0$ sont des entiers relatifs.
			
			\begin{enumerate}
				
				\item  Démontrer que $q$ divise le produit $np$.
				
				\item  En déduire que $q$ divise $n$.
				
			\end{enumerate}
			
			
			\item  Réciproquement, on suppose que $q$ divise $n$ et on souhaite trouver un couple $(x_0,y_0)$ d'entiers relatifs tels que $y_0=\frac mnx_0-\frac pq$. 
			
			\begin{enumerate}
				
				\item  On pose $n=qr$, où $r$ est un entier relatif non nul. Démontrer que l'on peut trouver deux entiers relatifs $u$ et $v$ tels que $qru-mv=1$.
				
				\item  En déduire qu'il existe un couple $(x_0,y_0)$ d'entiers relatifs tels que $y_0=\frac mn x_0-\frac pq$.
				
			\end{enumerate}
			
			
			\item  On donne l'algorithme suivant : 
			
			\begin{verbatim}
				Variables : 
				m,n,p,q entiers relatifs non nuls tels que pgcd(m,n)=pgcd(p,q)=1
				x entier naturel
				Algorithme
				Si q divise n alors
				x prend la valeur 0
				Tant que (mx/n - p/q) n'est pas un entier et (-mx/n-p/q) n'est pas un entier faire
				x prend la valeur x+1
				Fin Tant Que
				Si (mx/n-p/q) est entier alors Afficher x, mx/n-p/q
				Sinon Afficher -x,-mx/n-p/q
				Fin Si
				Sinon Afficher "Pas de solutions"
				Fin Si.
			\end{verbatim}
			
			Justifier que cet algorithme se termine. Que permet-il d'obtenir?
			
		\end{enumerate}
	\end{exo}
	\begin{corrige}
		\begin{enumerate}
			
			\item 
			
			\begin{enumerate}
				
				\item  Multiplions la relation $y_0=\frac mn x_0-\frac pq$ par $n$. On obtient 
				$$ny_0-mx_0=\frac{-pn}q.$$
				Mais $ny_0-mx_0$ est un entier relatif que nous notons $r$. Donc $q(-r)=pn$, ce qui dit bien que $q$ divise le produit $pn$.
				
				\item  Puisque $q$ et $p$ sont premier entre eux, ceci est une conséquence du théorème de Gauss.
				
			\end{enumerate}
			
			
			\item  
			
			\begin{enumerate}
				
				\item  Puisque $m$ et $n$ sont premiers entre eux, le théorème de Bézout nous dit qu'il existe deux entiers $u$ et $v$ tels que $nu-mv=1$. Remplaçant $n$ par $qr$, on obtient le résultat voulu.
				
				\item  Multiplions l'équation précédente par $\frac{-np}{q}=-pr$. On trouve
				$$\frac{-np}q=-pqr^2 u-(-prmv)=n(-pru)-m(-prv).$$ Posons alors $x_0=-prv$ et $y_0=-pru$. $x_0$ et $y_0$ sont bien deux entiers naturels. De plus, ils satisfont 
				$$\frac{-np}{q}=ny_0-mx_0$$
				équation qui est équivalente à 
				$$y_0=\frac mn x_0-\frac pq.$$
				Le point de coordonnées $(x_0,y_0)$ est donc bien un élément de $\Delta$.
				
			\end{enumerate}
			
			
			\item  L'algorithme permet de déterminer s'il y a un point à coordonnées entières sur la droite $\Delta$ et, le cas échéant, d'afficher les coordonnées d'un de ces points. Pour cela, il commence par étudier la condition nécessaire et suffisante d'appartenance d'un tel point en étudiant si $q$ divise $n$. Ensuite, si $q$ divise $n$, il regarde s'il existe un point à coordonnées entières sur $\Delta$ dont l'abscisse est 0, puis $\pm 1$, puis $\pm 2$. L'algorithme va terminer puisqu'on sait qu'un tel point existe et que toutes les valeurs successives possibles des abscisses sont prises. De plus, l'algorithme va produire un point à coordonnées entières $(x_0,y_0)$ sur $\Delta$ avec $|x_0|$ minimum.
			
		\end{enumerate}
	\end{corrige}
	
	\begin{exo}[{Calculs de pgcd}]
		Soient $n\in\mathbb Z$. Calculer les pgcd suivants :
		$$\begin{array}{lll}
			\mathbf 1.\ (n^2+n)\wedge (2n+1)&\quad\quad&\mathbf  2.\ (15n^2+8n+6)\wedge(30n^2+21n+13)\\
		\end{array}
		$$
	\end{exo}
	\begin{corrige}
		\begin{enumerate}
			
			\item  On remarque que $4(n^2+n)=(2n+1)^2-1$. Ainsi, si $d|(n^2+n)$ et $d|2n+1$, $d|1$. On a donc $(n^2+n)\wedge (2n+1)=1$.
			
			\item  Rappelons que si $a=bq+r$, on a $a\wedge b=b\wedge r$ et que $a\wedge b=a\wedge (a-b)$ (on écrit $a=b+(a-b)$).
			Ainsi, on trouve ici
			$$30n^2+21n+13=2(15n^2+8n+6)+5n+1,\ \ \ 15n^2+8n+6=3n(5n+1)+5n+6.$$
			On a donc
			\begin{eqnarray*}
				(30n^2+21n+13)\wedge (15n^2+8n+6)&=&(15n^2+8n+6)\wedge (5n+1)\\
				&=&(5n+1)\wedge (5n+6)\\
				&=&(5n+1)\wedge 5\\
				&=&1.
			\end{eqnarray*}
			
		\end{enumerate}
	\end{corrige}
	
	\begin{exo}[{Pile de livres}]
		J'ai une pile de livres. Lorsque je les range par 10, il m'en reste 3. Lorsque je les range par 17, il m'en reste 2. Quel est le nombre minimum de livres que je possède? Quels sont tous les nombres possibles?
	\end{exo}
	\begin{corrige}
		Soit $N$ le nombre de livres. L'information que l'on donne est que le reste dans la division euclidienne de $N$ par $10$ (respectivement par $17$) vaut $3$ (respectivement $2$). Ainsi, il existe $a$ et $b$ des entiers naturels tels que $N=10a+3$ et $N=17b+2$. Faisant la différence, on trouve que $(a,b)$ est une solution de l'équation de Bézout $17b-10a=1$. \\
		
		On résoud cette équation. Le couple $(a,b)=(5,3)$ est solution, et donc les couples solutions sont les couples $(a,b)$ avec $a=5+17k$ et $b=3+10k$, avec $k\in\mathbb N$ (il n'y a pas de $-$ car l'équation de Bézout comporte déjà un $-$). Ainsi, si $N$ est le nombre de livres, il existe $k\in\mathbb N$ (car le nombre de livres est positif!) tel que $N=53+170k$. \\
		
		Réciproquement, on vérifie que $53+170k= 10\times (17k+5)+3=17\times(10k+3)+2$. Le nombre minimal de livres possible est 53, les nombres possibles sont ceux de la forme $53+170k$ avec $k\in\mathbb N$. Bien sûr, le nombre minimal de livres possibles aurait pu être trouvé à l'aide d'un tableur.
	\end{corrige}
	
	\begin{exo}[{Pas de solutions rationnelles}]
		Montrer que l'équation $x^3-x^2+x+1=0$ n'a pas de solutions dans $\mathbb Q$.
	\end{exo}
	\begin{corrige}
		Admettons qu'il existe une solution rationnelle $x=p/q$, avec $p\wedge q=1$ et $q>0$. Alors, 
		remplaçant $x$ par $p/q$ dans l'équation et multipliant tout par $q^3$, on obtient :
		$$p^3-p^2q+pq^2+q^3=0.$$
		Puisque $q|p^2q+pq^2+q^3$, on obtient $q|p^3$ et donc $q=1$ puisque $p$ et $q$ sont
		premiers entre eux. En effectuant le même raisonnement avec $p$, on obtient $p=\pm 1$.
		Or $1$ et $-1$ ne sont pas solutions de l'équation. Il y a donc contradiction.
	\end{corrige}
	
	\begin{exo}[{Termes consécutifs d'une suite}]
		Soit $(u_n)$ la suite d'entiers définie par $u_0=14$ et $u_{n+1}=5u_n-6$. Démontrer que le pgcd de deux termes
		consécutifs de la suite est constant. Préciser sa valeur.
	\end{exo}
	\begin{corrige}
		De la relation $u_{n+1}=5u_n-6$, on tire que tout diviseur commun de $u_n$ et de $u_{n+1}$ divise aussi 6.
		Ainsi, le pgcd de $u_n$ et $u_{n+1}$ ne peut qu'être un élément de $\{1,2,3,6\}$. Par ailleurs, on va démontrer par
		récurrence sur $n$ que 2 divise $u_n$ et que $3$ ne divise pas $u_n$. C'est bien le cas pour $u_0=14$, et considérons $n\in\mathbb N$ tel que $2$ divise $u_n$ et $3$ ne divise pas $u_n$. Alors,  $2|5u_n$ et $2$ divise
		$6$, donc $2$ divise $u_{n+1}=5u_n-6$. De plus, si on avait $3|u_{n+1}$, alors $3|u_{n+1}+6$ et donc  $3|5u_{n}$. De plus, par le théorème de Gauss, puisque $3\wedge
		5=1$, on tire $3|u_{n}$, ce qui est une contradiction.\\
		
		Ainsi, pour tout entier $n$, on a prouvé que $2|u_n$ et $3$ ne divise pas $u_n$. Ainsi, si $d_n$ est le pgcd de $u_n$ et $u_{n+1}$, $2|d_n$ et $3$ ne divise pas $d_n$. Puisque l'on sait déjà que $d_n\in\{1,2,3,6\}$, on a nécessairement $d_n=2$.
	\end{corrige}
	
	\begin{exo}[{pgcd et somme ou produit prescrits}]
		\begin{enumerate}
			
			\item  Déterminer les couples d'entiers naturels de pgcd 18 et de somme 360.
			
			\item  Déterminer les couples d'entiers naturels de pgcd 18 et de produit 6480.
			
		\end{enumerate}
	\end{exo}
	\begin{corrige}
		\begin{enumerate}
			
			\item  Soit $(x,y)$ un tel couple d'entiers naturels. On écrit que $x=18x'$ et $y=18y'$ avec $x'\wedge y'=1$. Puisque $x+y=360$, on doit avoir $x'+y'=20$. Si $x'=1$, alors $y'=19$ et cette solution convient car $1$ et $19$ sont premiers entre eux. Si $x'=2$, alors $y'=18$ et cette solution ne convient pas car $2$ et $18$ ne sont pas premiers entre eux. On peut continuer ainsi à chercher de façon exhaustive toutes les solutions, et on trouve que les couples qui conviennent correspondent à $x'=1,3,7,9,11,13,17,19$. Les couples solutions sont donc :
			$$(18,342),\ (54,306)\, (126,234),\ (162,198)$$
			ainsi que les couples symétriques.
			
			\item  On procède exactement de la même façon, et cette fois on a l'équation $x'y'=20$. Les solutions qui vérifient en outre que $x'\wedge y'=1$ correspondent à $x'=1,4,5$ et $20$. Les couples solutions sont donc
			$$(18,360),\ (72,90),\ (90,72),\ (360,18).$$
			
		\end{enumerate}
	\end{corrige}
	
	\begin{exo}[{Équation avec pgcd et ppcm}]
		Trouver tous les couples d'entiers $(x,y)\in\mathbb N^2$ tels que $x\vee y+11(x\wedge y)=203$.
	\end{exo}
	\begin{corrige}
		On pose $d=x\wedge y$ et on écrit $x=dx'$, $y=dy'$ et $x'\wedge y'=1$. L'équation devient
		$$d(x'y'+11)=203=7\times 29.$$
		On a donc $d=1,7,29$ ou $203$ et on sépare les cas.
		
		\begin{itemize}
			
			\item  Les cas $d=29$ et $d=203$ ne conduisent pas à une solution, puisqu'on a alors $x'y'=7-11<0$ ou $x'y'=1-11<0$, ce
			qui est impossible puisque $x'$ et $y'$ sont positifs.
			
			\item  Pour $d=1$, on a $x'y'=192=2^6\times 3$. Puisque $x'\wedge y'=1$, on obtient que $(x',y')$ est un des couples
			$(192,1)$, $(3,2^6)$, $(2^6,3)$ ou $(1,192)$.
			
			\item  Pour $d=7$, on a $x'y'=18=2\times 3^2$. Puisque $x'\wedge y'=1$, on obtient que $(x',y')$ est un des couples
			$(18,1)$, $(2,9)$, $(9,2)$ ou $(1,18)$.
			
		\end{itemize}
		
		En remultipliant par $d$, on trouve que l'ensemble des solutions est constitué par les couples $(1,192)$, $(3,64)$,
		$(7,126)$, $(14,63)$ et leurs symétriques.
	\end{corrige}
	
	\begin{exo}[{Liens entre deux pgcd}]
		Soient $a,m,n\in\mathbb N^*$ avec $a\geq 2$ et $m\leq n$. On note $d=(a^n-1)\wedge (a^m-1)$ et $r$ le reste dans la division euclidienne de $n$ par $m$.
		
		\begin{enumerate}
			
			\item  Montrer que $a^n\equiv a^r\ [a^m-1]$.
			
			\item  En déduire que $d=(a^r-1)\wedge (a^m-1)$, puis que $d=a^{n\wedge m}-1$.
			
			\item  A quelle condition $a^m-1|a^n-1$?
			
		\end{enumerate}
	\end{exo}
	\begin{corrige}
		\begin{enumerate}
			
			\item  On écrit  $n=mq+r$ puis
			\begin{eqnarray*}
				a^n-a^r&=&a^r\left(a^{mq}-1\right)=a^r\left(\left(a^m\right)^q-1\right)\\
				&=&a^r\left(a^m-1\right)\left(a^{m(q-1)}+a^{m(q-2)}+\dots+1\right).
			\end{eqnarray*}
			Donc $a^m-1$ divise $a^n-a^r$, ce qui donne le résultat.
			
			\item  De la relation précédente, on tire
			$$\left(a^n-1\right)-\left(a^r-1\right)=Q\left(a^m-1\right)$$
			où $Q\in\mathbb N$. Ainsi, si $s$ divise $a^r-1$ et $a^m-1$, il divise $a^n-1$.
			Réciproquement si $s$ divise $a^n-1$ et $a^m-1$, alors il divise $a^r-1$.
			Ainsi, $d=(a^m-1)\wedge (a^r-1)$ (on utilise ici le même raisonnement
			que pour prouver que si $a=bq+c$, alors $a\wedge b=b\wedge c$). Mais,
			poursuivant les divisions euclidiennes successives, en notant $(r_k)$ la suite 
			des restes obtenus, on a $d=(a^{r_k}-1)\wedge (a^{r_{k+1}}-1)$, jusqu'à ce que
			$a^{r_{k+1}}=1$, c'est-à-dire $r_{k+1}=0$. On a donc $d=a^{r_{k_0}}-1$, où $r_{k_0}$ est le dernier reste non-nul 
			dans l'algorithme d'Euclide, c'est-à-dire que $r_{k_0}$ est  le pgcd de $n$ et $m$.
			
			\item  On a la suite d'équivalences :
			\begin{eqnarray*}
				a^m-1|a^n-1&\iff&d=a^m-1 \\
				&\iff&n\wedge m=m\\
				&\iff&m|n.
			\end{eqnarray*}
			
		\end{enumerate}
	\end{corrige}
	
	\subsection*{Nombres premiers et décomposition en produit de facteurs premiers}
	
	\begin{exo}[{Algorithme pour compter le nombre de nombres premiers inférieurs à un entier $n$}]
		\begin{enumerate}
			
			\item  Écrire une fonction $\textrm{divise}(p,q)$ d'argument deux entiers naturels non nuls $p$ et $q$ et renvoyant True si $p$ divise $q$, et False sinon.
			
			\item  Écrire une fonction $\textrm{estpremier}(p)$ d'argument un entier naturel $p$, renvoyant $1$ si $p$ est premier, et renvoyant $0$ sinon.
			
			\item  Écrire une fonction $\phi(n)$ d'argument un entier naturel $n$ et renvoyant le nombre de nombres premiers inférieurs ou égaux à $n$.
			
		\end{enumerate}
	\end{exo}
	\begin{corrige}
		\begin{enumerate}
			
			\item  On utilise le fait que $p$ divise $q$ si et seulement si le reste dans la division euclidienne de $q$ par $p$ est nul.
			
			\begin{verbatim}
				def divise(p,q):
				if ( (q%p) == 0):
				return True
				else:
				return False
			\end{verbatim}
			
			
			\item  Un entier naturel est premier si ses seuls diviseurs positifs sont $1$ et lui-même, et s'il n'est pas égal à $1$. On sépare donc le traitement de $1$, et pour les autres, on regarde si un des entiers entre $2$ et $p-1$ divise $p$.
			
			\begin{verbatim}
				def estpremier(p):
				if (p<2):
				return 0
				for i in range(2,p):
				if divise(i,p):
				return 0
				return 1
			\end{verbatim}
			
			
			\item 
			
			\begin{verbatim}
				def phi(n):
				total=0
				for m in range(2,n+1):
				total+=estpremier(m)
				return total
			\end{verbatim}
			
			
		\end{enumerate}
	\end{corrige}
	
	\begin{exo}[{Nombres de Fermat}]
		\begin{enumerate}
			
			\item  Soit $q$ un entier impair. Démontrer que, pour tout $x\in\mathbb R$,
			$$x^q+1=(x+1)(x^{q-1}-x^{q-2}+\dots+1).$$
			
			\item  Soit $m\in\mathbb N^*$ tel que $2^m+1$ soit premier. Montrer que $m=2^n$, où $n\in\mathbb N$.
			
		\end{enumerate}
	\end{exo}
	\begin{corrige}
		\begin{enumerate}
			
			\item  Il suffit de développer prudemment le membre de gauche et de constater que les termes se simplifient bien
			deux par deux. Remarquer qu'il faut que $q$ soit impair pour que le signe devant le dernier $1$ soit bien "$+$".
			
			\item  Soit $m\in\mathbb N^*$ tel que $2^m+1$ soit un nombre premier. Si $m$ n'est pas une puissance de $2,$ alors $m$ s'écrit $pq$ où $p=2^\alpha$ est un entier (avec $\alpha$ éventuellement égal à $0$),
			et $q$ est un nombre impair strictement supérieur  à $1$ (on obtient cette écriture en regroupant les nombres premiers impairs dans l'écriture en produit de facteurs premiers de $m$). On a alors :
			$$2^m+1=(2^{p})^q+1=(2^p+1)(2^{p(q-1)}-2^{p(q-2)}+\dots+1).$$
			Puisque $1<2^p+1<2^m+1$, on a obtenu une factorisation non triviale de $2^m+1$ qui n'est donc pas un nombre premier. C'est absurde, et donc $m$ est une puissance de $2.$
			
		\end{enumerate}
	\end{corrige}
	
	\begin{exo}[{Nombres de Mersenne}]
		Soient $a,n\geq 2$ des entiers. 
		
		\begin{enumerate}
			
			\item  Montrer que si $a^n-1$ est premier,
			alors $a=2$ et $n$ est premier.
			
			\item  On note $M_n=2^n-1$ le $n$-ième nombre de Mersenne. Vérifier que $M_{11}$ n'est pas premier.
			
		\end{enumerate}
	\end{exo}
	\begin{corrige}
		\begin{enumerate}
			
			\item  Montrons d'abord que $a=2$. Pour cela, on factorise $a^n-1$ sous la forme
			$$a^n-1=(a-1)(a^{n-1}+a^{n-2}+\dots+a+1).$$
			$a-1$ est un diviseur de $a^n-1$, distinct de $a^n-1$. Si on suppose que $a^n-1$ est premier,
			on a donc nécessairement $a-1=1$, c'est-à-dire $a=2$. De même, si $n=pq$, alors
			$$2^{pq}-1=(2^p)^q-1=(2^p-1)((2^p)^{q-1}+(2^p)^{q-2}+\dots+2^p+1).$$
			Donc $2^p-1$ divise $2^n-1$ qui est premier. On a donc ou bien $p=n$ (si $2^p-1=2^n-1$)
			ou bien $p=1$ (si $2^p-1=1$). Donc les seuls diviseurs de $n$ sont $1$ et $n$, i.e. $n$ est
			premier.
			
			\item  On a $M_{11}=2047=23\times 89$. On a donc $M_p$ premier implique $p$ premier,
			mais la réciproque est fausse.
			
		\end{enumerate}
	\end{corrige}
	
	\begin{exo}[{Une version faible du théorème de la progression arithmétique}]
		\begin{enumerate}
			
			\item  Démontrer que si un entier naturel est congru à $3$ modulo $4$, il admet un facteur premier congru à $3$ modulo $4$.
			
			\item  Démontrer qu'il existe une infinité de nombres premiers de la forme $4k+3$.
			
		\end{enumerate}
	\end{exo}
	\begin{corrige}
		\begin{enumerate}
			
			\item  Décomposons $n$ en produit de facteurs premiers : $n=p_1^{\alpha_1}\cdots p_r^{\alpha_r}$. Si $n$ n'admet aucun facteur premier congru à $3$ modulo $4$, ceci signifie que l'on a $p_i\equiv 2\ [4]$ ou $p_i\equiv 1\ [4]$ pour tout $i=1,\dots,r$ (le cas $p_i\equiv 0\ [4]$ est exclu puisque $p_i$ est premier).
			D'autre part, puisque $n$ est congru à $3$ modulo $4$, il est impair et aucun des $p_i$ n'est pair. On a donc pour tout $i=1,\dots,r$, $p_i\equiv 1\ [4]$ et donc $n=p_1^{\alpha_1}\cdots p_r^{\alpha_r}\equiv 1\ [4]$, une contradiction.
			
			\item  Supposons qu'il en existe simplement un nombre fini, notés $p_1,\dots,p_k$ et posons $n$ le
			produit de tous ces nombres :
			$$n=p_1\times\dots\times p_k=3\times 7\times 11\times 19\times \dots.$$
			Considérons alors $m=4n-1$. Alors $m$ est congru à 3 modulo 4.
			Par la question précédente, $m$ admet un diviseur premier congru à $3$ modulo $4$ et donc il existe $i\in\{1,\dots,k\}$ tel que $p_i|m=4n-1$. Écrivant $1=4n-m$, on trouve que $p_i|1$ puisque $p_i|m$ et $p_i|n$, ce qui est une contradiction.
			
		\end{enumerate}
	\end{corrige}
	
	\begin{exo}[{Intervalles sans nombres premiers}]
		Soit $q\geq 1$ un entier. Trouver un intervalle de longueur $q$ ne contenant pas de nombres premiers.
	\end{exo}
	\begin{corrige}
		Soit $n\geq 1$ et $I=\{n!+2,n!+3,\dots,n!+n\}$. Prenons un entier $m=n!+k$ dans 
		$I$. Alors $k$ est un diviseur strict de $m$, puisque $1<k<m$, et $k|k$, $k|n!$ puisque
		$k\leq n$, donc $k|m$. Ainsi, $I$ est un intervalle de $(n-1)$ entiers ne comportant
		pas de nombres premiers. Si on a choisi $n=q+1$, alors on a répondu au problème posé par l'énoncé.
	\end{corrige}
	
	\begin{exo}[{Divisibilité et puissance}]
		Soit $a,b\in\mathbb N^*$ tels que $a^2|b^2$. Démontrer que $a|b$.
	\end{exo}
	\begin{corrige}
		Décomposons $a$ et $b$ en produit de facteurs premiers :
		$$a=\prod_{i=1}^r p_i^{\alpha_i},\ b=\prod_{i=1}^r p_i^{\beta_i}.$$
		On a donc 
		$$a^2=\prod_{i=1}^r p_i^{2\alpha_i},\ b^2=\prod_{i=1}^r p_i^{2\beta_i}.$$
		La condition $a^2|b^2$ se traduit par, pour tout $i=1,\dots,r$, $2\alpha_i\leq 2\beta_i$. Mais ceci entraîne que $\alpha_i\leq\beta_i$ et donc que $a|b$.
	\end{corrige}
	
	\begin{exo}[{Puissances égales}]
		\begin{enumerate}
			
			\item  Montrer qu'un entier naturel qui est à la fois un carré et un cube est aussi un entier à la puissance $6$.
			
			\item  Généralisation : soient $a,b,p,q,n$ des entiers naturels avec $p\wedge q=1$ et $n=a^p=b^q$. Montrer qu'il existe 
			un entier naturel $c$ tel que $n=c^{pq}$.
			
		\end{enumerate}
	\end{exo}
	\begin{corrige}
		\begin{enumerate}
			
			\item  Détaillons d'abord ce que signifie la question. Soit $n$ un entier qui est à la fois un carré et un cube,
			i.e. $n=a^2$ et $n=b^3$.
			On veut prouver que c'est le carré d'un cube, c'est-à-dire qu'il existe un entier $c$ tel que $n=(c^3)^2=c^6$.
			Décomposons $n$, $a$ et $b$ en produit de facteurs premiers :
			\begin{eqnarray*}
				n&=&p_1^{n_1}\dots p_r^{n_r}\\
				a&=&p_1^{\alpha_1}\dots p_r^{\alpha_r}\\
				b&=&p_1^{\beta_1}\dots p_r^{\beta_r}\\
			\end{eqnarray*}
			Par unicité de la décomposition en produit de facteurs premiers, puisque $n=a^2$, on a $n_i=2\alpha_i$, et
			puisque $n=b^3$, on a $n_i=3\beta_i$ pour tout $i\in\{1,\dots,r\}$. Ainsi chaque $n_i$ est multiple de 2
			et de 3. C'est donc aussi un multiple de 6 et $n_i=6\gamma_i$. Posons alors 
			$$c=p_1^{\gamma_1}\dots p_r^{\gamma_r}.$$
			Alors $c^6=n$.
			
			\item  Décomposons $a$ et $b$ en produit de facteurs premiers, $a=p_1^{\alpha_1}\dots p_r^{\alpha_r}$ 
			et $b=p_1^{\beta_1}\dots p_r^{\beta_r}$. L'égalité $a^p=b^q$ se traduit par
			$p\alpha_j=q\beta_j$ pour tout $j=1,\dots,r$. On a donc $q|p\alpha_j$ et puisque $p\wedge q=1$,
			on obtient $q|\alpha_j$. Ecrivons $\alpha_j=q\gamma_j$ et reportons. On trouve
			$pq\gamma_j=q\beta_j$ soit $\beta_j=p\gamma_j$. Ainsi, si on pose $c=p_1^{\gamma_1}\dots p_r^{\gamma_r}$
			on a bien $n=c^{pq}$.
			
		\end{enumerate}
	\end{corrige}
	
	\begin{exo}[{Le produit est un carré parfait}]
		Soient $a$ et $b$ deux entiers premiers entre eux tels que leur produit $ab$ est un carré parfait.
		Montrer que $a$ et $b$ sont deux carrés parfaits.
	\end{exo}
	\begin{corrige}
		Décomposons $a,b$ et $ab$ en produit de facteurs premiers, $a=p_1^{\alpha_1}\dots p_r^{\alpha_r}$ 
		et $b=p_1^{\beta_1}\dots p_r^{\beta_r}$, $ab=p_1^{2\gamma_1}\dots p_r^{2\gamma_r}$ (on a utilisé
		que $ab$ est un carré parfait). Bien sûr, pour tout $j$, on a la relation $2\gamma_j=\alpha_j+\beta_j$. De plus, puisque
		$a\wedge b=1$, si $\alpha_j$ est non-nul, alors $\beta_j=0$, et réciproquement. On a alors ou bien $\alpha_j=2\gamma_j$ et $\beta_j=0$, ou
		bien $\beta_j=2\gamma_j$ et $\alpha_j=0$. Dans tous les cas, $\alpha_j$ et $\beta_j$ sont tous les deux pairs. Ainsi,
		$a$ et $b$ sont des carrés parfaits.
	\end{corrige}
	
	\begin{exo}[{Application au calcul de pgcd}]
		Soient $a,b,c\in\mathbb N^*$ et soit $n\in\mathbb N^*$.
		
		\begin{enumerate}
			
			\item  Démontrer que $c|ab\implies c|(a\wedge c)(b\wedge c)$.
			
			\item  Démontrer que $(a\wedge b)^n=a^n \wedge b^n$.
			
			\item  (Plus difficile) Calculer $(a^2+ab+b^2)\wedge ab$.
			
		\end{enumerate}
	\end{exo}
	\begin{corrige}
		\begin{enumerate}
			
			\item  On décompose $a,$ $b$ et $c$ en produit de facteurs premiers : 
			$$a=p_1^{\alpha_1}\dots p_r^{\alpha_r},\ b=p_1^{\beta_1}\dots p_r^{\beta_r},\ c=p_1^{\gamma_1}\dots p_r^{\gamma_r}$$
			(on autorise que les $\alpha_i,$ les $\beta_i$ et les $\gamma_i$ sont nuls car $a,$ $b$ et $c$ n'ont pas forcément les mêmes facteurs premiers).
			La condition $c|ab$ signifie que $\gamma_j\leq \alpha_j+\beta_j$ pour tout $j$. De plus, on a
			$$a\wedge c=p_1^{\min(\alpha_1,\gamma_1)}\dots p_r^{\min(\alpha_r,\gamma_r)}$$
			$$b\wedge c=p_1^{\min(\beta_1,\gamma_1)}\dots p_r^{\min(\beta_r,\gamma_r)}$$
			et donc on a $c|(a\wedge c)(b\wedge c)$ si et seulement si, pour tout $j$ de $1,\dots,r$, on a
			$$\gamma_j\leq \min(\alpha_j,\gamma_j)+\min(\beta_j,\gamma_j).$$
			Mais si $\min(\alpha_j,\gamma_j)=\gamma_j$ ou $\min(\beta_j,\gamma_j)=\gamma_j$, c'est trivial.
			Sinon on a
			$$\min(\alpha_j,\gamma_j)+\min(\beta_j,\gamma_j)=\alpha_j+\beta_j\geq\gamma_j$$
			par hypothèse. Ceci prouve le résultat.
			
			\item  Écrivons à nouveau
			$$a=p_1^{\alpha_1}\dots p_r^{\alpha_r},\ b=p_1^{\beta_1}\dots p_r^{\beta_r}$$
			et posons $\gamma_j=\min(\alpha_j,\beta_j)$. Alors 
			$$a\wedge b=p_1^{\gamma_1}\dots p_r^{\gamma_r}\implies (a\wedge b)^n=p_1^{n\gamma_1}\dots p_r^{n\gamma_r}.$$
			De même, 
			$$a^n=p_1^{n\alpha_1}\dots p_r^{n\alpha_r},\ b^n=p_1^{n\beta_1}\dots p_r^{n\beta_r}$$
			de sorte que 
			$$a^n\wedge b^n=p_1^{\delta_1}\cdots p_r^{\delta_r}$$
			avec $\delta_j=\min(n\alpha_j,n\beta_j)=n\min(\alpha_j,\beta_j)=n\gamma_j.$
			
			\item  Soit $d=(a^2+ab+b^2)\wedge ab$. On peut commencer par remarquer que 
			$d=(a^2+b^2)\wedge ab$. Décomposons $a$, $b$ et $d$ en produits de facteurs premiers :
			$$a=p_1^{\alpha_1}\dots p_r^{\alpha_r}, b=p_1^{\beta_1}\dots p_r^{\beta_r},\ d=p_1^{\gamma_1}\dots p_r^{\gamma_r}.$$
			Puisque $d|ab$, on a $\gamma_j\leq \alpha_j+\beta_j$ pour tout $j$. 
			Si $\alpha_j=\beta_j$, on en déduit que $\gamma_j\leq 2\min(\alpha_j,\beta_j)$. 
			Si maintenant $\alpha_j\neq \beta_j$, on peut supposer $\alpha_j<\beta_j$, alors 
			$$a^2+b^2=p_j^{2\alpha_j}(q+p_j^{2\beta_j-2\alpha_j}r)$$
			avec $q\wedge p_j=1$. En particulier, $(q+p_j^{2\beta_j-2\alpha_j}r)\wedge p_j=1$, et donc 
			$\gamma_j\leq 2\alpha_j=2\min(\alpha_j,\beta_j)$. Ceci démontre en particulier que $d|a^2\wedge b^2$.
			Réciproquement, on va prouver que $a^2\wedge b^2|ab$. Puisqu'il est clair qu'on a aussi $a^2\wedge b^2|a^2+b^2$,
			on pourra conclure que $d=a^2\wedge b^2.$ Mais 
			$$ab=p_1^{\alpha_1+\beta_1}\dots p_r^{\alpha_r+\beta_r}\textrm{ et }a^2\wedge b^2=p_1^{\min(2\alpha_1,2\beta_1)}\dots p_r^{\min(2\alpha_r,2\beta_r)}$$
			ce qui prouve ce que l'on voulait.
			
		\end{enumerate}
	\end{corrige}
	
	\begin{exo}[{Nombre de diviseurs d'un entier}]
		Soit $n$ un nombre entier, $n=p_1^{\alpha_1}\dots p_r^{\alpha_r}$
		sa décomposition en produit de facteurs premiers. On note $d(n)$ le nombre
		de diviseurs de $n$.
		
		\begin{enumerate}
			
			\item  Montrer que  $d(n)=\prod_{i=1}^r (\alpha_i+1)$.
			
			\item  Montrer que $n$ est un carré parfait si et seulement si $d(n)$ est impair.
			
			\item  Montrer que $\prod_{d|n}d=\sqrt{n}^{d(n)}$.
			
		\end{enumerate}
	\end{exo}
	\begin{corrige}
		\begin{enumerate}
			
			\item  Un diviseur de $n$ s'écrit de façon unique sous la forme $d=p_1^{\beta_1}\dots p_r^{\beta_r}$ avec
			$0\leq \beta_i\leq \alpha_i$ pour $i=1,\dots,r$. Notant $I_i=\{0,\dots,\alpha_i\}$, on a donc une bijection
			de $I_1\times\dots\times I_r$ sur $\{\textrm{diviseurs de }$n$\}$ donné par
			$$(\beta_1,\dots,\beta_r)\mapsto p_1^{\beta_1}\dots p_r^{\beta_r}.$$
			Les deux ensembles ont même cardinal, et donc $d(n)=\prod_{i=1}^r(\alpha_i+1)$.
			
			\item  Rappelons qu'un produit de nombre entiers est impair si et seulement si tous ses facteurs sont impairs.
			On en déduit
			\begin{eqnarray*}
				d(n)\textrm{ est impair}&\iff&\forall i\in\{1,\dots,r\}, \alpha_j+1\textrm{ est impair}\\
				&\iff&\forall i\in\{1,\dots,r\},\ \alpha_j\textrm{ est pair}\\
				&\iff&n\textrm{ est un carré parfait.}
			\end{eqnarray*}
			
			\item  L'idée est de regrouper ensemble les facteurs $d$ et $n/d$. Si le premier est inférieur à $\sqrt{n}$, le second
			est supérieur à $\sqrt{n}$. Précisément, si $n$ n'est pas un carré parfait :
			$$\prod_{d|n}d=\prod_{d|n,d<\sqrt{n}}d\times \frac{n}{d}=\prod_{d|n,d<\sqrt{n}} n=n^{d(n)/2}.$$
			Si $n$ est un carré parfait, $n=m^2$, alors
			$$\prod_{d|n}d=m\prod_{d|n,d<\sqrt{n}}d\times \frac{n}{d}=n^{1/2}n^{(d(n)-1)/2}=n^{d(n)/2}.$$
			
		\end{enumerate}
	\end{corrige}
	
	\begin{exo}[{Formule de Legendre}]
		Dans cet exercice, si $m$ est un entier et $q$ un nombre premier, $\nu_q(m)$ désigne le plus grand entier naturel $k$ tel que $q^k|m.$ Soit $n\geq 1$ un entier et $p\geq 2$ un nombre premier. Soit aussi $N$ tel que $p^N\geq n$.
		
		\begin{enumerate}
			
			\item  Soit $u\geq 1$ un entier. Combien-y-a-t-il de multiples de $u$ dans $\{1,\dots,n\}$?
			
			\item  Soit $k\geq 1$ et $A_k=\big \{l\in\{1,\dots,n\}:\ \nu_p(l)=k\big\}$.
			Calculer le cardinal de $A_k$.
			
			\item  En déduire la formule de Legendre : 
			$$\nu_p(n!)=\sum_{k=1}^N \left\lfloor \frac n{p^k}\right\rfloor.$$
			
			\item  Par combien de zéros se termine le nombre $1000!\ $?
			
		\end{enumerate}
	\end{exo}
	\begin{corrige}
		\begin{enumerate}
			
			\item  Un multiple de $u$ s'écrit $ku$. Soit $mu$ le plus grand multiple dans $\{1,\dots,n\}$. Alors $um\leq n<u(m+1)$ et donc
			$$m\leq \frac nu <m+1.$$
			Puisque $m$ est un entier, on reconnait $m=\left\lfloor \frac nu\right\rfloor$. Les multiples de $u$ dans $\{1,\dots,n\}$ étant les entiers $ku$, avec $1\leq k\leq m$, on déduit que leur nombre est égal à $\lfloor \frac nu\rfloor.$
			
			\item  Soit $l\in\{1,\dots,n\}$. Alors $\nu_p(l)=k$ si et seulement si $p^k$ divise $l$ mais $p^{k+1}$ ne divise pas $l$. D'après la question précédente, on déduit que 
			$$\card(A_k)=\left\lfloor \frac n{p^k}\right\rfloor-\left\lfloor \frac n{p^{k+1}}\right\rfloor.$$
			
			\item  On commence par écrire 
			\begin{align*}
				\nu_p(n!)&=\sum_{l=1}^n \nu_p(l)\\
				&=\sum_{k=1}^N k \card\left(\{l\in\{1,\dots,n\}:\ \nu_p(l)=k\}\right)
			\end{align*}
			où on a regroupé les éléments de $\{1,\dots,n\}$ suivant leur valuation $p$-adique. En utilisant le résultat de la question précédente,
			\begin{align*}
				\nu_p(n!)&=\sum_{k=1}^N k \left(\left\lfloor \frac n{p^k}\right\rfloor-\left\lfloor \frac n{p^{k+1}}\right\rfloor\right)\\
				&=\sum_{k=1}^N  k \left\lfloor \frac n{p^k}\right\rfloor-\sum_{k=1}^N k \left\lfloor \frac n{p^{k+1}}\right\rfloor\\
				&=\sum_{k=1}^N  k \left\lfloor \frac n{p^k}\right\rfloor-\sum_{k=2}^{N+1} (k-1) \left\lfloor \frac n{p^{k}}\right\rfloor\\
				&=\sum_{k=1}^N \left\lfloor \frac n{p^k}\right\rfloor-N \left\lfloor\frac{n}{p^{N+1}}\right\rfloor\\
				&=\sum_{k=1}^N \left\lfloor \frac n{p^k}\right\rfloor
			\end{align*}
			où la dernière ligne vient du fait que $p^{N+1}>n$.
			
			\item  Il suffit de calculer la plus grande puissance $N$ telle que $10^N|1000!$.
			Comme $10=2\times 5$, on obtient $N=\min(v_2(1000!),v_5(1000!))$. En utilisant la formule de Legendre, on trouve
			\begin{align*}
				v_5(1000!)&=\Big\lfloor\frac{1000}5\Big\rfloor+\Big\lfloor \frac{1000}{25}\Big\rfloor+\Big\lfloor \frac{1000}{125}\Big\rfloor+\Big\lfloor \frac{1000}{625}\Big\rfloor\\
				&=249\\
				v_2(1000!)&\geq \Big\lfloor\frac{1000}{2}\Big\rfloor=500.
			\end{align*}
			Ainsi, $1000!$ se termine par $249$ zéros.
			
		\end{enumerate}
	\end{corrige}
	
	\begin{exo}[{Produit des diviseurs d'un entier}]
		Déterminer tous les entiers naturels dont le produit des diviseurs (positifs) est égal à $45^{42}$.
	\end{exo}
	\begin{corrige}
		Remarquons d'abord que $45^{42}$ se décompose en facteurs premiers sous la forme $3^{84}5^{42}$. 
		Ainsi, un entier $n$ dont le produit des diviseurs est égal à $45^{42}$ est forcément de la forme $3^p5^q$. 
		Maintenant, les diviseurs d'un tel entier sont tous les nombres de la forme $3^k 5^l$ avec $0\leq k\leq p$ 
		et $0\leq l\leq q$. Ainsi, le produit de ces diviseurs est égal à
		\begin{eqnarray*}
			\prod_{k=0}^p \prod_{l=0}^q 3^k 5^l&=&\prod_{k=0}^p \left(3^{k(q+1)}\prod_{l=0}^q 5^l\right)\\
			&=&\prod_{k=0}^p \left(3^{k(q+1)}5^{\frac{q(q+1)}2}\right)\\
			&=&3^{\frac{p(p+1)(q+1)}2}5^{\frac{q(q+1)(p+1)}2}.
		\end{eqnarray*}
		$p$ et $q$ sont donc solutions du système
		$$\left\{
		\begin{array}{rcl}
			p(p+1)(q+1)&=&168\\
			q(q+1)(p+1)&=&84.
		\end{array}
		\right.$$
		Si on fait le quotient de ces deux équations, on obtient $p=2q$, ce qui donne encore
		$$q(q+1)(2q+1)=84.$$
		La seule solution est $q=3$ et donc $p=6$. Ainsi, il y a un unique entier naturel dont le produit des diviseurs est $45^{42}$. Il s'agit de
		$3^65^3$.
	\end{corrige}
	
	\subsection*{Congruences}
	
	\begin{exo}[{Puissances itérées}]
		\begin{enumerate}
			
			\item  Déterminer, suivant les valeurs de $n\in\mathbb N$, le reste de la division euclidienne de $2^n$ par $5$.
			
			\item  Quel est le reste de la division par 5 de $1357^{2013}$?
			
		\end{enumerate}
	\end{exo}
	\begin{corrige}
		\begin{enumerate}
			
			\item  La méthode pour ce type d'exercice est toujours la même et est très importante à savoir. On commence par rechercher le premier entier $k\geq 1$ tel que $2^k\equiv 1\ [5]$. On va ensuite raisonner modulo $k$. On trouve successivement :
			$$2^1\equiv 2\ [5],\quad 2^2\equiv 4\ [5],\quad 2^3\equiv 3\ [5],\quad 2^4\equiv 1\ [5].$$
			On va donc classer les entiers $n$ modulo $4$. En effet, si $n=4q+r$, alors sachant que 
			$2^{4q}\equiv 1^q\ [5]$ soit $2^{4q}\equiv 1\ [5]$, on trouve que 
			$$2^n\equiv 2^r\ [5].$$
			Ainsi, on obtient 
			
			\begin{itemize}
				
				\item  Si $n\equiv 0\ [4]$, alors $2^n\equiv 1\ [5]$;
				
				\item  Si $n\equiv 1\ [4]$, alors $2^n\equiv 2\ [5]$;
				
				\item  Si $n\equiv 2\ [4]$, alors $2^n\equiv 4\ [5]$;
				
				\item  Si $n\equiv 3\ [4]$, alors $2^n\equiv 3\ [5]$;
				
			\end{itemize}
			
			
			\item  On commence par effectuer la division euclidienne de $1357$ par $5$, et on trouve que $1357\equiv 2\ [5]$, d'où $1357^{2013}\equiv 2^{2013}\ [5]$. De plus, $2013\equiv 1\ [4]$. On en déduit que $1357^{2013}\equiv 2^1\equiv 2\ [5]$.
			
		\end{enumerate}
	\end{corrige}
	
	\begin{exo}[{Un entier divisible par $6$}]
		Démontrer que pour tout $n\in\mathbb Z$, $n(n+2)(7n-5)$ est divisible par $6$.
	\end{exo}
	\begin{corrige}
		Puisque $6=2\times 3$ avec $2\wedge 3=1$, il suffit de prouver que $n(n+2)(7n-5)$ est divisible par $2$ et est divisible par $3$. Commençons par la divisibilité par $2$. On va raisonner par disjonction de cas. Si $n\equiv 0\ [2]$, alors clairement $n(n+2)(7n+5)\equiv 0\ [2]$. Si $n\equiv 1\ [2]$, alors $7n-5\equiv 2\ [2]$ et donc $7n-5\equiv 0\ [2]$. A nouveau, on a $n(n+2)(7n-5)\equiv 0\ [2]$.\\
		
		Étudions maintenant la divisibilité par $3$.  On raisonne encore par disjonction de cas, le cas $n\equiv 0\ [3]$ se traite comme précédemment. Si $n\equiv 1\ [3]$, on conclut en utilisant que $n+2\equiv 0\ [3]$ et si $n\equiv 2\ [3]$, on conclut en remarquant que $7n-5\equiv 9\ [3]$ et donc $7n-5\equiv 0\ [3]$.
	\end{corrige}
	
	\begin{exo}[{Divisible par $30$.}]
		Soit $n\in\mathbb N$ et $a=n^5-n$. 
		
		\begin{enumerate}
			
			\item  Démontrer que $a$ est divisible par $5$.
			
			\item  En remarquant que $a=n(n-1)(n+1)(n^2+1)$, démontrer que $a$ est divisible par $2$ et par $3$.
			
			\item  Démontrer que $a$ est divisible par $30$.
			
		\end{enumerate}
	\end{exo}
	\begin{corrige}
		\begin{enumerate}
			
			\item  Le plus facile est de remarquer que, d'après le petit théorème de Fermat, $n^5\equiv n\ [5]$, ce qui est exactement le résultat demandé.
			
			\item  La vérification est immédiate. On remarque alors que, parmi $n$ et $n-1$, l'un des deux entiers est pair, et donc $2|a$. De la même façon, parmi les 3 entiers consécutifs $n-1$, $n$ et $n+1$, l'un au moins est divisible par $3$ et donc $3|a$.
			
			\item  Puisque $2\wedge 3=1$, puisque $2|a$ et $3|a$, on sait que $6|a$. Puisque $5\wedge 6=1$ et que $5|a$, $6|a$, on en déduit que $30|a$.
			
		\end{enumerate}
	\end{corrige}
	
	\begin{exo}[{Somme de trois cubes consécutifs}]
		Démontrer que la somme de trois cubes consécutifs est toujours divisible par 9.
	\end{exo}
	\begin{corrige}
		On développe $n^3+(n+1)^3+(n+2)^3$ :
		$$n^3+(n+1)^3+(n+2)^3=3n^3+9n^2+15n+9=9(n^2+1)+3n(n^2+5).$$
		Il suffit donc de démontrer que $9|3n(n^2+5)$ ou encore que $3|n(n^2+5)$.
		Si $n\equiv 0\ [3]$, c'est bien sûr vrai. Si $n\equiv 1\ [3]$, on a
		$n^2\equiv 1\ [3]$ et donc $n^2+5\equiv 6\equiv 0\ [3]$. Si $n\equiv 2\ [3]$,
		alors $n^2+5\equiv 4+5\equiv 9\equiv 0\ [3]$. Dans tous les cas,
		$3n(n^2+5)$ est divisible par 9.
	\end{corrige}
	
	\begin{exo}[{Deux équations avec des congruences}]
		\begin{enumerate}
			
			\item  Déterminer les entiers naturels $n$ tels que $5^n\equiv -1\ [13]$.
			
			\item  Déterminer les entiers naturels $n$ tels que 13 divise $5^{2n}+5^n$.
			
		\end{enumerate}
	\end{exo}
	\begin{corrige}
		\begin{enumerate}
			
			\item  On calcule les premiers termes et on trouve $5^0\equiv 1\ [13]$, $5^1\equiv 5\ [13]$, $5^2\equiv -1\ [13]$, $5^3\equiv -5\ [13]$, $5^4\equiv 1\ [13]$, $5^5\equiv 5\ [13]$, $5^6\equiv -1\ [13]$,... On voit clairement apparaître le cycle $1,5,-1,-5,1,5,-1$, ce qui nous incite à démontrer par récurrence sur $\mathcal p\in\mathbb N$ la propriété $\mathcal P(p)$ suivante :
			$$\mathcal P(p)="5^{4p}\equiv 1\ [13],\ 5^{4p+1}\equiv 5\ [13],\ 5^{4p+2}\equiv -1\ [13],\ 5^{4p+3}\equiv -5\ [13]".$$
			La propriété $\mathcal P(0)$ est vraie comme le montre le calcul précédent. Soit $p\in\mathbb N$ tel que $\mathcal P(p)$ est vraie, et prouvons $\mathcal P(p+1)$. Alors on a 
			$$5^{4(p+1)}=5^{4p}5^4\equiv 1\times 1\equiv 1\ [13].$$
			La démonstration pour les trois autres cas est exactement similaire. Donc $\mathcal P(p+1)$ est vraie.\\
			
			Ainsi, les entiers naturels solutions de $5^n\equiv -1\ [13]$ sont exactement les entiers de la forme $4p+2$, avec $p\in\mathbb N$.
			
			\item  Écrivons $5^{2n}+5^n=5^n(5^n+1)$. Si $13|5^{2n}+5^n$, puisque $5\wedge 13=1$, le théorème de Gauss assure que $13|5^{n}+1$, autrement dit que $5^n\equiv -1\ [13]$. D'après la question précédente, ceci est équivalent à dire que $n=4p+2$, avec $p\in\mathbb N$. Réciproquement, si $n=4p+2$ pour un certain $p\in\mathbb N$, on sait que $13|5^n+1$ et donc $13|5^n(5^n+1)=5^{2n}+5^n$. Les entiers $n$ solutions sont donc exactement ceux qui s'écrivent $4p+2$ avec $p\in\mathbb N$.
			
		\end{enumerate}
	\end{corrige}
	
	\begin{exo}[{Puissance et congruence}]
		Soit $a,b\in\mathbb Z$ et $n\geq 2$. On suppose que $a\equiv b\ [n]$. Démontrer que $a^n\equiv b^n\ [n^2]$.
	\end{exo}
	\begin{corrige}
		Puisque $a\equiv b\ [n]$, il existe $k\in\mathbb Z$ tel que $a=b+kn$. Mettons ceci à la puissance $n$. On en déduit que 
		\begin{eqnarray*}
			a^n&=&(b+kn)^n\\
			&=&\sum_{j=0}^n \binom n j b^{n-j}k^j n^j\\
			&=& b^n+\binom n1 b^{n-1}kn+\sum_{j=2}^n \binom n j b^{n-j}k^j n^j\\
			&=&b^n +b^{n-1}k n^2+n^2\sum_{j=2}^n \binom nj b^{n-j}k^j n^{j-2}.
		\end{eqnarray*}
		Puisque pour tout $j\geq 2$, $\binom nj b^{n-j}k^jn^{j-2}$ est un entier, on en déduit que $a^n=b^n+\ell n^2$ avec $\ell\in\mathbb Z$ et donc que $a^n\equiv b^n\ [n^2]$.
	\end{corrige}
	
	\begin{exo}[{La preuve par neuf}]
		Montrer que tout entier naturel est congru modulo 9 à la somme des chiffres de son écriture décimale.
		En déduire que, quels que soient les entiers naturels $x=\overline{a_n\dots a_0}$, $y=\overline{b_m\dots b_0}$
		et $z=\overline{c_p\dots c_0}$, si $xy=z$, alors $\left(\sum_{i=0}^n a_i\right)\left(\sum_{i=0}^m b_i\right)\equiv\left(\sum_{i=0}^p c_i\right)\ [9]$.
	\end{exo}
	\begin{corrige}
		Le point crucial est que $10$ est congru à 1 modulo 9. Ainsi, $10^2=10\times 10\equiv 1\times 1=1[9]$, et une récurrence facile montre
		que pour tout entier $k$, alors $10^k$ est congru à 1 modulo 9. Notons $x=\overline{a_n\dots a_0}$. Alors on a :
		\begin{eqnarray*}
			x&\equiv&10^n a_n+10^{n-1}a_{n-1}+\dots+10 a_1+a_0\ [9]\\
			&\equiv& a_n +a_{n-1}+\dots +a_0\ [9].
		\end{eqnarray*}
		La deuxième partie de l'exercice s'en déduit aisément. En effet, on écrit que
		$$xy=z\implies xy\equiv z\ [9]\implies \left(\sum_{i=0}^n a_i\right)\left(\sum_{i=0}^m b_i\right)\equiv\left(\sum_{i=0}^p c_i\right)\ [9].$$
	\end{corrige}
	
	\begin{exo}[{Résolution d'équations diophantiennes en raisonnant modulo un entier}]
		Résoudre les équations suivantes dans $\mathbb Z^2$ :
		
		\begin{enumerate}
			
			\item  $x^2-5y^2=3$ (on raisonnera modulo $5$);
			
			\item  $15x^2-8y^2=9$ (on raisonnera modulo $5$).
			
		\end{enumerate}
	\end{exo}
	\begin{corrige}
		\begin{enumerate}
			
			\item  On écrit l'équations modulo 5, et on trouve que $x^2\equiv 3\ [5]$. Or, si $x\equiv\pm 1\ [5]$,
			alors $x^2\equiv 1\ [5]$ et si $x\equiv \pm 2\ [5]$, alors $x^2\equiv 4\ [5]$ et donc l'équation
			$x^2=3\ [5]$ n'a pas de solutions. Ainsi, l'équation de départ n'a pas de solutions dans $\mathbb Z^2$.
			
			\item  On écrit l'équation modulo $5$, et on trouve $-8y^2\equiv 4\ [5]$
			et donc $2y^2\equiv 4\ [5]$ puisque $-8\equiv 2\ [5]$. Mais :
			
			\begin{itemize}
				
				\item  si $y\equiv 0\ [5]$, alors $2y^2\equiv 0\ [5]$;
				
				\item  si $y\equiv \pm 1\ [5]$, alors $2y^2\equiv 2\ [5]$;
				
				\item  si $y\equiv \pm 2\ [5]$, alors $2y^2\equiv 8\ [5]\equiv 3\ [5]$.
				
			\end{itemize}
			
			Dans tous les cas, il est impossible que $2y^2\equiv 4\ [5]$. Ainsi, l'équation n'admet pas de solutions.
			
		\end{enumerate}
	\end{corrige}
	
	\begin{exo}[{Divisible par 7}]
		Démontrer que, pour tout entier naturel $n$, $3^{2n+1} + 2^{4n+2}$ est divisible par 7.
	\end{exo}
	\begin{corrige}
		Notons, pour $n\in\mathbb N$, la propriété $\mathcal P(n)=$"$3^{2n+1} + 2^{4n+2}$ est divisible par $7$".
		Prouvons par récurrence que pour tout $n\in\mathbb N$, $\mathcal P(n)$ est vérifiée.\\
		
		Initialisation : On a $3^1+2^2=7$ qui est bien divisible par 7.\\
		
		Hérédité : Soit $n\in\mathbb N$ tel que $\mathcal P(n)$ est vraie. Remarquons que $3^2\equiv 2\ [7]$ et que $2^4\equiv 2\ [7]$. Alors, on écrit (modulo 7) :
		\begin{eqnarray*}
			3^{2n+3}+2^{4n+6}&\equiv&3^2 3^{2n+1}+2^4 2^{4n+2}\ [7]\\
			&\equiv& 2 \times 3^{2n+1}+2\times 2^{4n+2}\ [7]\\
			&\equiv& 2\times (3^{2n+1}+2^{4n+2})\ [7]\\
			&\equiv& 2\times 0\ [7]\\
			&\equiv& 0\ [7].
		\end{eqnarray*}
		Ainsi, $7$ divise $3^{2n+3}+2^{4n+6}$ et donc $\mathcal P(n+1)$ est vraie.\\
		
		En conclusion, par le principe de récurrence, $\mathcal P(n)$ est vraie pour tout $n\in\mathbb N$.
	\end{corrige}
	
	\begin{exo}[{Théorème de Pascal}]
		Soit $m$ un entier naturel non-nul, et soit $(r_i)$ la suite d'entiers définie
		par $r_0=1$ et $r_{i+1}$ est le reste de la division euclidienne de $10r_i$ par $m$.
		
		\begin{enumerate}
			
			\item  Démontrer que, pour tout entier naturel $a=\overline{a_n\dots a_0}$ en écriture décimale, on a
			$a\equiv\sum_{i=0}^n a_ir_i\ [m]$. 
			
			\item  En déduire des critères simples permettant de reconnaitre sur l'écriture décimale d'un réel s'il est ou non
			divisible par 3, par 9, par 10, par 11.
			
		\end{enumerate}
	\end{exo}
	\begin{corrige}
		\begin{enumerate}
			
			\item  Puisque $a$ s'écrit $a=\sum_{i=0}^n a_i 10^i$, il suffit de démontrer que pour 
			tout $i\geq 0$, $10^i\equiv r_i\ [m]$. On va démontrer ce résultat par récurrence sur $i$.
			La propriété est clairement vérifiée pour $i=0$. Supposons qu'elle est vraie au rang $i$,
			à savoir $r_i\equiv 10^i\ [m]$, et prouvons-la au rang $i+1$. On a 
			$10^{i+1}\equiv 10r_i\ [m]$ et par définition de $r_{i+1}$, on a aussi $r_{i+1}\equiv 10r_i\ [m]$.
			Ainsi, on a bien $10^{i+1}\equiv r_{i+1}\ [m]$, ce qui achève de prouver que, pour tout $i\geq 0$, on a
			$10^i\equiv r_i\ [m]$.
			
			\item  Pour chaque valeur proposée (3,6,9,10), on va étudier le comportement de la suite $r_i$ afin d'en déduire
			le critère de divisibilité proposé.
			
			\begin{itemize}
				
				\item  Pour $3$, on va prouver que $r_i$ est toujours égal à $1$. C'est vrai pour $i=0$, et si c'est vrai au rang $i$, alors
				$$10^{r_i}=10=3\times 3+1,$$
				et donc $r_{i+1}=1$. On a donc prouvé qu'un nombre est divisible par $3$ si et seulement si la somme de ses chiffres est divisible par 3.
				
				\item  Pour $9$, le raisonnement est similaire, et on prouve qu'un nombre est divisible par $9$ si la somme de ses chiffres est divisible par 9.
				
				\item  Pour $10$, on a $r_0=1$, $r_1=0$ et on prouve par une récurrence immédiate que $r_i=0$ pour $i\geq 1$. Un nombre est donc divisible par $10$ si et seulement si
				son dernier chiffre est nul (bien sûr, on le savait déjà!).
				
				\item  Pour $11$, on remarque que $r_0=1$, que $r_1=10$, puisque que $r_2=1$, $r_3=10,\dots$. On prouve donc par récurrence sur $k$ que
				$r_{2k}=1$ et $r_{2k+1}=10$. C'est vrai pour $k=0$, et si c'est vrai au rang $k$, alors
				$$10r_{2k+1}=100=9\times 11+1,$$
				ce qui prouve que $r_{2k+2}=1$. Il vient ensuite $r_{2k+3}=10$, ce qui prouve l'assertion par récurrence. Pour en déduire un critère
				de divisibilité commode, on remarque que $10\equiv -1\ [11]$. $a$ est donc divisible par $11$ si et seulement si
				$\sum_{k=0}^n (-1)^k a_k$ est divisible par 11.
				
			\end{itemize}
			
			
		\end{enumerate}
	\end{corrige}
	
	\begin{exo}[{Équations de congruence du premier degré}]
		Résoudre les équations suivantes :
		
		\begin{enumerate}
			
			\item  $5x\equiv 3\ [17]$;
			
			\item  $10x\equiv 6\ [34]$;
			
			\item  $10x \equiv 5\ [34]$.
			
		\end{enumerate}
	\end{exo}
	\begin{corrige}
		\begin{enumerate}
			
			\item  On commence par rechercher un entier $a$ tel que $5a\equiv 1\ [17]$ (on dit que $a$ est un inverse de $5$ modulo $17$). On cherche donc $a\in\mathbb Z$ tel qu'il existe $k\in\mathbb Z$ avec $5a=1+17k$. On cherche donc les coefficients d'une relation de Bézout. On peut bien sûr appliquer l'algorithme d'Euclide pour les déterminer, mais ici il est facile de constater que $5\times 7=1+17\times 2$. Autrement dit, $5 \times 7\equiv 1\ [17]$. \\
			
			Soit maintenant $x\in\mathbb Z$ une solution de l'équation. Alors en multipliant l'équation $5x\equiv 3\ [17]$ par $7$, on trouve $35x\equiv 21\ [17]$ et comme $35\equiv 1\ [17]$ et $21\equiv 4\ [17]$, on obtient finalement $x\equiv 4\ [17]$ c'est-à-dire qu'il existe $k\in\mathbb Z$ tel que $x=4+17k$. \\
			
			Réciproquement, si $x$ s'écrit $4+17k$, alors $5x=20+17\times (5k)$ et donc \begin{eqnarray*}
				5x&\equiv& 20 \ [17]\\
				& \equiv& 3\ [17].
			\end{eqnarray*}
			L'ensemble des solutions est donc $\{4+17k:\ k\in\mathbb Z\}$. 
			
			\item  On remarque que
			\begin{eqnarray*}
				10x\equiv 6\ [34]&\iff& \exists k\in\mathbb Z,\ 10x=6+34k\\
				&\iff&\exists k\in\mathbb Z,\ 5x=3+17k\\
				&\iff&5x\equiv 3\ [17]\\
				&\iff&\exists k\in\mathbb Z,\ x=4+17k.
			\end{eqnarray*}
			
			\item  Dans ce cas, l'équation n'admet pas de solutions. En effet, si $x$ est solution, alors il existe $k\in\mathbb Z$ tel que $10x=5+34k$ soit $5=10x-34k$. Mais ceci est impossible car $2$ divise le membre de droite de l'égalité, mais pas le membre de gauche. 
			
		\end{enumerate}
	\end{corrige}
	
	\begin{exo}[{Un système de congruences}]
		Résoudre le système suivant, d'inconnue $x\in\mathbb Z$ : 
		$$\left\{
		\begin{array}{rcl}
			x&\equiv&1\ [5]\\
			x&\equiv&2\ [11].
		\end{array}\right.$$
	\end{exo}
	\begin{corrige}
		Soit $x$ une solution du système. Puisque $x\equiv 1\ [5]$, il existe $k\in\mathbb Z$ tel que $x=1+5k$. Reportant dans la deuxième équation, on trouve $1+5k\equiv 2\ [11]$ soit $5k\equiv 1\ [11]$. Multiplions ceci par $2$. Puisque $10\equiv -1\ [11]$, on obtient 
		$$-k\equiv 2\ [11]\iff k\equiv -2\ [11].$$
		Il existe donc $\ell\in\mathbb Z$ tel que $k=-2+11\ell$. Finalement, on a prouvé que si $x$ est solution, alors $x=1-10+55\ell=-9+55\ell$, $\ell\in\mathbb Z$.\\
		
		Réciproquement, si $x=-9+55\ell$, $\ell\in\mathbb Z$, on écrit $x=1+5\times (-2+11\ell)$ puis $x=2+11\times(-1+5\ell)$ pour prouver que $x$ est bien solution du système.
	\end{corrige}
	
	\begin{exo}[{Nombres palindromes}]
		Un nombre palindrome est un nombre qui se lit indifféremment de gauche à droite ou de droite à gauche.
		Par exemple, 2002, 12321 sont des nombres palindromes. Prouver qu'un nombre palindrome ayant un nombre pair de chiffres est divisible par 11
	\end{exo}
	\begin{corrige}
		Écrivons un tel nombre $n=a_1\dots a_p a_p\dots a_1$, c'est-à-dire 
		\begin{eqnarray*}
			n&=&10^{2p-1}a_1+10^{2p-2}a_2+\dots+10^p a_p+10^{p-1}a_{p}+10 a_2+a_1\\
			&=&\sum_{k=1}^p (10^{k-1}+10^{2p-k})a_k.
		\end{eqnarray*}
		Remarquons ensuite que $10=-1\textrm{ (mod }11)$, et donc $10^k=(-1)^k\textrm{ (mod }11)$. On obtient donc
		\begin{eqnarray*}
			10^{k-1}+10^{2p-k}&=&(-1)^{k-1}+(-1)^{2p-k}\textrm{ (mod }11)\\
			&=&(-1)^{k-1}+(-1)^k\textrm{ (mod }11)\\
			&=&0\textrm{ (mod }11)
		\end{eqnarray*}
		Ainsi, chaque terme $(10^{k-1}+10^{2p-k})a_k$ est congru à 0 modulo 11. Il en est de même pour $n$.
	\end{corrige}
	
	\begin{exo}[{Une équation diophantienne}]
		Le but de l'exercice est de déterminer tous les couples d'entiers $(m,n)\in\mathbb N^2$ tels que $2^m-3^n=1$.
		
		\begin{enumerate}
			
			\item  Déterminer le reste de la division euclidienne de $3^n$ par $8$ (on distinguera les cas $n$ pair et $n$ impair).
			
			\item  En déduire que si $(m,n)\in\mathbb N^2$ est solution de $2^m-3^n=1$, alors $m\leq 2$.
			
			\item  Déterminer toutes les solutions de l'équation.
			
		\end{enumerate}
	\end{exo}
	\begin{corrige}
		\begin{enumerate}
			
			\item  On remarque que $3^2=9\equiv 1\ [8]$. En utilisant les propriétés sur les congruences, on trouve que, si $n=2k$ est pair, 
			$$3^{2k}=(3^2)^k \equiv 1^k\ [8]\equiv 1\ [8]$$
			et si $n=2k+1$ est impair, 
			$$3^{2k+1}=3^{2k}3\equiv 1\times 3\ [8]\equiv 3\ [8].$$
			Le reste de la division euclidienne de $3^n$ par 8 est donc égal à 1 si $n$ est pair et à 3 si $n$ est impair.
			
			\item  Supposons qu'il existe une solution avec $m\geq 3$. Alors $2^m=2^{m-3}2^3=82^{m-3}\equiv 0\ [8]$.
			On en déduit que 
			$$2^m-3^n\equiv 1\ [8]\implies -3^n\equiv 1\ [8].$$
			Mais ceci est impossible, puisque si $n$ est pair, $-3^n\equiv -1\ [8]$ et $-1$ n'est pas congru à $1$ modulo $8$, 
			et si $n$ est impair, $-3^n\equiv -3\ [8]\equiv 5\ [8]$ et $5$ n'est pas congru à $1$ modulo $8$.
			
			\item  Pour $m=0$, on a $2^m-3^n=1\iff -3^n=0$, ce qui n'est jamais le cas. Si $m=1$, on a $2^m-3^n=1\iff 3^n=1\iff n=0$.
			Le couple $(1,0)$ est donc solution. Si $m=2$, alors $2^m-3^n=1\iff 3^n=3\iff n=1$. Le couple $(2,1)$ est solution.
			Donc les solutions sont les couples $(1,0)$ et $(2,1)$.
			
		\end{enumerate}
	\end{corrige}
	
	\begin{exo}[{Somme de deux carrés divisible par 7}]
		Soit $a$ et $b$ deux entiers tels que $a^2+b^2$ soit divisible par $7$. Démontrer que $a$ et $b$ sont divisibles par $7$.
	\end{exo}
	\begin{corrige}
		Commençons par écrire un tableau décrivant les carrés modulo 7.
		$$\begin{array}{c|c|c|c|c|c|c|c}
			x\ [7]&0&1&2&3&4&5&6\\
			\hline
			x^2\ [7]&0&1&4&2&2&4&1\\
		\end{array}$$
		On distingue alors 4 cas :
		
		\begin{itemize}
			
			\item  Si $a^2=0$ modulo 7, alors $a^2+b^2=0$ modulo 7 si et seulement si $b=0$ modulo 7 : c'est le cas où $a$ et $b$ sont tous les deux divisibles par 7.
			
			\item  Si $a^2=1$ modulo 7, pour que $7|a^2+b^2$, il faudrait que $b^2=6$ modulo 7, ce qui est impossible.
			
			\item  Si $a^2=2$ modulo 7, pour que $7|a^2+b^2$, il faudrait que $b^2=5$ modulo 7, ce qui est impossible. 
			\item  Si $a^2=4$ modulo 7, pour que $7|a^2+b^2$, il faudrait que $b^2=3$ modulo 7, ce qui est impossible. 
			
		\end{itemize}
		
		Ainsi la seule possibilité est bien que $a$ et $b$ soient divisibles par 7.
	\end{corrige}
	
	\begin{exo}[{Coefficients binomiaux}]
		\begin{enumerate}
			
			\item  Soit $n\geq 1$. Montrer que $(n+1)|\binom{2n}{n}$.
			
			\item  Soit $p\geq 2$ premier. Montrer que $p|\binom{p}{k}$ pour $k\in\{1,\dots,p-1\}$.
			
			\item  En déduire une preuve du petit théorème de Fermat : si $n\geq 1$ et $p$ est premier, $n^p\equiv n\ [p]$.
			
		\end{enumerate}
	\end{exo}
	\begin{corrige}
		\begin{enumerate}
			
			\item  Il est clair que $(n+1)\binom{2n}{n+1}=n\binom{2n}{n}$. Maintenant,
			comme $(n+1)\wedge n=1$, on en déduit que $(n+1)|\binom{2n}n$.
			
			\item  Puisque $\binom{p}{k}$ est entier, on sait que $k!(p-k)! | p!=p\times(p-1)!$.
			Maintenant, puisque $p$ est premier, on sait aussi que $k!\wedge p=1$ et $(p-k)!\wedge p=1$. Par le théorème de Gauss,
			$k!(p-k)! |(p-1)!$ et donc $\frac{(p-1)!}{k!(p-k)!}$ est un entier. Autrement dit, $p|\binom{p}{k}$.
			
			\item  Si $n=1$, le résultat est trivial. Supposons le résultat établi au rang $n$ et prouvons le au rang $n+1$.
			On a
			$$(n+1)^p=n^p+1+\sum_{k=1}^{p-1}\binom{p}{k}n^{k}\equiv n^p+1\ [p]$$
			d'après la question précédente. Par hypothèse de récurrence,
			$$(n+1)^p\equiv n+1\ [p].$$
			
		\end{enumerate}
	\end{corrige}
	
	\subsection*{Exercices complémentaires}
	
	\begin{exo}[{Divisibilité}]\ \textcolor{red!70!black}{$\star$}
		Montrer que pour tout $n\in\Z$, $6$ divise $n^3-n$, et que $30$ divise $n^5-n$.
	\end{exo}
	\begin{corrige}
		$n^3-n=(n-1)n(n+1)$ est un produit de trois entiers consécutifs, donc divisible par $2$ et par $3$, donc par $6$. Ensuite $n^5-n=(n^3-n)(n^2+1)$ est divisible par $6$. Modulo $5$ : si $n\equiv0,\pm1$, alors $5\mid n(n^2-1)$ ; si $n\equiv\pm2$, alors $n^2+1\equiv5\equiv0$. Donc $5$ divise $n^5-n$. Comme $6$ et $5$ sont premiers entre eux, $30\mid n^5-n$.
	\end{corrige}
	
	\begin{exo}[{Équation de Bézout}]\ \textcolor{red!70!black}{$\star$}
		\begin{enumerate}
			\item À l'aide de l'algorithme d'Euclide, calculer $\pgcd(35,27)$ et trouver une relation de Bézout.
			\item Résoudre dans $\Z^2$ l'équation $35x+27y=4$.
		\end{enumerate}
	\end{exo}
	\begin{corrige}
		\begin{enumerate}
			\item $35=1\cdot27+8$, $27=3\cdot8+3$, $8=2\cdot3+2$, $3=1\cdot2+1$ : $\pgcd(35,27)=1$. En remontant : $1=3-2=3-(8-2\cdot3)=3\cdot3-8=3(27-3\cdot8)-8=3\cdot27-10\cdot8=13\cdot27-10\cdot35$. Donc $35\cdot(-10)+27\cdot13=1$.
			\item Une solution particulière est $(x_0,y_0)=(-40,52)$ (car $35\times(-40)+27\times52=-1400+1404=4$). Si $(x,y)$ est solution, $35(x-x_0)=-27(y-y_0)$ ; comme $\pgcd(35,27)=1$, le lemme de Gauss donne $27\mid x-x_0$. Les solutions sont $(x,y)=(-40+27k,\ 52-35k)$, $k\in\Z$.
		\end{enumerate}
	\end{corrige}
	
	\begin{exo}[{Congruences}]\ \textcolor{red!70!black}{$\star$}
		\begin{enumerate}
			\item Déterminer le reste de la division euclidienne de $3^{2026}$ par $7$.
			\item Résoudre dans $\Z$ la congruence $5x\equiv 3\ [11]$.
			\item Montrer que l'équation $x^2+y^2=4z+3$ n'a pas de solution entière.
		\end{enumerate}
	\end{exo}
	\begin{corrige}
		\begin{enumerate}
			\item $3^6\equiv1\ [7]$ (petit théorème de Fermat) et $2026=6\times337+4$, donc $3^{2026}\equiv3^4=81\equiv4\ [7]$. Le reste est $4$.
			\item $5\times9=45\equiv1\ [11]$, donc $x\equiv9\times3=27\equiv5\ [11]$. Ensemble des solutions : $\{5+11k,\ k\in\Z\}$.
			\item Les carrés sont congrus à $0$ ou $1$ modulo $4$, donc $x^2+y^2\in\{0,1,2\}$ modulo $4$, jamais $3$ ; or $4z+3\equiv3\ [4]$.
		\end{enumerate}
	\end{corrige}
	
	\begin{exo}[{Valuation $p$-adique}]\ \textcolor{red!70!black}{$\star$}
		Soit $p$ premier et $n\in\N^*$. Montrer que l'exposant de $p$ dans la décomposition de $n!$ est $\dis\sum_{k\geq 1}\left\lfloor \frac{n}{p^k}\right\rfloor$. Par combien de zéros se termine l'écriture décimale de $100!$ ?
	\end{exo}
	\begin{corrige}
		Chaque entier $m\leq n$ contribue $v_p(m)=\#\{k\geq1:\ p^k\mid m\}$ à $v_p(n!)=\sum_{m=1}^n v_p(m)$. Il y a $\lfloor n/p^k\rfloor$ multiples de $p^k$ entre $1$ et $n$, d'où $v_p(n!)=\sum_{k\geq1}\lfloor n/p^k\rfloor$. Pour $100!$ : $v_5=20+4=24$ et $v_2\geq50$, donc $100!$ se termine par $24$ zéros.
	\end{corrige}
	\section{Groupes, anneaux, corps}
	
	\subsection*{Groupes}
	
	\begin{exo}[{Un exemple de groupes}]
		On définit, pour $(x,y)$ et $(x',y')$ dans $\mathbb R^*\times\mathbb R$, 
		$$(x,y)\star (x',y')=(xx',xy'+y).$$
		
		\begin{enumerate}
			
			\item  Démontrer que $(\mathbb R^*\times \mathbb R,\star)$ est un groupe. Est-il commutatif?
			
			\item  Simplifier $(x,y)^n$ pour tout $(x,y)\in\mathbb R^*\times\mathbb R$ et tout $n\in\mathbb N^*$.
			
		\end{enumerate}
	\end{exo}
	\begin{corrige}
		\begin{enumerate}
			
			\item  On n'a pas affaire à une loi "classique" et donc on ne peut pas démontrer qu'on a un sous-groupe d'un groupe connu. Il faut donc vérifier à la main les trois propriétés d'un groupe ainsi que le fait qu'il s'agit bien d'une loi interne.
			
			\begin{itemize}
				
				\item  Si $(x,y)$ et $(x',y')$ sont dans $\mathbb R^*\times \mathbb R$, alors $xx'\in\mathbb R^*$ et $xy'+y\in\mathbb R$ et donc il s'agit bien d'une loi interne.
				
				\item  La loi $\star$ est associative : en effet, si $(x,y)$, $(x',y')$ et 
				$(x'',y'')$ sont dans $\mathbb R^*\times\mathbb R$, alors d'une part
				\begin{align*}
					\big( (x,y)\star (x',y')\big)\star (x'',y'')&=
					(xx',xy'+y)\star (x'',y'')\\
					&=(xx'x'',xx'y''+xy'+y)
				\end{align*}
				et d'autre part
				\begin{align*}
					(x,y)\star \big((x',y')\star (x'',y'')\big)&=
					(x,y)\star (x'x'',x'y''+y')\\
					&=(xx'x'',x(x'y''+y')+y)\\
					&=(xx'x'',xx'y''+xy'+y).
				\end{align*}
				
				\item  La loi possède un élément neutre, qui est $(1,0)$. Il est en effet facile de vérifier que 
				$$(x,y)\star (1,0)=(1,0)\star (x,y)=(x,y).$$
				
				\item  Tout élément $(x,y)$ possède un inverse. Expliquons comment le trouver. Il n'est pas très difficile de remarquer qu'il doit s'écrire sous la forme $(1/x,a)$. De
				$$(x,y)\star (1/x,a)=(1,xa+y)$$
				on voit qu'on doit avoir $xa+y=0$ et donc $a=-y/x$. On vérifie alors que $(1/x,-y/x)$ est un inverse de $(x,y)$ : 
				$$(x,y)\star (1/x,-y/x)=(1/x,-y/x)\star (x,y)=(1,0).$$
				
			\end{itemize}
			
			Le groupe n'est pas commutatif. En effet on a 
			$$(1,1)\star (2,1)=(2,2)$$ 
			alors que 
			$$(2,1)\star (1,1)=(2,3).$$
			
			\item  On remarque que 
			$$(x,y)^2=(x^2,xy+y)\textrm{ et }(x,y)^3=(x^3,x^2y+xy+y).$$
			On prouve alors par récurrence sur $n\in\mathbb N^*$ que
			$$(x,y)^n=(x^n,x^{n-1}y+x^{n-2}y+\cdots+xy+y).$$
			
		\end{enumerate}
	\end{corrige}
	
	\begin{exo}[{Sous-groupes ou non?}]
		Dans les questions suivantes, déterminer si la partie $H$ est un sous-groupe du groupe $G$.
		
		\begin{enumerate}
			
			\item  $G=(\mathbb Z,+)$; $H=\{\textrm{nombres pairs}\}$.
			
			\item  $G=(\mathbb Z,+)$; $H=\{\textrm{nombres impairs}\}$.
			
			\item  $G=(\mathbb R,+)$; $H=[-1,+\infty[$.
			
			\item  $G=(\mathbb R^*,\times)$; $H=\mathbb Q^*$.
			
			\item  $G=(\{\textrm{bijections de $E$ dans $E$}\},\circ)$; $H=\{f\in G;\ f(x)=x\}$ où $E$ est un ensemble et $x\in E$.
			
			\item  $G=(\{\textrm{bijections de $E$ dans $E$}\},\circ)$; $H=\{f\in G;\ f(x)=y\}$ où $E$ est un ensemble et $x,y\in E$ avec $x\neq y$.
			
		\end{enumerate}
	\end{exo}
	\begin{corrige}
		\begin{enumerate}
			
			\item  $H$ est un sous-groupe de $G$. En effet, $0\in H$, si $x,y\in H$, alors $-x$ et $x+y$ sont deux entiers pairs
			et donc $-x\in H$, $x+y\in H$. Le théorème de caractérisation des sous-groupes nous dit que $H$ est un sous-groupe de $G$.
			
			\item  $0\notin H$ et donc $H$ n'est pas un sous-groupe de $G$.
			
			\item  $2\in H$ et $-2\notin H$ : $H$ n'est pas un sous-groupe de $G$.
			
			\item  $1\in H$. Si $x=p/q$ et $y=p'/q'$ sont deux rationnels non-nuls, alors $1/x=q/p$ et $x\times y=\frac{p\times p'}{q\times q'}$
			sont deux rationnels non nuls. $H$ est un sous-groupe de $G$.
			
			\item  $Id_E$, l'élément neutre de $G$, est élément de $H$. De plus, si $f,g\in H$, alors 
			$$f(x)=x\implies f^{-1}(f(x))=f^{-1}(x)\implies f^{-1}(x)=x$$
			et
			$$f\circ g(x)=f(g(x))=f(x)=x.$$
			Ainsi, $f^{-1}$ et $f\circ g$ sont éléments de $H$, et $H$ est un sous-groupe de $G$.
			
			\item  L'élément neutre de $G$, $Id_E$, n'est pas élément de $H$ qui n'est donc pas un sous-groupe de $G$.
			
		\end{enumerate}
	\end{corrige}
	
	\begin{exo}[{Quelques exemples de sous-groupes}]
		Démontrer pour chaque question que $H$ est un sous-groupe de $G$.
		
		\begin{enumerate}
			
			\item  $G=(\mathbb C^*,\times)$ et $H=\{z\in \mathbb C^*:\ \exists n\in\mathbb N^*,\ z^n=1\}.$
			
			\item  $G=(\mathbb R^*,\times)$ et $H=\{a+b\sqrt 2:\ a,b\in\mathbb Q,\ (a,b)\neq (0,0)\}$.
			
		\end{enumerate}
	\end{exo}
	\begin{corrige}
		\begin{enumerate}
			
			\item  On commence par vérifier que $H\subset G$ (c'est-à-dire que $0\notin H$). On vérifie ensuite que $1\in H$ puisque $1^1=1$. Soit $z_1,z_2\in H$. Alors il existe $n_1\in\mathbb N$ et $n_2\in\mathbb N$ tels que $z_1^{n_1}=1$ et $z_2^{n_2}=1$ (attention! on n'a pas forcément $n_1=n_2$). Mais alors 
			$$(z_1z_2)^{n_1n_2}=(z_1^{n_1})^{n_2}(z_2^{n_2})^{n_1}=1^{n_2}1^{n_1}=1.$$
			Ainsi, $z_1z_2\in H$. De plus,
			$$\left(\frac 1z_1\right)^{n_1}=\frac{1}{z_1^{n_1}}=1$$
			et donc $1/z_1\in H$. Ainsi, $H$ est bien un sous-groupe de $G$.
			
			\item  $1=1+0\sqrt 2\in H$ (attention! la condition $(a,b)\neq (0,0)$ signifie qu'on ne peut pas avoir $a=0$ ET $b=0$ en même temps, mais il est tout à fait possible de prendre $a=1$ et $b=0$). Si $x=a+b\sqrt 2$ et $y=c+d\sqrt 2$ sont éléments de $H$, alors 
			$$\frac 1x=\frac{1}{a+b\sqrt 2}=\frac{a-b\sqrt 2}{a^2-2b^2}=\frac{a}{a^2-2b^2}+\frac{-b}{a^2-2b^2}\sqrt 2$$
			est lui aussi un élément de $H$ (on n'a $a^2-2b^2\neq 0$ si $(a,b)\neq (0,0)$ car $\sqrt 2$ est irrationnel). De même, $xy\in H$ puisque 
			$$xy=(ac+2bd)+(ad+bc)\sqrt 2$$
			est un élément de $H$. Remarquons que l'on ne peut pas avoir $ac+2bd=0$ et $ad+bc=0$, sinon on aurait $xy=0$ et donc ou bien $x=0$ ou bien $y=0$
			ce qui n'est pas le cas. Ainsi $H$ est un sous-groupe de $G$.
			
		\end{enumerate}
	\end{corrige}
	
	\begin{exo}[{Quelques sous-groupes usuels}]
		Soit $(G,\cdot)$ un groupe. Démontrer que les parties suivantes sont des sous-groupes de $G$ :
		
		\begin{enumerate}
			
			\item  $C(G)=\{x\in G;\ \forall y\in G, xy=yx\}$ ($C(G)$ s'appelle le centre de $G$);
			
			\item  $aHa^{-1}=\{aha^{-1};\ h\in H\}$ où $a\in G$ et $H$ est un sous-groupe de $G$.
			
			\item  On suppose de plus que $G$ est commutatif. On dit que $x$ est un élément de torsion de $G$ s'il existe $n\in\mathbb N^*$ tel que $x^n=e$. Démontrer que l'ensemble des éléments de torsion de $G$ est un sous-groupe de $G$.
			
		\end{enumerate}
	\end{exo}
	\begin{corrige}
		Il suffit, pour chaque cas, d'appliquer le théorème de caractérisation des sous-groupes.
		
		\begin{enumerate}
			
			\item  $e$ est élément de $C(G)$ car $ey=ye=y$ pour tout $y\in G$. Soient $x_1,x_2\in C(G)$. Alors, 
			pour tout $y\in G$, on a 
			$$x_1x_2y=x_1(x_2y)=(x_1y)x_2=yx_1x_2$$
			et donc $x_1x_2\in C(G)$. Enfin, si $x\in C(G)$, alors pour tout $y\in G$,
			$$xy=yx\implies xyx^{-1}=yxx^{-1}=y\implies x^{-1}xyx^{-1}=x^{-1}y\implies yx^{-1}=x^{-1}y$$
			où on a multiplié à droite puis à gauche par $x^{-1}$. On en déduit que $x^{-1}\in C(G)$ qui est donc
			un sous-groupe de $G$.
			
			\item  Puisque $H$ est un sous-groupe de $G$, $e\in H$ et donc $aea^{-1}\in aHa^{-1}$. Mais $aea^{-1}=e$ et donc $e\in aHa^{-1}$.
			Soient $x=aha^{-1}$ et $y=ah'a^{-1}$ deux éléments de $aHa^{-1}$ avec donc $h,h'\in H$.
			On a 
			$$xy=aha^{-1}ah'a^{-1}=ahh'a^{-1}\in aHa^{-1}$$
			puisque $hh'\in H$ ($H$ est un sous-groupe de $G$). Enfin, si on choisit $h'=h^{-1}$, le calcul précédent montre que 
			$$xy=yx=e$$
			et donc $x^{-1}=y\in aHa^{-1}$
			puisque $h^{-1}\in H$. $aHa^{-1}$ est donc bien un sous-groupe de $G$.
			
			\item  Notons $T$ l'ensemble des éléments de torsion de $G$. On a $e^1=e$, donc $e\in T$. De plus, si $x,y\in T$, avec respectivement $x^n=e$ et $y^m=e$, il suffit de remarquer que 
			$$(y^{-1})^m=(y^m)^{-1}=e^{-1}=e$$
			puis d'utiliser le fait que $x$ et $y^{-1}$ commutent pour prouver que
			$$(xy^{-1})^{nm}=(x^n)^m((y^{-1})^m)^n=e.$$
			Ainsi, $xy^{-1}$ est élément de $T$, et $T$ est bien un sous-groupe de $G$.
			
		\end{enumerate}
	\end{corrige}
	
	\begin{exo}[{Sous-groupe d'une courbe}]
		Montrer que $H=\{x+y\sqrt 3;\ x\in\mathbb N,\ y\in\mathbb Z,\ x^2-3y^2=1\}$ est un sous-groupe de $(\mathbb R_+^*,\times)$.
	\end{exo}
	\begin{corrige}
		La première chose à remarquer est que $H\subset \mathbb R_+^*$. Pour $x+y\sqrt 3\in H$, puisque $x^2-3y^2>0$ et $x\in\mathbb N$, on a $x>\sqrt 3|y|$ et donc $x+y\sqrt 3>0$. On remarque ensuite que $1=1+0\sqrt 3$ est bien un élément de $H$. Soient $a=x+y\sqrt 3$ et $b=u+v\sqrt 3$ deux éléments de $H$. Alors :
		$$(x+y\sqrt 3)(u+v\sqrt 3)=(xu+3yv)+\sqrt 3(xv+yu).$$
		On remarque ensuite que 
		\begin{eqnarray*}
			(xu+3yv)^2-3(xv+yu)^2&=&x^2u^2+9y^2v^2-3x^2v^2-3y^2u^2\\
			&=&x^2(u^2-3v^2)+3y^2(3v^2-u^2)\\
			&=&x^2-3y^2\\
			&=&1.
		\end{eqnarray*}
		De plus, il est clair que $xu+3yv$ et $xv+yu$ sont éléments de $\mathbb Z$. Il reste à voir que 
		$xu+3yv$ est élément de $\mathbb N$. Mais c'est clair car $x\geq \sqrt 3 |y|$ et $u\geq \sqrt 3 |v|$. Ainsi, $ab\in H$.\\
		
		Démontrons finalement que $H$ est bien stable par passage à l'inverse. On a 
		$$\frac 1a=\frac 1{x+y\sqrt 3}=\frac {x-y\sqrt 3}{x^2-3y^2}=x-y\sqrt 3\in H$$
		puisque $x^2-3y^2=1$. Ainsi, $H$ est bien un sous-groupe de $(\mathbb R_+^*,\times)$.
	\end{corrige}
	
	\begin{exo}[{Sous-groupe engendré par une partie}]
		Dans cet exercice, $G$ désigne un groupe.
		
		\begin{enumerate}
			
			\item  Soit $(H_i)_{i\in I}$ une famille quelconque de sous-groupes de $G$.
			Démontrer que $\bigcap_{i\in I}H_i$ est un sous-groupe de $G$.
			
			\item  Soit $X$ une partie de $G$. On note $\langle X\rangle$ l'intersection de tous les sous-groupes de $G$ contenant $X$. Démontrer que $\langle X\rangle$
			est le plus petit sous-groupe de $G$ contenant $X$.
			
			\item  Démontrer que 
			$$\langle X\rangle=\left\{x_1^{\veps_1}\cdots x_n^{\veps_n}:\ n\in\mathbb N,\ x_i\in X, \veps_i=\pm 1\textrm{ pour }i=1,\dots,n\right\}$$
			(avec la convention qu'un produit vide vaut $1_G$).
			
		\end{enumerate}
	\end{exo}
	\begin{corrige}
		\begin{enumerate}
			
			\item  Notons $H=\bigcap_{i\in I}H_i$. 
			
			\begin{itemize}
				
				\item   $1_G\in H$ puisque $1_G$ appartient à tous les $H_i$. 
				
				\item  Soit $x,y\in H$. Pour tout $i\in I$, $x\in H_i$ et $y\in H_i$. Puisque $H_i$ est un sous-groupe de $G$, on a $xy\in H_i$ et comme c'est vrai pour tout $i\in I$, $xy\in H$. 
				
				\item  Soit $x\in H$. Pour tout $i\in I$, $x\in H_i$ et donc $x^{-1}\in H_i$ puisque $H_i$ est un sous-groupe de $G$. Puisque c'est vrai pour tout $i\in I$, $x^{-1}\in H$. 
				
			\end{itemize}
			
			Finalement, on a prouvé que $H$ est un sous-groupe de $G$. 
			
			\item  Notons $(H_i)_{i\in I}$ la famille des sous-groupes de $G$ qui contiennent $X$, de sorte que $\langle X\rangle=\bigcap_{i\in I}H_i$. Alors, d'après la question précédente, $\langle X\rangle$ est bien un sous-groupe de $G$. Puisque $X\subset H_i$ pour tout $i\in I$, $X\subset \langle X\rangle$. Puisque $\langle X\rangle$ est contenu dans tout sous-groupe de $G$ contenant $X$, $\langle X\rangle$ est bien le plus petit sous-groupe de $G$ contenant $X$.
			
			\item  Posons
			$$H=\left\{x_1^{\veps_1}\cdots x_n^{\veps_n}:\ n\in\mathbb N,\ x_i\in X, \veps_i=\pm 1\textrm{ pour }i=1,\dots,n\right\}.$$
			On commence par remarquer que $H$ est un sous-groupe de $G$. En effet, 
			
			\begin{itemize}
				
				\item  $1_G\in H$ (cas du produit vide)
				
				\item  si $x,y\in H$, on écrit $x=x_1^{\veps_1}\cdots x_n^{\veps_n}$ et $y=y_1^{\nu_1}\cdots y_m^{\nu_m}$. Alors 
				$$xy=x_1^{\veps_1}\cdots x_n^{\veps_n}y_1^{\nu_1}\cdots y_m^{\nu_m}$$
				est bien dans $H$.
				
				\item  si $x\in H$, on écrit $x=x_1^{\veps_1}\cdots x_n^{\veps_n}$ d'où
				$$x^{-1}=x_n^{-\veps_n}\cdots x_1^{-\veps_1}$$
				qui est bien élément de $H$ puisque $-\veps_i\in\{-1,1\}$.
				
			\end{itemize}
			
			On remarque ensuite que $X\subset H$ (cas du produit avec un seul élément). On obtient donc $\langle X\rangle \subset H$. Reste à démontrer l'inclusion réciproque. Pour cela, on va prouver par récurrence sur $n\in\mathbb N$ la propriété suivante :
			$$\mathcal P_n="\textrm{tout }x\in G\textrm{ qui s'écrit }x=x_1^{\veps_1}\cdots x_n^{\veps_n}\textrm{ avec }x_i\in X,\ \veps_i=\pm 1\textrm{ appartient à }\langle X\rangle".$$
			Initialisation : soit $x\in G$ tel que $x=x_1$ ou $x=x_1^{-1}$ avec $x_1\in X$. Alors si $x=x_1$, $x\in \langle X\rangle$. Si $x=x_1^{-1}$, alors, comme $x_1\in\langle X\rangle$ et que $\langle X\rangle$ est un sous-groupe de $G$, on a $x=x_1^{-1}\in \langle X\rangle$.\\
			
			Hérédité : soit $n\in\mathbb N^*$ tel que $\mathcal P_n$ est vraie. Soit $x=x_1^{\veps_1}\dots x_n^{\veps_n}x_{n+1}^{\veps_{n+1}}.$ Alors
			par hypothèse de récurrence, $x_1^{\veps_1}\dots x_n^{\veps_n}\in \langle X\rangle$ et comme dans l'initialisation, $x_{n+1}^{\veps_{n+1}}\in\langle X\rangle$. Puisque $\langle X\rangle$ est un groupe, 
			$$x_1^{\veps_1}\dots x_n^{\veps_n}x_{n+1}^{\veps_{n+1}}\in\langle X\rangle.$$
			On a donc prouvé par récurrence que pour tout $n\in\mathbb N^*$, $\mathcal P_n$ est vraie. Ceci signifie exactement que $H\subset\langle X\rangle$. On a donc prouvé, par double inclusion, que $H=\langle X\rangle$.
			
		\end{enumerate}
	\end{corrige}
	
	\begin{exo}[{Intersection de deux sous-groupes}]
		Soit $G$ un groupe et $H_1,H_2$ deux sous-groupes de $G$. Démontrer que $H_1\cap H_2$ est un sous-groupe de $G$.
	\end{exo}
	\begin{corrige}
		On va vérifier les hypothèses du théorème de caractérisation des sous-groupes.
		
		\begin{itemize}
			
			\item  On sait que $1_G\in H_1$ (puisque $H_1$ est un sous-groupe de $G$) et que $1_G\in H_2$ (puisque $H_2$ est un sous-groupe de $G$). Donc $1_G\in H_1\cap H_2$.
			
			\item  Soit $x,y\in H_1\cap H_2$. Alors, $x\in H_1,$ $y\in H_1$ et donc puisque $H_1$ est un sous-groupe de $G$, $xy\in H_1$. De la même façon, on prouve que $xy\in H_2$ et donc que $xy\in H_1\cap H_2$.
			
			\item  Soit $x\in H_1\cap H_2$. Alors, $x\in H_1$ et puisque $H_1$ est un sous-groupe de $G$, $x^{-1}\in H_1$. De même, $x^{-1}\in H_2$ et finalement $x^{-1}\in H_1\cap H_2$.
			
		\end{itemize}
	\end{corrige}
	
	\begin{exo}[{Produit de groupe et sous-groupe du produit}]
		Un sous-groupe d'un groupe produit est-il nécessairement produit de deux sous-groupes?
	\end{exo}
	\begin{corrige}
		Non, ce n'est pas le cas. Prenons $\mathbb Z^2=\mathbb Z\times \mathbb Z$ et $H=\{(x,x);\ x\in\mathbb Z\}$. $H$ est clairement un sous-groupe de $\mathbb Z^2$, et $H$ ne s'écrit pas $H=A\times B$, sinon on aurait $A=B=\mathbb Z$ ce qui n'est pas le cas.
	\end{corrige}
	
	\begin{exo}[{Union de deux sous-groupes}]
		Soit $G$ un groupe et $H,K$ deux sous-groupes de $G$. Démontrer que $H\cup K$ est un sous-groupe de $G$ si et seulement si $H\subset K$ ou $K\subset H$.
	\end{exo}
	\begin{corrige}
		Si $H\subset K$, alors $H\cup K=K$ qui est un sous-groupe de $G$. De même, si $K\subset H$, $H\cup K=H$ qui est un sous-groupe de $G$.\\
		
		Supposons maintenant que $H\cup K$ est un sous-groupe de $G$ et que ni $H\subset K$, ni $K\subset H$. Alors on peut trouver $x\in H\backslash K$ et $y\in K\backslash H$. Puisque $H\cup K$ est un groupe et que $x,y\in H\cup K$, on a $xy\in H\cup K$. Mais si $xy\in H$, alors $y=x^{-1}(xy)$
		est le produit de deux éléments de $H$, qui est un sous-groupe de $G$, et donc $y\in H$ ce qui est une contradiction. On obtient de même une contradiction dans l'autre cas possible $xy\in K$. L'hypothèse de départ est donc fausse, et on a bien $H\subset K$ ou $K\subset H$.
	\end{corrige}
	
	\begin{exo}[{Produit de deux sous-groupes}]
		Soit $(G,\cdot)$ un groupe et $A$, $B$ deux sous-groupes de $G$.
		On note $AB=\{ab;\ a\in A,\ b\in B\}$. Montrer que $AB$ est un sous-groupe de $G$ si
		et seulement si $AB=BA$.
	\end{exo}
	\begin{corrige}
		Supposons d'abord que $AB=BA$. Alors $AB$ est un sous-groupe de $G$ car :
		
		\begin{enumerate}
			
			\item  $e\in AB$, car $e=ee$ avec $e\in A$ et $e\in B$ (ce sont des sous-groupes);
			
			\item  $AB$ est stable par passage au produit. En effet, si $x=ab\in AB$ et $y=a'b'\in AB$, alors
			$xy=aba'b'$. Or, $ba'$ est un élément de $BA$, c'est donc aussi un élément de $AB$ et donc
			$ba'=a''b''$ avec $a''\in A$ et $b''\in B$. On en déduit que
			$$xy=aa''b''b'\in AB$$
			puisque $aa''\in A$ et $b''b'\in B$.
			
			\item  $AB$ est stable par passage à l'inverse. En effet, si $x=ab\in AB$, alors $x^{-1}=b^{-1}a^{-1}$
			est élément de $BA$ et $BA=AB$.
			
		\end{enumerate}
		
		Réciproquement, supposons que $AB$ est un sous-groupe de $G$ et prouvons que $AB=BA$. Soit d'abord $x=ab\in AB$.
		Alors $x^{-1}=b^{-1}a^{-1}\in AB$ puisque $AB$ est un groupe et donc $b^{-1}a^{-1}=a'b'$ avec $a'\in A$ et $b'\in B$. On passe à l'inverse :
		$$ab=b'^{-1}a'^{-1}\in BA.$$
		Pour l'autre inclusion, considérons $y=ba\in BA$. Alors $y^{-1}=a^{-1}b^{-1}\in AB$, et donc $y=(y^{-1})^{-1}\in AB$ puisque $AB$ est un groupe. .
	\end{corrige}
	
	\begin{exo}[{Théorème de Lagrange}]
		Soit $(G,\cdot)$ un groupe fini et $H$ un sous-groupe de $G$.
		
		\begin{enumerate}
			
			\item  Montrer que pour tout $a\in G$, $H$ et $aH=\{ah;\ h\in H\}$ ont le même nombre d'éléments.
			
			\item  Soient $a,b\in G$. Démontrer que $aH=bH$ ou $aH\cap bH=\varnothing$.
			
			\item  En déduire que le cardinal de $H$ divise le cardinal de $G$.
			
		\end{enumerate}
	\end{exo}
	\begin{corrige}
		\begin{enumerate}
			
			\item  Soit $f:H\to aH$ définie par $f(h)=ah$. Il s'agit clairement d'une surjection de $H$ sur $aH$. De plus, si $ah_1=ah_2$, alors $h_1=h_2$ car $a$ est inversible, et donc $f$ est aussi injective. $f$ est donc une bijection de $H$ sur $aH$; ces deux ensembles ont le même nombre d'éléments.
			
			\item  Supposons que $aH\cap bH\neq \varnothing$ et prouvons que $aH=bH$. Par symétrie, il suffit de prouver que $aH\subset bH$. Soit $x\in aH\cap bH$, $x=ah_1=bh_2$. Prenons $y=ah\in aH$. Alors 
			$a=bh_2h_1^{-1}$ et donc $y=b h_2h_1^{-1}h\in bH$.
			
			\item  La réunion des ensembles $aH$ est clairement égale à $G$ (si $x\in G$, il est dans $xH$). On ne garde que les $aH$ deux à deux disjoints et par les deux questions précédentes, on réalise ainsi une partition de $G$ avec des ensembles qui ont tous le même cardinal, à savoir le cardinal de $H$. Si $k$ est le nombre d'ensembles nécessaires pour réaliser cette partition, on a 
			$$k\card(H)=\card(G)$$
			et donc le cardinal de $H$ divise celui de $G$.
			
		\end{enumerate}
	\end{corrige}
	
	\begin{exo}[{Des propriétés bien connues}]
		Traduire en termes de morphismes de groupes les propriétés bien connues suivantes (dont le domaine de validité a volontairement été omis) :
		
		\begin{enumerate}
			
			\item  $\ln(xy)=\ln(x)+\ln(y)$;
			
			\item  $|zz'|=|z||z'|$;
			
			\item  $\sqrt{xy}=\sqrt{x}\sqrt{y}$;
			
			\item  $e^{x+y}=e^xe^y$;
			
		\end{enumerate}
	\end{exo}
	\begin{corrige}
		\begin{enumerate}
			
			\item  La fonction $\ln$ est un morphisme du groupe $(\mathbb R_+^*,\times)$ dans le groupe $(\mathbb R,+)$.
			
			\item  La fonction $|\cdot|$ est un morphisme de groupes de $(\mathbb C^*,\times)$ dans lui-même, ou de $(\mathbb C^*,\times)$ dans $(\mathbb R^*,\times)$,
			ou de $(\mathbb C^*,\times)$ dans $(\mathbb R_+^*,\times)$.
			Attention, même si la propriété est vraie pour $z=0$, il faut exclure $0$ (car il n'est pas inversible pour la multiplication).
			
			\item  La fonction $\sqrt{}$ est un morphisme du groupe $(\mathbb R_+^*,\times)$ dans lui-même.
			
			\item  La fonction $\exp$ est un morphisme de groupe de $(\mathbb R,+)$ dans $(\mathbb R_+^*,\times)$.
			
		\end{enumerate}
	\end{corrige}
	
	\begin{exo}[{Exponentielle complexe}]
		Justifier que $\exp$ est un morphisme de $(\mathbb C,+)$ dans $(\mathbb C^*,\cdot)$. Quel est son image? Son noyau?
	\end{exo}
	\begin{corrige}
		On sait que, pour tous $z,w\in \mathbb C$, on a
		$$\exp(z+w)=\exp(z)\exp(w).$$
		Ceci signifie exactement que $\exp$ est un morphisme de $(\mathbb C,+)$ dans $(\mathbb C^*,\cdot)$ (la fonction exponentielle ne prend jamais la valeur zéro). De plus, soit $w\in\mathbb C^*, w=re^{i\theta}$ avec $r>0$. Soit $a=\ln r$ et posons $z=a+i\theta$. Alors
		$$\exp(z)=\exp(a)\exp(i\theta)=re^{i\theta}=w.$$
		L'exponentielle est un morphisme surjectif de $(\mathbb C,+)$ dans $(\mathbb C^*,\cdot)$. Déterminons son noyau. Si $\exp(z)=1$,  posons $z=x+iy$. Alors 
		$$\exp(z)=\exp(x)\exp(iy)=1=1\exp(i0).$$
		Ceci est équivalent à $x=0$ et il existe $k\in \mathbb Z$ tel que $y=k2\pi$. On a donc
		$$\ker(\exp)=\{2ik\pi;\ k\in\mathbb Z\}.$$
	\end{corrige}
	
	\begin{exo}[{Morphismes de $\mathbb Z$ dans $\mathbb Z$}]
		Déterminer tous les morphismes de $(\mathbb Z,+)$ dans lui-même. Lesquels sont injectifs? surjectifs?
	\end{exo}
	\begin{corrige}
		Soit $f$ un morphisme de $(\mathbb Z,+)$. Prouvons par récurrence que pour tout $n\geq 1$, on a
		$f(n)=nf(1)$. C'est vrai pour $n=1$, et si c'est vrai pour $n$, alors
		$$f(n+1)=f(n)+f(1)=nf(1)+f(1)=(n+1)f(1).$$
		De plus, pour $n\leq 0$, on a $-n\geq 0$ et donc $f(-n)=-nf(1)$. On en déduit :
		$$0=f(0)=f(n+(-n))=f(n)+f(-n)=f(n)-nf(1).$$
		Ainsi, on a toujours $f(n)=nf(1)$, quel que soit $n\in\mathbb Z$.\\
		
		Caractérisons maintenant les morphismes surjectifs. Supposons donc que $f$ est surjectif.
		Tout élément de $\mathbb Z=f(\mathbb Z)$ est un multiple de $f(1)$.
		Or, les seuls éléments de $\mathbb Z$ qui divisent tous les autres entiers sont $1$ et $-1$.
		On en déduit que $f(1)=1$ ou $f(1)=-1$, et donc que $f(n)=n$ ou $f(n)=-n$. 
		Réciproquement, ces deux applications sont clairement des morphismes surjectifs de $(\mathbb Z,+)$.\\
		
		Déterminons enfin les morphismes injectifs. Soit $f$ un morphisme et $n\in \ker(f)$. Alors $f(n)=nf(1)=0$. Si $f(1)\neq 0$, alors $f(n)=0\iff n=0$ et $f$ est injectif, et si $f(1)=0$, alors $f$ n'est pas injectif. Donc tous les morphismes de $(\mathbb Z,+)$ dans $(\mathbb Z,+)$ sont injectifs sauf l'application identiquement nulle.
	\end{corrige}
	
	\begin{exo}[{Morphismes de $(\mathbb Q,+)$ dans $(\mathbb Z,+)$}]
		Déterminer tous les morphismes de groupes de $(\mathbb Q,+)$ dans $(\mathbb Z,+)$.
	\end{exo}
	\begin{corrige}
		Soit $f$ un tel morphisme de groupe. On va commencer par démontrer que pour $p$ et $q$ des entiers naturels, on a 
		$f(p)=pf(1)$ et $f\left(\frac 1q\right)=\frac{1}{q}f(1)$. La première des deux propriétés se démontre aisément par récurrence.
		Pour la deuxième, on écrit que 
		$$f(1)=f\left(\frac 1q+\cdots+\frac 1q\right)=qf\left(\frac 1q\right).$$
		Notons ensuite $a=f(1)$.  Alors $a/q$ est un entier pour tout entier $q$, et donc $a=0$. On en déduit que $f(p)=f(1/q)=0$ pour tous les entiers $p$ et $q$,
		puis que $f(p/q)=pf(1/q)=0$. Finalement, on trouve que $f$ est le morphisme nul.
	\end{corrige}
	
	\begin{exo}[{Morphisme entre groupes de torsion et groupes sans torsion}]
		Dans un groupe $(G,\cdot)$, un élément $x$ est dit de torsion s'il existe $n\geq 1$ tel que $x^n=e$. On dit que $G$ est de torsion si tous ses éléments sont de torsion.
		On dit que $G$ est sans torsion si son seul élément de torsion est l'élément neutre.
		Soit $G_1$ un groupe de torsion et $G_2$ un groupe sans torsion. Déterminer tous les morphismes de groupe de $G_1$ dans $G_2$.
	\end{exo}
	\begin{corrige}
		Remarquons d'abord que le morphisme qui à tout $x$ de $G_1$ associe l'élément neutre $e$ de $G_2$ convient.
		Réciproquement, soit $x$ un élément de $G_1$ et notons $y=f(x)$. Puisque $x$ est un élément de torsion,
		il existe $n\geq 1$ tel que $x^n=e$. Mais alors, 
		$$y^n=(f(x))^n=f(x^n)=f(e)=e.$$
		Comme $G_2$ est sans torsion, on a $y=e$ et donc $f(x)=e$ pour tout $x\in G_1$.
	\end{corrige}
	
	\begin{exo}[{Groupes non isomorphes}]
		Démontrer que les groupes multiplicatifs $(\mathbb R^*,\cdot)$ et $(\mathbb C^*,\cdot)$ ne sont pas isomorphes.
	\end{exo}
	\begin{corrige}
		Supposons que ces deux groupes sont isomorphes et soit $f$ un isomorphisme de $(\mathbb C^*,\cdot)$ dans $(\mathbb R^*,\cdot)$. Posons $a=f(i)$. Alors
		$$f(i^4)=a^4=1$$
		et donc $a^2=1$ puisque $a^2>0$. D'où $1=a^2=f(i^2)=f(-1)$ et $1=f(1)$. $f$ ne peut pas être injectif, on a obtenu une contradiction.
	\end{corrige}
	
	\begin{exo}[{Groupes non isomorphes}]
		Démontrer que les groupes $(\mathbb R^*,\times)$ et $(\mathbb Q^*,\times)$ ne sont pas isomorphes.
	\end{exo}
	\begin{corrige}
		Procédons par l'absurde et supposons qu'il existe $f:(\mathbb R^*,\times)\to(\mathbb Q^*,\times)$ un isomorphisme. En particulier, il existe $a\in\mathbb R^*$ tel que $f(a)=2$. Si $a>0$, alors il existe un réel $b$ tel que $a=b^2$. On a alors $2=f(a)=f(b)\times f(b)$ et donc $\sqrt 2=\pm f(b)$. Mais ceci contredit que $\sqrt 2$ est irrationnel. \\
		
		Si $a<0$, alors il existe un réel $b$ tel que $a=-b^2$. On a alors $2=f(a)=f(-1)f(b)\times f(b)$. Mais $f(-1)^2=f(-1\times -1)=f(1)=1$ et donc $f(-1)=\pm 1$. Comme $f$ est un isomorphisme et $f(1)=1$, on a $f(-1)=-1$. On obtient alors $2=-f(b)^2$, ce qui est impossible puisque $2>0$ et $-f(b)^2<0$. \\
		
		Dans les deux cas, on a obtenu une contradiction et l'existence d'un tel isomorphisme est impossible. On aurait aussi pu utiliser des arguments de cardinalité : $\mathbb Q^*$ est dénombrable, et $\mathbb R^*$ ne l'est pas. Donc il ne peut pas exister de bijection entre ces deux ensembles.
	\end{corrige}
	
	\begin{exo}[{Automorphisme intérieur}]
		Soit $(G,\cdot)$ un groupe. Pour $a\in G$, on note $\tau_a:G\to G$ défini par $\tau_a(x)=axa^{-1}$. 
		
		\begin{enumerate}
			
			\item  Démontrer que $\tau_a$ est un endomorphisme de $G$.
			
			\item  Vérifier que, pour tous $a,b\in G$, $\tau_a\circ \tau_b=\tau_{ab}$.
			
			\item  Montrer que $\tau_a$ est bijective et déterminer son inverse.
			
			\item  En déduire que $\Theta=\{\tau_a;\ a\in G\}$ muni du produit de composition est un groupe.
			
		\end{enumerate}
	\end{exo}
	\begin{corrige}
		\begin{enumerate}
			
			\item  Il suffit d'appliquer la définition : pour tous $x,y\in G$, on a 
			$$\tau_a(xy)=axya^{-1}=axa^{-1}aya^{-1}=\tau_a(x)\tau_a(y).$$
			
			\item  Soit $x\in G$. On a 
			$$\tau_a\circ\tau_b(x)=\tau_a(bxb^{-1})=abxb^{-1}a^{-1}$$
			tandis que 
			$$\tau_{ab}(x)=abx(ab)^{-1}=abxb^{-1}a^{-1}.$$
			On a donc $\tau_a\circ\tau_b=\tau_{ab}$.
			
			\item  Soit $a\in G$. On pourrait prouver que $\tau_a$ est injectif en calculant son noyau, puisqu'il est surjectif, mais c'est plus facile d'appliquer la question précédente. Avec $b=a^{-1}$, elle donne
			$$\tau_a\circ \tau_{a^{-1}}=\tau_{aa^{-1}}=\tau_e=Id_G$$
			et en inversant le rôle joué par $a$ et $b$, on a aussi
			$$\tau_{a^{-1}}\circ\tau_a=Id_G.$$
			Ainsi, $\tau_a$ est inversible d'inverse $\tau_{a^{-1}}$.
			
			\item  On va prouver que $\Theta=\{\tau_a;\ a\in G\}$ est un sous-groupe de $(S_G,\circ)$. Il est non-vide parce qu'il contient $\tau_e$. Si $\tau_a,\tau_b\in\Theta$, alors 
			$$(\tau_a)^{-1}=\tau_{a^{-1}}\in\Theta$$
			et 
			$$\tau_a\circ\tau_b=\tau_{ab}\in\Theta.$$
			Ainsi, $(\Theta,\circ)$ est bien un sous-groupe de $(S_G,\circ)$.
			
		\end{enumerate}
	\end{corrige}
	
	\begin{exo}[{Somme des valeurs}]
		Soit $f$ un morphisme non constant d'un groupe fini $(G,\cdot)$ dans $(\mathbb C^*,\cdot)$. Calculer
		$\sum_{x\in G}f(x)$.
	\end{exo}
	\begin{corrige}
		Puisque $f$ n'est pas constante, il existe $a\in G$ tel que $f(a)\neq 1$. Maintenant, l'application $x\mapsto ax$ est une permutation de $G$ : en effet, pour tout $y\in G$, il existe un unique $x\in G$ tel que $y=ax$ ($x$ est égal à $a^{-1}y$). On en déduit que
		$$\sum_{x\in G}f(ax)=\sum_{x\in G}f(x).$$
		Mais d'autre part, puisque $f$ est un morphisme de groupes, on a aussi
		$$\sum_{x\in G}f(ax)=\sum_{x\in G}f(a)f(x)=f(a)\sum_{x\in G}f(x).$$
		Ainsi, il vient
		$$\big(f(a)-1\big)\times \sum_{x\in G}f(x)=0.$$
		Puisque $f(a)\neq 1$, on en déduit que $\sum_{x\in G}f(x)=0.$
	\end{corrige}
	
	\subsection*{Anneaux}
	
	\begin{exo}[{Anneau de matrices}]
		Soit $A$ l'ensemble des matrices s'écrivant $\begin{pmatrix}a&b\\0&a\end{pmatrix}$ avec $a$ et $b$ des entiers relatifs. 
		
		\begin{enumerate}
			
			\item  Démontrer que $A$ est un anneau pour les lois d'addition et de produits de matrices.
			
			\item  Déterminer les éléments inversibles de $A$. 
			
		\end{enumerate}
	\end{exo}
	\begin{corrige}
		\begin{enumerate}
			
			\item  Il suffit de démontrer que $A$ est un sous-anneau de $\mathcal M_2(\mathbb R).$ Pour cela, on remarque que
			
			\begin{itemize}
				
				\item  $I_2\in A$ (choisir $a=1$ et $b=0$).
				
				\item  Si $M=\begin{pmatrix}a&b\\0&a\end{pmatrix}$ et $M'=\begin{pmatrix}a'&b'\\0&a'\end{pmatrix}$ sont dans $A$, alors 
				$$M-M'=\begin{pmatrix}a-a'&b-b'\\0&a-a'\end{pmatrix}$$
				et 
				$$MM'=\begin{pmatrix}aa'&ab'+ba'\\0&aa'\end{pmatrix}$$
				sont dans $A.$
				
			\end{itemize}
			
			
			\item  Soit $M=\begin{pmatrix}a&b\\0&a\end{pmatrix}$ inversible dans $A$. Alors il existe $M'=\begin{pmatrix}a'&b'\\0&a'\end{pmatrix}$ tel que $MM'=I_2$. On obtient $aa'=1$, et comme $a$ et $a'$ sont dans $\mathbb Z$, on doit avoir $a=\pm 1$. Dans ce cas, on aura $a'=a$. De plus, on doit aussi avoir 
			$ab'+ba'=0$ ce qui entraîne $b'=-b$.\\
			
			Réciproquement, soit $M\in A$ qui s'écrit
			$$M=\begin{pmatrix}a&b\\0&a\end{pmatrix}$$
			avec $a=\pm 1$.
			Posons 
			$$M'=\begin{pmatrix}a&-b\\0&a\end{pmatrix}.$$
			Alors 
			$$MM'=\begin{pmatrix}a^2&ab-ba\\0&a^2\end{pmatrix}=I_2.$$
			Ainsi, $M'$ est inversible. On a donc démontré que les éléments inversibles de $A$ sont ceux pour lesquels $a=1$ ou $a=-1.$
			
		\end{enumerate}
	\end{corrige}
	
	\begin{exo}[{Anneau des entiers de Gauss}]
		On appelle ensemble des entiers de Gauss noté $\mathbb Z[i]$ l'ensemble des nombres complexes qui s'écrivent $a+ib$, avec $a$ et $b\in\mathbb Z.$
		
		\begin{enumerate}
			
			\item  Démontrer que $\mathbb Z[i]$ est un anneau.
			
			\item  Pour tout nombre complexe $z$, on note $N(z)=z\bar z.$
			
			\begin{enumerate}
				
				\item  Démontrer que, pour tous nombres complexes $z$ et $z'$, $N(z)N(z')=N(zz').$
				
				\item  Démontrer que, pour tout entier de Gauss $z$, $N(z)$ est un entier naturel.
				
				\item  Soit $z$ un entier de Gauss inversible. Déduire des questions précédentes que $N(z)=1$.
				
				\item  Quels sont les éléments inversibles de $\mathbb Z[i]$?
				
			\end{enumerate}
			
			
		\end{enumerate}
	\end{exo}
	\begin{corrige}
		\begin{enumerate}
			
			\item  On va prouver que $\mathbb Z[i]$ est un sous-anneau de $\mathbb C.$
			En effet,
			
			\begin{itemize}
				
				\item  $1=1+0i\in\mathbb Z[i];$
				
				\item  Soit $z=a+ib$ et $z'=a'+ib'\in\mathbb Z[i].$ Alors 
				$$z-z'=(a-a')+i(b-b')\in\mathbb Z[i]$$
				(puisque $a-a'$ et $b-b'\in\mathbb Z$) et 
				$$zz'=(aa'-bb')+i(ab'+a'b)\in\mathbb Z[i]$$
				(puisque $aa-bb'$ et $ab'+a'b\in\mathbb Z$).
				
			\end{itemize}
			
			
			\item  
			\begin{enumerate}
				
				\item  On remarque que $N(z)=|z|^2.$ Puisque $|zz'|=|z|\cdot |z'|,$ en mettant au carré cette égalité, on a le résultat demandé.
				
				\item  Si $z=a+ib$, alors $N(z)=a^2+b^2$ et $a^2,b^2$ sont des entiers naturels, donc $N(z)$ aussi.
				
				\item  Soit $z$ un entier de Gauss inversible et soit $z'$ son inverse. Alors on sait que $zz'=1$ et donc $N(z)\times N(z')=1$. Or le produit de deux entiers naturels est égal à $1$ si et seulement si ces deux entiers sont égaux à $1$. Donc $N(z)=1.$
				
				\item  Soit $z=a+ib\in\mathbb Z[i]$ inversible. Alors $N(z)=1$, et donc $a^2+b^2=1.$ Or, puisque $a^2$ et $b^2$ sont des entiers naturels, ceci n'est possible que dans quatre cas : $(a,b)=(1,0),$ $(a,b)=(-1,0),$ $(a,b)=(0,1),$
				et $(a,b)=(0,-1).$ Réciproquement, il est facile de vérifier que $1$, $-1$, $i$ et $-i$ sont inversibles dans $\mathbb Z[i],$ d'inverse respectif $1$, $-1$, $-i$ et $i$. Les éléments inversibles de $\mathbb Z[i]$ sont donc $1$, $-1$, $i$ et $-i.$
				
			\end{enumerate}
			
			
		\end{enumerate}
	\end{corrige}
	
	\begin{exo}[{Endomorphisme de groupes}]
		Soit $(G,+)$ un groupe commutatif. On note $\textrm{End}(G)$ l'ensemble des endomorphismes de $G$ sur lequel on définit la loi $+$ par $f+g:G\to G,\ x\mapsto f(x)+g(x)$. Démontrer que $(\textrm{End}(G),+,\circ)$ est un anneau.
	\end{exo}
	\begin{corrige}
		On remarque d'abord que $+$ et $\circ$ sont bien des lois de composition interne sur $\textrm{End}(G)$. Ensuite, on vérifie tous les points de la définition d'un anneau.
		
		\begin{enumerate}
			
			\item  $(\textrm{End}(G),+)$ est un groupe commutatif. En effet, la loi $+$ est associative et commutative, l'application $0_G:G\to G,\ g\mapsto 0$ est un élément neutre pour la loi $+$, et tout élément $f\in\textrm{End}(G)$ admet un inverse $-f:G\to G,\ x\mapsto -f(x)$.
			
			\item  La loi $\circ$ est associative.
			
			\item  La loi $\circ$ possède un élément neutre, qui est l'application identité.
			
			\item  La loi $\circ$ est distributive par rapport à la loi $+$ : pour tous $f,g,h\in\textrm{End}(G)$ et tout $x\in G$, 
			$$((f+g)\circ h)(x)=(f+g)(h(x))=f(h(x))+g(h(x))=(f\circ h+g\circ h)(x)$$
			et 
			$$(f\circ (g+h))(x)=f((g+h)(x))=f(g(x)+h(x))=f(g(x))+f(h(x))=(f\circ g+f\circ h)(x).$$
			
		\end{enumerate}
		
		Ainsi, $(\textrm{End}(G),+,\circ)$ est un anneau.
	\end{corrige}
	
	\begin{exo}[{Rationnels à dénominateur impair}]
		Soit $\displaystyle A=\left\{\frac mn;\ m\in\mathbb Z,\ n\in 2\mathbb N+1\right\}$ (c'est-à-dire que $A$ est l'ensemble des rationnels à dénominateur impair). Démontrer que $(A,+,\times)$ est un anneau. Quels sont ses éléments inversibles?
	\end{exo}
	\begin{corrige}
		On va démontrer que $A$ est un sous-anneau de $(\mathbb Q,+,\times)$. Pour cela, soient $x=\frac mn$  et $y=\frac{m'}{n'}\in A$. Alors : 
		$$x-y=\frac{mn'-m'n}{nn'}\textrm{ et }xy=\frac{mm'}{nn'}.$$
		Comme $nn'$, produit de deux nombres impairs, est impair, et que $A$ est non vide puisqu'il contient $1$, on en déduit que $A$ est bien un sous-anneau de $(\mathbb Q,+,\times)$. \\
		
		Déterminons ensuite les inversibles de $A$. Soit $x=\frac mn\in A$ inversible, et soit $y=\frac{m'}{n'}\in A$ tel que $xy=1$. On en déduit que $mm'=nn'$. En particulier, $m$ est nécessairement impair. Réciproquement, si $x=\frac{m}n$ avec $m$ impair, alors $y=\frac nm$ est dans $A$ (si jamais $m<0$, il suffit d'écrire $y=\frac{-n}{-m}$ pour vérifier qu'il est bien dans $A$), et $xy=1$. Ainsi, les inversibles de $A$ sont les éléments $\frac mn$ avec $m\in\mathbb Z$, $n\in\mathbb N^*$, et $m,n$ impairs.
	\end{corrige}
	
	\begin{exo}[{Décimaux}]
		Soit $\mathbb D$ l'ensemble des nombres décimaux, 
		$$\mathbb D=\left\{\frac{n}{10^k};\ n\in\mathbb Z, k\in\mathbb N\right\}.$$
		Démontrer que $(\mathbb D,+,\times)$ est un anneau. Quels sont ses éléments inversibles?
	\end{exo}
	\begin{corrige}
		On va prouver que $(\mathbb D,+,\times)$ est un sous-anneau de $(\mathbb Q,+,\times)$. On remarque d'abord que $\mathbb D\subset\mathbb Q$, puis que $1\in\mathbb D$. De plus, soient $x=\frac n{10^k}$ et $y=\frac{m}{10^l}$ deux éléments de $\mathbb D$. Alors 
		$$x-y=\frac{n10^l-m10^k}{10^{k+l}}\textrm{ et }xy=\frac{nm}{10^{k+l}}$$
		sont clairement des éléments de $\mathbb D$, et $(\mathbb D,+,\times)$ est bien un sous-anneau de $(\mathbb Q,+,\times)$.\\
		
		Déterminons ensuite les inversibles de $(\mathbb D,+,\times)$. Soit $x=\frac{n}{10^k}$ inversible, d'inverse $y=\frac{m}{10^l}$. Alors 
		$$xy=1\iff nm=10^{k+l}.$$
		On en déduit que les seuls diviseurs premiers de $n$ sont $2$ et $5$, autrement dit que $n$ s'écrit $\pm 2^p5^q$ pour $p,q\in\mathbb N$. Réciproquement, soit $x=\frac{\pm 2^p5^q}{10^k}$ et montrons que $x$ est inversible dans $\mathbb D$. Posons $y=\frac{\pm 10^k}{2^p5^q}$. Il suffit de vérifier que $y$ est élément de $\mathbb D$. Mais on peut aussi écrire
		$$y=\frac{\pm 10^k 2^q 5^p}{2^{p+q}5^{p+q}}=\frac{\pm 10^k 2^q5^p}{10^{p+q}}\in\mathbb D.$$
		Ainsi, les inversibles de $(\mathbb D,+,\times)$ sont les éléments $\frac{\pm 2^p5^q}{10^k}$, avec $p,q,k\in\mathbb N$.
	\end{corrige}
	
	\begin{exo}[{Centre d'un anneau}]
		Soit $A$ un anneau. On appelle centre de $A$ et l'on note $C(A)$ l'ensemble des éléments $a\in A$ tels que, pour tout $b\in A$, $ab=ba.$ Démontrer que $C(A)$ est un sous-anneau de $A$.
	\end{exo}
	\begin{corrige}
		Il suffit de vérifier le théorème de caractérisation des sous-anneaux.
		
		\begin{itemize}
			
			\item  $1\in C(A)$, puisque $1a=a1=a$ pour tout $a\in A.$
			
			\item  Soit $a,a'\in C(A)$ et soit $b\in A$. Alors 
			$$(a-a')b=ab-a'b=ba-ba'=b(a-a')$$ 
			(on a utilisé deux fois la distributivité) et donc $a-a'\in C(A).$
			
			\item  Soit $a,a'\in C(A)$ et soit $b\in A$. Alors 
			$$(aa')b=a(a'b)=a(ba')=(ab)a'=(ba)a'=b(aa')$$
			(on a utilisé plusieurs fois l'associativité) et donc $aa'\in C(A).$
			
		\end{itemize}
	\end{corrige}
	
	\begin{exo}[{Anneau des fonctions de $\mathbb R$ dans $\mathbb R$}]
		L'anneau des fonctions de $\mathbb R$ dans $\mathbb R$ est-il intègre?
	\end{exo}
	\begin{corrige}
		Non, cet anneau n'est pas intègre. Il suffit de trouver deux fonctions non identiquement nulles dont le produit est nulle. Par exemple, si $f=1$ sur $]-\infty,0]$ et $0$ sur $]0,+\infty[$, et si $g=0$ sur $]-\infty,0]$ et $1$ sur $]0,+\infty[$, alors ni $f$ ni $g$ n'est nulle et pourtant $fg=0.$
	\end{corrige}
	
	\begin{exo}[{Éléments nilpotents}]
		Un élément $x$ d'un anneau $A$ est dit nilpotent s'il existe un entier $n\geq 1$ tel que $x^n=0$.
		On suppose que $A$ est commutatif, et on fixe $x,y$ deux éléments nilpotents.
		
		\begin{enumerate}
			
			\item  Montrer que $xy$ est nilpotent.
			
			\item  Montrer que $x+y$ est nilpotent.
			
			\item  Montrer que $1_A-x$ est inversible.
			
			\item  Dans cette question, on ne suppose plus que $A$ est commutatif. Soit $u,v\in A$ tels que $uv$ est nilpotent. Montrer que $vu$ est nilpotent.
			
		\end{enumerate}
	\end{exo}
	\begin{corrige}
		Soient $n,m$ tels que $x^n=0$ et $y^m=0$.
		
		\begin{enumerate}
			
			\item  Puisque $x$ et $y$ commutent, on a $(xy)^n=x^ny^n=0\times y^n=0$.
			
			\item  Remarquons d'abord que pour $p\geq n$, on a $x^p=x^{p-n}x^n=0$.
			D'après la formule du binôme, $(x+y)^{n+m}=\sum_{k=0}^{n+m}\binom{n+m}{k}x^ky^{n+m-k}$.
			Mais, pour $k\geq n$, $x^k=0\implies x^ky^{n+m-k}=0$. D'autre part, pour $k<n$,
			on a $n+m-k\geq m$ et donc $y^{n+m-k}=0\implies x^k y^{n+m-k}=0$. Ainsi, $(x+y)^{n+m}=0$. On pourrait
			même se contenter de prendre la puissance $n+m-1$.
			
			\item  L'idée est d'utiliser l'identité remarquable (toujours valable dans un anneau)
			$$1-x^p=(1-x)(1+x+\dots+x^{p-1}).$$
			Si on l'applique pour $p=n$, alors on obtient
			$$1=(1-x)(1+x+\dots+x^{n-1})$$
			ce qui implique que $1-x$ est inversible d'inverse $1+x+\dots+x^{n-1}$.
			
			\item  Soit $n\geq 1$ tel que $(uv)^n=0$. Alors
			$$(vu)^{n+1}=v(uv)^nu=v\times0\times u=0.$$
			Ainsi, $vu$ est nilpotent.
			
		\end{enumerate}
	\end{corrige}
	
	\begin{exo}[{Anneau de Boole}]
		On dit qu'un anneau $A$ est un anneau de Boole si, pour tout $x\in A$, $x^2=x$. On fixe $A$ un tel anneau.
		
		\begin{enumerate}
			
			\item  Démontrer que, pour tout $x\in A$, $x=-x$.
			
			\item  Montrer que $A$ est commutatif.
			
			\item  On suppose que $A$ est intègre. Montrer que $A$ contient exactement deux éléments.
			
		\end{enumerate}
	\end{exo}
	\begin{corrige}
		\begin{enumerate}
			
			\item  On applique la propriété à l'élément $x+x$. Il vient 
			$$x+x=(x+x)^2=x^2+x^2+x^2+x^2=x+x+x+x.$$
			Après simplification, on trouve $x+x=0$, soit $x=-x$.
			
			\item  Soient $x,y\in A$. On doit prouver $xy=yx$. Appliquons la propriété à l'élément $x+y$. On a
			$$(x+y)=(x+y)^2=x^2+y^2+xy+yx=x+y+xy+yx.$$
			Après simplification, on trouve $xy+yx=0$ soit $xy=-yx$, soit $xy=yx$ en appliquant le résultat
			de la question précédente.
			
			\item  Soit $x\in A.$ On sait que $x^2-x=0$. Mais $x^2-x=x(x-1).$ Par intégrité de $A,$ on trouve $x=0$ ou $x=1.$ Ainsi, $A$ possède exactement deux éléments, $0$ et $1.$
			
		\end{enumerate}
	\end{corrige}
	
	\subsection*{Corps}
	
	\begin{exo}[{Q(i)}]
		Montrer que $\mathbb Q(i)=\{a+ib:\ a,b\in\mathbb Q\}$ est un corps.
	\end{exo}
	\begin{corrige}
		On va commencer par démontrer que $\mathbb Q(i)$ est un sous-anneau de $\mathbb C.$ Pour cela, on remarque que 
		
		\begin{itemize}
			
			\item  $1\in\mathbb Q(i)$
			
			\item  si $z=a+ib$ et $z'=a'+ib'\in\mathbb Q(i),$ alors
			$$z-z'=(a-a')+i(b-b')\in\mathbb Q(i)$$
			et 
			$$zz'=(a+ib)(a'+ib')=(aa'-bb')+i(ab'+a'b)\in\mathbb Q(i).$$
			
		\end{itemize}
		
		Ensuite, soit $z=a+ib\in\mathbb Q(i),$ $z\neq 0.$ Alors 
		$$\frac1z=\frac{a-ib}{a^2+b^2}=\frac{a}{a^2+b^2}+i\frac{-b}{a^2+b^2}\in\mathbb Q(i).$$
		Ainsi, $1/z$ est dans $\mathbb Q(i)$ et donc tout élément non nul de $\mathbb Q(i)$ admet un inverse dans $\mathbb Q(i)$. Ceci achève de prouver que $\mathbb Q(i)$ est un corps.
	\end{corrige}
	
	\begin{exo}[{$\mathbb Q[\sqrt d]$}]
		Soit $d\in\mathbb N$ tel que $\sqrt d\notin \mathbb Q$. On note 
		$$\mathbb Q[\sqrt d]=\{a+b\sqrt d;\ (a,b)\in\mathbb Q ^2\}.$$
		Démontrer que $(\mathbb Q[\sqrt d],+,\times)$ est un corps.
	\end{exo}
	\begin{corrige}
		On va démontrer qu'il s'agit d'un sous-corps de $(\mathbb R,+,\times)$. Remarquons d'abord qu'il est bien contenu dans $\mathbb R$ et que $0,1\in \mathbb Q[\sqrt d]$. Soient $x,y\in\mathbb Q[\sqrt d]$. 
		On les écrit $x=a+b\sqrt d$ et $y=a'+b'\sqrt d$. Alors :
		\begin{eqnarray*}
			x-y&=&(a-a')+(b-b')\sqrt d\\
			xy&=&(aa'+dbb')+(ab'+a'b)\sqrt d
		\end{eqnarray*}
		ce qui prouve que $x-y$ et $xy\in\mathbb Q[\sqrt d]$. D'autre part, si $x\neq 0$, alors 
		$$\frac 1x=\frac{1}{a+b\sqrt d}=\frac{a-b\sqrt d}{a^2-b^2d}=\frac{a}{a^2-b^2 d}-\frac{b}{a^2-b^2 d}\sqrt d\in\mathbb Q[\sqrt d]$$
		et donc $\frac 1x\in \mathbb Q[\sqrt d]$. Remarquons qu'il était possible de multiplier par la quantité conjuguée qui est non-nulle car $\sqrt d\notin \mathbb Q$.  Finalement, on a bien prouvé que $(\mathbb Q[\sqrt d],+,\times)$ est un sous-corps de $(\mathbb R,+,\times)$.
	\end{corrige}
	
	\begin{exo}[{Un corps de matrices}]
		Soit $\mathcal C=\left\{\begin{pmatrix}a&-b\\b&a\end{pmatrix}:\ (a,b)\in\mathbb R^2\right\}.$ 
		
		\begin{enumerate}
			
			\item  Démontrer que $(\mathcal C,+,\times)$ est un corps.
			
			\item  Démontrer que $\mathcal C$ est isomorphe à $\mathbb C$.
			
		\end{enumerate}
	\end{exo}
	\begin{corrige}
		\begin{enumerate}
			
			\item  On va commencer par démontrer que $\mathcal C$ est un sous-anneau de $\mathcal M_2(\mathbb R)$. Pour cela, on remarque que 
			
			\begin{itemize}
				
				\item  la matrice $I_2$ est bien élément de $\mathcal C$ (prendre $a=1$, $b=0$).
				
				\item  Soit $A,A'$ deux matrices de $\mathcal C$, écrivons les 
				$$A=\begin{pmatrix}a&-b\\b&a\end{pmatrix}, A'=\begin{pmatrix}a'&-b'\\b'&a'\end{pmatrix}.$$
				Alors 
				$$A+A'=\begin{pmatrix}(a+a')&-(b+b')\\(b+b')&(a+a')\end{pmatrix}\in\mathcal C$$
				et 
				$$A\times A'=\begin{pmatrix}aa'-bb'&-(ab'+a'b)\\ab'+a'b&aa'-bb'\end{pmatrix}\in\mathcal C.$$
				
			\end{itemize}
			
			De plus, $(\mathcal C,+,\times)$ est bien un anneau commutatif car le calcul précédent montre facilement que $A\times A'=A'\times A.$ Pour conclure, il faut encore démontrer que si on choisit $A\in\mathcal C$, $A\neq 0,$ alors $A$ est inversible et $A^{-1}\in\mathcal C$. Notons toujours 
			$$A=\begin{pmatrix}a&-b\\b&a\end{pmatrix}.$$
			Il est facile de voir que $A$ est inversible (par exemple par ce que son déterminant est $a^2+b^2\neq 0$). De plus, l'inverse de $A$ est 
			$$A^{-1}=\frac{1}{a^2+b^2}\begin{pmatrix}
				a&b\\
				-b&a
			\end{pmatrix}$$
			et cette matrice est bien dans $\mathcal C$ (choisir $a'=a/(a^2+b^2)$ et $b'=-b/(a^2+b^2)$). Ainsi, on a bien démontré que $\mathcal C$ est un corps.
			
			\item  Considérons $\phi(A)=a+ib$, où on a toujours noté 
			$$A=\begin{pmatrix}a&-b\\b&a\end{pmatrix}.$$
			On a bien $\phi(I_2)=1$. De plus, soit
			$$A=\begin{pmatrix}a&-b\\b&a\end{pmatrix}, A'=\begin{pmatrix}a'&-b'\\b'&a'\end{pmatrix}.$$
			Alors d'après le calcul de $A+A'$ fait ci-dessus,
			$$\phi(A+A')=(a+a')+i(b+b')=\phi(A)+\phi(A')$$
			et de même, d'après le calcul de $A\times A',$
			$$\phi(AA')=(aa-bb')+i(ab'+a'b)=(a+ib)\times(a'+ib')=\phi(A)\phi(A').$$
			Ainsi, $\phi$ est bien un morphisme d'anneaux (ou de corps). C'est même un isomorphisme : si $A\in \ker(\phi)$, $\phi(A)=0\iff A=0$ et donc 
			$\ker(\phi)=\{0\}$. De plus, $\textrm{Im}(\phi)=\mathbb C$, 
			puisque si $z=a+ib\in\mathbb C$, $z=\phi(A)$ où $A=\begin{pmatrix}a&-b\\b&a\end{pmatrix}.$ Ainsi, les corps $\mathcal C$ et $\mathbb C$ sont isomorphes.
			
		\end{enumerate}
	\end{corrige}
	
	\begin{exo}[{Anneau intègre fini}]
		Soit $A$ un anneau intègre (commutatif) fini. Démontrer que $A$ est un corps.
	\end{exo}
	\begin{corrige}
		Fixons $a\in A$ non nul et considérons le morphisme de groupes $(A,+)\to (A,+),\ x\mapsto ax$. Alors ce morphisme de groupes est injectif, car son noyau est réduit à $\{0_A\}$ puisque $A$ est intègre. Puisque $A$ est fini, ce morphisme est nécessairement bijectif, et donc il existe $x\in A$ tel que $ax=1_A$. Par commutativité de $A$, on a aussi $xa=1_A$ et donc $a$ admet un inverse : $A$ est un corps.\\
		
		Remarquons que l'on peut se passer de l'hypothèse que $A$ est commutatif, par exemple en faisant le même raisonnement avec $x\mapsto xa$, et en prouvant que l'inverse à droite et l'inverse à gauche coïncident.
	\end{corrige}
	
	\subsection*{Exercices complémentaires}
	
	\begin{exo}[{Caractérisation des sous-groupes}]\ \textcolor{red!70!black}{$\star$}
		Soit $(G,\cdot)$ un groupe et $H\subset G$. Montrer que $H$ est un sous-groupe de $G$ si et seulement si $H\neq\emptyset$ et $\forall (x,y)\in H^2,\ xy^{-1}\in H$. Montrer que l'intersection de deux sous-groupes est un sous-groupe, et que la réunion de deux sous-groupes est un sous-groupe si et seulement si l'un contient l'autre.
	\end{exo}
	\begin{corrige}
		($\Rightarrow$) Un sous-groupe contient $e$ et est stable par inverse et produit. ($\Leftarrow$) $H\neq\emptyset$ : soit $x\in H$, alors $e=xx^{-1}\in H$ ; pour $y\in H$, $y^{-1}=ey^{-1}\in H$ ; pour $x,y\in H$, $xy=x(y^{-1})^{-1}\in H$.
		
		L'intersection de deux sous-groupes contient $e$ et est stable par $(x,y)\mapsto xy^{-1}$. Réunion : si $H\subset K$ alors $H\cup K=K$ est un sous-groupe. Réciproquement, supposons $H\cup K$ sous-groupe avec $H\not\subset K$ et $K\not\subset H$ : soit $h\in H\setminus K$ et $k\in K\setminus H$. Alors $hk\in H\cup K$. Si $hk\in H$, $k=h^{-1}(hk)\in H$ : absurde ; si $hk\in K$, $h=(hk)k^{-1}\in K$ : absurde.
	\end{corrige}
	
	\begin{exo}[{Le groupe $\mathcal S_3$}]\ \textcolor{red!70!black}{$\star$}
		\begin{enumerate}
			\item Dresser la liste des éléments de $\mathcal S_3$ et donner l'ordre de chacun.
			\item Le groupe $\mathcal S_3$ est-il commutatif ?
			\item Déterminer tous les sous-groupes de $\mathcal S_3$.
		\end{enumerate}
	\end{exo}
	\begin{corrige}
		\begin{enumerate}
			\item $\mathcal S_3$ contient l'identité (ordre $1$), les trois transpositions $(1\,2),(1\,3),(2\,3)$ (ordre $2$) et les deux $3$-cycles $(1\,2\,3),(1\,3\,2)$ (ordre $3$).
			\item Non : $(1\,2)\circ(1\,3)=(1\,3\,2)$ alors que $(1\,3)\circ(1\,2)=(1\,2\,3)$.
			\item L'ordre d'un sous-groupe divise $6$. Ordre $1$ : $\{\id\}$. Ordre $2$ : engendré par un élément d'ordre $2$, soit $\langle(1\,2)\rangle$, $\langle(1\,3)\rangle$, $\langle(2\,3)\rangle$. Ordre $3$ : engendré par un élément d'ordre $3$, soit $\{\id,(1\,2\,3),(1\,3\,2)\}$ (les deux $3$-cycles engendrent le même sous-groupe). Ordre $6$ : $\mathcal S_3$. Au total, $6$ sous-groupes.
		\end{enumerate}
	\end{corrige}
	
	\begin{exo}[{Morphismes}]\ \textcolor{red!70!black}{$\star$}
		Montrer que $\exp:(\R,+)\to(\R_+^*,\times)$ est un isomorphisme de groupes. Déterminer le noyau et l'image du morphisme $z\mapsto |z|$ de $(\C^*,\times)$ dans $(\R^*,\times)$ et celui de $x\mapsto e^{ix}$ de $(\R,+)$ dans $(\C^*,\times)$.
	\end{exo}
	\begin{corrige}
		$\exp(x+y)=\exp x\exp y$ : c'est un morphisme ; il est bijectif de $\R$ sur $\R_+^*$ de réciproque $\ln$, donc c'est un isomorphisme.
		
		$z\mapsto|z|$ est un morphisme car $|zw|=|z||w|$ ; son noyau est le groupe $\mathbb U$ des complexes de module $1$ et son image est $\R_+^*$. $x\mapsto e^{ix}$ est un morphisme car $e^{i(x+y)}=e^{ix}e^{iy}$ ; son noyau est $2\pi\Z$ et son image est $\mathbb U$.
	\end{corrige}
	
	\begin{exo}[{Anneaux $\Z[i]$ et $\Z/n\Z$}]\ \textcolor{red!70!black}{$\star$}
		\begin{enumerate}
			\item Montrer que $\Z[i]=\{a+ib\mid a,b\in\Z\}$ est un sous-anneau de $\C$. Déterminer ses éléments inversibles.
			\item Montrer que $\Z/n\Z$ est un corps si et seulement si $n$ est premier.
			\item Déterminer l'inverse de $7$ dans $\Z/26\Z$.
		\end{enumerate}
	\end{exo}
	\begin{corrige}
		\begin{enumerate}
			\item $0,1\in\Z[i]$, et $\Z[i]$ est stable par différence et par produit : $(a+ib)(c+id)=(ac-bd)+i(ad+bc)$. Si $z=a+ib$ est inversible, il existe $w\in\Z[i]$ avec $zw=1$, donc $|z|^2|w|^2=1$ avec $|z|^2=a^2+b^2$ et $|w|^2$ entiers positifs : $a^2+b^2=1$, d'où $z\in\{1,-1,i,-i\}$. Réciproquement ces quatre éléments sont inversibles ($i\cdot(-i)=1$). 
			\item Si $n$ est premier, tout $\bar k\neq0$ est premier avec $n$ donc inversible (Bézout) : c'est un corps. Si $n=ab$ avec $1<a,b<n$, alors $\bar a\bar b=0$ avec $\bar a,\bar b\neq0$ : diviseurs de zéro, donc pas un corps.
			\item $7\times15=105=4\times26+1$ : l'inverse de $\bar7$ est $\overline{15}$.
		\end{enumerate}
	\end{corrige}
	\section{Polynômes et fractions rationnelles}
	
	\subsection*{Pour réviser}
	
	\begin{exo}[{Expression du reste}]
		Soit $P\in \mathbb K[X]$, soit $a,b\in\mathbb K$ avec $a\neq b$.
		
		\begin{enumerate}
			
			\item  Soit $R$ le reste de la division euclidienne de $P$ par $(X-a)(X-b)$. Exprimer $R$ en fonction de $P(a)$ et de $P(b)$.
			
			\item  Soit $R$ le reste de la division euclidienne de $P$ par $(X-a)^2$. Exprimer $R$ en fonction de $P(a)$ et de $P'(a)$.
			
		\end{enumerate}
	\end{exo}
	\begin{corrige}
		Dans les deux cas, on remarquera que  $R$ est de degré au plus $1$ et s'écrit donc $R(X)=\alpha X+\beta$. 
		
		\begin{enumerate}
			
			\item  Évaluons la relation 
			$$P(X)=(X-a)(X-b)Q(X)+\alpha X+\beta$$
			en $a$ et en $b$. On trouve le système 
			$$\left\{
			\begin{array}{rcl}
				\alpha a +\beta &=&P(a)\\
				\alpha b +\beta &=&P(b).
			\end{array}\right.$$
			On en déduit alors facilement que 
			$$\alpha=\frac{P(a)-P(b)}{a-b}\textrm{ et }\beta=\frac{aP(b)-bP(a)}{a-b}.$$
			
			\item  Évaluons la relation
			$$P(X)=(X-a)^2Q(X)+\alpha X+\beta$$
			au point $a$. On trouve $P(a)=a\alpha+\beta$. Dérivons maintenant la relation précédente :
			$$P'(X)=2(X-a)Q(X)+(X-a)^2Q'(X)+\alpha.$$
			On évalue à nouveau en $a$ et on trouve que 
			$$\alpha=P'(a).$$
			En revenant à la première équation, on en déduit que $\beta=P(a)-aP'(a).$
			
		\end{enumerate}
	\end{corrige}
	
	\begin{exo}[{Racines rationnelles}]
		Soit $P(X)=a_nX^n+\dots+a_0$ un polynôme à coefficients dans $\mathbb Z$, avec
		$a_n\neq 0$ et $a_0\neq 0$. On suppose que $P$ admet une racine rationnelle $p/q$
		avec $p\wedge q=1$. Démontrer que $p|a_0$ et que $q|a_n$. Le polynôme $P(X)=X^5-X^2+1$
		admet-il des racines dans $\mathbb Q$?
	\end{exo}
	\begin{corrige}
		On écrit que $P(p/q)=0$ et on met tout au même dénominateur en multipliant par $q^n$. On trouve
		$$a_np^n +a_{n-1}p^{n-1}q+\dots+a_1pq^{n-1}+a_0q^n=0.$$
		On commence par isoler $a_0q^n$ et on trouve que
		$$p(a_np^{n-1}+a_{n-1}p^{n-2}q+\dots+a_1q^{n-1})=-a_0q^n.$$
		En particulier, $p|a_0q^n$. Puisque $p\wedge q=1$, on en déduit que $p|a_0$.
		De même, en isolant $a_np^n$, on trouve
		$$q(a_{n-1}p^{n-1}+\dots+a_0q^{n-1})=-a_np^n,$$
		soit $q|a_n p^n$, soit, puisque $p\wedge q=1$, $q|a_n$.\\
		
		Par conséquent, si le polynôme 
		$X^5-X^2+1$ admet une racine rationnelle $p/q$, alors $p|1$ et $q|1$, et
		donc $p=\pm 1$ et $q=\pm 1$. Autrement dit, les seules racines rationnelles possibles
		sont $1$ et $-1$. Or, elles ne sont pas racines de $P$. Donc $P$ n'admet pas 
		de racines rationnelles.
	\end{corrige}
	
	\begin{exo}[{Somme des racines}]
		Soit $n\in\mathbb N^*$ et $P\in\mathbb C_n[X]$. On note, pour $p<n$, $u_p$ la somme des racines de $P^{(p)}$. Démontrer que
		$u_0,\dots,u_{n-1}$ forme une progression arithmétique.
	\end{exo}
	\begin{corrige}
		Écrivons 
		$$P(X)=a_nX^n+a_{n-1}X^{n-1}+\dots+a_0.$$
		Alors, par les relations coefficients/racines, on sait que la somme des racines de $P$ vaut
		$u_0=-a_{n-1}/a_n$. Plus généralement, on a 
		$$P^{(p)}(X)=n(n-1)\dots(n-p+1)a_nX^{n-p}+(n-1)\dots(n-p)a_{n-1}X^{n-p-1}+\dots+p!a_p,$$
		de sorte que
		$$u_p=-\frac{(n-1)\dots(n-p)}{n(n-1)\dots(n-p+1)}\times\frac{a_{n-1}}{a_n}=-\frac{(n-p)}{n}\times\frac{a_{n-1}}{a_n}.$$
		On a donc
		$$u_{p+1}-u_p=-\left(\frac{n-p-1}{n}-\frac{n-p}n\right)\times\frac{a_{n-1}}{a_n}=\frac{a_{n-1}}{na_n}.$$
		On obtient bien une progression arithmétique de raison $\frac{a_{n-1}}{na_n}$.
	\end{corrige}
	
	\begin{exo}[{Décomposer!}]
		Décomposer en produits d'irréductibles de $\mathbb R[X]$ les polynômes suivants :
		$$\begin{array}{lllll}\mathbf{1.}\ \ X^4+1&\quad&\mathbf{2.}\ X^8-1&\quad&\mathbf{3.}\ (X^2-X+1)^2+1
		\end{array}$$
	\end{exo}
	\begin{corrige}
		\begin{enumerate}
			
			\item  On commence par chercher les racines complexes pour factoriser dans $\mathbb C[X]$, puis on regroupe les
			racines complexes conjuguées. 
			\begin{eqnarray*}
				X^4+1&=&(X-e^{i\pi/4})(X-e^{3i\pi/4})(X-e^{5i\pi/4})(X-e^{7i\pi/4})\\
				&=&\big((X-e^{i\pi/4})(X-e^{7i\pi/4})\big)\big((X-e^{3i\pi/4})(X-e^{5i\pi/4})\big)\\
				&=&(X^2-\sqrt 2X+1)(X^2+\sqrt 2X+1).
			\end{eqnarray*}
			Les deux polynômes de degré 2 que l'on obtient n'ont pas de racines réelles,
			ils sont donc irréductibles dans $\mathbb R[X]$.
			
			\item  On commence par utiliser une identité remarquable, puis la réponse à la question précédente :
			\begin{eqnarray*}
				X^8-1&=&(X^4-1)(X^4+1)\\
				&=&(X^2-1)(X^2+1)(X^2-\sqrt 2X+1)(X^2+\sqrt 2X+1)\\
				&=&(X-1)(X+1)(X^2+1)(X^2-\sqrt 2X+1)(X^2+\sqrt 2X+1).
			\end{eqnarray*}
			
			\item  On commence par factoriser le polynôme dans $\mathbb C[X]$ en remarquant qu'il s'agit alors d'une différence de
			deux carrés :
			$$(X^2-X+1)^2+1=(X^2-X+1)^2-i^2=(X^2-X+1-i)(X^2-X+1+i).$$
			On factorise alors chacun des polynômes de degré 2 dans $\mathbb C$, par exemple en calculant leur
			discriminant ou en remarquant que $i$ (resp. $-i$) sont des racines évidentes. On trouve :
			$$(X^2-X+1)^2+1=(X+i)(X-1-i)(X-i)(X-1+i).$$
			En regroupant les termes conjugués, on trouve finalement :
			$$(X^2-X+1)^2+1=(X^2+1)(X^2-2X+2).$$
			
		\end{enumerate}
	\end{corrige}
	
	\begin{exo}[{pgcd dans deux corps différents}]
		Soit $\mathbb K\subset\mathbb L$ deux corps. On considère $P,Q\in\mathbb K[X]$.
		
		\begin{enumerate}
			
			\item  Démontrer que $P$ et $Q$ sont premiers entre eux dans $\mathbb K[X]$ si et seulement s'ils sont premiers entre eux dans $\mathbb L[X]$.
			
			\item  Plus généralement, démontrer que le pgcd de $P$ et de $Q$ vus comme polynômes de $\mathbb K[X]$ est égal au pgcd de $P$ et de $Q$ vus comme polynôme de $\mathbb L[X]$.
			
		\end{enumerate}
	\end{exo}
	\begin{corrige}
		\begin{enumerate}
			
			\item  Si $P$ et $Q$ sont premiers entre eux dans $\mathbb K[X]$, alors d'après le théorème de Bézout, il existe $U,V\in\mathbb K[X]\subset \mathbb L[X]$ tels que $UP+VQ=1$. Mais alors, d'après l'autre partie du théorème de Bézout, on a aussi que $P$ et $Q$ sont premiers entre eux dans $\mathbb L[X]$.\\
			
			Supposons maintenant que $P$ et $Q$ sont premiers entre eux dans $\mathbb L[X]$. S'ils n'étaient pas premiers entre eux dans $\mathbb K[X]$, il existerait un polynôme $D\in\mathbb K[X]$ avec $\deg(D)\geq 1$ divisant à la fois $P$ et $Q$. Mais comme $D\in\mathbb L[X]$, ceci contredirait le fait que $P$ et $Q$ sont premiers entre eux dans $\mathbb L[X]$.
			
			\item  Notons $P\wedge_{\mathbb K}Q$ le pgcd dans $\mathbb K[X]$ et $P\wedge_{\mathbb L}Q$ le pgcd dans $\mathbb L[X]$. On peut procéder par récurrence sur $\deg(P)+\deg(Q)$. Si les deux polynômes sont constants, c'est clair! (leur pgcd est 1). Sinon, notons $A$ et $R$ respectivement le quotient et le reste dans la division euclidienne de $P$ par $Q$ dans $\mathbb K[X]$ (on peut toujours supposer $\deg(P)\geq \deg(Q)$). Mais alors, par unicité dans la division euclidienne, $A$ et $R$ sont aussi le quotient et le reste dans la division euclidienne de $P$ par $Q$ dans $\mathbb L[X]$. On a donc, d'après l'algorithme d'Euclide, 
			$$P\wedge_{\mathbb L}Q=Q\wedge_{\mathbb L}R$$
			et 
			$$P\wedge_{\mathbb K}Q=Q\wedge_{\mathbb K}R.$$
			Puisque $\deg(R)+\deg(Q)<\deg(P)+\deg(Q)$, l'hypothèse de récurrence nous donne que
			$$Q\wedge_{\mathbb K}R=Q\wedge_{\mathbb L}R$$
			et donc finalement que
			$$P\wedge_{\mathbb K}Q=P\wedge_{\mathbb L}Q.$$
			
		\end{enumerate}
	\end{corrige}
	
	\begin{exo}[{Un calcul de somme}]
		\begin{enumerate}
			
			\item  Décomposer en éléments simples la fraction rationnelle $\displaystyle\frac{1}{X(X+1)(X+2)}$.
			
			\item  En déduire la limite de la suite $(S_n)$ suivante : $\displaystyle S_n=\sum_{k=1}^n \frac{1}{k(k+1)(k+2)}$.
			
		\end{enumerate}
	\end{exo}
	\begin{corrige}
		\begin{enumerate}
			
			\item  On sait que la fraction rationnelle admet une décomposition en éléments simples de la forme
			$$\frac{a}{X}+\frac{b}{X+1}+\frac{c}{X+2}.$$
			Par les techniques usuelles (identification, multiplication par $X$ et faire $X=0$,...), on trouve
			$$\frac{1}{X(X+1)(X+2)}=\frac{1/2}X-\frac{1}{X+1}+\frac{1/2}{X+2}.$$
			
			\item  Utilisant la décomposition en éléments simples précédente, il vient 
			\begin{eqnarray*}
				S_n&=&\sum_{k=1}^n \frac{1/2}{k}-\sum_{k=1}^n\frac{1}{k+1}+\sum_{k=1}^n \frac{1/2}{k+2}\\
				&=&\sum_{k=1}^n \frac{1/2}{k}-\sum_{k=2}^{n+1}\frac{1}{k}+\sum_{k=3}^{n+2} \frac{1/2}{k}\\
				&=&\frac12+\frac14-\frac{1}2-\frac{1}{n+1}+\frac{1/2}{n+1}+\frac{1/2}{n+2}.
			\end{eqnarray*}
			On en déduit immédiatement que $(S_n)$ converge vers $\frac14$.
			
		\end{enumerate}
	\end{corrige}
	
	\begin{exo}[{Un calcul de somme}]
		Soit $P\in\mathbb R[X]$ un polynôme de degré $n\geq 1$ possédant $n$ racines distinctes $x_1,\dots,x_n$ non-nulles.
		
		\begin{enumerate}
			
			\item  Décomposer en éléments simples la fraction rationnelle $\displaystyle \frac1{XP(X)}$.
			
			\item  En déduire que $\displaystyle\sum_{k=1}^n \frac{1}{x_k P'(x_k)}=\frac{-1}{P(0)}$.
			
		\end{enumerate}
	\end{exo}
	\begin{corrige}
		\begin{enumerate}
			
			\item  On écrit la forme à priori de la décomposition en éléments simples qui est ici
			$$\frac{1}{XP(X)}=\frac{\alpha_0}{X}+\sum_{k=1}^n \frac{\alpha_k}{X-x_k}.$$
			On calcule $\alpha_0$ en multipliant tout par $X$ et en faisant tendre $X$ vers $0,$ et on trouve
			$\alpha_0=\frac{1}{P(0)}$. Pour $k\geq 1$, on multiplie tout par $X-x_k$ et on fait tendre $X$ vers $x_k$.
			On trouve cette fois
			$$\alpha_k=\lim_{x\to x_k} \frac{x-x_k}{x_k P(x)}=\lim_{x\to x_k} \frac{x-x_k}{x_k (P(x)-P(x_k))}=\frac{1}{x_k P'(x_k)}.$$
			La décomposition en éléments simples est donc 
			$$\frac{1}{XP(X)}=\sum_{k=1}^n \frac{1}{x_k P'(x_k)}\times\frac{1}{X-x_k}+\frac{1}{P(0)}\times\frac{1}{X}.$$
			
			\item  On va multiplier par $X$ cette fraction rationnelle, et on va étudier sa limite en $+\infty$. On trouve d'une part
			$$\lim_{x\to+\infty}\frac x{xP(x)}=0$$
			et d'autre part
			$$\lim_{x\to +\infty}\sum_{k=1}^n\frac{1}{x_k P'(x_k)}\times\frac{x}{x-x_k}+\frac{1}{P(0)}\times\frac{x}{x}=\sum_{k=1}^n \frac{1}{x_k P'(x_k)}+\frac1{P(0)}.$$
			Identifiant ces deux égalités, on trouve le résultat voulu!
			
		\end{enumerate}
	\end{corrige}
	
	\begin{exo}[{Enveloppe convexe des zéros}]
		Soit $P\in\mathbb C_n[X]$ admettant $n$ racines simples $\alpha_1,\dots,\alpha_n$.
		Soient $A_1,\dots,A_n$ les points du plan complexe d'affixe respectives $\alpha_1,\dots,\alpha_n$.
		
		\begin{enumerate}
			
			\item  Décomposer la fraction rationnelle $P'/P$ en éléments simples.
			
			\item  Soit $\beta$ une racine de $P'$, et soit $B$ son image dans le plan complexe. Déduire de la question précédente que 
			$$\sum_{j=1}^n \frac{1}{\beta-\alpha_j}=0.$$
			
			\item  En déduire que $B$ est un barycentre de la famille de points $(A_1,\dots,A_n)$, avec des coefficients positifs.
			Interpréter géométriquement cette propriété.
			
		\end{enumerate}
	\end{exo}
	\begin{corrige}
		\begin{enumerate}
			
			\item   On va étudier séparément les parties polaires relatives à chaque racine. 
			On peut factoriser $P$ en $P(X)=(X-\alpha_j) Q(X)$, soit en dérivant 
			$P'(X)=Q(X)+(X-\alpha_j) Q'(X)$. On trouve alors
			$$\frac{P'(X)}{P(X)}=\frac{1}{X-\alpha_j}+\frac{Q'(X)}{Q(X)}.$$
			Or, $\alpha_j$ n'est pas racine de $Q$, donc $Q'/Q$ n'admet pas $\alpha_j$ pour pôle et $\frac{1}{X-\alpha_j}$ est la partie
			polaire de $P'/P$ relative à $a$. En résumé, la décomposition en éléments simples recherchée est
			$$\frac{P'}{P}=\sum_{j=1}^p \frac{1}{X-\alpha_j}.$$
			
			\item  Il suffit d'évaluer l'équation précédente en $\beta$.
			
			\item  On multiplie par la quantité conjuguée et on trouve
			$$\sum_{j=1}^n \frac{\bar\beta-\bar\alpha_j}{|\beta-\alpha_j|^2}=0.$$
			Prenant le conjugué de cette expression, on trouve :
			$$\left(\sum_{j=1}^n \frac{1}{|\beta-\alpha_j|^2}\right)\beta=\sum_{j=1}^n \frac{1}{|\beta-\alpha_j|^2}\alpha_j,$$
			ce qui correspond bien au résultat souhaité. On vient donc de prouver que toute racine de $P'$ est dans l'enveloppe convexe des racines de $P$.
			Ce résultat s'appelle le théorème de Lucas, il est aussi valide si les racines de $P$ ne sont pas simples. La preuve
			est similaire, si ce n'est que la décomposition en éléments simples de $P'/P$ est plus difficile à obtenir.
			C'est un bon exercice!
			
		\end{enumerate}
	\end{corrige}
	
	\subsection*{Pour progresser}
	
	\begin{exo}[{Tout polynôme positif est somme de deux carrés}]
		Soit $P\in\mathbb R[X]$ non constant tel que $P(x)\geq 0$ pour tout réel $x$.
		
		\begin{enumerate}
			
			\item  Montrer que le coefficient dominant de $P$ est positif et que les racines réelles de $P$ sont de multiplicité paire.
			
			\item  Montrer qu'il existe un polynôme $C\in\mathbb C[X]$ tel que $P=C\overline{C}$.
			
			\item  En déduire qu'il existe $A$ et $B$ dans $\mathbb R[X]$ tels que $P=A^2+B^2$.
			
		\end{enumerate}
	\end{exo}
	\begin{corrige}
		\begin{enumerate}
			
			\item  Notons $P(x)=a_nx^n+\cdots$. Alors $\lim_{x\to+\infty}P(x)=\lim_{x\to+\infty}a_nx^n$ et cette limite vaut $+\infty$ si $a_n>0$ et $-\infty$ si $a_n<0$.
			Si la limite est $-\infty$, alors on ne peut pas avoir $P(x)\geq 0$ au voisinage de $+\infty$ et donc $a_n>0$. D'autre part, soit $a$ une racine de $P$ de multiplicité $m$.
			On sait que $P(x)\sim_{a}\lambda (x-a)^m$. Si $m$ est impair, on en déduit que $P$ s'annule en changeant de signe en $a$. Ceci contredit que $P(x)\geq 0$ pour tout $x$
			proche de $a$. Ainsi, les racines de $P$ sont de multiplicité paire.
			
			\item  On décompose $P$ en produits de facteurs d'irréductibles : on a
			$$P(X)=\lambda \prod_{i=1}^s (X-a_i)^{2m_i}\prod_{j=1}^t (X^2+c_j X+d_j)^{q_j}$$
			où les $a_i$ sont les racines réelles de $P$ (elles sont de multiplicité paire!), 
			où les $X^2+c_j X+d_j$ sont des polynômes du second degré à discriminant négatif. Les racines de ces polynômes
			sont des nombres complexes conjugués $\mu_j$, $\overline{\mu_j}$. En particulier, on a 
			$$X^2+c_jX+d_j=(X-\mu_j)\overline{(X-\mu_j)}.$$
			Posons donc
			$$C(X)=\sqrt{\lambda} \prod_{i=1}^s (X-a_i)^{m_i} \prod_{j=1}^t (X-\mu_j)^{q_j}.$$
			On a bien $P=C\overline C$. 
			
			\item  Décomposons $C$ en partie réelle et en partie imaginaire : $C(X)=A(X)+iB(X)$ avec $A,B\in\mathbb R[X]$.
			D'après la question précédente, on a $P=C\overline C=A^2+B^2$.
			
		\end{enumerate}
	\end{corrige}
	
	\begin{exo}[{Polynôme réciproque}]
		On dit qu'un polynôme $P\in\mathbb C[X]$ de degré $n$ est réciproque s'il s'écrit $P=a_nX^n+\dots+a_0$ avec
		$a_k=a_{n-k}$ pour tout $k$ dans $\{0,\dots,n\}$. 
		
		\begin{enumerate}
			
			\item  Soit $P\in\mathbb C[X]$ de degré $n$. Démontrer que $P$ est réciproque si et seulement si $P(X)=X^n P\left(\frac 1X\right)$.
			
			\item  Montrer qu'un produit de polynômes réciproques est réciproque.
			
			\item  On suppose que $P$ et $Q$ sont réciproques et que $Q|P$. Démontrer que $\frac PQ$ est réciproque.
			
			\item  Soit $P\in\mathbb C[X]$ un polynôme réciproque.
			
			\begin{enumerate}
				
				\item  Démontrer que si $\alpha$ est une racine de $P$, alors $\alpha\neq 0$ et $\alpha^{-1}$ est une racine de $P$. 
				
				\item  Démontrer que si $1$ est une racine de $P$, alors sa multiplicité est supérieure ou égale à $2$.
				
				\item  Démontrer que si le degré de $P$ est impair, alors $-1$ est racine de $P$.
				
				\item  Démontrer que si $P$ est de degré pair et si $-1$ est une racine de $P$, alors sa multiplicité est supérieure ou égale à $2$.
				
			\end{enumerate}
			
			
			\item  Démontrer que tout polynôme réciproque de $\mathbb C[X]$ de degré $2n$ se factorise en 
			$$P=a_{2n}(X^2+b_1X+1)\dots(X^2+b_n X+1).$$
			Que peut-on dire si le degré de $P$ est impair?
		\end{enumerate}
	\end{exo}
	\begin{corrige}
		\begin{enumerate}
			
			\item  Soit $P=a_nX^n+\dots+a_0$, alors
			$$X^n P\left(\frac 1X\right)=a_0X^n+\dots+a_n.$$
			Ainsi, si $P$ est réciproque, on a bien $X^n P(1/X)=P(X)$. Réciproquement, si $X^nP(1/X)=P(X)$, alors on a nécessairement
			$a_0=a_n$, $a_1=a_{n-1}$, etc... Donc $P$ est réciproque.
			
			\item  Soient $P$ et $Q$ réciproques, de degrés respectifs $n$ et $m$. Alors 
			$$X^n P(1/X)=P(X)\textrm{ et }X^m Q(1/X)=Q(X).$$
			On en déduit que 
			$$X^{n+m}(PQ)(1/X)=X^n P(1/X)X^m Q(1/X)=P(X)Q(X)=(PQ)(X).$$
			Ainsi, d'après la question précédente, $PQ$ est réciproque.
			
			\item  Le raisonnement est complètement identique, en utilisant le quotient au lieu du produit!
			
			\item  
			
			\begin{enumerate}
				
				\item  Puisque $P$ est réciproque, $a_0=a_n\neq 0$ et donc $P(0)=a_0\neq 0$. D'autre part, si $\alpha$ est racine de $P$, alors la relation
				$P(\alpha)=\alpha^n P(\alpha^{-1})$ prouve que $\alpha^{-1}$ est aussi racine de $P$.
				
				\item  Dérivons la relation de la première question. On trouve, pour tout $x\neq 0$, 
				$$P'(x)=nx^{n-1}P(1/x)-x^{n-2}P'(1/x).$$
				On évalue en 1, et on trouve
				$$P'(1)=-P'(1)$$ et donc $P'(1)=0$. 
				On en déduit que $1$ est racine au moins double.
				
				\item  On utilise encore le résultat de la première question, et on remarque que $P(-1)=-P(-1)$ puisque le degré
				de $P$ est impair. Donc $P(-1)=0$. 
				
				\item  On raisonne exactement comme deux questions plus haut.
				
			\end{enumerate}
			
			
			\item  On va procéder par récurrence sur $n$, le cas $n=1$ étant trivial. Supposons donc que le résultat a été démontré pour tout polynôme réciproque
			de degré $2n$, et prouvons-le pour un polynôme réciproque $P$ de degré $2n+2$. Soit $\alpha$ une racine de $P$. Alors, on sait que 
			$\alpha\neq 0$ et que $\alpha^{-1}$ est aussi racine de $P$. Si $\alpha\neq 1,-1$, $\alpha^{-1}\neq \alpha$ et on peut factoriser
			$P$ par $(X-\alpha)(X-\alpha^{-1})$. Or, il est facile de vérifier que $(X-\alpha)(X-\alpha^{-1})$ s'écrit $(X^2+b_{n+1}X+1)$. D'autre part,
			si $\alpha=1$ ou $\alpha=-1$, alors $\alpha$ est racine de multiplicité au moins deux, et on peut factoriser par $(X-\alpha)^2$. Un tel 
			polynôme s'écrit encore $(X^2+b_{n+1}X+1)$. Donc, dans tous les cas, en notant $Q=X^2+b_{n+1}X+1$, on a $Q|P$ et
			$P$, $Q$ réciproques. On en déduit que $\frac PQ$ est réciproque, de degré $2n$, donc par l'hypothèse de récurrence s'écrit
			$$\frac PQ=a_{2n+2}(X^2+b_1X+1)\dots(X^2+b_n X+1).$$
			On remultiplie par $Q$, et on a bien prouvé que le résultat est vrai au rang $n+1$.\\
			
			Si maintenant $P$ est réciproque de degré impair $2n+1$, alors $-1$ est racine de $P$ et $P$ se factorise par le polynôme réciproque $Q=X+1$. 
			Donc $\frac PQ$ est réciproque de degré pair $2n$, donc s'écrit $a_{2n+1}(X^2+b_1X+1)\dots(X^2+b_n X+1).$
			Ainsi, tout polynôme réciproque de degré impair $2n+1$ se factorise en 
			$$P=a_{2n+1}(X+1)(X^2+b_1X+1)\dots(X^2+b_n X+1).$$
			
		\end{enumerate}
	\end{corrige}
	
	\begin{exo}[{Irréductibilité par réduction}]
		Soit $p$ un nombre premier. Pour tout polynôme $P\in\mathbb Z[X]$, on note $\bar P$ son image (par réduction modulo $p$ de ses coefficients) dans $(\mathbb Z/p \mathbb Z)[X]$.
		
		\begin{enumerate}
			
			\item  Soit $P\in\mathbb Z[X]$ unitaire. Montrer que si $\bar P$ est irréductible dans $(\mathbb Z/p\mathbb Z)[X]$, alors $P$ est irréductible dans $\mathbb Z[X]$. La réciproque est-elle vraie?
			
			\item  Démontrer que le polynôme $X^4+X+1$ est irréductible dans $\mathbb Z[X]$.
			
		\end{enumerate}
	\end{exo}
	\begin{corrige}
		\begin{enumerate}
			
			\item  On procède par contraposition en démontrant que si $P$ est réductible dans $\mathbb Z[X]$, alors $\bar P$ est réductible dans $(\mathbb Z/p\mathbb Z)[X]$.
			Supposons en effet $P=AB$ avec $A,B\in\mathbb Z[X]$ non-inversibles. Écrivons $P(X)=X^n+\cdots$, $A(X)=aX^p+\cdots$ et $B(X)=bX^q+\cdots$, $a,b\in\mathbb Z$. Alors $ab=1$ et $p+q=n$. Quitte à changer $a$ en $-a$, on peut supposer $a=1$ et $b=1$. Ceci entraine que $p\geq 1$ et $q\geq 1$. Mais alors, on a $\bar P=\bar A\bar B$ avec $\bar A(X)=\bar 1 X^p+\cdots$ et $\bar B(X)=\bar 1 X^q+\cdots$. Les polynômes $\bar A$ et $\bar B$ ne sont donc pas inversibles et $\bar P$ est bien réductible dans $(\mathbb Z/p\mathbb Z)[X]$.\\
			
			La réciproque est fausse. Le polynôme $X^2+2$ est irréductible dans $\mathbb Z[X]$ (car il est irréductible dans $\mathbb R[X]$). Son image dans $(\mathbb Z/2\mathbb Z)[X]$ est $X^2$, qui est réductible car $X^2=X\times X$.
			
			\item  D'après ce qui précède, il suffit de démontrer que le polynôme $X^4+X+\bar 1$ est irréductible dans $(\mathbb Z/2\mathbb Z)[X]$. S'il ne l'est pas, ou bien il admet une racine dans $\mathbb Z/2\mathbb Z$, ou bien il se factorise en produit de deux polynômes de degré 2. Il est facile de vérifier que ni $\bar 0$ ni $\bar 1$ ne sont racines du polynôme. Donc s'il réductible, il s'écrit 
			$$X^4+X+\bar 1=(X^2+aX+b)(X^2+cX+d)$$
			(on peut effectivement toujours se ramener à des polynômes unitaires en multipliant l'un des polynômes par une constante et en divisant l'autre par la même constante). Le terme constant donne $bd=1$ ce qui implique encore $b=d=\bar 1$. Le terme de degré $1$ donne alors $a+c=\bar 1$ tandis que le terme de degré $3$ donne $a+c=\bar 0$. Il y a contradiction. Le polynôme $X^4+X+1$ est bien irréductible dans $\mathbb Z[X]$.
			
		\end{enumerate}
	\end{corrige}
	
	\begin{exo}[{Critère d'irréductibilité d'Eisenstein}]
		Si $P\in\mathbb Z[X]$, on appelle contenu de $P$, et on note $c(P)$, le pgcd des coefficients de $P$. 
		
		\begin{enumerate}
			
			\item  Soit $P,Q\in\mathbb Z[X]$ et $p$ un nombre premier. On suppose que $p$ divise tous les coefficients de $PQ$. Montrer que $p$ divise tous les coefficients de $P$ ou tous les coefficients de $Q$.
			
			\item  Soit $P,Q\in\mathbb Z[X]$ et $R(X)=\frac{PQ}{c(P)c(Q)}\in\mathbb Z[X]$. Démontrer que $c(R)=1$. En déduire que l'on a $c(PQ)=c(P)c(Q)$.
			
			\item  Soit $Q$ un polynôme de $\mathbb Z[X]$. On suppose que $Q$ n'est pas irréductible dans $\mathbb Q[X]$. Démontrer qu'il existe deux polynômes $A$ et $B$ de $\mathbb Z[X]$ tels que $Q=AB$, avec $\deg(A)<\deg(Q)$ et $\deg(B)<\deg(Q)$.
			
			\item  Soit $A(X)=a_n X^{n}+\dots+a_1X+a_0\in\mathbb Z[X]$. On suppose qu'il existe un nombre premier $p$ tel que 
			$$p|a_k,\textrm{ pour tout }0\leq k\leq n-1,\ p\not\mid a_n,\ p^2\not\mid a_0.$$
			Démontrer que $A$ est irréductible dans $\mathbb Q[X]$. 
			
			\item  Démontrer qu'il existe dans $\mathbb Q[X]$ des polynômes irréductibles de tout degré $n\geq 1$.
			
		\end{enumerate}
	\end{exo}
	\begin{corrige}
		\begin{enumerate}
			
			\item  Il est possible de faire une preuve directe par l'absurde. Le plus simple, néanmoins est de raisonner dans $\mathbb Z/p\mathbb Z[X]$. Pour $R\in\mathbb Z[X]$, notons en effet $\bar R$ son projeté dans $\mathbb Z/p\mathbb Z[X]$. Alors, on a $\overline{PQ}=\bar P\bar Q$. De plus, puisque $p$ divise tous les coefficients de $PQ$, on a $\overline{PQ}=0$. Puisque $\mathbb Z/p\mathbb Z$ est intègre, ceci implique que $\bar P=0$ ou $\bar Q=0$. La première éventualité signifie que $p$ divise tous les coefficients de $P$, la seconde que $p$ divise tous les coefficients de $Q$.
			
			\item  Si $c(R)\neq 1$, il existe un nombre premier $p$ qui divise tous les coefficients de $R$. Notons $P_1=P/c(P)$ et $Q_1=Q/c(Q)$. Alors $R=P_1Q_1$ et donc d'après la première question, $p$ divise tous les coefficients de $P_1$ ou tous les coefficients de $Q_1$. Ceci contredit la définition du contenu. Puisque $1=c(R)=\frac{c(PQ)}{c(P)c(Q)}$, on obtient bien que $c(PQ)=c(P)c(Q)$.
			
			\item  Remarquons d'abord qu'on peut se ramener à $c(Q)=1$ (quitte ensuite à multiplier par $c(Q)$). Factorisons  $Q=CD$ dans $\mathbb Q[X]$. Soit $\alpha$ et $\beta$ de sorte que $C_1=\alpha C$ et $D_1=\beta D$ soit éléments de $\mathbb Z[X]$. Alors $\alpha \beta Q=C_1D_1$. Utilisant le résultat de la question précédente, on a donc
			$\alpha\beta=c(C_1)c(D_1)$. Mais alors 
			$$Q=\frac{C_1 D_1}{\alpha\beta}=\frac{C_1D_1}{c(C_1)c(D_1)}=\frac{C_1}{c(C_1)}\times \frac {D_1}{c(D_1)}.$$
			On a le résultat voulu en posant $A=\frac{C_1}{c(C_1)}$ et $B=\frac{D_1}{c(D_1)}$.
			
			\item  Supposons que $A$ n'est pas irréductible dans $\mathbb Q[X]$. Alors d'après la question précédente,
			$A$ s'écrit $BC$, avec $1<\deg(B)<\deg(A)=n$. Projetons l'égalité $A=BC$ dans $\mathbb Z/p\mathbb Z$. Il vient 
			$$\overline{a_n}X^n =\overline{B}\times\overline{C}.$$
			De plus, puisque $\deg(B)+\deg(C)=\deg(A)$, que $\deg(\bar B)\leq \deg(B)$, $\deg(\bar C)\leq \deg(C)$
			et $\overline{a_n}\neq \bar 0$, $\deg(\bar B)=\deg(B)$ et $\deg(\bar C)=\deg(C)$. Par unicité de la réduction en produits d'irréductibles dans $\mathbb Z/p\mathbb Z[X]$, on a $\bar B=\overline{b_k}X^k$ et 
			$\bar C=\overline{c_l}X^l$. En particulier, $\overline{b_0}=\overline{c_0}=0$, c'est-à-dire $p|b_0$ et $p|c_0$. Puisque $a_0=b_0c_0$, on a $p^2|a_0$, une contradiction.
			
			\item  Les polynômes de degré $1$ sont irréductibles. Pour les polynômes de degré $n\geq 2$, il suffit de considérer par exemple $X^n-2$, auquel on peut très facilement appliquer le critère d'Eisenstein.
			
		\end{enumerate}
	\end{corrige}
	
	\begin{exo}[{Polynômes et fonctions polynomiales}]
		Soit $\mathbb K$ un corps commutatif. On note $\mathcal F(\mathbb K)$ l'ensemble des fonctions polynomiales de $\mathbb K$ dans lui-même. Pour tout polynôme $P\in\mathbb K[X]$, on note $\tilde P$ la fonction polynomiale associée.
		
		\begin{enumerate}
			
			\item  On suppose $\mathbb K$ infini. Vérifier que le morphisme de $\mathbb K$-algèbres $P\in\mathbb K[X]\mapsto \tilde P\in\mathcal F(K)$ est injectif. Que peut-on en déduire sur $\mathbb K[X]$  et $\mathcal F(\mathbb K)$?
			
			\item  Combien y-a-t-il d'éléments dans $(\mathbb Z/2\mathbb Z)[X]$ et dans $\mathcal F(\mathbb Z/2\mathbb Z)$?
			
			\item  On suppose désormais jusqu'à la fin de l'exercice que $\mathbb K$ est de cardinal fini $q$. Soit $P\in\mathbb K[X]$ et $R$ le reste dans la division euclidienne de $P$ par $\prod_{a\in \mathbb K}(X-a)$. Démontrer que l'on a $\tilde P=\tilde R$. 
			
			\item  Démontrer que $\mathcal F(\mathbb K)$ et $\mathbb K_{q-1}[X]$ sont isomorphes.
			
		\end{enumerate}
	\end{exo}
	\begin{corrige}
		\begin{enumerate}
			
			\item  Si $\tilde P=0$, alors le polynôme $P$ admet une infinité de racines. C'est donc nécessairement le polynôme nul. On en déduit que $\mathbb K[X]$ et
			$\mathcal F(\mathbb K)$ sont isomorphes.
			
			\item  Il est d'abord clair que $(\mathbb Z/2\mathbb Z)[X]$ est infini, la suite de polynômes $(X^n)$ étant par exemple une suite de polynômes distincts de $\mathbb Z/2\mathbb Z$. Pour chercher le cardinal de $\mathcal F(\mathbb Z/2\mathbb Z)$, on commence par remarquer qu'il est majoré par 4, qui est le nombre d'applications de $\mathbb Z/2\mathbb Z$ dans $\mathbb Z/2\mathbb Z$ (deux choix pour l'image de $\bar 0$, deux choix pour l'image de $\bar 1$). Mais, $f_1(x)=\bar 0$, $f_2(x)=\bar 1$, $f_3(x)=x$, $f_4(x)=x+\bar 1$ sont quatre fonctions polynomiales qui sont toutes différentes comme le montre le tableau suivant :
			\begin{eqnarray*}
				f_1(\bar 0)=\bar 0&&f_1(\bar 1)=\bar 0\\
				f_2(\bar 0)=\bar 1&&f_2(\bar 1)=\bar 1\\
				f_3(\bar 0)=\bar 0&&f_3(\bar 1)=\bar 1\\
				f_4(\bar 0)=\bar 1&&f_4(\bar 1)=\bar 0.
			\end{eqnarray*}
			Ainsi, le cardinal de $\mathcal F(\mathbb Z/2\mathbb Z)$ est égal à 4. Cet exemple montre toute la différence, pour les corps finis, entre polynôme et fonction polynomiale. Les polynômes $X$ et $X^2$ de $(\mathbb Z/2\mathbb Z)[X]$ sont différents (ils n'ont pas le même degré...). Les fonctions polynomiales $x$ et $x^2$ sont égales (les images de chaque point sont égales). C'est cette différence que le reste de l'exercice se propose d'explorer.
			
			\item  Notons $Q$ le polynôme $Q(X)=\prod_{a\in\mathbb K}(X-a)$ et $P=AQ+R$. Alors, on a 
			$$\tilde P(x)=\tilde A(x)\tilde Q(x)+\tilde R(x)=\tilde R(x)$$
			puisque $\tilde Q(x)=0$ pour tout $x\in\mathbb K$.
			
			\item  Considérons l'application $\phi:\mathbb K_{q-1}[X]\to \mathcal F(\mathbb K),\ R\mapsto \tilde R$. C'est clairement une application linéaire. Elle est injective : si $R_1,R_2\in\mathbb K_{q-1}[X]$ sont tels que $\tilde R_1=\tilde R_2$, alors $R_1-R_2$ admet au moins $q$ racines et est de degé inférieur ou égal à $q-1$: c'est donc le polynôme nul. De plus, si $f\in\mathcal F(\mathbb K)$, alors il existe $P\in\mathbb K[X]$ tel que $\tilde P=f$. Si on note $R$ le reste de la division euclidienne de $P$ par $\prod_{a\in \mathbb K}(X-a)$, alors $\phi(P)=\tilde R=\tilde P=f$ ce qui montre que $\phi$ est surjective. Ainsi, $\phi$ est un isomorphisme de $\mathbb K_{q-1}[X]$ sur $\mathcal F(\mathbb K)$.
			
		\end{enumerate}
	\end{corrige}
	
	\subsection*{Généralités}
	
	\begin{exo}[{Racine carrée de $X$}]
		Démontrer qu'il n'existe pas de fraction rationnelle $F$ tel que $F^2=X$.
	\end{exo}
	\begin{corrige}
		Si $F$ est une fraction rationnelle, alors $\deg(F^2)=2\deg (F)$.
		Or, $\deg(X)=1$, et l'équation $2n=1$ n'a pas de solutions dans $\mathbb Z$.
	\end{corrige}
	
	\begin{exo}[{Pôles et racines}]
		Soient $p$ et $q$ deux entiers naturels premiers entre eux. Déterminer les racines
		et les pôles de $(X^p-1)/(X^q-1)$, en précisant leur ordre de multiplicité.
	\end{exo}
	\begin{corrige}
		Les racines de $X^p-1$ (resp. $X^q-1$) sont les racines $p$-ièmes de l'unité
		(resp. $q$-ièmes) et elles sont toutes simples. Mais il peut y avoir des simplifications
		entre le numérateur et le dénominateur et on doit déterminer les racines communes.
		Soit donc $\omega$ une racine de $X^p-1$ et de $X^q-1$. Alors $\omega^p=\omega^q=1$.
		Mais, puisque $p$ et $q$ sont premiers entre eux, on peut d'après le théorème de Bézout
		trouver deux entiers $u,v\in\mathbb Z$ tels que $pu+qv=1$. Alors
		$$\omega=\omega^{pu+qv}=(\omega^p)^u(\omega^q)^v=1.$$
		Réciproquement, $1$ est racine commune des deux polynômes. On a donc prouvé que les racines de la fraction rationnelle
		sont les racines $p$-ièmes de l'unité autre que 1. Ses pôles sont les racines $q$-ièmes de l'unité, autre que 1.
		Les multiplicités respectives sont à chaque fois égales à 1.
	\end{corrige}
	
	\begin{exo}[{Fraction rationnelle à valeurs rationnelles sur les entiers}]
		Soit $R(X)=\frac{P(X)}{Q(X)}$ une fraction rationnelle de $\mathbb R(X)$ avec $P\wedge Q=1$ et telle que $R(n)\in\mathbb Q$ pour une infinité d'entiers  $n\in\mathbb N$. On veut démontrer que $R(x)=\frac{P_1(X)}{Q_1(X)}$ où $P_1,Q_1\in\mathbb Z[X]$. On note $\omega(R)=\deg(P)+\deg(Q)$.
		
		\begin{enumerate}
			
			\item  Démontrer le résultat si $\omega(R)=0$.
			
			\item  Soit $d\geq 0$. On suppose que le résultat est vrai pour toute fraction rationnelle $R$ tel que $\omega(R)\leq d$ et on souhaite le prouver pour toute fraction rationnelle telle que $\omega(R)=d+1$. On fixe donc $R=P/Q$ une fraction rationnelle telle que $\omega(R)=d+1$ et $R(n)\in\mathbb Q$ pour tout
			$n\in\mathbb N$.
			
			\begin{enumerate}
				
				\item  Pourquoi peut-on supposer que $\deg(P)\geq \deg(Q)$?
				
				\item  Soit $n\in\mathbb N$ tel que $Q(n)\neq 0$ et $P(n)/Q(n)\in\mathbb Q$. On pose 
				$$P_0(X)=\frac{P(X)Q(n)-Q(X)P(n)}{Q(n)(X-n)}.$$
				Justifier que $P_0\in\mathbb R[X]$ avec $\deg(P_0)<\deg(P)$.
				
				\item  Conclure.
				
			\end{enumerate}
			
			
		\end{enumerate}
	\end{exo}
	\begin{corrige}
		\begin{enumerate}
			
			\item  Si $\omega(R)=0$, alors $R$ est le quotient de deux réels $R=p/q$. L'hypothèse entraîne que $R$ est à valeurs dans $\mathbb Q$, c'est-à-dire que $p/q$ est rationnel. On peut donc écrire $p/q=p_1/q_1$ avec $p_1,q_1\in\mathbb Z$.
			
			\item  
			
			\begin{enumerate}
				
				\item  Si $\deg(P)<\deg(Q)$, il suffira de prouver le résultat à la fraction rationnelle $1/R$ qui vérifie les mêmes hypothèses que $R$, si ce n'est que cette fois le degré de son numérateur sera supérieur à celui de son dénominateur.
				
				\item  Considérons $S$ le polynôme $S(X)=P(X)Q(n)-Q(X)P(n)$. Alors, puisque $\deg(P)\geq \deg(Q)$, $\deg(S)\leq \deg(P)$. De plus, $S(n)=0$. Ainsi, $S$ se factorise en $S(X)=(X-n)S_1(X)$ avec $\deg(S_1)<\deg(P)$. Mais alors, $P_0=\frac{S_1}{Q(n)}$ est bien un élément de $\mathbb R[X]$ de degré inférieur strict au degré de $P$.
				
				\item  Considérons $T(X)=\frac{P_0(X)}{Q(X)}$. Alors $\omega(T)=\deg(P_0)+\deg(Q)<\omega(R)$. De plus, si $m\neq n$ est un entier tel que $R(m)\in\mathbb Q$, on a, 
				$$T(m)=\frac{P(m)Q(n)-Q(m)P(n)}{Q(m)Q(n)(m-n)}=\frac{P(m)}{Q(m)(n-m)}-\frac{P(n)}{Q(n)(n-m)}\in\mathbb Q.$$
				Ainsi, $T(m)\in\mathbb Q$ pour une infinité d'entiers $m$. D'après l'hypothèse de récurrence, 
				$T=\frac{P_2}{Q_2}$ 
				où $P_2,Q_2\in \mathbb Z[X]$. On remonte ensuite à $R$ en remarquant que 
				$$R(X)=(X-n)T(X)+\frac{P(n)}{Q(n)}.$$
				Puisque $P(n)/Q(n)$ est rationnel, il s'écrit encore $c/d$ avec $c,d\in\mathbb Z$.
				Ainsi, mettant tout au même dénominateur, on a 
				$$R(X)=\frac{P_1(X)}{Q_1(X)}$$
				où
				$$P_1(X)=d(X-n)P_2(X)+c$$
				et $$Q_1(X)=dQ_2(X)$$
				qui sont bien des polynômes à coefficients entiers.
				
			\end{enumerate}
			
			
		\end{enumerate}
	\end{corrige}
	
	\subsection*{Décomposition en éléments simples}
	
	\begin{exo}[{Tous les cas possibles}]
		Décomposer sur $\mathbb R$ les fractions rationnelles suivantes :
		$$\begin{array}{lll}
			\displaystyle\mathbf{1.}\quad\frac{X^2+2X +5}{X^2-3X+2}&\quad\quad\mathbf{2.}\quad \displaystyle\frac{X^2+3X+1}{(X-1)^2(X-2)}
			&\quad\quad\mathbf{3.}\quad \displaystyle \frac 1{X^4-1}
		\end{array}$$
	\end{exo}
	\begin{corrige}
		\begin{enumerate}
			
			\item  
			\begin{itemize}
				
				\item  On commence par calculer la partie entière de la fraction rationnelle : 
				$$X^2+2X+5=1(X^2-3X+2)+5X+3.$$
				
				\item  On factorise le dénominateur $$X^2-3X+2=(X-1)(X-2).$$
				
				\item  On écrit la décomposition en éléments simples a priori : 
				$$\frac{X^2+2X+5}{X^2-3X+2}=1+\frac a{X-1}+\frac{b}{X-2}.$$
				
				\item  On calcule les coefficients $a$ et $b$ : posons $P(X)=X^2+2X+5$, $Q(X)=X^2-3X+2$, $Q'(X)=2X-3.$ 
				$$a=\frac{P(1)}{Q'(1)}=\frac{8}{-1}=-8$$
				$$b=\frac{P(2)}{Q'(2)}=\frac{13}{1}=13.$$
				
				\item  Ainsi, 
				$$ \frac{X^2+2X+5}{X^2-3X+2}=1-\frac {8}{X-1}+\frac{13}{X-2}.$$
				
			\end{itemize}
			
			
			\item  
			\begin{itemize}
				
				\item  Puisque le degré du dénominateur est strictement supérieur au degré du numérateur, la partie entière est nulle.
				
				\item  On a donc 
				$$\frac{X^2+3X+1}{(X-1)^2(X-2)}=\frac{a}{X-2}+\frac{b}{X-1}+\frac{c}{(X-1)^2}.$$
				
				\item  Calculons $a$. On pose $P(X)=X^2+3X+1$, $Q(X)=(X-1)^2(X-2)$, $Q'(X)=2(X-1)(X-2)+(X-1)^2$ et 
				$$a=\frac{P(2)}{Q'(2)}=\frac{11}{1}=11.$$
				
				\item  Calculons $c$. On multiplie par $(X-1)^2$ et on évalue en $X=1$.
				\begin{eqnarray*}
					\frac{X^2+3X+1}{X-2}&=&\frac{11(X-1)^2}{X-2}+b(X-1)+c\\ 
					\frac{5}{-1}&=& c\implies c=-5.
				\end{eqnarray*}
				On a donc 
				$$\frac{X^2+3X+1}{(X-1)^2(X-2)}=\frac{11}{X-2}+\frac{b}{X-1}+\frac{-5}{(X-1)^2}.$$
				
				\item  Pour calculer $b$, on multiplie par $X$ et on fait tendre $X$ vers l'infini.
				$$\frac{X(X^2+3X+1)}{(X-1)^2(X-2)}=\frac{11X}{X-2}+\frac{bX}{X-1}+\frac{-5X}{(X-1)^2}.$$
				Quand $X\to+\infty,$
				$$1= 11+ b+ 0 \implies b=-10.$$
				Conclusion :
				$$\frac{X^2+3X+1}{(X-1)^2(X-2)}=\frac{11}{X-2}-\frac{10}{X-1}-\frac{5}{(X-1)^2}.$$
				
			\end{itemize}
			
			
			\item  
			\begin{itemize}
				
				\item  La partie entière est nulle.
				
				\item  On factorise le dénominateur : 
				$$X^4-1=(X^2-1)(X^2+1)=(X-1)(X+1)(X^2+1).$$
				
				\item  On écrit 
				$$\frac1{X^4-1}=\frac a{X-1}+\frac{b}{X+1}+\frac{cX+d}{X^2+1}.$$
				
				\item  On pose $P(X)=1$, $Q(X)=X^4-1$ de sorte que $Q'(X)=4X^3$ et
				$$a=\frac{P(1)}{Q'(1)}=\frac14$$
				$$b=\frac{P(-1)}{Q'(-1)}=\frac{-1}{4}.$$
				
				\item  Pour calculer $c$ et $d$, on effectue d'abord la décomposition en éléments neutres sur $\mathbb C$ et on regroupe  les termes conjugués.
				
				\item  Décomposons la fraction rationnelle sur $\mathbb C.$
				$$X^4-1=(X-1)(X+1)(X^2+1)=(X-1)(X+1)(X-i)(X+i).$$
				On écrit 
				$$\frac 1{X^4-1}=\frac{a}{X-1}+\frac{b}{X+1}+\frac{e}{X-i}+\frac{f}{X+i}.$$
				On pose $P(X)=1$, $Q(X)=X^4-1$, $Q'(X)=4X^3.$  On a déjà vu $a=1/4$, $b=-1/4$ et on a 
				$$e=\frac{P(i)}{Q'(i)}=\frac{1}{4i^3}=\frac{1}{-4i}=\frac{i}4$$
				$$f=\frac{P(i)}{Q'(-i)}=\frac{1}{4(-i)^3}=\frac{1}{4i}=\frac{-i}4.$$
				Ainsi, 
				$$\frac 1{X^4-1}=\frac{1/4}{X-1}-\frac{1/4}{X+1}+\frac{i/4}{X-i}-\frac{i/4}{X+i}.$$
				
				\item  On termine en passant de $\mathbb C$ à $\mathbb R$.
				\begin{eqnarray*}
					\frac 1{X^4-1}&=&\frac{1/4}{X-1}-\frac{1/4}{X+1}+\frac{i/4}{X-i}-\frac{i/4}{X+i}\\ 
					&=&\frac{1/4}{X-1}-\frac{1/4}{X+1}+\frac 14\cdot \frac{i(X+i)-i(X-i)}{(X-i)(X+i)}\\ 
					&=&\frac{1/4}{X-1}-\frac{1/4}{X+1}+\frac 14\cdot \frac{-2}{X^2+1}\\ 
					&=&\frac{1/4}{X-1}-\frac{1/4}{X+1}-\frac{1/2}{X^2+1}.
				\end{eqnarray*}
				
			\end{itemize}
			
			
		\end{enumerate}
	\end{corrige}
	
	\begin{exo}[{Pôles simples}]
		Décomposer en éléments simples les fractions rationnelles suivantes : 
		$$\begin{array}{lll}
			\displaystyle\mathbf{1.}\quad\frac{1}{X^3-X}&\quad\quad\mathbf{2.}\quad  \displaystyle \frac{X^3}{(X-1)(X-2)(X-3)}
		\end{array}$$
	\end{exo}
	\begin{corrige}
		\begin{enumerate}
			
			\item  La partie entière est nulle, et le dénominateur se factorise en $X(X-1)(X+1)$.
			Multipliant par $X$ et faisant $X=0$, on trouve la partie polaire relativement à $X-0$, et ainsi de suite...
			On trouve finalement
			$$\frac{1}{X^3-X}=\frac{-1}{X}+\frac{1/2}{X-1}+\frac{1/2}{X+1}.$$
			
			\item  En appliquant exactement la même méthode, on trouve que la décomposition en éléments simples est
			$$1+\frac{1}{2(X-1)}-\frac{8}{X-2}+\frac{27}{2(X-3)}.$$
			
		\end{enumerate}
	\end{corrige}
	
	\begin{exo}[{Pôles multiples}]
		Décomposer en éléments simples les fractions rationnelles suivantes : 
		$$\begin{array}{lll}
			\displaystyle\mathbf{1.}\quad\frac{2X^2+1}{(X^2-1)^2}&
			\quad\quad\mathbf{2.}\quad\displaystyle\frac{X^3+1}{(X-1)^3}
		\end{array}$$
	\end{exo}
	\begin{corrige}
		\begin{enumerate}
			
			\item  Le dénominateur se factorise en $(X^2-1)^2=(X-1)^2(X+1)^2$. Il faut donc étudier la partie polaire relative à +1 et à -1.
			Commençons par étudier la partie polaire relative à $-1$. La fraction rationnelle s'écrit sous la forme
			$$\frac{2X^2+1}{(X-1)^2(X+1)^2}=\frac{\lambda_1}{X+1}+\frac{\lambda_2}{(X+1)^2}+G(X),$$
			où $G(X)=\frac{\mu_1}{X-1}+\frac{\mu_2}{(X-1)^2}$ est une fraction rationnelle dont $-1$ n'est pas un pôle. Multipliant cette égalité
			par $(X+1)^2$ et faisant $X=-1$, on trouve $\lambda_2=3/4$. On calcule ensuite 
			$$\frac{2X^2+1}{(X-1)^2(X+1)^2}-\frac{3/4}{(X+1)^2}=\frac{(5X+1)/4}{(X+1)(X-1)^2}=\frac{\lambda_1}{X+1}+G(X).$$
			On multiplie cette fois par $X+1$, et on fait $X=-1$, et on trouve $\lambda_1=-1/4$. Pour étudier la partie polaire relative à $1$,
			on peut procéder de la même façon ou simplement remarquer que la fraction rationnelle est paire. On en déduit que
			$$\frac{2X^2+1}{(X-1)^2(X+1)^2}=\frac{-1/4}{X+1}+\frac{3/4}{(X+1)^2}+\frac{1/4}{X-1}+\frac{3/4}{(X-1)^2}.$$
			
			\item  La partie entière de cette fraction rationnelle est égale à 1, et on a 
			$$\frac{X^3+1}{(X-1)^3}=1+\frac{3X^2-3X+2}{(X-1)^3}=1+\frac{a}{(X-1)^3}+\frac{b}{(X-1)^2}+\frac{c}{X-1}.$$
			Pour trouver $a$, on multiplie par $(X-1)^3$ et on fait $X=1$. On trouve
			$$a=2.$$
			Pour trouver $b$, on soustrait $\frac2{(X-1)^3}$, et on trouve
			\begin{eqnarray*}
				\frac{3X^2-3X+2}{(X-1)^3}-\frac 2{(X-1)^3}&=&\frac{3X}{(X-1)^2}\\
				&=&\frac{b}{(X-1)^2}+\frac c{X-1}.
			\end{eqnarray*}
			Multipliant par $(X-1)^2$ et faisant $X=1$, on trouve
			$$b=3.$$
			Finalement, on retranche encore $\frac{3}{(X-1)^2}$ de sorte que
			\begin{eqnarray*}
				\frac{3X}{(X-1)^2}-\frac{3}{(X-1)^2}&=&\frac{3}{X-1}\\
				&=&\frac{c}{X-1}.
			\end{eqnarray*}
			On a donc $c=3$. Finalement, la décomposition en éléments simples recherché est
			$$1+\frac{2}{(X-1)^3}+\frac{3}{(X-1)^2}+\frac{3}{X-1}.$$
			
		\end{enumerate}
	\end{corrige}
	
	\begin{exo}[{Pôle multiple et facteur irréductible de degré $2$}]
		Décomposer en éléments simples sur $\mathbb R$ la fraction rationnelle suivante : 
		$$\frac{X^4+1}{(X+1)^2(X^2+1)}$$
	\end{exo}
	\begin{corrige}
		On commence par réaliser la décomposition en éléments simples sur $\mathbb C$.
		Les pôles sont $-1$ (double), $i$ et $-i$. La fraction rationnelle étant à coefficients réels, 
		les parties polaires sont conjuguées. On a donc
		$$\frac{X^4+1}{(X+1)^2(X^2+1)}=1+\frac{a}{(X+1)^2}+\frac{b}{X+1}+\frac{c}{X-i}+\frac{\bar c}{X+i}.$$
		Multipliant par $X-i$ et faisant $X=i$, on trouve
		$$c=\frac{i^4+1}{(i+i)(i+1)^2}=-\frac12.$$
		De même, on a
		$$a=\frac{1+1}{1+1}=1.$$
		En retranchant $\frac1{(X+1)^2}$, on trouve 
		$$\frac{X^4-X^2}{(X+1)^2(X^2+1)}=\frac{X^2(X-1)}{(X+1)(X^2+1)}=1+\frac{b}{X+1}+\frac c{X-i}+\frac{\bar c}{X+i}.$$
		Multipliant par $X+1$ et faisant $X=-1$, on trouve
		$$b=-1.$$
		Finalement,
		$$\frac{X^4+1}{(X+1)^2(X^2+1)}=1+\frac{1}{(X+1)^2}-\frac{1}{X+1}-\frac{1/2}{X-i}-\frac{1/2}{X+i}.$$
		Finalement, en regroupant les deux derniers termes, on trouve la décomposition sur $\mathbb R$ :
		$$\frac{X^4+1}{(X+1)^2(X^2+1)}=1+\frac{1}{(X+1)^2}-\frac{1}{X+1}-\frac{X}{X^2+1}.$$
	\end{corrige}
	
	\begin{exo}[{Avec paramètres}]
		Décomposer en éléments simples les fractions rationnelles suivantes :
		$$\begin{array}{lll}
			\displaystyle\mathbf{1.}\quad \frac{1}{X^n-1}&
			\displaystyle\quad\quad\mathbf{2.}\quad\frac{X^{n-1}}{X^n-1}&
		\end{array}$$
	\end{exo}
	\begin{corrige}
		\begin{enumerate}
			
			\item  Les pôles de $1/(X^n-1)$ sont les racines $n$-ièmes de l'unité, c'est-à-dire les complexes $x_k=e^{2ik\pi/n}$, $k=0,\dots,n-1$.
			Chaque pôle est simple, la partie polaire correspondante est donc de la forme $\frac{c_k}{X-x_k}$ avec $c_k=\frac{1}{P'(x_k)}=\frac{1}{nx_k^{n-1}}$.
			Or, $x_k^{n-1}=\frac{x_k^n}{x_k}=\frac1{x_k}=e^{-2ik\pi/n}$. La décomposition en éléments simples recherché vaut donc
			$$\frac1{X^n-1}=\frac1n\sum_{k=0}^{n-1}\frac{e^{2ik\pi/n}}{X-e^{2ik\pi/n}}.$$
			
			\item  Les racines de $X^n-1$ sont les complexes $\omega_k=e^{2ik\pi/n}$, $0\leq k\leq n-1$. 
			La fraction rationnelle admet donc $n$ pôles, qui sont tous simples. Sa décomposition en éléments simples a pour forme
			$$\frac{X^{n-1}}{X^n-1}=\sum_{k=0}^{n-1}\frac{\alpha_k}{X-\omega_k}$$
			avec $\alpha_k=\frac{\omega_k^{n-1}}{n\omega_k^{n-1}}=\frac{1}{n}$. La décomposition en éléments simples est
			$$\frac{X^{n-1}}{X^n-1}=\frac1n\sum_{k=0}^{n-1}\frac{1}{X-e^{2ik\pi/n}}.$$
			
		\end{enumerate}
	\end{corrige}
	
	\begin{exo}[{Avec paramètres}]
		Décomposer en éléments simples les fractions rationnelles suivantes :
		$$\begin{array}{lll}
			\displaystyle\mathbf{1.}\quad \frac{X^m}{(X-1)^n}&
			\displaystyle\quad\quad\mathbf{2.}\quad\frac{1}{X(X+1)\cdots (X+n)}&
		\end{array}$$
	\end{exo}
	\begin{corrige}
		\begin{enumerate}
			
			\item  On écrit $X=(X-1)+1$ et on utilise la formule du binôme. Il vient 
			$$X^m=\big((X-1)+1\big)^m =\sum_{k=0}^m \binom mk (X-1)^k.$$
			On distingue alors deux cas pour l'écriture de la décomposition en éléments simples : si $m\geq n,$ alors 
			$$\frac{X^m}{(X-1)^n}=\sum_{k=n}^m \binom mk (X-1)^{k-n}+\sum_{k=0}^{n-1}\frac{\binom mk}{(X-1)^{n-k}},$$
			le premier terme étant la partie entière et le second la partie polaire.
			Si $m<n$, alors on écrit simplement
			$$\frac{X^m}{(X-1)^n}=\sum_{k=0}^{m}\frac{\binom mk}{(X-1)^{n-k}}.$$
			
			\item  On peut écrire 
			$$\frac1{X(X+1)\cdots (X+n)}=\sum_{k=0}^n\frac{a_k}{X+k}.$$
			Reste à déterminer $a_k$ pour $k=0,\dots,n$. Le plus simple ici, comme le dénominateur est déjà factorisé, est de multiplier par $X+k$ et de faire tendre $X$ vers $-k$. On obtient alors 
			$$a_k=\frac{1}{(-k)(-k+1)\cdots(-k+(k-1)) (-k+(k+1))\cdots (-k+n)}=\frac{(-1)^k}{k!(n-k)!}.$$
			
		\end{enumerate}
	\end{corrige}
	
	\begin{exo}[{Avec paramètres}]
		Décomposer en éléments simples la fraction rationnelle suivante :
		$$\frac{1}{(X-1)(X^n-1)}.$$
	\end{exo}
	\begin{corrige}
		La fraction admet 1 comme pôle double et $\omega_k=e^{2i\pi k/n}$, $k=1,\dots,n-1$ comme pôles simples. On commence par calculer le coefficient devant $(X-1)^2$. On écrit que la fraction est égale à 
		$$\frac{1}{(X-1)^2(1+X+\dots+X^{n-1})}=\frac{\lambda}{(X-1)^2}+Q(X),$$
		où $1$ n'est pas pôle double de $Q$. Multipliant par $(X-1)^2$ et faisant tendre $X$ vers $1,$ on en déduit que
		$\lambda=\frac{1}{n}.$ Le coefficient devant $\frac1{X-\omega_k}$ est lui obtenu en dérivant le dénominateur, et en remplaçant par $\omega_k$. On trouve :
		$$\frac{1}{\omega_k^n-1+(\omega_k-1)n\omega_k^{n-1}}=\frac{\omega_k}{n(\omega_k-1)}.$$
		Enfin, pour calculer le terme en $\frac1{X-1}$, on écrit que
		$$\frac{1}{(X-1)^2(1+X+\dots+X^{n-1})}-\frac{1/n}{(X-1)^2}=\frac{\mu}{(X-1)}+R(X),$$
		où $1$ n'est plus un pôle de $R$. On a alors
		\begin{eqnarray*}
			\frac{1}{(X-1)^2(1+X+\dots+X^{n-1})}-\frac{1/n}{(X-1)^2}\\
			=\frac{n-(1+X+\dots+X^{n-1})}{n(X-1)^2(1+X+\dots+X^{n-1})}\\
			=-\frac{(X-1)+(X^2-1)+\dots+(X^{n-1}-1)}{n(X-1)^2(1+X+\dots+X^{n-1})}\\
			=-\frac{1+(1+X)+(1+X+X^2)+\dots+(1+X+\dots+X^{n-2})}{n(X-1)(1+X+\dots+X^{n-1})}.
		\end{eqnarray*}
		Multipliant par $X-1$ et faisant $X=1$, on trouve
		$$\mu=-\frac{1+2+\dots+n-1}{n^2}=\frac{1-n}{2n}.$$
		Finalement, la décomposition en éléments simples de la fraction rationnelle est donc :
		$$\frac{1-n}{2n(X-1)}+\frac{1}{n(X-1)^2}+\sum_{k=1}^{n-1}\frac{\omega_k}{n(\omega_k-1)}\frac{1}{X-\omega_k}.$$
	\end{corrige}
	
	\subsection*{Applications}
	
	\begin{exo}[{Calcul de primitives}]
		Déterminer une primitive des fonctions suivantes :
		$$\begin{array}{lll}
			\mathbf{1.}\ \displaystyle x\mapsto \frac{1}{1-x^2}\textrm{ sur }]1,+\infty[&\quad\quad&\mathbf{2.}\ \displaystyle x\mapsto\frac x{x^3-7x+6}\textrm{ sur }]2,+\infty[.
		\end{array}$$
	\end{exo}
	\begin{corrige}
		On va réaliser à chaque fois une décomposition en éléments simples.
		
		\begin{enumerate}
			
			\item  On peut factoriser $1-X^2$ en $(1-X)(1+X)$ et on sait qu'il existe des réels $a$ et $b$ de sorte que 
			$$\frac{1}{1-X^2}=\frac{a}{X-1}+\frac{b}{X+1}.$$
			Notons $Q(X)=1-X^2$, de sorte que $Q'(X)=-2X$. On a $a=1/Q'(1)=-1/2$ et $b=1/Q'(-1)=1/2.$ Ainsi, 
			$$\frac{1}{1-X^2}=\frac{-1/2}{X-1}+\frac{1/2}{X+1}.$$
			Ainsi, une primitive de $x\mapsto \frac1{1-x^2}$ sur $]1,+\infty[$ est la fonction 
			$$x\mapsto \frac 12\ln(x+1)-\frac 12\ln(x-1).$$
			
			\item  On peut factoriser $X^3-7X+6$ en $(X+3)(X-1)(X-2)$ (pour cela, on remarque auparavant que $1$ est racine évidente par exemple). On en déduit qu'il existe trois réels $a$, $b$ et $c$ tels que 
			$$\frac{X}{X^3-7X+6}=\frac{a}{X+3}+\frac{b}{X-1}+\frac{c}{X-2}.$$
			Posons $P(X)=X$ et $Q(X)=X^3-7X+6$, de sorte que $Q'(X)=3X^2-7.$ Alors on a 
			$$a=\frac{P(-3)}{Q'(-3)}=-\frac3{20},\ b=\frac{P(1)}{Q'(1)}=-\frac 14,\ 
			c=\frac{P(2)}{Q(2)}=\frac 25.$$
			On a donc 
			$$\frac{x}{x^3-7x+6}=-\frac{3}{20(x+3)}-\frac1{4(x-1)}+\frac2{5(x-2)}.$$
			Une primitive de cette fonction sur $]2,+\infty[$ est donc donnée par 
			$$x\mapsto -\frac3{20}\ln(x+3)-\frac 14\ln(x-1)+\frac 25\ln(x-2).$$
			
		\end{enumerate}
	\end{corrige}
	
	\begin{exo}[{Calcul d'intégrale}]
		Soit $f(x)=\frac{5x^2+21x+22}{(x-1)(x+3)^2}$, $x\in ]1,+\infty[$.
		
		\begin{enumerate}
			
			\item  Décomposer en éléments simples la fraction rationnelle
			$$\frac{5X^2+21X+22}{(X-1)(X+3)^2}$$
			en éléments simples.
			
			\item  En déduire la primitive de $f$ sur $]1,+\infty[$ qui s'annule en 2.
			
		\end{enumerate}
	\end{exo}
	\begin{corrige}
		\begin{enumerate}
			
			\item  On écrit la décomposition a priori :
			$$\frac{5X^2+21X+22}{(X-1)(X+3)^2}=\frac{a}{X-1}+\frac{b}{X+3}+\frac{c}{(X+3)^2}.$$
			Pour déterminer $a$, on multiplie par $X-1$ et on fait tendre $X$ vers $1$ : 
			on trouve 
			$$a=\frac{5+21+22}{4^2}=\frac{48}{16}=3.$$
			Pour déterminer $c$, on multiplie par $(X+3)^2$ et on trouve
			$$c=\frac{45-63+22}{-4}=-1.$$
			Pour déterminer $b$, on multiplie par $X$ et on fait tendre $X$ vers $+\infty$. On a donc 
			$$\frac{3X}{X-1}+\frac {bX}{X+3}-\frac{X}{(X+3)^2}=\frac{5X^3+21X^2+22X}{(X-1)(X+3)^2}.$$
			En faisant tendre $X$ vers $+\infty$, on obtient 
			$$3+b=5\iff b=2.$$
			
			\item  On a donc obtenu, pour $x>1,$ $$f(x)=\frac{3}{x-1}+\frac{2}{x+3}-\frac1{(x+3)^2}.$$
			On intègre chacun des éléments simples de la décomposition précédente, en tenant compte du fait que l'on travaille sur l'intervalle $]1,+\infty[$. Les primitives de $f$ sur cet intervalle sont donc les fonctions
			$$F(x)=3\ln(x-1)+2\ln(x+3)+\frac{1}{x+3}+d.$$
			La primitive qui s'annule en $2$ et celle pour laquelle $d$ vérifie l'équation 
			$$3\ln(1)+2\ln 5+\frac 15+d=0.$$
			La primitive de $f$ sur l'intervalle $]1,+\infty[$ qui s'annule en 2 est donc la fonction $F$ définie par
			$$F(x)=3\ln(x-1)+2\ln(x+3)+\frac{1}{x+3}-2\ln 5-\frac 15.$$
			
		\end{enumerate}
	\end{corrige}
	
	\begin{exo}[{Calcul d'intégrale}]
		Pour $x>0,$ on pose 
		$$f(x)=\frac{x^4+x^3+4x^2+2x+2}{x^3+x}.$$
		
		\begin{enumerate}
			
			\item  Décomposer en éléments simples la fraction rationnelle
			$$\frac{X^4+X^3+4X^2+2X+2}{X^3+X}.$$
			
			\item  En déduire la valeur de $\int_1^2 f(x)dx.$
			
		\end{enumerate}
	\end{exo}
	\begin{corrige}
		\begin{enumerate}
			
			\item  On commence par effectuer la division euclidienne de $X^4+X^3+4X^2+2X+2$ par $X^3+X$. On trouve que le quotient vaut $X+1$ et le reste vaut $3X2+X+2$. On a donc 
			\begin{align*}
				\frac{X^4+X^3+4X^2+2X+2}{X^3+X}&=X+1+\frac{3X^2+X+2}{X^3+X}\\
				&=X+1+\frac{3X^2+X+2}{X(X-i)(X+i)}\\
				&=X+1+\frac{a}{X}+\frac{b}{X-i}+\frac{\bar b}{X+i}.
			\end{align*}
			Posons $P(X)=3X^2+X+2$ et $Q(X)=X^3+X$. Alors $a=P(0)/Q'(0)=2,$
			$b=P(i)/Q'(i)=\frac12-\frac i2.$
			On obtient
			\begin{align*}
				\frac{X^4+X^3+4X^2+2X+2}{X^3+X}&=X+1+\frac{2}{X}+\frac{1}2\frac{1-i}{X-i}+\frac 12\frac{1+i}{X+i}\\
				&=X+1+\frac 2X+\frac{X+1}{X^2+1}.
			\end{align*}
			
			\item  On calcule l'intégrale en cherchant une primitive de chacun des termes précédents. On écrit encore
			$$\frac{X+1}{X^2+1}=\frac{X}{X^2+1}+\frac{1}{X^2+1}.$$
			On obtient donc
			\begin{align*}
				\int_1^2 f(x)dx&=\left[\frac{x^2}2+x+2\ln(x)+\frac{1}2\ln(x^2+1)+\arctan(x)\right]_1^2\\
				&=\frac 52+\frac 32\ln(2)+\frac 12\ln(5)+\arctan(2)-\frac \pi4.
			\end{align*}
			
		\end{enumerate}
	\end{corrige}
	
	\begin{exo}[{Dérivée $n$-ème d'une fraction rationnelle}]
		Soit $n\geq 1$. Déterminer la dérivée $n$-ème de la fonction suivante  :
		$$f(x)=\frac{1}{x^2-3x+2}.$$
	\end{exo}
	\begin{corrige}
		On va commencer par décomposer en éléments simples la fraction rationnelle
		$$\frac{1}{X^2-3X+2}.$$
		On peut factoriser $X^2-3X+2$ en $(X-1)(X-2).$ On a donc
		$$\frac{1}{X^2-3X+2}=\frac{a}{X-1}+\frac b{X-2}.$$
		Notons $P(X)=X^2-3X+2$ de sorte que $P'(X)=2X-3$. Alors $a=\frac1{P'(1)}=-1$ 
		et $b=\frac 1{P'(2)}=1.$ Ainsi, on a
		$$f(x)=\frac{1}{x-2}-\frac 1{x-1}.$$
		D'autre part, il est facile de voir (par exemple par récurrence) que la dérivée $n$-ème de $x\mapsto \frac 1x$ est $x\mapsto \frac{(-1)^n n!}{x^{n+1}}.$
		On en déduit finalement que 
		$$f^{(n)}(x)=\frac{(-1)^n n!}{(x-2)^{n+1}}-\frac{(-1)^n n!}{(x-1)^{n+1}}.$$
	\end{corrige}
	
	\begin{exo}[{Un calcul de somme}]
		\begin{enumerate}
			
			\item  Décomposer en éléments simples la fraction rationnelle $\displaystyle\frac{1}{X(X+1)(X+2)}$.
			
			\item  En déduire la limite de la suite $(S_n)$ suivante : $\displaystyle S_n=\sum_{k=1}^n \frac{1}{k(k+1)(k+2)}$.
			
		\end{enumerate}
	\end{exo}
	\begin{corrige}
		\begin{enumerate}
			
			\item  On sait que la fraction rationnelle admet une décomposition en éléments simples de la forme
			$$\frac{a}{X}+\frac{b}{X+1}+\frac{c}{X+2}.$$
			Par les techniques usuelles (identification, multiplication par $X$ et faire $X=0$,...), on trouve
			$$\frac{1}{X(X+1)(X+2)}=\frac{1/2}X-\frac{1}{X+1}+\frac{1/2}{X+2}.$$
			
			\item  Utilisant la décomposition en éléments simples précédente, il vient 
			\begin{eqnarray*}
				S_n&=&\sum_{k=1}^n \frac{1/2}{k}-\sum_{k=1}^n\frac{1}{k+1}+\sum_{k=1}^n \frac{1/2}{k+2}\\
				&=&\sum_{k=1}^n \frac{1/2}{k}-\sum_{k=2}^{n+1}\frac{1}{k}+\sum_{k=3}^{n+2} \frac{1/2}{k}\\
				&=&\frac12+\frac14-\frac{1}2-\frac{1}{n+1}+\frac{1/2}{n+1}+\frac{1/2}{n+2}.
			\end{eqnarray*}
			On en déduit immédiatement que $(S_n)$ converge vers $\frac14$.
			
		\end{enumerate}
	\end{corrige}
	
	\begin{exo}[{Un calcul de somme}]
		Soit $P\in\mathbb R[X]$ un polynôme de degré $n\geq 1$ possédant $n$ racines distinctes $x_1,\dots,x_n$ non-nulles.
		
		\begin{enumerate}
			
			\item  Décomposer en éléments simples la fraction rationnelle $\displaystyle \frac1{XP(X)}$.
			
			\item  En déduire que $\displaystyle\sum_{k=1}^n \frac{1}{x_k P'(x_k)}=\frac{-1}{P(0)}$.
			
		\end{enumerate}
	\end{exo}
	\begin{corrige}
		\begin{enumerate}
			
			\item  On écrit la forme à priori de la décomposition en éléments simples qui est ici
			$$\frac{1}{XP(X)}=\frac{\alpha_0}{X}+\sum_{k=1}^n \frac{\alpha_k}{X-x_k}.$$
			On calcule $\alpha_0$ en multipliant tout par $X$ et en faisant tendre $X$ vers $0,$ et on trouve
			$\alpha_0=\frac{1}{P(0)}$. Pour $k\geq 1$, on multiplie tout par $X-x_k$ et on fait tendre $X$ vers $x_k$.
			On trouve cette fois
			$$\alpha_k=\lim_{x\to x_k} \frac{x-x_k}{x_k P(x)}=\lim_{x\to x_k} \frac{x-x_k}{x_k (P(x)-P(x_k))}=\frac{1}{x_k P'(x_k)}.$$
			La décomposition en éléments simples est donc 
			$$\frac{1}{XP(X)}=\sum_{k=1}^n \frac{1}{x_k P'(x_k)}\times\frac{1}{X-x_k}+\frac{1}{P(0)}\times\frac{1}{X}.$$
			
			\item  On va multiplier par $X$ cette fraction rationnelle, et on va étudier sa limite en $+\infty$. On trouve d'une part
			$$\lim_{x\to+\infty}\frac x{xP(x)}=0$$
			et d'autre part
			$$\lim_{x\to +\infty}\sum_{k=1}^n\frac{1}{x_k P'(x_k)}\times\frac{x}{x-x_k}+\frac{1}{P(0)}\times\frac{x}{x}=\sum_{k=1}^n \frac{1}{x_k P'(x_k)}+\frac1{P(0)}.$$
			Identifiant ces deux égalités, on trouve le résultat voulu!
			
		\end{enumerate}
	\end{corrige}
	
	\begin{exo}[{Somme de l'inverse des dérivées aux racines}]
		Soit $P(X)=\prod_{k=1}^{n}(X-x_k)\in\mathbb R_n[X]$ un polynôme scindé à racines simples de degré $n\geq 2$.
		
		\begin{enumerate}
			
			\item  Décomposer en éléments simples $1/P$.
			
			\item  En déduire la valeur de $\sum_{k=1}^n \frac1{P'(x_k)}$.
			
		\end{enumerate}
	\end{exo}
	\begin{corrige}
		\begin{enumerate}
			
			\item  Les racines de $P$ étant simples, on a 
			$$\frac1{P(X)}=\sum_{k=1}^n \frac{\lambda_k}{X-x_k}.$$
			De plus, pour tout $k=1,\dots,n$,
			$$\lambda_k=\lim_{x\to x_k}\frac{x-x_k}{P(x)}=\frac 1{P'(x_k)}.$$
			
			\item  Multipliant par $X$, on a pour tout $x\in\mathbb R$,
			$$\frac x{P(x)}=\sum_{k=1}^n \frac 1{P'(x_k)}\times \frac x{x-x_k}.$$
			Faisons tendre $x$ vers $+\infty$. Puisque $\deg(P)\geq 2$, on obtient
			$$\sum_{k=1}^n \frac1{P'(x_k)}=0.$$
			
		\end{enumerate}
	\end{corrige}
	
	\begin{exo}[{Polynôme et dérivé}]
		\begin{enumerate}
			
			\item  Décomposer en éléments simples la fraction $\frac{P'}P$, où $P$ est un polynôme de $\mathbb C[X]$.
			
			\item  En déduire les polynômes $P\in\mathbb C[X]$ tels que $P'|P$.
			
		\end{enumerate}
	\end{exo}
	\begin{corrige}
		\begin{enumerate}
			
			\item  On va étudier séparément les parties polaires relatives à chaque racine. Soit donc $a$ une racine de $P$,
			de multiplicité $m$. Alors on peut factoriser $P$ en $P(X)=(X-a)^m Q(X)$, soit en dérivant 
			$P'(X)=m(X-a)^{m-1}Q(X)+(X-a)^m Q'(X)$. On trouve alors
			$$\frac{P'(X)}{P(X)}=\frac{m}{X-a}+\frac{Q'(X)}{Q(X)}.$$
			Or, $a$ n'est pas racine de $Q$, donc $Q'/Q$ n'admet pas $a$ pour pôle et $\frac{m}{X-a}$ est la partie
			polaire de $P'/P$ relative à $a$. En résumé, si $P(X)=\lambda(X-a_1)^{m_1}\dots (X-a_p)^{m_p}$, alors on trouve
			$$\frac{P'}{P}=\sum_{i=1}^p \frac{m_i}{X-a_i}.$$
			
			\item  Si $P'|P$, alors $P(X)=\lambda(X-a)P'$ (car $\deg P'=\deg P-1$) et $P'/P=\frac{1/\lambda}{X-a}$. 
			Par unicité de la décomposition en éléments simples et le résultat de la question précédente, 
			ceci entraîne que $a$ est l'unique racine de $P$, et
			donc que $P(X)=\lambda (X-a)^m$. Réciproquement, les polynômes de cette forme sont tels que $P'|P$.
			
		\end{enumerate}
	\end{corrige}
	
	\begin{exo}[{Enveloppe convexe des zéros}]
		Soit $P\in\mathbb C_n[X]$ admettant $n$ racines simples $\alpha_1,\dots,\alpha_n$.
		Soient $A_1,\dots,A_n$ les points du plan complexe d'affixe respectives $\alpha_1,\dots,\alpha_n$.
		
		\begin{enumerate}
			
			\item  Décomposer la fraction rationnelle $P'/P$ en éléments simples.
			
			\item  Soit $\beta$ une racine de $P'$, et soit $B$ son image dans le plan complexe. Déduire de la question précédente que 
			$$\sum_{j=1}^n \frac{1}{\beta-\alpha_j}=0.$$
			
			\item  En déduire que $B$ est un barycentre de la famille de points $(A_1,\dots,A_n)$, avec des coefficients positifs.
			Interpréter géométriquement cette propriété.
			
		\end{enumerate}
	\end{exo}
	\begin{corrige}
		\begin{enumerate}
			
			\item   On va étudier séparément les parties polaires relatives à chaque racine. 
			On peut factoriser $P$ en $P(X)=(X-\alpha_j) Q(X)$, soit en dérivant 
			$P'(X)=Q(X)+(X-\alpha_j) Q'(X)$. On trouve alors
			$$\frac{P'(X)}{P(X)}=\frac{1}{X-\alpha_j}+\frac{Q'(X)}{Q(X)}.$$
			Or, $\alpha_j$ n'est pas racine de $Q$, donc $Q'/Q$ n'admet pas $\alpha_j$ pour pôle et $\frac{1}{X-\alpha_j}$ est la partie
			polaire de $P'/P$ relative à $a$. En résumé, la décomposition en éléments simples recherchée est
			$$\frac{P'}{P}=\sum_{j=1}^p \frac{1}{X-\alpha_j}.$$
			
			\item  Il suffit d'évaluer l'équation précédente en $\beta$.
			
			\item  On multiplie par la quantité conjuguée et on trouve
			$$\sum_{j=1}^n \frac{\bar\beta-\bar\alpha_j}{|\beta-\alpha_j|^2}=0.$$
			Prenant le conjugué de cette expression, on trouve :
			$$\left(\sum_{j=1}^n \frac{1}{|\beta-\alpha_j|^2}\right)\beta=\sum_{j=1}^n \frac{1}{|\beta-\alpha_j|^2}\alpha_j,$$
			ce qui correspond bien au résultat souhaité. On vient donc de prouver que toute racine de $P'$ est dans l'enveloppe convexe des racines de $P$.
			Ce résultat s'appelle le théorème de Lucas, il est aussi valide si les racines de $P$ ne sont pas simples. La preuve
			est similaire, si ce n'est que la décomposition en éléments simples de $P'/P$ est plus difficile à obtenir.
			C'est un bon exercice!
			
		\end{enumerate}
	\end{corrige}
	
	\subsection*{Exercices complémentaires}
	
	\begin{exo}[{Division euclidienne}]\ \textcolor{red!70!black}{$\star$}
		Effectuer la division euclidienne de $A=X^5+3X^3-X+2$ par $B=X^2+1$ dans $\R[X]$, puis dans $\C[X]$. Que remarque-t-on ?
	\end{exo}
	\begin{corrige}
		$X^5+3X^3-X+2=(X^2+1)(X^3+2X)+(-3X+2)$. Vérification : $(X^2+1)(X^3+2X)=X^5+3X^3+2X$. Le quotient est $X^3+2X$ et le reste $-3X+2$, dans $\R[X]$ comme dans $\C[X]$ : par unicité de la division euclidienne, le résultat ne dépend pas du corps dans lequel on se place.
	\end{corrige}
	
	\begin{exo}[{Racines multiples}]\ \textcolor{red!70!black}{$\star$}
		Pour quelles valeurs de $a\in\R$ le polynôme $P=X^3-3X+a$ admet-il une racine double ? Factoriser alors $P$.
	\end{exo}
	\begin{corrige}
		$P$ a une racine double $x$ si et seulement si $P(x)=P'(x)=0$. Or $P'=3X^2-3$ s'annule en $\pm1$. $P(1)=a-2=0$ si $a=2$ ; $P(-1)=a+2=0$ si $a=-2$. Pour $a=2$ : $P=(X-1)^2(X+2)$. Pour $a=-2$ : $P=(X+1)^2(X-2)$.
	\end{corrige}
	
	\begin{exo}[{Factorisations}]\ \textcolor{red!70!black}{$\star$}
		Factoriser $X^4+1$ dans $\C[X]$ puis dans $\R[X]$. Même question pour $X^6-1$.
	\end{exo}
	\begin{corrige}
		$X^4+1=(X-\omega)(X-i\omega)(X+\omega)(X+i\omega)$ avec $\omega=e^{i\pi/4}=\frac{1+i}{\sqrt2}$. Dans $\R[X]$, en regroupant les racines conjuguées : $X^4+1=(X^2+1)^2-2X^2=(X^2+\sqrt2X+1)(X^2-\sqrt2X+1)$.
		
		$X^6-1=\prod_{k=0}^5\left(X-e^{ik\pi/3}\right)$ dans $\C[X]$ ; dans $\R[X]$ : $X^6-1=(X^3-1)(X^3+1)=(X-1)(X+1)(X^2+X+1)(X^2-X+1)$.
	\end{corrige}
	
	\begin{exo}[{Éléments simples}]\ \textcolor{red!70!black}{$\star$}
		Décomposer en éléments simples dans $\R(X)$ :
		\begin{enumerate}
			\item $\dfrac{1}{X^2-1}$ ;
			\item $\dfrac{1}{X(X^2+1)}$ ;
			\item $\dfrac{X^3+1}{X(X-1)^2}$.
		\end{enumerate}
	\end{exo}
	\begin{corrige}
		\begin{enumerate}
			\item $\dfrac{1}{X^2-1}=\dfrac12\left(\dfrac1{X-1}-\dfrac1{X+1}\right)$.
			\item $\dfrac{1}{X(X^2+1)}=\dfrac1X-\dfrac{X}{X^2+1}$.
			\item Division : $X^3+1=1\cdot X(X-1)^2+(2X^2-X+1)$, donc $\dfrac{X^3+1}{X(X-1)^2}=1+\dfrac aX+\dfrac b{X-1}+\dfrac c{(X-1)^2}$ avec $a=1$ (en $0$), $c=2$ (en $1$) et $a+b=2$ (coefficient de $X^2$), donc $b=1$ : \quad $1+\dfrac1X+\dfrac1{X-1}+\dfrac2{(X-1)^2}$.
		\end{enumerate}
	\end{corrige}
	
\end{document}
